\documentclass[11pt,a4paper]{article}
\usepackage[margin=25mm]{geometry}
\usepackage{amsmath,amssymb,amsthm,mathtools}
\usepackage{booktabs,longtable,array,enumitem}
\usepackage[expansion=false]{microtype}
\usepackage{xcolor}
\usepackage[colorlinks=true,linkcolor=blue!45!black,citecolor=blue!45!black,urlcolor=blue!45!black]{hyperref}
\usepackage{fancyhdr}
\pdfmapfile{+cm.map}
\pdfmapfile{+cmextra.map}
\pdfmapfile{+symbols.map}
\pdfmapfile{+latxfont.map}
\pdfmapfile{+euler.map}
\setlength{\emergencystretch}{3em}
\setlength{\parskip}{4pt}
\setlength{\parindent}{0pt}
\setlist{itemsep=3pt,topsep=5pt,leftmargin=*}
\pagestyle{fancy}\fancyhf{}\fancyhead[L]{\small GCT\_NSRH: checked reductions and open gates}\fancyfoot[C]{\thepage}
\setlength{\headheight}{14pt}
\newcommand{\Q}{\mathbb Q}\newcommand{\Z}{\mathbb Z}\newcommand{\C}{\mathbb C}\newcommand{\N}{\mathbb N}\newcommand{\F}{\mathbb F}
\newcommand{\R}{\mathbb R}
\newcommand{\A}{\mathbb A}\newcommand{\PP}{\mathbb P}\newcommand{\GL}{\mathrm{GL}}
\newcommand{\Hom}{\operatorname{Hom}}\newcommand{\Spec}{\operatorname{Spec}}\newcommand{\Tr}{\operatorname{Tr}}\newcommand{\rank}{\operatorname{rank}}\newcommand{\Sym}{\operatorname{Sym}}\newcommand{\Res}{\operatorname{Res}}\newcommand{\IC}{\operatorname{IC}}\newcommand{\Gr}{\operatorname{Gr}}\newcommand{\Rep}{\operatorname{Rep}}\newcommand{\Frob}{\operatorname{Frob}}\newcommand{\End}{\operatorname{End}}
\newcommand{\id}{\operatorname{id}}\newcommand{\et}{\mathrm{\acute et}}\newcommand{\boxt}{\boxtimes}
\newcommand{\mathscr}{\mathcal}
\renewcommand{\S}{\ensuremath{\mathsection}}
\renewcommand{\textbullet}{\ensuremath{\bullet}}
\newcommand{\fid}[1]{\textsf{#1}}\newcommand{\gate}[1]{\textsf{#1}}
\newtheoremstyle{named}{8pt}{5pt}{\itshape}{} {\bfseries}{.}{.5em}{}
\theoremstyle{named}\newtheorem{finding}{Finding}[section]
\theoremstyle{definition}\newtheorem{definition}[finding]{Definition}
\theoremstyle{remark}\newtheorem{remark}[finding]{Remark}
\title{Purity from path sums and the nonstandard GCT positivity problem\\[4pt]\large Checked reductions, explicit barriers, and the missing realization}
\author{Research run \texttt{GCT\_NSRH\_20261121}}
\date{21 November 2026 (Europe/Berlin)}
\begin{document}\maketitle
\begin{abstract}
We investigate whether Alp\"oge's stipulated path-quadric sign mechanism transfers to nonstandard geometric complexity theory. Iterated Parseval, finite spectral visibility and unbounded tensor powers give an upper weight bound; specified duality gives equality. We extend the criterion to signed Euler multiplicities and a controlled countable spectrum. An explicit normalized Gaussian construction now reproduces every effective finite coefficient trace, gives actual opposite-sign path cohomology, and embeds its End/Hom data fully faithfully in perverse sheaves. For that construction, however, the two pointwise bounds are equivalent to the coefficient-purity assertion itself.

The supplied GCT IV, VI, VII and VIII versions propose a nonstandard-RH program without specifying finite-field coefficient objects or weights. We locate their exact algebraic targets and test concrete realizations. In rank three, an explicit finite-flat center and regular matrices retain the actual algebra, its 64-dimensional tensor module, both actions, its bar and a positive joint basis. Relative smooth and coherent projective-line models retain the collision data. A Cayley construction and its Gaussian image give exact pure Weil and perverse End/Hom realizations at fixed parameters; a coherent Gaussian convolution carries the actual pure tensor maps through all degrees.

The stronger comparisons have proved limits. The normalized all-extension word trace has no finite Weil-spectrum realization; its stack realization has unbounded dual spectrum. Algebraic spectral retuning cannot retain the actual whole relative collision order as its full regular centralizer. Under parameter retuning, a fixed composition-preserving trace comparison already fails on two actual generator words. A concrete commutant projector has an unavoidable pole in VIII's designated global form. Positive adjoint realizations kill the collision radical. Every unnormalized whole-tensor word trace in this rank-three model nevertheless has an explicit projective motive with hard Lefschetz; a specified summand obstructs the stronger positive representation-kernel model. A concrete graded-kernel category realizes the actual rank-three relations and split Grothendieck algebra, with a perverse arithmetic lift; its summand completion and adjoints require a larger category, so this does not identify the designated canonical basis.

Ordinary and motivic Satake supply pure Tate classical multiplicity objects without selecting the designated nonstandard basis. Primitive Grassmannian motives and compressed witnesses give two-row rectangular results. A known Kahle--Micha{\l}ek theorem refutes the attached homogeneous Ehrhart PH1; separate encoding obstructions delimit replacement complexity claims. Every checked finding has its full proof and a separate Grothendieck-framework translation with exact source and archival locators. Independently bounded local coefficient systems, designated arithmetic and basis comparisons, and strict orbit-closure multiplicity inequalities remain open.
\end{abstract}
\begin{center}\fbox{\parbox{.91\linewidth}{\textbf{Research status: work in progress.} The conditional theorems and counterexamples below are proved and checked. No nonstandard RH for the designated GCT objects, general canonical-basis positivity, general positive plethysm rule, or complexity separation is proved. The final part records the exact open gates.}}
\end{center}
\newpage\setcounter{tocdepth}{1}{\small\tableofcontents}\newpage

\part{The exact question and its sources}
\section{Scope, premises, and conventions}\label{sec:scope}
The main question is a transfer question. We accept the RH and Dirichlet-GRH results in the supplied Alp\"oge material as stipulated premises, as requested. We do not re-referee those results. Their acceptance does not supply a theorem relating an automorphic or finite-field trace system to a nonstandard quantum canonical basis.

The relevant finite-field description is in \cite[printed p.~3, footnote~3, attached to the paragraph before the timeline; also step~3]{Alpoge}. Timeline step~5 itself concerns a von Mangoldt functional and continues onto p.~4. We use ``the step-5 mechanism'' in the prompt's sense, but retain this more precise locator. The note names a path-quadric identity, middle-cohomological signs, iterated Parseval, and the absence of cancellation of inverse roots. It does not state a standalone theorem with all the geometric hypotheses used below.

The supplied \texttt{ALPOGE\_Main.lean} consists of five exported wrappers, including the RH wrapper \texttt{riemannHypothesis}. Its body is \texttt{rhOfAll}; the imported module is \texttt{Supplements.Kernels.\allowbreak MainTheoremsAll}. The supplied \texttt{ALPOGE\_Statements.lean} defines the zeta/Dirichlet zero and nonvanishing propositions. Neither file contains a nonstandard quantum-group, path-concentration, or theta/GCT transfer theorem. The imported development is not in this ZIP. The supplied scan records zero occurrences for the listed prohibited constructs; this run did not repeat that scan or execute the missing development. These are access and locator facts, not objections to the stipulated theorem. No Masullo result is used: none is identified or supplied in this prompt.

We write $v$ for the \emph{formal quantum parameter}, $q$ for a fixed prime power, and $Q=q^n$ for an extension-field cardinality. In the GCT sources the quantum parameter is usually called $q$. No substitution $v=Q$ or $v=Q^{1/2}$ constructs a Weil structure. All weights in the purity conclusion are integers. The Frobenius convention is geometric Frobenius, so $\Q_\ell(-1)$ has eigenvalue $q$ over $\F_q$.

``Finding'' means a deduction established for this investigation, including elementary lemmas in the precise form needed. It is not a claim of priority over the mathematical literature. Known research results are cited, not presented as new findings. A source statement checked in context is not thereby an independently checked proof of the entire cited theorem.

\section{What IV, VI, VII and VIII actually specify}\label{sec:sources}
\subsection{Versions and the missing RH statement}
The attached IV is the revised Blasiak--Mulmuley--Sohoni paper, arXiv cs/0703110v4 (2013), not merely the 2007 version referred to by VII and VIII. It proves nonstandard Schur--Weyl duality and the two-row crystal result. VI is arXiv 0704.0229v4 (2009); VII and VIII are 0709.0749v2 and 0709.0751v2 (2008). Rendered title-page dates produced by an undated \verb|\maketitle| are not used as version dates.

\textbf{Source finding A01.} In these four supplied versions there is no formulated nonstandard RH specifying finite-field objects, Frobenius actions, cohomological degrees and weights. This is an exact-source conclusion, not a theorem about what other work might contain. The relevant locators are:

{\small\begin{longtable}{p{.13\linewidth}p{.35\linewidth}p{.44\linewidth}}
\toprule Work&Exact locator&Content and status\\\midrule\endhead
IV&\S1.1, Problem 1.1; \S\S1.2--1.5; \S19&Positive-rule and basis program; the revised text has no finite-field purity statement.\\
VI&Introduction, Figure 1.1; \S1.7, Observation 1.7.1 and Figure 1.2&The proposed RH/extensions connection is marked by question marks; standard canonical positivity motivates the program.\\
VI&\S4.2, Proposition 4.2.1; \S4.2.2, Question 4.2.6 (PH0??); \S4.3&Asks for compatible canonical bases and analogous positivity. No nonstandard weight theorem is stated.\\
VII&\S1 after Conjectures 1.3--1.4; \S4.1, equations (9)--(10), Theorem 4.4; \S4.2, Theorem 4.5&Nonstandard minors and Laplace coefficients need interpretations analogous to perverse-sheaf constructions.\\
VIII&\S1, pp.~2--3; \S2.3 after Conjecture 2.5&An extension of perverse sheaves and finite-field RH is proposed as a way to understand minors and globalization.\\
VIII&\S2.4.2 after Conjecture 2.7; \S2.4.3 after Conjecture 2.9&Topological interpretations, Hard Lefschetz, and a theory for noncommutative, nonflat objects are suggested; the general case includes correction terms.\\\bottomrule
\end{longtable}}
These statements are supported by complete reading of their contexts and matching numbered PDF statements, not merely a search for the phrase ``Riemann hypothesis.'' Their appropriate conclusion is that a precise realization/purity package must first be specified.

\subsection{The algebraic objects}
Let $H$ be a complex connected classical reductive group, $X=V_\mu(H)$, $G=\GL(X)$, and $\rho:H\to G$ the representation homomorphism. The generalized plethysm multiplicity is
\begin{equation}\label{eq:plethysm}
a_{\lambda,\mu}^{\pi}=\dim\Hom_H\bigl(V_\pi,\Res_\rho V_\lambda(G)\bigr).
\end{equation}
For Kronecker coefficients, $H=\GL(V)\times\GL(W)$, $X=V\otimes W$, and $\rho(g,h)=g\otimes h$.

VII, Problem~1.1 and Theorems~1.2, 2.1 and 3.1, constructs a nonstandard matrix-coordinate bialgebra
\[
 S=\mathcal O(M_v^H(X)),\qquad \mathcal A_r^H=(S_r)^*,
\]
with a quantized map from $H_v$ to the associated nonstandard quantum object. Its matrix generators satisfy the projector relation
\[
 P^{+,H}_{X,X}(\mathbf u\otimes\mathbf u)
 =(\mathbf u\otimes\mathbf u)P^{+,H}_{X,X}
\]
(VII, \S2, equations (4)--(6)). The projector is constructed from the braiding on $X_v\otimes X_v$. The determinant may vanish, so the strict term may be quantum semigroup; an antipode is not automatic. The dual nonstandard symmetric-group deformation is $\mathcal B_r^H(v)$ (VII, \S5).

These are noncommutative algebras over fields containing $\Q(v)$, sometimes an algebraic extension. Their bar involution is not arithmetic Frobenius. Complete reducibility or a compact real form is not a Weil structure. Nonflatness and noncommutativity have different roles: nonflat commutative schemes still have \(\ell\)-adic sheaves, whereas a noncommutative coordinate algebra is not the coordinate ring of a scheme in the EGA/SGA sense.

General plethysm duality and reciprocity are VII, Conjectures~1.3--1.4 and refined Conjecture~6.1. They would give
\begin{equation}\label{eq:comparison}
a_{\lambda,\mu}^{\pi}=\sum_\alpha m_\lambda^\alpha n_\pi^\alpha,
\end{equation}
where $n_\pi^\alpha$ is an $H_v$ constituent multiplicity and $m_\lambda^\alpha$ is a Specht multiplicity after specialization at $v=1$. In revised IV, Theorems~1.5/12.1 establish the Kronecker duality, and \S1.4, equation~(11), gives the corresponding Kronecker decomposition. It would be incorrect to leave that particular comparison conjectural merely because VII predates the revision.

\subsection{The actual positivity targets}\label{sec:targets}
The basis constructions themselves have separate prerequisites. VIII, Conjecture~2.2, concerns compatible local crystal bases; Question~2.3's ordinary balanced-triple condition fails in examples; Conjecture~2.5 proposes a unique minimum-degree globalization. Conjectures~2.6 and 3.3 concern cells, Conjectures~3.1--3.2 the dual globalization, and Conjecture~3.7 the \emph{opposite} basis's quasi-subcells and specialization. VIII's Proposition~2.4 contains an omitted proof and assumes the local-crystal conjecture; we do not record it as independently proved.

\paragraph{Kronecker, coordinate side.}
VIII, Conjecture~2.7, asks that expansion coefficients of $CB(\mathcal A_r^H)$ in the external canonical basis $CB(Z)$, $Z=\End(X_v^{\otimes r})$, be polynomials with nonnegative \emph{rational} coefficients (equation~(9)). Multiplication and comultiplication constants have the form
\begin{equation}\label{eq:kroneckerpositive}
 \pm(v-v^{-1})^a f(v),\qquad a\in\Z_{\ge0},
\end{equation}
where $f$ is bar-invariant, coefficientwise nonnegative rational and unimodal. The same is proposed for $CB(S)$.

\paragraph{Kronecker, dual side.}
VIII, Conjecture~3.4, applies to structure constants and external expansion coefficients of $E(r)$ and $E^{\mathrm{opp}}(r)$. Its form is
\begin{equation}\label{eq:dualpositive}
 \pm(v^{1/2}-v^{-1/2})^a f(v),
\end{equation}
with $f$ nonnegative rational, bar-invariant and unimodal in the half-power variable. Thus raw coefficientwise positivity is not the stated assertion even in the Kronecker case.

\paragraph{General plethysm.}
VIII, Conjectures~2.8 and 3.5, concern normalized coefficients of minimal polynomials of algebraic structure/expansion coefficients. The displayed normalizations are equations~(10) and (14):
\[
 (-1)^e(v-v^{-1})^{e'}c(v),\qquad
 (-1)^e(v^{1/2}-v^{-1/2})^{e'}g(v).
\]
The nonzero normalized polynomial is required to be positive at $v=1$ (called saturated there). Conjectures~2.9 and 3.6 ask for a dominant positive unimodal part and a small correction, for example
\[
 c=c^0+c^1,\qquad
 \frac{\|c^1\|_2}{\|c^0\|_2}
 \le\frac1{\operatorname{poly}(\langle\mu\rangle,\langle\pi\rangle,\langle r\rangle)},
 \qquad c^0(1)\gg |c^1(1)|.
\]
Conjecture~3.6 is explicitly informal. A theorem target must fix the quantitative meaning of the polynomial and the asymptotic comparison. VIII, \S5.4.1, Figure~30, displays negative entries in a normalized example; these are cited source computations, not a recomputation of the whole algebra in this run.

There is also a sign-convention issue in ``coefficients of the minimal polynomial'': the usual monic polynomial of 1 is $T-1$, while the source later refers to trace and norm without the conventional sign on trace. We record this ambiguity rather than use it to disprove the intended conjecture. ``Frobenius automorphisms'' in VIII, \S5.4.1, refers to algebraic conjugations over $\Q(v)$, not to finite-field Weil Frobenius.

\subsection{Positivity hypotheses in the flip}
VI, Hypothesis~1.6.4 (PH1), is substantially stronger than numerical nonnegativity. For each triple there should be a rational polytope $P_{\lambda,\mu}^{\pi}$ whose Ehrhart function is
\[
 \widetilde a_{\lambda,\mu}^{\pi}(t)=a_{t\lambda,\mu}^{t\pi},
\]
with $\mu$ fixed (VI, \S1.6, definitions preceding Theorem~1.6.1). It must admit a description $Ax\le b$ with $A$ independent of $\lambda,\pi$ and $b$ homogeneous linear in them; ambient dimension and entry bitlengths are polynomially bounded as specified there; membership and a separating hyperplane are computable in polynomial time. The number of inequalities need not be polynomial.


\paragraph{The attached homogeneous PH1 is already false.}
Kahle--Micha{\l}ek \cite[Theorem~5 and Remark~6, p.~4]{KM}
prove that the fixed-inner-weight stretching function
\[
 \phi(s)=a_{(4s),(3)}^{(7s,5s)}
       =\frac{s+r(s)}3,\qquad
 (r(0),\ldots,r(5))=(3,-1,1,0,2,-2),
\]
with $r$ periodic modulo six, is not the Ehrhart function of a
rational polytope. Its degree is one, but $\phi(1)=0$ and
$-\phi(-1)=1$, contradicting the interior-versus-total inequality
in Ehrhart reciprocity. Their p.~2 explicitly identifies the target
as VI, Hypothesis~1.6.4. By the rectangular identification on their
p.~6, the same function is
\[
 g((4s,4s,4s),(4s,4s,4s),(7s,5s)).
\]
This is a cited prior refutation, not a new finding of this run.
Their pp.~3--4 distinguish the later 2011 formulation allowing
inhomogeneous families $Ax\le sb+c$. The cited theorem does not
exclude every dimension of that later formulation. We therefore
do not leave the old homogeneous assertion as an open gate, or
equate its failure with a refutation of every later positivity
program. Finding~\fid{F34} separately tests the compact dimension
requirement on a family with a genuine positive simplex model.

VI, Hypothesis~1.6.5 (SH), asks for a polynomial bound on the saturation index in $\dim G+\operatorname{ht}\lambda+\operatorname{ht}\pi$ and asymptotic generic strict saturation. Hypothesis~1.6.6 (PH2) bounds the positivity index; Hypothesis~1.6.7 (PH3) bounds the modular index of a positive rational Hilbert-series form. Strict saturation in every case is not the attached revised hypothesis.

The conditional implication to relaxed PHflip is VI, Theorem~1.6.8. A printed discrepancy must be retained: its part~(a) displays $a_{b\lambda,b\mu}^{b\pi}$ although the stretching definition and its stated proof through PH1 keep $\mu$ fixed. We use only the fixed-$\mu$ consequence directly justified by the displayed Ehrhart model and record the all-three-scaling formulation as requiring clarification. This editorial point is not needed for our main obstruction. More fundamentally, VIII, \S6, explicitly asks for additional effective-label properties. An abstract canonical basis or an abstract positive count does not prove PH1, SH, or the existence of a lower-bound obstruction.

\section{The Kazhdan--Lusztig template, with weights fixed}\label{sec:KL}
Let $G/\F_Q$ be split reductive, $B$ a Borel subgroup, $C_w=BwB/B$ and $X_w=\overline{C_w}$. Put
\[
 \IC_w^u=j_{!*}\Q_\ell[\ell(w)][-\ell(w)]
\]
with the usual extension to $X_w$; the notation records the unshifted convention, constant in degree zero on $C_w$. Then
\begin{equation}\label{eq:KL}
 P_{x,w}(T)=\sum_{j\ge0}\dim\mathcal H^{2j}(\IC_w^u)_x\,T^j.
\end{equation}
Odd stalk groups vanish, and each eigenvalue in degree $2j$ is $Q^j$. This is the Tate/parity statement, stronger than an absolute-value bound. Exact primary locators for this form are Haines \cite[Theorems 2.0.1, 5.0.4; Proposition 3.0.2; \S4]{Haines}. The original KL geometric identification is \cite[Theorem 4.3]{KL}.

In the modern proof, a Bott--Samelson resolution is proper, its fibers admit affine pavings, and the decomposition theorem isolates the intersection complex. The affine paving provides even/Tate cohomology. The compatibility with Weil structures must be checked: a geometric splitting need not itself be Frobenius-equivariant. Haines \S4 discusses precisely this issue. His introduction also notes that the original KL proof does not require the decomposition theorem; the latter is the requested modern template, not a claim about historical logical necessity.

Deligne's upper-weight theorem is \cite[Theorem 3.3.1; Corollaries 3.3.4--3.3.6]{WeilII}. BBD intermediate extension preserves purity (5.3.2), pure perverse sheaves are geometrically semisimple (5.3.8), and pure complexes split geometrically (5.4.5); relative Hard Lefschetz is 5.4.10 \cite{BBD}. None of these statements identifies an arbitrary pure trace with the dimension polynomial \eqref{eq:KL}. Arithmetic Frobenius semisimplicity is distinct: BBD, pp.~138--139, explicitly discusses pure unipotent Jordan-block sheaves.

{\small\begin{longtable}{p{.25\linewidth}p{.29\linewidth}p{.38\linewidth}}
\toprule Standard ingredient&Nonstandard datum supplied&Unfilled role\\\midrule\endhead
Flag schemes, Schubert strata&$H_v$, nonstandard FRT bialgebra; support-defined cells&An actual space/category with those strata and stalk functors.\\
Distinguished $\IC_w$ and its normalized class&Candidate minimum-degree bases&Objects whose classes are the \emph{specified} bases; existence and uniqueness.\\
Weil structure and closed-point Frobenius&Algebraic bar involution&Arithmetic models, traces, degrees and integer weights.\\
Proper resolutions/convolution&Algebra multiplication and comultiplication&Correspondences or functors with a decomposition theorem.\\
Tate parity and effective multiplicities&Signed rational forms in VIII&Positive multiplicity objects, signs, factors and denominators.\\
Verdier duality and Hard Lefschetz&Bar symmetry/unimodality conjectures&Compatible duality and Lefschetz isomorphisms.\\
Crystal/path labels&Conjectural cells; proved two-row tableaux&Effective labels, specialization and polynomial-time verification in general.\\\bottomrule
\end{longtable}}
For standard quantum groups, quiver and flag constructions supply related geometric realizations; it would be misleading to assert that every quantum canonical-basis constant is a stalk of one Schubert variety.

\part{Checked findings}
\section{The sign mechanism and its exact strength}\label{sec:main}
This section contains the main proved reduction. Its hypotheses explicitly retain the information that the phrase ``middle cohomology'' compresses.

\begin{finding}[\fid{F01}: path identity and Parseval estimate]\label{f:path}
Let $V=\F_Q^d$, $d\ge1$, with a nondegenerate bilinear pairing, and let $\psi$ be a nontrivial additive character. For $k\ge2$ and any function $t:V\to\C$, put
\[
 M=\sum_{x\in V}t(x),\quad
 P_k(x_1,\ldots,x_k)=\sum_{i=1}^{k-1}\langle x_i,x_{i+1}\rangle,
 \quad N_k=\sum_{P_k=0}\prod_{i=1}^k t(x_i).
\]
Then $QN_k=M^k+R_k$, where
\[
 R_k=\sum_{h\in\F_Q^\times}\sum_{x_1,\ldots,x_k\in V}
       \prod_i t(x_i)\,\psi(hP_k),
\]
and
\begin{equation}\label{eq:parseval}
 |R_k|\le (Q-1)Q^{d(k-1)/2}\|t\|_2^2\|t\|_\infty^{k-2}.
\end{equation}
No nonnegativity assumption on $t$ is required.
\end{finding}
\begin{proof}
Orthogonality gives $\sum_h\psi(hz)=Q\,1_{z=0}$; separating $h=0$ gives the identity. For $h\ne0$, set $K_h(x,y)=\psi(h\langle x,y\rangle)$ and $D=\operatorname{diag}(t)$. Nondegeneracy gives
\[
 (K_hK_h^*)_{x,x'}=\sum_y\psi(h\langle x-x',y\rangle)
 =Q^d1_{x=x'},
\]
so $\|K_h\|=Q^{d/2}$. The inner sum is
$t^{\mathsf T}K_hD K_h\cdots D K_ht$, with $k-1$ kernels and $k-2$ diagonal factors. The norm of the row functional $t^{\mathsf T}$ is $\|t\|_2$, even for complex $t$ (it is the Hermitian dual of $\bar t$). Cauchy--Schwarz and submultiplicativity give the bound for one $h$; sum the $Q-1$ terms. For $k=2$ there are no diagonal factors.
\end{proof}
This refines Fable suggestion~3's scalar toy in the exact form needed here. It is an open-path matrix coefficient, not in general a trace moment: tracing the power closes the path into a cycle.

\begin{finding}[\fid{F02}: finite spectral no-cancellation]\label{f:spectral}
For distinct nonzero complex numbers $\alpha_j$ and nonzero constants $c_j$, let $p_n=\sum_jc_j\alpha_j^n$. For a nonempty finite spectrum,
\[
 \limsup_{n\to\infty}|p_n|^{1/n}=\max_j|\alpha_j|.
\]
Consequently $|p_n|\le Cn^A R^n$ for all $n\ge1$, $R>0$, implies $|\alpha_j|\le R$. The same holds for any subexponential prefactor.
\end{finding}
\begin{proof}
Near zero,
\[
 \sum_{n\ge1}p_nz^n=\sum_j\frac{c_j\alpha_jz}{1-\alpha_jz}.
\]
Every $1/\alpha_j$ is a genuine pole, since the other summands are regular there. The radius of convergence is $1/\max|\alpha_j|$. Cauchy--Hadamard proves the equality; the bound gives radius at least $1/R$. The empty case is vacuous. Coincident spectral values must first be grouped; the conclusion concerns their \emph{nonzero} total coefficients.
\end{proof}
Positive multiplicities are one sufficient way to ensure that regrouping does not erase a root. Eigenvalues themselves need not be positive, or real. The traces of powers do not see nilpotent Jordan off-diagonal entries, so semisimplicity is not required for this lemma.

\begin{finding}[\fid{F03}: tensor-path purity criterion]\label{f:criterion}
Fix a prime power $q>1$, an integer $a$, $d\ge1$, and $B>0$. For every $n\ge1$ let $Q=q^n$ and $t_n:\F_Q^d\to\C$ satisfy $|t_n(x)|\le BQ^{a/2}$. Form $M_n,N_{k,n},R_{k,n}$ as in \fid{F01}. Assume:
\begin{enumerate}[label=(S\arabic*)]
\item A fixed finite multiset $\mathcal A$ of nonzero algebraic numbers, with positive integer multiplicities, and $\epsilon\in\{1,-1\}$ satisfy
\[
 M_n=\epsilon\sum_{\alpha\in\mathcal A}\alpha^n\qquad(n\ge1).
\]
\item For an unbounded set of integers $k\ge2$ there is, for each $k$, a fixed finite multiset $\mathcal B_k$ of nonzero algebraic numbers, possibly empty, with positive integer multiplicities such that
\[
 q^nN_{k,n}=-\epsilon^k\sum_{\beta\in\mathcal B_k}\beta^n\qquad(n\ge1).
\]
\end{enumerate}
Then $|\alpha|\le q^{(a+d)/2}$ for every $\alpha\in\mathcal A$.

For the all-conjugate assertion, require for every embedding $\sigma$ of the algebraic spectral data a corresponding trace system $t_{n,\sigma}$ satisfying these hypotheses with spectra $\sigma(\mathcal A),\sigma(\mathcal B_k)$ and a constant $B_\sigma$ independent of $n$. If in addition
\begin{equation}\label{eq:duality}
 \mathcal A=q^{a+d}/\mathcal A
 \quad\text{as multisets},
\end{equation}
then every eigenvalue has every complex absolute value $q^{(a+d)/2}$: the spectrum is pure of weight $a+d$.
\end{finding}
\begin{proof}
The pointwise bound gives $\|t_n\|_2^2\le B^2Q^{a+d}$. Equation~\eqref{eq:parseval} therefore gives
\begin{equation}\label{eq:overhead}
 |R_{k,n}|\le B^kQ^{k(a+d)/2+d/2+1}.
\end{equation}
By the exact path identity and (S1)--(S2),
\begin{equation}\label{eq:positivespectrum}
 -\epsilon^kR_{k,n}=
 \sum_{\beta\in\mathcal B_k}\beta^n+
 \sum_{(\alpha_1,\ldots,\alpha_k)\in\mathcal A^k}
 (\alpha_1\cdots\alpha_k)^n.
\end{equation}
After grouping equal roots all coefficients are positive integers. In particular the diagonal tuple contributes $\alpha^k$ and cannot be cancelled. For each \emph{fixed} $k$, \fid{F02} and \eqref{eq:overhead} give
\[
 |\alpha|^k\le q^{k(a+d)/2+d/2+1},\qquad
 |\alpha|\le q^{(a+d)/2+(d/2+1)/k}.
\]
Let $k$ tend to infinity through the unbounded set. This order of limits is essential: constants such as $B^k$ disappear in the $n$th-root limit before $k$ varies. Applying the same argument to each $\sigma$ gives all-conjugate upper bounds. Since $a+d$ is an integer, $q^{a+d}$ is rational and fixed by every $\sigma$. Apply the upper bound to $q^{a+d}/\sigma(\alpha)$, which occurs by \eqref{eq:duality}, to obtain the reverse inequality.
\end{proof}
The purity conclusion is conditional on a \emph{specified Tate duality}. An early draft allowed an arbitrary algebraic dualizing scalar; the independent check found that its conjugates could have different absolute values. The final theorem uses the integer-weight scalar $q^{a+d}$ and the correction is recorded in the audit.

\begin{finding}[\fid{F04}: controlled background terms]\label{f:errors}
Replace (S2) by
\[
 q^nN_{k,n}=-\epsilon^k\sum_{\beta\in\mathcal B_k}\beta^n+E_{k,n}.
\]
Suppose $|E_{k,n}|\le C_kn^{A_k}q^{n(k(a+d)/2+\eta_k)}$ for all $n$, with $\eta_k/k\to0$ along the unbounded set of $k$. Then \fid{F03}'s upper bound, and its purity conclusion under the same all-embedding duality hypotheses, still hold, provided that the modified identity and error estimate also hold for every corresponding trace system $t_{n,\sigma}$, with constants $C_{k,\sigma},A_{k,\sigma}$ and $\eta_{k,\sigma}/k\to0$ allowed to depend on $\sigma$.
\end{finding}
\begin{proof}
The positive spectral sum in \eqref{eq:positivespectrum} is now $\epsilon^k(E_{k,n}-R_{k,n})$. Its exponential rate for fixed $k$ is at most
$q^{k(a+d)/2+\max(d/2+1,\eta_k)}$. Apply \fid{F02} and take the $k$th root. Since $\max(d/2+1,\eta_k)/k\to0$, the same limit follows. Duality supplies the reverse bound as before.
\end{proof}
This is a sufficient error budget, not a necessity theorem for every possible proof. It permits suitable removed Tate or boundary terms; it does not justify discarding uncontrolled terms.

\begin{finding}[\fid{F17}: finite extensions and a separate dual system]\label{f:extensiondual}
The upper-weight conclusion of \fid{F03}--\fid{F04} still holds if, for each allowed $k$, the opposite-sign identity and its error estimate are known only for $n=e_kr$, all $r\ge1$, where the positive integer $e_k$ may depend on $k$.

Self-duality can also be replaced by a separate dual coefficient system. Suppose the first spectrum $\mathcal A$ has the proved upper weight $a+d$, a second spectrum $\mathcal A^D$ has the proved upper weight $d-a$, and a perfect Frobenius-equivariant pairing identifies
\[
 \mathcal A^D=\{q^d/\alpha:\alpha\in\mathcal A\}
\]
with multiplicities. Then $\mathcal A$ is pure of weight $a+d$, provided both upper bounds hold under every embedding.
\end{finding}
\begin{proof}
For fixed $k$, apply the spectral lemma in the variable $r$. The positive multiset now contains $\alpha^{e_kk}$; equal powered roots still have positive multiplicities, so none disappears. Its exponential bound is
\[
 (q^{e_k})^{k(a+d)/2+d/2+1}
\]
(or the corresponding maximum with $\eta_k$). Taking the $(e_kk)$th root gives exactly the same bound on $|\alpha|$ as before, and the unbounded-$k$ limit is unchanged. Constants may depend on $k$ and $e_k$. This uses full multiples of $e_k$, not an arbitrary infinite subsequence such as the odd degrees in \fid{F05}(d).

For the second assertion apply the dual upper bound to $q^d/\alpha$:
\[
 |\alpha|\ge\frac{q^d}{q^{(d-a)/2}}=q^{(d+a)/2}.
\]
Combine it with the first upper bound. Since $q^d$ is rational, the same calculation holds under every embedding.
\end{proof}
This refinement allows a geometric construction to choose a finite field of definition separately for each path length. The upper weight $d-a$ for the dual system is an independent hypothesis: the original pointwise upper weight $a$ alone does not give the required upper estimate for the dual coefficient system. The refinement also distinguishes Verdier duality with a dual coefficient system from an assertion that a single compact-support module is self-dual. To use the latter option geometrically, the necessary compact-support/ordinary-cohomology comparison still has to be proved; it is not supplied by dimension alone.

\begin{finding}[\fid{F24}: signed Euler spectra at even path lengths]\label{f:signed}
Keep the pointwise setup and path identity of \fid{F03}, and put $w=a+d$.
Replace its input sign hypothesis by
\[
 M_n=\sum_{\alpha\in A}c_\alpha\alpha^n\quad(n\ge1),
 \qquad c_\alpha\in\mathbb R\setminus\{0\},
\]
where $A$ is finite, consists of distinct nonzero roots, and is nonempty.
Thus $A$ is the \emph{visible spectrum}, after grouping equal roots and deleting
zero coefficients. Suppose that for an unbounded set of \emph{even} $k\ge2$,
\[
 q^nN_{k,n}=-\sum_{\beta\in B_k}b_{k,\beta}\beta^n+E_{k,n},
 \qquad b_{k,\beta}>0,
\]
for every $n\ge1$, with the error bound of \fid{F04}. Then every
$\alpha\in A$ satisfies $|\alpha|\le q^{w/2}$. The positive coefficients
$b_{k,\beta}$ may be real; positive integer multiplicities suffice.

The all-embedding conclusion holds when each embedding has the corresponding
trace system, real transformed input coefficients, and error bounds. Integer
Euler multiplicities meet the coefficient condition. A reciprocal relation
$A=q^w/A$ among the visible supports, or a separately bounded visible dual
system as in \fid{F17}, supplies purity of the visible spectrum.

If the identities are available only at $n=e_kr$, the same conclusion holds
provided that the largest original radius remains visible after powering
by $e_k$, for an unbounded set of the allowed even $k$. This extra condition
cannot be omitted for signed spectra.
\end{finding}
\begin{proof}
Let $R=\max_{\alpha\in A}|\alpha|$ and
$A_R=\{\alpha:|\alpha|=R\}$. In the group algebra of the multiplicative group
generated by $A_R$, form
\[
 F=\sum_{\alpha\in A_R}c_\alpha[\alpha],\qquad
 F^k=\sum_\rho d_{k,\rho}[\rho],
\]
grouping equal products. The coefficients $d_{k,\rho}$ are real and
\[
 \sum_\rho d_{k,\rho}
   =\left(\sum_{\alpha\in A_R}c_\alpha\right)^k\ge0
\]
because $k$ is even. Moreover $F^k\ne0$. Otherwise evaluating
$[\alpha]\mapsto\alpha^n$ would give
$(\sum_{A_R}c_\alpha\alpha^n)^k=0$ for every $n\ge1$.
The Vandermonde matrix for $n=1,\ldots,|A_R|$ is invertible, contradicting
the nonzero coefficients on distinct nonzero roots. This argument uses
nonnilpotence of $F$, not a false assertion that a group algebra with
torsion has no zero divisors. Since a nonzero finite collection of
nonpositive real numbers has negative sum, at least one $d_{k,\rho}$ is
strictly positive.

Every root in $F^k$ has modulus $R^k$. A product containing a smaller-radius
root has smaller modulus, so it cannot cancel the chosen positive
coefficient. The path identity gives
\[
 M_n^k+\sum_{\beta\in B_k}b_{k,\beta}\beta^n=E_{k,n}-R_{k,n}.
\]
Adding the positive $B_k$ spectrum cannot cancel the chosen top-radius
coefficient either. The left side therefore has a surviving root of modulus
$R^k$. Its other coefficients need not be positive. Applying \fid{F02}
and \eqref{eq:overhead} for fixed $k$ gives
\[
 R^k\le q^{kw/2+\max(d/2+1,\eta_k)}.
\]
Take the $k$th root and the unbounded even-$k$ limit. The all-embedding and
dual arguments are the same as in \fid{F03} and \fid{F17}, but only visible
roots can be used in a reciprocal pairing.

For the finite-extension assertion, group first:
\[
 M_{er}=\sum_\gamma d_{e,\gamma}\gamma^r,
 \qquad d_{e,\gamma}=\sum_{\alpha^e=\gamma}c_\alpha.
\]
If the nonzero grouped spectrum is nonempty, put
$R_e=\max_{d_{e,\gamma}\ne0}|\gamma|^{1/e}$. Applying the preceding
argument in $r$ with $e=e_k$ gives
\[
 R_{e_k}\le q^{w/2+\max(d/2+1,\eta_k)/k}.
\]
Thus $R_{e_k}=R$ along an unbounded allowed set suffices. This equality is
automatic if the top-radius coefficients have one common sign, or if no
ratio of two distinct top-radius roots is a root of unity. It is not
automatic for general signed coefficients. Each embedding needs its own
top-radius visibility condition.
\end{proof}

\paragraph{Two exact limits of the refinement.}
First, take $a=0$, $d=1$, and let $t_n$ be the constant function $0$ for even
$n$ and $2$ for odd $n$. Then $|t_n|\le2$ and
\[
 M_n=q^n-(-q)^n.
\]
Both visible roots have modulus $q>q^{1/2}$. For every even $k$ choose
$e_k=2$. At every required degree $n=2r$ all counts vanish, so the
opposite-sign hypothesis holds with empty $B_k$ and zero error. The top
spectrum disappeared under powering. This refutes the earlier draft
extension without visibility, while leaving the positive-multiset
statement \fid{F17} intact. Geometrically the pointwise sequence itself
is the alternating trace of the constant line minus the unramified line
with Frobenius $-1$; its failure is not an exotic nonalgebraic sequence.

Second, the real-coefficient condition cannot simply be dropped for the
prescribed unbounded even set. Let $a=0$, $d=1$, and $t_n(x)=i$ for every
$x,n$. Then $M_n=iq^n$ has the root $q>q^{1/2}$ with coefficient $i$.
For $k=2m\ge6$ with $k\equiv2\pmod4$, put $a_j=x_{2j-1}$ and $b_j=x_{2j}$. The path
polynomial is
\[
 P_k=\sum_{j=1}^{m}a_jb_j+
       \sum_{j=1}^{m-1}a_{j+1}b_j.
\]
Its coefficient matrix is triangular with diagonal entries one, so it is
equivalent to a nondegenerate bilinear pairing. Its zero count over
$\F_Q$ is $Q^m+(Q^m-1)Q^{m-1}$, including characteristic two.
Since $i^k=-1$,
\[
 QN_{k,n}=-Q^k-Q^{k/2+1}+Q^{k/2}.
\]
Take $B_k=\{q^k,q^{k/2+1}\}$ and $E_{k,n}=Q^{k/2}$; the error has
$\eta_k=0$. Every other stated hypothesis holds for all $n$. This is a
counterexample to that deleted hypothesis, not to the integer-Euler
version or to possible other criteria with phase control.

\paragraph{What the stronger theorem removes.}
For a bounded constructible Weil complex, alternating cohomological
multiplicities are integers. The input need not be supported in one parity
if the even-path output has the prescribed negative visible spectrum and
the pointwise bound holds. Identical roots that already cancel between
degrees remain invisible: for example
$\Q_\ell(-1)\oplus\Q_\ell(-1)[-1]$ has zero trace in every extension.
Consequently \fid{F24} strictly weakens the input parity gate; it does not
infer bounds for erased cohomology or provide the path-output sign.


\subsection{Countable trace completions: an upper bound and a duality obstruction}
The finite-spectrum hypothesis has a useful extension, but its duality
requirement becomes restrictive. Call a visible weighted spectrum
$\mathcal A=\{(\alpha,c_\alpha):\alpha\in A\}$ \emph{admissible} if
$A\subset\C^*$ is bounded, its roots are distinct, every
$\{\alpha:|\alpha|\ge\epsilon\}$ is finite, all $c_\alpha\ne0$, and
for some nonnegative integer $N$,
\[
 S_N=\sum_{\alpha\in A}|c_\alpha|\,|\alpha|^N<\infty.
\]
Equal roots are grouped and zero coefficients removed before imposing
these conditions. The empty spectrum represents zero.

\begin{finding}[\fid{F32}: admissible countable spectra and visible reciprocal duality]\label{f:countable}
For a nonempty admissible spectrum, the maximum radius $R$ is attained
and $u_n=\sum c_\alpha\alpha^n$ converges absolutely for $n\ge N$.
It satisfies
\[
 \limsup_{n\to\infty}|u_n|^{1/n}=R.
\]
An eventual sequence has at most one admissible visible expansion.
Finite sums, tensor products, and integer power substitutions preserve
admissibility after grouping, although power substitution may erase the
entire largest-radius part.

The signed even-path criterion of \fid{F24} remains true when its
input visible spectrum and each positive path-output spectrum are
admissible countable spectra. The identities and bounds need hold only
for all sufficiently large extension degrees, with thresholds depending
on the path length. The input coefficients must still be real, and the
finite-extension and all-embedding versions retain the same explicit
visibility and transformed-hypothesis requirements.

If $A$ is admissible and $\kappa\ne0$, and every $\kappa/\alpha$ is
visible in a bounded dual support, then $A$ is finite. In particular,
an infinite completion of this kind cannot have bounded visible
reciprocal duality. If the nonzero grouped coefficients are integers,
absolute convergence of one nonnegative moment already implies the
boundedness and local-finiteness conditions above.
\end{finding}
\begin{proof}
Choose a root $\alpha_0$. The finite set of roots of modulus at least
$|\alpha_0|$ contains a largest one; it is a global maximum $R>0$.
For $n\ge N$, absolute summation gives $|u_n|\le S_NR^{n-N}$.
The generating function of the shifted sequence is
\[
 H_N(T)=\sum_{m\ge0}u_{N+m}T^m
       =\sum_{\alpha\in A}\frac{c_\alpha\alpha^N}{1-\alpha T}
       \quad (|T|<1/R).
\]
The equality follows from the double absolute bound
$S_N/(1-R|T|)$. Its right side converges locally normally away from
$\{1/\alpha\}$. Indeed, on a compact set with $|T|\le L$, roots
of modulus at most $1/(2\max(1,L))$ have denominator modulus at
least $1/2$, giving a normally convergent tail by $2S_N$; only
finitely many other roots remain. The pole set is locally finite.
At $T=1/\alpha$ the residue is
$-c_\alpha\alpha^{N-1}\ne0$; all other terms are holomorphic in a
neighborhood. Thus the Taylor series has radius exactly $1/R$.
Cauchy--Hadamard, and the harmless fixed shift by $N$, give the
asserted exponential rate. To prove uniqueness, subtract two proposed
expansions at a common sufficiently large moment index. Their grouped
difference is admissible. If nonempty its exponential rate is positive,
contradicting an eventually zero sequence.

For a tensor power of order $k$, a product of modulus at least
$\epsilon$ has every factor of modulus at least
$\epsilon/R^{k-1}$. Only finitely many factor tuples can therefore
produce any nonzero product, or any product outside a prescribed disk.
After grouping to coefficients $d_{k,\rho}$,
\[
 \sum_\rho |d_{k,\rho}|\,|\rho|^N\le S_N^k.
\]
This proves tensor closure, including all multiplicative relations
among roots. Finite sums use a common moment index and the triangle
inequality. For the substitution $n=er$, group the at most $e$ roots
with a given $e$th power and choose $N'$ with $eN'\ge N$; boundedness
and the moment estimate prove admissibility. Cancellation on the top
circle is possible, as $q^n-(-q)^n$ already shows.

For completeness, consider the countable signed path argument directly.
Write $w=a+d$ and let $A_R$ be the finite top circle. For every allowed
even $k$, the finite group-algebra element
$F=\sum_{\alpha\in A_R}c_\alpha[\alpha]$ has
$F^k\ne0$: otherwise evaluation $[\alpha]\mapsto\alpha^n$ for
all positive integers $n$, followed by Vandermonde invertibility,
would force every coefficient to vanish. Its coefficient sum is
$(\sum c_\alpha)^k\ge0$, so at least one of its real grouped
coefficients is positive. The corresponding root has modulus $R^k$.
No product using a smaller-radius root reaches that modulus, and the
positive path-output spectrum cannot cancel its coefficient.

Write the assumed path output as
$q^nN_{k,n}=-\sum_\beta b_{k,\beta}\beta^n+E_{k,n}$, with
$b_{k,\beta}>0$ and
$|E_{k,n}|\le C_kn^{s_k}q^{n(kw/2+\eta_k)}$,
$\eta_k/k\to0$. The exact path identity gives
\[
 M_n^k+\sum_\beta b_{k,\beta}\beta^n=E_{k,n}-R_{k,n},\qquad
 |R_{k,n}|\le B^kq^{n(kw/2+d/2+1)}.
\]
The left side is admissible by closure and visibly contains the root
just found. Applying the proved spectral-radius formula at fixed $k$
gives
\[
 R^k\le q^{kw/2+\max(\eta_k,d/2+1)}.
\]
Taking $k$th roots and then unbounded even $k$ proves
$R\le q^{w/2}$. Constants depending on $k$ and polynomial factors
in $n$ disappeared in the preceding fixed-$k$ limit. If only
$n=e_kr$ are used, the top circle must survive grouping for an unbounded
allowed set of $k$; then the same estimate has the factor $e_k$ on
both sides and yields the same conclusion. For each field embedding,
admissibility, real coefficients, estimates and visibility must be
verified for the transformed system separately.

For the reciprocal assertion let the dual support have radius bound
$R_D<\infty$. It is positive when $A\ne\varnothing$, and
$|\alpha|\ge |\kappa|/R_D$ for every $\alpha\in A$.
Local finiteness away from zero makes $A$ finite. No convergence
assumption on the dual support is needed for this deduction.
Consequently the two independent bounds used in \fid{F17}, together
with visible reciprocity $q^d/A$, still give purity, but now also
force the original visible support to be finite.

Finally suppose the coefficients are nonzero integers and the moment
sum converges. For $N=0$ every term is at least one, so the support is
finite. For $N\ge1$,
\[
 \#\{\alpha:|\alpha|\ge\epsilon\}\le S_N\epsilon^{-N},\qquad
 |\alpha|\le S_N^{1/N}.
\]
These prove the two support conditions. The requirement $N\ge0$
matters if boundedness was not separately given: roots $2^j$ with
coefficients one have a convergent moment at index $-1$.
\end{proof}

\paragraph{The natural completion of the actual normalized trace.}
The trace in \fid{F26}, after $v=q^{\pm n/2}$, is
\[
 h_n=24+\frac{16q^{-2n}}{(1+q^{-n})^4}
 =24+16\sum_{j\ge0}(-1)^j\binom{j+3}{3}
                   (q^{-(j+2)})^n.
\]
The binomial expansion is absolutely convergent for every $n\ge1$.
It gives an admissible integer Euler spectrum
$A=\{1,q^{-2},q^{-3},\ldots\}$, with
$S_1=24+16q^{-2}/(1-q^{-1})^4$. Its radius is one, consistent with
an upper weight-zero bound. This observation does not construct the
pointwise and path-output identities required by the theorem.
For every fixed $\kappa\ne0$, the reciprocal support
$\{\kappa,\kappa q^2,\kappa q^3,\ldots\}$ is unbounded.
With these nonzero integer coefficients its nonnegative moments do
not converge. Uniqueness excludes replacing the same eventual trace
by a different admissible finite or pure reciprocal expansion.
Thus this completion repairs the finite-spectrum requirement for an
upper-bound argument but cannot supply bounded visible reciprocal
duality for this actual word trace. Regularized traces, non-atomic
spectra, duals outside this convergence class, and other comparison
dictionaries are not excluded.

\subsection{What the hypotheses mean geometrically}\label{sec:tracegeom}
Let $\mathcal F$ be a constructible Weil sheaf on $\A^d/\F_q$. Its trace functions $t_n$ have the required pointwise estimate if stalk ranks are bounded by $B$ and all stalk eigenvalues have weight at most $a$. Let
\[
 H_k=\{(x_1,\ldots,x_k):P_k(x_1,\ldots,x_k)=0\}
 \subset(\A^d)^k.
\]
The character kernel is represented by an Artin--Schreier sheaf pulled back along $hP_k$; the coefficient system on $H_k$ is the restriction of $\mathcal F^{\boxt k}$. The Grothendieck trace formula identifies the sums with alternating traces of $R\Gamma_c$.

If $R\Gamma_c(\A^d_{\bar{\F}_q},\mathcal F)$ is concentrated in degree $d$, then $\epsilon=(-1)^d$. If the relevant cohomology of $H_k$ is concentrated in degree $kd-1$, its sign is
\[
 (-1)^{kd-1}=-\bigl((-1)^d\bigr)^k.
\]
Multiplication by $q^n$ is the twist $(-1)$. These facts give precisely (S1)--(S2), or \fid{F04} when background terms have been controlled. A suitable duality of the coefficient system and cohomology supplies \eqref{eq:duality}.

Being a hypersurface of dimension $kd-1$ does \emph{not} prove this concentration. Singularities, exceptional fibers, compact-support boundary groups and Tate contributions are separate obligations. An affine variety also need not identify $H_c^d$ with its dual without an additional comparison. This is the information the proposed sign route still needs. Finding~\ref{f:signed} removes the input common-parity requirement at even lengths, but retains an explicit negative visible path-output spectrum. Its spectral core can be supplied geometrically or by a rigorous algebraic/category substitute; our argument does not prove that schemes are the only possible source of it.

\subsection{Countermodels to omitted hypotheses}
\begin{finding}[\fid{F05}: precise failures of abbreviated sign arguments]\label{f:countermodels}
The following general implications are false.
\begin{enumerate}[label=(\alph*)]
\item Path identity plus Parseval implies the expected middle weight.
\item Bounds for alternating traces control every individual cohomological eigenvalue without non-cancellation.
\item An upper-weight bound implies purity of that weight.
\item Bounds on an arbitrary infinite set of extension degrees suffice.
\item Nonnegative weights and Parseval make the Fourier operator positive.
\item Punctual purity implies arithmetic Frobenius semisimplicity.
\end{enumerate}
\end{finding}
\begin{proof}
(a) Take $t_n=1$ on $\F_{q^n}$, $a=0,d=1$. Then $M_n=q^n$, with eigenvalue $q>\sqrt q$, while \fid{F01} holds for all $k,n$. The actual cohomology is $H_c^2(\A^1)=\Q_\ell(-1)$, not the assumed middle degree.

(b) On $\Spec\F_q$, the complex $\Q_\ell(-1)\oplus\Q_\ell(-1)[-1]$ has a $q$-eigenvalue in degrees 0 and 1. All alternating traces vanish, and its virtual zeta function is 1.

(c) The eigenvalue 1 satisfies the weight-at-most-1 bound but has actual weight zero.

(d) The spectrum $\{R,-R\}$ has trace zero for every odd $n$, with $R$ arbitrarily large. Our lemma uses every extension degree; we assert no general necessity theorem excluding other suitably controlled subsequences.

(e) For $t=1$ on $\F_3$, the normalized Fourier matrix has diagonal entry $e^{2\pi i/3}/\sqrt3$ at $x=1$ and hence is not Hermitian. Over $\F_2$ the normalized matrix is $2^{-1/2}\left(\begin{smallmatrix}1&1\\1&-1\end{smallmatrix}\right)$, with eigenvalues $1,-1$. Both preserve the Euclidean norm.

(f) The matrix $\left(\begin{smallmatrix}1&1\\0&1\end{smallmatrix}\right)$ has pure weight-zero eigenvalues and trace 2 for every power, but is not diagonalizable. It defines a continuous $\ell$-adic representation through $\widehat\Z\to\Z_\ell$ and the associated unipotent matrices. This is the standard arithmetic-versus-geometric distinction also discussed in BBD after 5.3.8.
\end{proof}
In particular Fable suggestion~3's proposed identification of Fourier-after-multiplication with a positive operator is false in that literal generality. This does not refute the valid norm estimate or an independently constructed different positive operator.

\section{From purity to positivity: what is additionally required}\label{sec:positivity}
\begin{finding}[\fid{F06}: weights, positive values, and unimodality are different]\label{f:phase}
Exact purity does not imply a positive Frobenius trace. Even positivity at every positive numerical specialization, together with bar invariance, does not imply coefficientwise positivity. Pure graded pieces with a duality need not have unimodal dimensions.
\end{finding}
\begin{proof}
On $\Spec\F_q$, the rank-one Weil sheaf with Frobenius $-1$ is pure of weight zero but has trace $(-1)^n$ over $\F_{q^n}$. Next
\[
 h(v)=v^2-1+v^{-2}\ge1\quad(v\in\mathbb R\setminus\{0\})
\]
by $v^2+v^{-2}\ge2$, but its central coefficient is negative and $h(v^{-1})=h(v)$. Since $h(1)=1$, a factorization $\pm(v-v^{-1})^a f(v)$ with $a\ge0$ and nonnegative Laurent $f$ would force $a=0$, positive sign and $f=h$, which is impossible. Finally take $M_0=\Q_\ell$, $M_1=0$, $M_2=\Q_\ell(-2)$, with grade $j$ assigned weight $2j$ and the perfect pairing $M_0\otimes M_2\to\Q_\ell(-2)$. Its dimension sequence is $(1,0,1)$, not unimodal. Every degree-one operator $L$ has $L^2=0$, so it cannot satisfy the isomorphisms of \fid{F07} for $D=2$. This fixes the grading convention; it is not the usual degree-two Lefschetz grading on the cohomology of $\PP^1$.
\end{proof}
The last example is an abstract graded Weil object, not the cohomology of a smooth projective variety to which Hard Lefschetz already applies. It isolates precisely the extra theorem that such geometry would supply.

\begin{finding}[\fid{F07}: effective categorification and Lefschetz suffice]\label{f:category}
Let $\mathscr C$ be a split semisimple additive category with grading shift $\langle1\rangle$ and an additive product. Suppose its designated simple representatives satisfy finite decompositions
\[
 L_a*L_b\simeq\bigoplus_{c,j}L_c\langle j\rangle^{\oplus m_{ab}^{c,j}}.
\]
If an isomorphism from its graded split Grothendieck ring to the target algebra sends $[L_a]$ to the \emph{specified} basis element $b_a$ and grading shift to $v$, then its structure constants are
$\sum_jm_{ab}^{c,j}v^j\in\N[v,v^{-1}]$.

For a normalized multiplicity space $M=\bigoplus_{j=0}^D M_j$, suppose there is an operator $L:M_j\to M_{j+1}$ such that
$L^{D-2j}:M_j\to M_{D-j}$ is an isomorphism for $j\le D/2$. Then $\sum_j\dim M_jv^j$ is symmetric and unimodal.
\end{finding}
\begin{proof}
Apply the split Grothendieck relation to the displayed direct sum and compare coefficients under the basis-preserving isomorphism. Multiplicities of actual direct summands are nonnegative integers. For the second assertion the displayed isomorphisms give symmetry. If $j<D/2$ and $Lx=0$, then $L^{D-2j}x=0$, so $x=0$. Thus $\dim M_j\le\dim M_{j+1}$ up to the middle; symmetry gives the other half.
\end{proof}
For VIII this must be applied to the correctly \emph{normalized} multiplicity polynomial, with the signs, $(v-v^{-1})$ factors and positive rational denominators justified separately. The grading step must match the variable in which unimodality is asserted. An arbitrary derived $K_0$ realization is insufficient because $[K[1]]=-[K]$.

\begin{finding}[\fid{F08}: pure geometry and positive Gram determinants do not fix a basis]\label{f:twopoints}
There is a smooth projective finite \(\F_q\)-scheme, an integral basis of effective $K_0$ classes, and a positive definite Gram form, such that multiplication has a negative structure constant in that basis.
\end{finding}
\begin{proof}
Take $X=\Spec\F_q\sqcup\Spec\F_q$. Its $K_0$ ring is $\Z^2$ with componentwise product. The realization $H^0_{\et}(X_{\bar{\F}_q},\Q_\ell)=\Q_\ell^2$ has identity Frobenius and is pure of weight zero. Put
\[
 b_1=(1,1),\qquad b_2=(1,2).
\]
They form an integral basis, since the determinant is 1, and both are actual vector bundles, with the displayed ranks on the two components. The natural real pairing has Gram matrix
\[
 \begin{pmatrix}2&3\\3&5\end{pmatrix},
\]
whose principal minors are $2,5,1>0$. Nevertheless
\[
 b_2^2=(1,4)=3b_2-2b_1.
\]
All assertions follow directly. Exterior powers are also defined componentwise: $\lambda_t(r,s)=((1+t)^r,(1+t)^s)$. Thus passing to the ordinary geometric $\lambda$-ring framework does not remove this counterexample.
\end{proof}
This disproves a general bridge from purity/determinants/effectivity to coefficients of a chosen basis. It does not disprove the positivity of any designated GCT canonical basis.

\section{An obstruction in the actual nonstandard Hecke inclusion}\label{sec:ambient}
\begin{finding}[\fid{F09}: the inherited positive cone has rank one]\label{f:hecke}
Let $R=\Z[v,v^{-1}]$, $R_+=\N[v,v^{-1}]$, $\delta=v+v^{-1}$, and let $H=R1\oplus Rc$ with $c^2=\delta c$. In $H\otimes_RH$, put
\[
 P=c\otimes c+(c-\delta)\otimes(c-\delta),\qquad A=R[P].
\]
Equip the ambient tensor algebra with its positive cone $C_+$ in the basis
$1,c\otimes1,1\otimes c,c\otimes c$. Then
\begin{equation}\label{eq:cone}
 A=R1\oplus RP,\qquad P^2=\delta^2P,\qquad A\cap C_+=R_+1.
\end{equation}
Consequently no basis of the rank-two algebra $A$ consists entirely of ambient-positive elements, although its internal basis $\{1,P\}$ has positive multiplication.
\end{finding}
\begin{proof}
Set $b=c-\delta$. Then $cb=0$ and $b^2=-\delta b$, so the cross terms vanish and
$P^2=\delta^2(c\otimes c+b\otimes b)=\delta^2P$. Also
\[
 P=\delta^2\,1-\delta(c\otimes1+1\otimes c)+2c\otimes c.
\]
This proves that $1,P$ span $A$ and are independent. For $a+bP$, its four ambient coefficients are
\[
 a+\delta^2b,\quad-\delta b,\quad-\delta b,\quad2b.
\]
If all lie in $R_+$, then $2b\in R_+$ implies coefficientwise $b\ge0$. A nonzero such $b$ makes $\delta b$ nonzero and nonnegative, contradicting $-\delta b\ge0$. Thus $b=0$, and then $a\in R_+$. The converse is immediate. Finally $\delta^2=v^2+2+v^{-2}$ is positive, proving the intrinsic positivity assertion.
\end{proof}
This is the actual rank-one nonstandard generator of revised IV, \S1.4, equation~(8), and Definition~11.1, equation~(93). We fix the consistent earlier \S3.1 convention
$(T-v)(T+v^{-1})=0$, $c=C'_s=T+v^{-1}$ and $C_s=T-v=c-\delta$.
The opening of \S11.1 interchanges the exponents while retaining the idempotent assertion; this printed normalization slip is recorded rather than silently used. The proof above is unambiguous and also works over $\Q[v,v^{-1}]$ with the rational positive cone.

\paragraph{Categorical consequence.}
Suppose the usual graded tensor Hecke category has split $K_0$ identified with $H\otimes H$, with actual objects giving the ambient positive cone. Every object in an additive graded subcategory whose character lies in $A$ has class in $A\cap C_+=R_+1$. Its Grothendieck image therefore spans at most $R1$, not $A$. Hence the full nonstandard algebra cannot be obtained as such an effective subcategory through this fixed inclusion.

This is a proved obstruction to a specific natural route. It allows a different category, a different heart or character map, or a derived realization with signs. It is not a counterexample to VIII's signed-normalized conjecture and does not assert that $P$ must be a designated canonical-basis vector. The contrast between intrinsic positivity and failure of inherited positivity is exactly why a new realization is needed for this approach.

\begin{finding}[\fid{F10}: an algebraic replacement and its limits]\label{f:pivots}
Let $R=\Q[v,v^{-1}]$ and $R_+=\Q_{\ge0}[v,v^{-1}]$. Suppose an $R$-algebra $B$ has a positive basis $(e_j)$ and an $R$-subalgebra $A$ has a basis $(b_i)$ with expansions $b_i=\sum_jp_{ji}e_j$, $p_{ji}\in R_+$. If each $b_i$ has an exclusive coordinate $j_i$ with coefficient $r_iv^{m_i}$, $r_i>0$, and all other $b_h$ have zero coefficient there, then $A$ has positive structure constants in this basis. Disjoint positive supports, without this pivot normalization, do not suffice, even with bar invariance and generic semisimplicity. Disjoint supports only after specialization are weaker still.
\end{finding}
\begin{proof}
Write $b_sb_t=\sum_i m_{st}^i b_i$. Its ambient coefficients are positive. The exclusive coordinate is $r_iv^{m_i}m_{st}^i$, so division by the positive monomial preserves the cone and proves $m_{st}^i\in R_+$.

For the second assertion put
\[
 p=v+v^{-1},\quad h=v^2-1+v^{-2},\quad c=ph=v^3+v^{-3},
\]
and take $B=R\times R[t]/(t^2-c)$ with basis
$1=(1,1)$, $x=(0,t)$, $y=(0,1)$. It has positive multiplication
$x^2=cy$, $xy=x$, $y^2=y$. The free subalgebra with basis $1,x,z=py$ has disjoint positive ambient supports but
\[
 x^2=hz,\quad xz=px,\quad z^2=pz.
\]
The coefficient $h$ is negative in its middle term, and \fid{F06} shows it also fails the stated signed-factor positivity form. The product presentation proves associativity, and $t^2-c$ is separable over $\Q(v)$, so the generic algebra is semisimple. Since $p$ is invertible over $\Q(v)$, the generic subalgebra equals the generic ambient algebra and is also semisimple. All the displayed data are bar-invariant.

Finally in the positive algebra $\Q[v][x]/(x^2-x)$, the basis $b_0=1$, $b_1=x+v$ has nonnegative external coefficients and disjoint singleton supports modulo $v$, but direct expansion gives
\[
 b_1^2=(1+2v)b_1-v(1+v)b_0.
\]
This proves the last assertion.
\end{proof}
The statement provides a possible algebraic route that does not require new finite-field geometry. VIII, Conjecture~2.2, asserts disjoint supports in $L_\alpha/vL_\alpha$, not the global exclusive monomial pivots of \fid{F10}. Nor is raw positivity necessarily appropriate for all of VIII's signed targets. The missing pivot/normalization theorem cannot be obtained merely by repeating the local support assertion.

\begin{finding}[\fid{F11}: two limits of the Grothendieck-ring shortcut]\label{f:lambda}
Neither virtual $\lambda$-ring identities nor reduction of a noncommutative algebra modulo $p$ supplies the missing effective positive cone and arithmetic Frobenius.
\end{finding}
\begin{proof}
In $K_0(\Rep\mathbb G_m)=\Z[z,z^{-1}]$, the line character $z$ satisfies
\[
 \lambda_t(z-1)=\frac{1+zt}{1+t}
 =1+(z-1)t+(1-z)t^2+\cdots.
\]
Thus $\lambda^2(z-1)=1-z$, outside the effective cone $\N[z,z^{-1}]$. Also
$\gamma_t(z-1)=\lambda_{t/(1-t)}(z-1)=1+(z-1)t$. The augmentation/\(\gamma\)-filtration bookkeeping supplies no additional positivity order.

For the Frobenius assertion take $A=e_{12}$ and $B=e_{21}$ in $M_2(\F_p)$. Their $p$th powers vanish, whereas $(A+B)^2=I$, hence $(A+B)^p=I$ for $p=2$ and $A+B$ for odd $p$. The map $x\mapsto x^p$ is not additive. Therefore the automatic absolute Frobenius of a commutative scheme is not obtained by the same formula on a noncommutative coordinate algebra.
\end{proof}
A separately constructed quantum or categorical Frobenius remains possible. Also braided tensor categories do not automatically carry the symmetric exterior-power structure required by SGA~6.

\section{The theta route: exact positive consequences}\label{sec:theta}
\subsection{Periods and multiplicities are different data}
Alp\"oge's step~4 (p.~3) asserts positivity of a functional via a determinant. Step~7 (p.~4) describes determinants containing GGP periods, interlacing values 0 or 1, and controlled failures of interlacing. Step~8 describes iterated lifts between totally definite unitary groups and a gain of multiplicity. Page~1 and the introductory paragraph on p.~3 explain this as the analogue of Stepanov order of vanishing, not an assertion about the dimension of a Kronecker Hom space. The supplied Lean wrappers contain no theorem identifying the two kinds of multiplicity.

The following two findings state exactly what a positive projector or a chain of branching rules can give. They also show why merely finding a positive determinant is not the missing representation-theoretic formula.

\begin{finding}[\fid{F12}: positive projectors already realize numerical multiplicities]\label{f:projector}
For a finite group $H$ and a finite-dimensional complex unitary representation $U$, put
\[
 P_H=\frac1{|H|}\sum_{h\in H}\rho(h).
\]
It is the orthogonal projector onto $U^H$, with $\Tr P_H=\rank P_H=\dim U^H$ and nonnegative principal minors in an orthonormal basis. For partitions of $D$,
\[
 g(\lambda,\mu,\nu)=\Tr P_{S_D}
 \quad\text{on }[\lambda]\otimes[\mu]\otimes[\nu].
\]
For $\lambda\vdash dn$ and $\dim V\ge\ell(\lambda)$,
\[
 a_\lambda(d[n])=\Tr P_{H_{d,n}}\quad\text{on }[\lambda],\qquad
 H_{d,n}=(S_n)^d\rtimes S_d\subset S_{dn},
\]
where $a_\lambda(d[n])$ is the multiplicity in $\Sym^d(\Sym^nV)$.
\end{finding}
\begin{proof}
Inversion gives $P_H^*=P_H$, and counting each element $|H|$ times in the double sum gives $P_H^2=P_H$. Its image is fixed by $H$, and it is the identity on invariant vectors. Thus it is the claimed orthogonal projector. In an orthonormal basis its principal submatrices are Gram matrices of the vectors $P_He_i$, so their determinants are squared exterior norms. Its spectrum consists of $\dim U^H$ ones and the remaining zeros.

Specht modules are self-dual and can be made unitary by averaging a form; the triple invariant-space description gives the Kronecker assertion. For plethysm, invariants under the within-block permutations give $(\Sym^n V)^{\otimes d}$, and under block permutations give $\Sym^d(\Sym^nV)$. Apply invariants to Schur--Weyl decomposition
\[
 V^{\otimes dn}=\bigoplus_\lambda S_\lambda V\otimes[\lambda]
\]
and compare the $S_\lambda V$ multiplicity. This is also the projector explicitly used in BIP, equation~(4.1), with Proposition~3.3.
\end{proof}
This is a proof of the useful reformulation, not a new numerical nonnegativity result. BIP, Proposition~4.5, already gives projected tableaux spanning the highest-weight space; its Theorem~4.7 expresses evaluations as sums of signed column contributions. Independence and a compatible basis remain separate.

For any spanning vectors $v_i$ of a multiplicity space, their Gram matrix satisfies
\[
 \sum_{|I|=r}\det G_I=\sum_{|I|=r}\bigl\|\bigwedge_{i\in I}v_i\bigr\|^2.
\]
This sum is positive exactly when the rank is at least $r$. Determinant positivity therefore restates the rank question; it does not by itself determine which minors are nonzero or provide an efficient basis-selection rule. We make no claim that every such rank question is intrinsically hard.

\begin{finding}[\fid{F13}: when an interlacing chain counts a target multiplicity]\label{f:branching}
Let $\mathscr C_r,\ldots,\mathscr C_0$ be split semisimple linear categories of finite-length objects and let $F_i:\mathscr C_i\to\mathscr C_{i-1}$ be additive exact functors. Suppose
\[
 F_i S_i(a)\simeq\bigoplus_b S_{i-1}(b)^{\oplus b_i(a,b)},\qquad b_i(a,b)\in\{0,1\},
\]
with finite sums. Then the multiplicity of $S_0(c)$ in $F_1\cdots F_rS_r(a)$ is the number of allowed paths from $a$ to $c$:
\[
 \sum_{a_{r-1},\ldots,a_1}
 b_r(a,a_{r-1})\cdots b_1(a_1,c).
\]
If the allowed steps are specified by integral interlacing inequalities and every intermediate coordinate is bounded by the fixed endpoints, these paths are the integer points of the resulting polytope.
\end{finding}
\begin{proof}
Apply each additive functor to the previous finite direct-sum decomposition. The multiplicities multiply along paths and add over intermediate simples. Semisimplicity identifies the final coefficient with the desired multiplicity. Every product of 0/1 entries is 1 exactly for an allowed path. Sending a path to its array of intermediate integer weights is a bijection with the integer solutions of the inequalities. Boundedness makes the polyhedron a polytope.
\end{proof}
Ordinary restriction along consecutive block inclusions $\GL_r\supset\GL_{r-1}$ has the interlacing rule $\lambda_j\ge\mu_j\ge\lambda_{j+1}$; see \cite[Theorem 2.1, equation (2.2), and the patterns (2.4)]{Molev}. The Kronecker map $(g,h)\mapsto g\otimes h$ and the plethysm map $g\mapsto S_\mu(g)$ are different maps. A theta lift between dual pairs is not automatically their composite restriction functor.

A small object mismatch is instructive. $S_{(d^r)}(\C^r)=\det^d$ is one-dimensional and has a unique usual interlacing pattern, whereas the rectangular Specht module is a different representation. For example
\[
 [(2,2)]\otimes[(2,2)]=[(4)]\oplus[(2,2)]\oplus[(1^4)].
\]
To check this without a table lookup, the permutation action of $S_4$ on the three pairings of four letters is the direct sum of the trivial representation and $[(2,2)]$. The latter has character $(2,0,2,-1,0)$ on cycle types $(1^4),(2,1,1),(2,2),(3,1),(4)$. Its pointwise square equals this character plus the trivial and sign characters, proving the identity. A correctly identified chain could still count higher multiplicities; the example only prevents substitution of a determinant character for a rectangular Specht module.

\subsection{A lattice obstruction to ordinary interlacing alone}
\begin{finding}[\fid{F18}: homogeneous difference constraints and a rectangular parity ray]\label{f:saturation}
Let $u$ range over an integral label set closed under positive-integer scaling. Fix a number of coordinates and a finite list of inequalities
\[
 P(u)=\{x\in\mathbb R^s:x_i-x_j\le b_{ij}(u)\},\qquad x_0=0,
\]
where every specified $b_{ij}(u)$ is integral and $b_{ij}(Nu)=Nb_{ij}(u)$ for every $N\ge1$. Assume $P(u)$ is bounded, and put $h(u)=\#(P(u)\cap\Z^s)$. Then
\[
 h(Nu)>0\quad\Longleftrightarrow\quad h(u)>0.
\]
The same property holds for a fixed finite sum of such counting functions. Thus these models cannot count all two-row rectangular Kronecker coefficients. On the particular ray with all three partitions $(d,d)$, the established two-row formula gives
\[
 g((d,d),(d,d),(d,d))=\boldsymbol1_{2\mid d}
 =\#\bigl(d\{1/2\}\cap\Z\bigr)\qquad(d\ge1).
\]
This ray has an exact rational Ehrhart model despite its failure of the displayed ordinary support property.
\end{finding}
\begin{proof}
Homogeneity gives $P(Nu)=NP(u)$. If $x_i-x_j\le b$ with $b\in\Z$, then
$\lfloor x_i\rfloor\le\lfloor x_j+b\rfloor=\lfloor x_j\rfloor+b$.
Hence coordinatewise flooring, including $x_0=0$, takes every real point of $P(u)$ to an integral one. If $z\in P(Nu)\cap\Z^s$, apply this to $z/N$. Conversely an integral $x\in P(u)$ gives $Nx\in P(Nu)\cap\Z^s$. Positivity of a fixed finite sum is positivity of at least one summand, so the same argument applies. Boundedness is only needed to make these finite counts.

For $u=((1,1),(1,1),(1,1))$, the triple tensor product is the sign representation of $S_2$, so $g(u)=0$. The explicit $S_4$ character calculation following \fid{F13} gives $g(2u)=1$: the class sizes are $(1,6,3,8,6)$ and the character $(2,0,2,-1,0)$ has cubic average $(8+3\cdot8-8)/24=1$. Thus an exact model of the stated form on all two-row rectangles is impossible.

For the whole ray use IV, \S18.3, equation~(172), printed p.~128. In its column convention, $(d,d)=[d,0]$ as a two-row label and $[0,0,d,0]$ as a four-row label. The coefficient for the second rectangle is the constant term of that equation's generating function. Only $k=0$ and $j=0$ in its first sum survive; the second sum is empty. The remaining condition is that $n_3,n_2,n_1$ are all even and nonnegative. Substituting $(n_3,n_2,n_1)=(0,d,0)$ gives exactly the displayed parity formula. The single rational point $\{1/2\}$ has an integral point after dilation by $d$ exactly when $d$ is even.
\end{proof}
The difference-constraint result includes ordinary interlacing with integral endpoint rows. Its exclusions are essential: extra congruences, sum equations, nonhomogeneous endpoint shifts, a changing constraint list or a changed lattice can evade it. For example $2z=d$ retains precisely the parity that ordinary integral differences lose. The example concerns all two-row rectangles, not a model restricted only to IP's large-parameter region.

Nor does it contradict VI's definition of ``strictly saturated'' (Definition~1.2.4, printed p.~14). That definition tests each nonzero residue polynomial separately; here $f_{\rm even}=1$ and $f_{\rm odd}=0$, so its saturation index is zero. We use the displayed support property to avoid conflating the two notions. IV, Theorem~18.2 and equations~(163)--(164), explicitly retain congruences and remove columns in pairs. A successful theta comparison must preserve these arithmetic data or equivalent information; the source does not identify them with a particular GGP local sign.

In Grothendieck's language, the tensor group-scheme map $\GL_{2,\Z}\times\GL_{2,\Z}\to\GL_{4,\Z}$ induces, over $\Q$, a positive restriction map on representation $K_0$. The coefficient of $[\det^d\boxtimes\det^d]$ in the restriction of $[S_{(d,d)}(V\otimes W)]$ is this parity function. Thus an effective $K_0$ homomorphism does not by itself give the ordinary support property. The proof above is already the exact integral-lattice dictionary; its group-scheme framework is SGA~3 I, 4.7.2--3, with cote~90, scan~17, as an action/comodule anchor.

\paragraph{Exact remaining theta gate.}
Construct functors and endpoint identifications making Alp\"oge's lift chain the target restriction, or identify its period spaces with the projector images of \fid{F12}. To obtain canonical-basis positivity, additionally identify grading, operations, and the specified basis as in \fid{F07}. For a \(\#P\) rule the path encodings and admissibility tests must have polynomial bounds; for PH1 the stronger polyhedral specifications in VI must also hold. None of these identifications is supplied by steps~4, 7 or~8.

% Add bibliography key BlasiakRep for Jonah Blasiak,
% Representation theory of the nonstandard Hecke algebra,
% arXiv:1201.2209v2 (2012), https://arxiv.org/abs/1201.2209v2.
\subsection{An actual nonstandard tower: specialization fails on lines and succeeds after mixing}
There is already a nonstandard tower to which the generic path mechanism
applies. Blasiak \cite[Theorem~5.1, p.~14; Proposition~5.2, pp.~14--15;
Definition~6.1, p.~23]{BlasiakRep} proves that the two-row nonstandard
Temperley--Lieb algebra over $K=\Q(v)$ is split semisimple, that its
successive restrictions are multiplicity free, and hence that its simple
modules have seminormal path lines. These are established source results,
not new consequences of purity. The source distinguishes these lines
from a basis compatible with classical decomposition at $v=1$; its
Remark~6.2 also gives a non-multiplicity-free example outside the
two-row quotient. In particular, its simple module at \emph{tower rank
four}
\[
 M_K=M_{(3,1),K}\otimes M_{(2,2),K}
\]
has dimension six and restriction
\begin{equation}\label{eq:towerrestriction}
 M_K|_{A_{3,K}}=N(\delta)\oplus N(D-2)\oplus\epsilon_+
                  \oplus\epsilon_-.
\end{equation}
Here the rank-three notation is that of \fid{F19}/\fid{F21}.
The following comparison computes exactly what happens to its path lines
in the natural ordinary Specht tensor lattice.

\begin{finding}[\fid{F30}: regular path mixing and a projector base-change defect]\label{f:tower}
Let $R=\Q[v,v^{-1}]_{(v-1)}$, $k=R/(v-1)=\Q$, and fix
\[
 M_R=\mathsf S_{(3,1),R}\otimes_R\mathsf S_{(2,2),R},
\]
with the specified nonstandard action. Write $A_R$ for the actual
nonstandard algebra acting at tower rank four. Its specialization has
\[
 M_1\simeq[(3,1)]\oplus[(2,1,1)].
\]
The two generic path lines through the penultimate summand $N(\delta)$
lie in a regular saturated rank-two plane $E_R$. No vector on either
line which is diagonally rescaled to be regular with nonzero reduction
can specialize into one of these two classical constituents.

Nevertheless a regular non-diagonal change of the six path representatives
in this same lattice does produce a specialization-compatible basis.
It mixes only two pairs of corresponding rank-two path labels, by
unit rescaling followed by the blocks
\[
 \begin{pmatrix}1&1\\1&-1\end{pmatrix}.
\]
Finally
\[
 \End_{A_R}(M_R)=R,\qquad \End_{A_1}(M_1)=k\times k,
\]
and the natural base-change map is $a\mapsto(a,a)$. Thus the two
classical projectors do not lift to endomorphisms commuting with the
whole nonstandard deformation, although the stated regular mixed basis
exists.
\end{finding}
\begin{proof}
\par\smallskip\noindent\textbf{A fixed saturated plane.}\quad
The ordinary parabolic Hecke algebra $H_3\subset H_4$ has trivial
projector
\[
 p_+=\frac{\sum_{w\in S_3}v^{\ell(w)}T_w}
              {\sum_{w\in S_3}v^{2\ell(w)}}.
\]
The denominator specializes to six and is a unit of $R$. The numerator
is multiplied by $v$ by each simple generator on either side; squaring
it therefore multiplies it by the displayed denominator. Hence $p_+$
is a regular central idempotent of $H_3$. Ordinary branching gives
\[
 M_{(3,1)}|_{H_3}=M_{(3)}\oplus M_{(2,1)},\qquad
 M_{(2,2)}|_{H_3}=M_{(2,1)}.
\]
It follows that
\[
 E_R=M_R(p_+\otimes1)
\]
is a free rank-two direct summand, stable under $A_{3,R}$, and its
generic fiber is exactly the $N(\delta)$ summand in
\eqref{eq:towerrestriction}. Because the projector is regular,
\begin{equation}\label{eq:towersaturated}
 E_R=E_K\cap M_R.
\end{equation}
Indeed an element of the intersection is fixed by $p_+\otimes1$
and so lies in its $R$-image. Every regular nonzero reduction of a
rescaled generic path vector in $E_K$ therefore lies in $E_1$.

\par\smallskip\noindent\textbf{The special decomposition and its class sum.}\quad
At $v=1$, $P_i=2(1+s_i)$ for the diagonal symmetric-group action,
so the specialized algebra image generates that action. On the five
cycle types $(1^4),(2,1,1),(2,2),(3,1),(4)$, the characters are
\[
 \chi_{31}=(3,1,-1,0,-1),\qquad
 \chi_{22}=(2,0,2,-1,0),\qquad
 \chi_{211}=(3,-1,-1,0,1).
\]
Their product satisfies
\[
 \chi_{31}\chi_{22}=(6,0,-2,0,0)=\chi_{31}+\chi_{211},
\]
which proves the stated classical decomposition in characteristic zero.
Let $U$ be the sum-zero subspace of $\mathbb R^4$, and $W$ the
sum-zero subspace of the permutation representation on the three
pairings of four letters. These are orthogonal models for $[31]$
and $[22]$. The $S_3$ fixing the fourth letter has invariant line
in $U$ generated by
\[
 w=(1,1,1,-3),\qquad \langle w,w\rangle=12.
\]
Thus $E_1=\mathbb Rw\otimes W$. Let $p$ project orthogonally onto
$\mathbb Rw$, put $\Pi=p\otimes1$, and let
\[
 Z=\sum_{1\le i<j\le4}(ij)
\]
act diagonally. For $i,j\le3$, the transposition fixes $w$ and
$p\rho_U((ij))p=p$. For an external transposition $(i4)$ its inner
product with $w$ is $-4$, giving
\[
 p\rho_U((i4))p=-\tfrac13p.
\]
The action on the three pairings is the quotient $S_4\to S_3$ with
Klein-four kernel. The three internal transpositions map bijectively
to its three transpositions, and so do the three external ones.
The sum of those three matrices in the natural three-point
permutation module is the all-ones matrix; it vanishes on $W$.
Consequently both transposition sums on $W$ are zero and
\begin{equation}\label{eq:towercompression}
 \Pi Z\Pi=p\otimes(0-\tfrac13\,0)=0.
\end{equation}

The class sum $Z$ acts on $[31]$ and $[211]$ by $+2$ and $-2$:
its scalar is six times the transposition character divided by the
dimension, namely $6(\pm1)/3$. If a nonzero $x\in E_1$ belonged
to either constituent, then $\Pi Z\Pi x=\pm2x$, contradicting
\eqref{eq:towercompression}. Thus $E_1$ meets each constituent only
in zero. Together with \eqref{eq:towersaturated}, this proves the
diagonal-normalization obstruction. All identities in this real
model have rational coefficients, so the conclusion holds over $k$.

\par\smallskip\noindent\textbf{Regular complementary path planes.}\quad
Put $J_R=M_R((1-p_+)\otimes1)$. It is the complementary regular
rank-four summand. The three central block idempotents
$z_+,z_-,z_{\rm mat}$ of the actual rank-three algebra in \fid{F19}
are defined over $R$, without dividing by the colliding-root
difference. They commute with the ordinary parabolic projector.
Set
\[
 L_{+,R}=J_Rz_+,\quad L_{-,R}=J_Rz_-,\quad F_R=J_Rz_{\rm mat}.
\]
These direct summands are free over the local ring $R$, with generic
types $\epsilon_+,\epsilon_-,N(D-2)$ and ranks $1,1,2$. Hence
\[
 M_R=L_{+,R}\oplus L_{-,R}\oplus E_R\oplus F_R.
\]
This separates $E_R$ from $J_R$ by the ordinary projector first;
it does not try to separate the two colliding matrix branches by a
singular nonstandard central projector.

At specialization the four $S_3$ types are respectively trivial,
sign, standard, standard. The first-factor orthogonal complement
$U_0=w^\perp\subset U$ is standard, so
$J_1=U_0\otimes W=1\oplus\mathrm{sgn}\oplus[21]$ and $F_1$
is its unique standard summand. Write these specialized spaces as
$L_+,L_-,E,F$. Since $Z$ commutes with $S_3$, it sends $E$ into
$E\oplus F$; there is no homomorphism from the standard module
to either scalar type. Equation~\eqref{eq:towercompression} removes
the $E$ component. Also $Z^2=4$ by the classical decomposition.
It follows that
\[
 Z(E)=F,\qquad Z(F)=E,\qquad J:=Z/2:E\longrightarrow F
\]
is an isomorphism with inverse $Z/2$. Its graph planes
\[
 G_+=\{x+Jx:x\in E\},\qquad G_-=\{x-Jx:x\in E\}
\]
have $Z$-eigenvalues $+2$ and $-2$. The classical restrictions
$[31]|_{S_3}=1\oplus[21]$ and
$[211]|_{S_3}=\mathrm{sgn}\oplus[21]$ therefore identify the
actual constituents as
\begin{equation}\label{eq:towergraphs}
 [31]=L_+\oplus G_+,\qquad [211]=L_-\oplus G_-.
\end{equation}

\par\smallskip\noindent\textbf{A two-pair mixing repair.}\quad
The idempotent $e=P_1/D$ from tower rank two is regular. Each of $E_R,F_R$
splits by $e,1-e$ into two rank-one free direct summands. Choose
bases $x_{+,R},x_{-,R}$ of these two lines in $E_R$, and
$f_{+,R},f_{-,R}$ in $F_R$. Their generic lines are exactly the
four tower paths through the two standard penultimate labels.
At specialization $e=(1+s_1)/2$. Since $J$ commutes with $s_1$,
it carries each line of $E$ to the matching line of $F$. Rescale
each $f_{\sigma,R}$ by a constant unit of $\Q$ so that its
reduction is $Jx_\sigma$.
For each $\sigma=+,-$, replace the two regular vectors by
\[
 g_{\sigma,+}=x_{\sigma,R}+f_{\sigma,R},\qquad
 g_{\sigma,-}=x_{\sigma,R}-f_{\sigma,R}.
\]
The displayed two-by-two block has determinant $-2$, a unit in
$R$, so these four vectors, together with bases of $L_{+,R}$
and $L_{-,R}$, form an $R$-basis. Their reductions lie in the
two graph planes in \eqref{eq:towergraphs}. This proves the
regular specialization-compatible basis claim in the original
fixed lattice. It uses non-diagonal mixing across distinct generic
penultimate summands, precisely the freedom left open by the first
part of the finding.

\par\smallskip\noindent\textbf{Why the classical projectors still do not lift.}\quad
By the cited split simplicity,
$\End_{A_K}(M_K)=K\,\id$. Any element of $\End_{A_R}(M_R)$
therefore becomes a scalar $\lambda\id$ generically. In an
$R$-basis its diagonal entries belong to $R$, forcing $\lambda\in R$.
Conversely all such scalars commute with the action. This proves
$\End_{A_R}(M_R)=R$. At specialization the image is the diagonal
$k[S_4]$ action with two pairwise nonisomorphic absolutely simple
constituents, each once, so its endomorphism ring is $k\times k$.
The base-change map sends a scalar to $(a,a)$; it is injective with
one-dimensional cokernel, and neither nontrivial classical projector
lies in its image.

The three-vector groups in the mixed basis span direct summands of
the underlying free $R$-module, whose reductions are the two
classical constituents. They are not $A_R$-stable generically,
which would contradict generic simplicity. Thus the constructive
basis comparison makes no false claim of a nonstandard-equivariant
projector lift.
\end{proof}

\paragraph{Exact scope and the coherent dictionary.}
The excluded operation is diagonal normalization of this tower's
path lines in the specified ordinary Specht tensor lattice; the
positive operation is the explicit regular mixing just constructed.
No arbitrary-lattice obstruction, VIII canonical-basis identification,
uniform general specialization algorithm or positive structure-constant
rule is inferred. The known path labels count generic simple dimensions,
not classical Kronecker multiplicities by themselves. Their verification
is polynomial in the expanded tower rank, not automatically in its
binary logarithm.


\section{Two further geometric routes and their limits}\label{sec:other}
\begin{finding}[\fid{F14}: what an ambient pure realization would prove]\label{f:sub}
A Frobenius-stable subquotient of a pure algebraic Weil module of weight $w$ is pure of weight $w$. A stable subquotient of its endomorphism module, with Frobenius acting by conjugation, is pure of weight zero.
\end{finding}
\begin{proof}
A basis adapted to a stable subspace makes Frobenius block upper triangular. The characteristic polynomial factors into those of the subspace and quotient, so their eigenvalues are selected from the original spectrum. This works under every embedding. Eigenvalues on $V\otimes V^*$ are ratios $\alpha/\beta$, of absolute value 1; apply the first assertion. No diagonalizability is required.
\end{proof}
Thus if the desired nonstandard multiplicity objects were already realized as such stable subquotients, a new spectral RH would not be needed for them. Stability, grading, effective classes, and a basis-preserving realization would still need proof. Declaring an arbitrary vector space to have identity Frobenius proves an artificial weight-zero assertion and identifies no GCT coefficient.

\begin{finding}[\fid{F15}: orbit models exist, but their ordinary purity does not determine Hilbert multiplicities]\label{f:orbit}
An integral polynomial tensor $f$ in a finite free representation $W_\Z$ of $\GL_{N,\Z}$ has a finite-type affine orbit-closure model over $\Z$ whose complex fiber is the usual reduced orbit closure. Ordinary pure cohomology of a projective scheme does not determine its polarized coherent Hilbert function or its equivariant coherent character.
\end{finding}
\begin{proof}
The orbit map induces $\Z[W]\to\Z[\GL_N]$. Its kernel $I$ defines the finite-type scheme $\Spec(\Z[W]/I)$. The quotient is a subring of $\Z[\GL_N]$ and is torsion-free, hence flat over $\Z$. Since $\C$ is flat over $\Z$, tensoring preserves the kernel. The resulting image ring is a subring of the reduced domain $\C[\GL_N]$, so the complex fiber is the reduced scheme-theoretic orbit closure. No corresponding equality after arbitrary reduction modulo $p$ is asserted; kernel formation can fail to commute with nonflat base change. For an explicit example, let $\mathbb G_m$ scale $f=p$ in $\A^1_\Z$. The map $\Z[x]\to\Z[t,t^{-1}]$, $x\mapsto pt$, is injective, so this model is $\A^1_\Z$. Reducing $f$ first gives zero and an orbit closure consisting of one point, whereas the special fiber of the model is $\A^1_{\F_p}$.

For the second assertion fix $X=\PP^1_{\F_q}$ and $L_r=\mathcal O(r)$, $r>0$. Monomials of degree $rd$ in two variables form a basis of $H^0(X,L_r^d)$, so its dimension is $rd+1$. The underlying scheme and its ordinary/intersection cohomology are independent of $r$: they are $\Q_\ell$ in degree 0 and $\Q_\ell(-1)$ in degree 2 (the projective-space calculation, \cite[Example 16.3, p.~109]{Milne}). Thus the pure Frobenius data do not determine the polarized Hilbert function. Finally, on a point with trivial $\mathbb G_m$ action, a weight-$r$ linearization of $\mathcal O$ gives character $z^{rd}$ on its $d$th tensor power, while the underlying Weil cohomology remains the same. This proves the character assertion.
\end{proof}
A projective integral model of the homogeneous orbit closure can likewise be formed by projectivization. For a reduced equidimensional projective fiber $X/\F_q$ of dimension $e$ with a dense smooth open $j:U\hookrightarrow X$, the unshifted complex $\IC_X^u=j_{!*}\Q_\ell[e][-e]$ is pure of weight zero (BBD~5.3.2), and $IH^i(X)=H^i(X_{\bar{\F}_q},\IC_X^u)$ is pure of weight $i$ by proper pushforward and duality (BBD, Remark~5.4.9, printed p.~143). For a nonreduced fiber pass to its reduction for this \(\ell\)-adic assertion. This is already ordinary finite-field mathematics. The GCT multiplicities instead belong to the graded coherent representation $\C[X]_d$, or to coherent sections with the required saturation and degree hypotheses. The polarization, linearization and comparison object cannot be omitted. Riemann--Roch is consistent with the example: it explicitly retains the varying Chern character of $L_r^d$.

Hilbert/Quot schemes represent the relevant coherent moduli under their usual hypotheses (FGA, Expos\'e~221, Theorems~3.1--3.2); they do not identify their intersection cohomology with the desired multiplicity spaces. Constructing an affine space from an already known multiplicity module encodes its unknown dimension in a scheme, but does not compute that dimension or produce a canonical indexing set. This route was pursued and reduces to an additional coherent-to-constructible comparison, not to an unproved ordinary Weil theorem.

\begin{finding}[\fid{F16}: a positive rank-one model exists with a different action]\label{f:rankonemodel}
The abstract rank-one algebra of \fid{F09} has a positive realization inside the same ambient algebra after changing the inclusion. Put $E=c\otimes c$. Then $R1\oplus RE$ is a subalgebra, and $P\mapsto E$ gives an $R$-algebra isomorphism from $R1\oplus RP$. The coefficient $\delta^2$ is the normalized Tate cohomology polynomial of $\PP^1\times\PP^1$. This isomorphism does not preserve the original ambient action relevant to GCT.
\end{finding}
\begin{proof}
We have $E^2=\delta^2E$, and $1,E$ are independent in the ambient tensor basis. Hence both algebras have the presentation $R[z]/(z^2-\delta^2z)$, with the respective chosen generators. The new basis $1,E$ is ambient-positive. The even cohomology of $\PP^1\times\PP^1$ is $\Q_\ell$, $\Q_\ell(-1)^{\oplus2}$ and $\Q_\ell(-2)$ in degrees 0, 2 and 4, by the projective-space calculation and K\"unneth. Its polynomial, normalized by $v^{-2}$ with $v^2$ counting a Tate degree, is
\[
 v^{-2}(1+2v^2+v^4)=v^{-2}+2+v^2=\delta^2.
\]
At $v=1$, let $s$ be the nontrivial element of $S_2$, so $c=1+s$. Then
\[
 P=2(1+s\otimes s),\qquad E=(1+s)\otimes(1+s).
\]
On the sign tensor sign representation, $P$ acts by 4 and $E$ by 0. Therefore the isomorphism changes the given ambient representation and cannot be substituted into the Kronecker comparison without a new compatibility theorem.
\end{proof}
This is a positive result for the abstract rank-one Hecke case and a precise limit of that result. Purity here is degree by degree: the total ungraded cohomology has several weights, not one weight-zero Weil module. The Tate-grade dimensions are $(1,2,1)$, whereas consecutive integer powers of $v$ give $(1,0,2,0,1)$; no assertion about VIII's precise unimodality normalization follows without fixing the grading step. This is neither an identification of VIII's full specified bases nor a new consequence of the sign-only theorem. In motive language the coefficient is realized by the Tate object $\mathbf1\oplus\mathbb L^{\oplus2}\oplus\mathbb L^2$, with a grading shift; the missing datum is the action comparison, not its purity.

\subsection{Derived lifts and a positive model of the actual rank-one action}
\begin{finding}[\fid{F27}: what changing parity or the effective cone accomplishes]\label{f:derived}
Let $R=\Z[v,v^{-1}]$, $\delta=v+v^{-1}$,
$H=R[c]/(c^2-\delta c)$, and let $P$ be the actual rank-one element
of \fid{F09}. In a graded additive monoidal categorification whose
split Grothendieck ring has the named effective basis
$1,c_1,c_2,c_1c_2$, the following hold.
\begin{enumerate}
\item $P$ has an explicit bounded two-term lift, but no bounded complex
with Euler class $P^k$ ($k\ge1$) has all its nonzero terms in a single
homological parity. An overall sign or internal grading shift does not
alter this obstruction.
\item The explicit zero-differential lift has tensor-power homotopy
amplitude $k$. Separately, if a nonzero bounded complex in an abelian tensor category
with faithful exact vector-space fiber functor has extreme nonzero
cohomology degrees $p,r$, its $k$th tensor power has amplitude $k(r-p)$.
Mixed cohomological parity persists in every tensor power.
\item A different positive cone does preserve the actual rank-one
character actions. The faithful map
\[
 H\longrightarrow R^2,\qquad a+bc\longmapsto(a+b\delta,a)
\]
becomes an isomorphism after inverting $\delta$. Its tensor square sends
\[
 P\longmapsto(\delta^2,0,0,\delta^2)
\]
on the four labels $++,+-,-+,--$. Thus the missing sign-tensor-sign
action in \fid{F16} is restored. In the even form $u=v^2$, the two
nonzero values are actual positive Tate classes
$u^{-1}+2+u$.
\end{enumerate}
These statements do not identify the new character cone with VIII's
canonical basis or construct the negative path-output spectrum of
\fid{F24}.
\end{finding}
\begin{proof}
The already proved quadratic relation gives
\[
 P^k=\delta^{2k-2}P
 =\delta^{2k}\,1-\delta^{2k-1}(c_1+c_2)
                  +2\delta^{2k-2}c_1c_2.
\]
Its four named coordinates have both signs for every $k$.
Let $B_0,B_1,B_2,B_3$ represent $1,c_1,c_2,c_1c_2$, with $B_0$ the
tensor unit. Internal shift $(1)$ represents $v$ and is distinct
from homological shift $[1]$. Set
\[
 X_+=B_3^{\oplus2}\oplus B_0(2)\oplus B_0^{\oplus2}\oplus B_0(-2),
\]
\[
 X_-=B_1(1)\oplus B_1(-1)\oplus B_2(1)\oplus B_2(-1).
\]
The complex $K^0=X_+$, $K^1=X_-$ with zero differential has class
$P$. Any complex with terms only in even degrees has effective Euler
class, and one with terms only in odd degrees has the negative of an
effective class. Neither can equal the displayed mixed-sign class of
$P^k$. Homotopy equivalence preserves Euler class, so it cannot repair
this; overall signs and grading monomials preserve the mixed pattern.

The zero-differential $K^{\otimes k}$ has nonzero terms in all degrees
$0,\ldots,k$. Use unit summands from $X_+$ and a $B_1$ summand from
each chosen $X_-$; $c_1^j=\delta^{j-1}c_1\ne0$ proves nonvanishing
when $j\ge1$. If a zero-differential complex is homotopy equivalent
to $D$, the composite maps back to it equal the identity degreewise,
because its homotopy correction $dh+hd$ is zero. Hence its degree-$i$
term is a retract of $D^i$; no nonzero extreme can disappear. This
proves the amplitude claim for this particular lift, not for every
complex with its Euler class.

For the separate cohomological assertion, apply the faithful exact
tensor fiber functor. Over a field the K\"unneth decomposition of
$H^j(K^{\otimes k})$ is the direct sum of
$H^{i_1}(K)\otimes\cdots\otimes H^{i_k}(K)$ with sum of degrees $j$.
The extreme summands are the nonzero powers of $H^p$ and $H^r$.
Opposite-parity nonzero groups in degrees $i,j$ give nonzero summands
in degrees $ki$ and $(k-1)i+j$, of opposite parity. Exactness and
faithfulness detect these assertions in the original category.
Shifts, Tate twists and finite base-field extensions do not reduce
the amplitude. This argument does not assume an arbitrary convolution
or pushforward is exact for a chosen heart.

For the character model, the map evaluates $c$ at $\delta$ and zero.
Its kernel is zero since $R$ is a domain and $\delta\ne0$; its image
is the pairs whose difference is in $\delta R$. Localizing makes it
surjective. The four-coordinate tensor map is injective by successive
evaluation of $1,c_1,c_2,c_1c_2$, and direct evaluation of $P$ gives
the stated tuple. At $v=1$, $c_i=1+s_i$ and $P=2(1+s_1s_2)$,
so both equal-sign characters have eigenvalue four and both others
have eigenvalue zero. Every module over the localized algebra splits
by its four actual orthogonal character idempotents. It becomes a
quasi-coherent sheaf on four copies of the localized base, a vector
bundle when its corresponding modules are finite projective.

For an integer-Tate formulation set $b=vc$ and use
\[
 H_{\rm even}=\Z[u,u^{-1}][b]/(b^2-(1+u)b),\qquad
 P=u^{-1}\{b_1b_2+(b_1-1-u)(b_2-1-u)\}.
\]
Scalar extension recovers the original algebra. On each rational
character point, represent $b_i$ by $\mathbf1\oplus\mathbf1(-1)$
if its sign is $+$, and by zero otherwise. The element $P$ is then
represented on each diagonal point by
\[
 T=\mathbf1(1)\oplus\mathbf1^{\oplus2}\oplus\mathbf1(-1),
 \qquad \Tr(F^n\mid T)=q^{-n}+2+q^n,
\]
and by zero off the diagonal. Tensoring with this object on the
two supported points gives the corresponding positive operation on
the character Grothendieck group. It has weights $-2,0,0,2$;
it is not a single pure weight-zero representation.
\end{proof}

There is an actual Hecke category satisfying the first hypothesis:
the product of the two rank-one Soergel categories. For one factor
take $A=\Q[x]$, $\deg x=2$, $s(x)=-x$, and
$B=A\otimes_{A^s}A(1)$. The rank-two free decomposition over
$A^s=\Q[x^2]$ gives $B\otimes_A B=B(1)\oplus B(-1)$.
The split Grothendieck/KL identification is the known result
\cite[\S3.5, Theorem~3.5, pp.~1107--1108; Theorem~1.1, p.~1093]{EW}.
The bounded category used above is its homotopy category; no
unspecified perverse $t$-exactness was invoked.

The character map is also a finite normalization of affine schemes:
$R^2$ is the integral closure of $H$ in its total fraction ring,
and $R^4$ is that of the tensor-square algebra. Indeed the separate
branch idempotents are integral; conversely an element integral over
these finite $R$-algebras is integral over the normal ring $R$ and
has each coordinate in $R$. The map is an isomorphism on
$\delta\ne0$, but it is not an integral identification of the named
KL effective basis with the full character basis.

The finite-spectrum lemma of \fid{F24} shows why tensoring the
derived input does not itself manufacture the required opposite sign:
every nonempty visible real Euler spectrum has a positive top-radius
coefficient in each even tensor power. Input shifts contribute sign
$(-1)^{kr}=1$ there. An extra odd shift of the output reverses a sign
formally, but its identification with path-hypersurface cohomology
still requires the trace identity. Full scalar localization also
cannot be silently realized by finite Weil traces:
$(1+q^n)^{-1}$ is excluded by \fid{F26}. A particular localized
idempotent may have its denominator cancel in character coordinates,
as $b/(1+u)\mapsto(1,0)$ does. The obstruction concerns the specified
non-Laurent all-extension trace, not every object in a localized
category.


\subsection{The actual rank-three algebra has a coherent scheme model}
\begin{finding}[\fid{F19}: a finite-flat center, its collision, and the lost radical]\label{f:rankthree}
Put $\delta=v+v^{-1}$, $D=\delta^2$, and take either
\[
 R=\Q[v,v^{-1}]_{(v-1)}
 \quad\text{or}\quad
 R_0=\Z[1/2,v,v^{-1},\delta^{-1},(D-1)^{-1},(D-2)^{-1}].
\]
Write $R$ for the chosen ring in the formulas below. Let $A$ be the actual nonstandard Hecke algebra $\check H_3$ over this ring, with its specified generators $P_1,P_2$. Define
\[
 z_1=\frac{D-1}{D},\qquad z_2=\frac{4(D-1)}{D^2},\qquad
 C=R[z]/((z-z_1)(z-z_2)).
\]
Then there is an explicit algebra isomorphism
\begin{equation}\label{eq:r3model}
 A\simeq R\times R\times M_2(C).
\end{equation}
On the two scalar blocks $(P_1/D,P_2/D)$ is respectively $(1,1)$ and $(0,0)$; on the matrix block it is
\[
 e=\begin{pmatrix}1&0\\0&0\end{pmatrix},\qquad
 f=\begin{pmatrix}1-z&1\\z(1-z)&z\end{pmatrix}.
\]
In particular $A$ is free of rank ten. At $v=1$ over a field $k$ of characteristic different from 2 and 3,
\[
 A_1\simeq k\times k\times M_2(k[\varepsilon]/\varepsilon^2).
\]
The actual specialization map to the ambient tensor group algebra has kernel the four-dimensional square-zero Jacobson radical and image the diagonal $k[S_3]$.

Over the local characteristic-zero ring, the generic center branches collide to order two:
\begin{equation}\label{eq:r3collision}
 z_1-z_2=\frac{D-1}{D^2}(v-v^{-1})^2.
\end{equation}
The algebra is also the global coherent endomorphism algebra of an explicitly given vector bundle on the ordinary commutative scheme
$S=\Spec R\sqcup\Spec R\sqcup\Spec C$.
\end{finding}
\begin{proof}
The source input is the complete presentation in IV, Theorem~11.16, equations~(111)--(112), printed p.~65. With $x=P_1$, $y=P_2$, $a_1=\delta$ and $a_2=D-2$, it is
\[
 x^2=Dx,\quad y^2=Dy,\quad
 x(yx-a_1^2)(yx-a_2^2)=y(xy-a_1^2)(xy-a_2^2).
\]
Set $e=x/D$, $f=y/D$. These are idempotents, and $z=(e-f)^2=e+f-ef-fe$ is central: multiplying on either side by $e$ gives $e-efe$, and multiplying on either side by $f$ gives $f-fef$. Consequently
\[
 efe=(1-z)e,\qquad fef=(1-z)f.
\]
Expansion of the displayed braid relation gives
\[
 D^5(z-z_1)(z-z_2)(e-f)=0.
\]
For example, its two sides expand using $xyx=D^3(1-z)e$ and $xyxyx=D^5(1-z)^2e$, and the corresponding expressions in $f$; the two factors are $D^2(1-z)-a_i^2$. Thus the presentation is equivalent to the idempotent relations and $g(z)(e-f)=0$, where $g(z)=(z-z_1)(z-z_2)$.

Multiplication by $e-f$ gives $zg(z)=0$. Since $g(0)=z_1z_2$ is a unit, the central Chinese remainder decomposition splits the algebra into its $z=0$ and $g(z)=0$ blocks. In the first, $g(0)(e-f)=0$ forces $e=f$, giving $R[e]/(e^2-e)=R\times R$.

In the second block the central ring is $C$. Every word is a $C$-linear combination of $1,e,f,ef$, by removing repeated idempotents, shortening alternating triples with the preceding identities, and replacing $fe$ by $e+f-ef-z$. The matrices in the statement satisfy those identities. Both $z$ and $1-z$ are units in $C$: their evaluations at the two roots are units, since
\[
 1-z_1=1/D,\qquad 1-z_2=(D-2)^2/D^2.
\]
Equivalently their norms in the quadratic algebra are units, which gives their inverses. Now
\[
 ef(1-e)=E_{12},\qquad (1-e)fe=z(1-z)E_{21}.
\]
Together with $e$ and $1-e$ these generate every matrix unit. The matrices $1,e,f,ef$ are also independent: the lower-left coordinate of a relation first kills its $f$ coefficient; the upper-right then kills its $ef$ coefficient; the lower-right kills its scalar coefficient; the upper-left kills the remaining $e$ coefficient. This proves the isomorphism and freeness.

At $v=1$, $D=4$ and $z_1=z_2=3/4$, giving the dual-number description. Its ideal
$J=0\times0\times M_2(\varepsilon k)$ has square zero and semisimple quotient $k\times k\times M_2(k)$, so it is exactly the Jacobson radical and has dimension four. In the original ambient inclusion,
\[
 P_i=2(1+s_i\otimes s_i).
\]
These generate the diagonal group algebra, since $s_i\otimes s_i=P_i/2-1$. Its dimension is six and it is semisimple in the stated characteristics. The image of $J$ is a nilpotent ideal of that semisimple image and hence is zero; dimension then makes $J$ the entire kernel.

Equation~\eqref{eq:r3collision} follows from $D-4=(v-v^{-1})^2$. Evaluation at the two roots identifies
\[
 C=\{(u_1,u_2)\in R^2:u_1-u_2\in(z_1-z_2)R\}.
\]
Indeed $a+bz$ evaluates to a pair whose difference is $b(z_1-z_2)$; conversely such a pair gives this $b$ and $a=u_1-bz_1$. Injectivity follows because $R$ is a domain and the two generic roots differ. At $v=1$, the factor multiplying $(v-1)^2$ in \eqref{eq:r3collision} is a unit over the local characteristic-zero ring. The generic central idempotents divide by $z_1-z_2$ and do not extend through this collision; the finite algebra $C$ over this local ring has local special fiber and is itself local, so has no nontrivial idempotent.

Finally take the trivial vector bundle $E$ of ranks 1, 1 and 2 on the three components of $S$. Its algebra of global coherent endomorphisms is exactly $R\times R\times M_2(C)$. The displayed matrices retain the named nonstandard generators. This proves the coherent realization.
\end{proof}
This strengthens the construction beyond an abstract matching-dimension model. It also corrects a possible flatness misreading: the family remains rank ten; the \emph{specialized inclusion} loses injectivity. IV, Remark~11.9 (p.~61), already distinguishes abstract specialization from ambient image; Proposition~11.14 (pp.~64--65) gives the generic simple dimensions $1,1,2,2$ and the merging of the two two-dimensional representations.

Over a finite field of good characteristic at $v=1$, $S$ has two reduced rational points and one double rational point. Its geometric \(\ell\)-adic cohomology is that of three rational points, since a nilpotent thickening does not change the \emph{\'etale topos} (\cite[Proposition~59.45.4, tag~03SI]{StacksNil}). Thus it is pure in degree zero and has dimension three; this purity does not see the coherent radical. Replacing the three vector-bundle ranks by constant \(\ell\)-adic sheaves gives a chosen endomorphism algebra of dimension six on the reduced points, not a faithful realization of the ten-dimensional coherent algebra.

A full endomorphism algebra of a finite semisimple object over an algebraically closed coefficient field is a product of matrix algebras; it cannot be $A_1$. An algebra acting as a \emph{subalgebra} of such an endomorphism algebra can be nonsemisimple, so that broader construction is not excluded. The coherent bundle above models the algebra; \fid{F21} below supplies the actual right tensor module. Neither construction identifies VIII's based left/right structure or formal quantum grading with Frobenius weights. Those remain precise comparison tasks. In Grothendieck's language the result lives in finite-flat schemes and coherent $\End(E)$, not in an unproved RH; EGA~I's affine scheme formalism, SGA~3's module/comodule language, and cote~90, scan~17, supply the framework, while the complete algebra proof is given here.

\subsection{The actual tensor action and its coherent deformation}
\begin{finding}[\fid{F21}: the specified rank-three tensor lattice is realized]\label{f:r3tensor}
Use the ring $R_0$, algebra $A$, center $C$, roots $z_1,z_2$ and named
generators of \fid{F19}. Let $V=W=R_0^2$ and give
$M=(V\otimes W)^{\otimes3}$ the actual right tensor-monomial action of
GCT~IV, \S6.1, p.~33, and Definition~11.1, equation~(93). Write
$\epsilon_+$ for $P_i=D$, $\epsilon_-$ for $P_i=0$, and
\[
 N(a):\quad P_1=\begin{pmatrix}D&0\\a&0\end{pmatrix},\qquad
 P_2=\begin{pmatrix}0&a\\0&D\end{pmatrix}.
\]
These are right row matrices, as in IV Proposition~11.14,
equation~(110). There is an action-preserving isomorphism over $R_0$,
\begin{equation}\label{eq:r3tensor}
 M\simeq\epsilon_+^{20}\oplus\epsilon_-^4
       \oplus N(\delta)^{16}\oplus N(D-2)^4.
\end{equation}
In particular every displayed change of basis is regular at $v=1$ in
characteristic different from $2,3$; no colliding-root difference is inverted.

Put $L_i=C/(z-z_i)$ and
$S=\Spec R_0\sqcup\Spec R_0\sqcup\Spec C$. For
\[
 E=(\mathcal O,\mathcal O,\mathcal O_C^2),\qquad
 \mathcal F=(\mathcal O^{20},\mathcal O^4,
             \widetilde L_1^{16}\oplus\widetilde L_2^4),
\]
the coherent model realizes this same action as
\[
 M\simeq\Gamma(S,\mathcal Hom(E,\mathcal F)),\qquad A=\Gamma(S,\End E),
\]
with right action by precomposition. It is faithful before specialization
and kills the four-dimensional radical after $v=1$.
\end{finding}
\begin{proof}
Let $L=R_0^2$. The cited ordinary tensor action, with the consistent
normalization $c_i=T_i+v^{-1}$, is
\[
 \mathbf v_{\mathbf k}c_i=
 \begin{cases}
 v\mathbf v_{\mathbf k}+\mathbf v_{\mathbf k s_i},&k_i<k_{i+1},\\
 \delta\mathbf v_{\mathbf k},&k_i=k_{i+1},\\
 v^{-1}\mathbf v_{\mathbf k}+\mathbf v_{\mathbf k s_i},&k_i>k_{i+1}.
 \end{cases}
\]
On $(112,121,211)$ it gives
\[
 c_1=\begin{pmatrix}\delta&0&0\\0&v&1\\0&1&v^{-1}\end{pmatrix},\qquad
 c_2=\begin{pmatrix}v&1&0\\1&v^{-1}&0\\0&0&\delta\end{pmatrix}.
\]
The row-basis matrix
\[
 B_3=\begin{pmatrix}v^2&v&1\\-1-v^{-2}&v&1\\
                         -v^{-1}&-v^{-2}&\delta\end{pmatrix}
\]
has determinant $(D-1)^2$. Multiplying gives
\[
 B_3c_1B_3^{-1}=\begin{pmatrix}\delta&0&0\\0&\delta&0\\0&1&0\end{pmatrix},
 \quad
 B_3c_2B_3^{-1}=\begin{pmatrix}\delta&0&0\\0&0&1\\0&0&\delta\end{pmatrix}.
\]
Thus this block is $\epsilon\oplus U$, with
\[
 c_1|_U=\begin{pmatrix}\delta&0\\1&0\end{pmatrix},\qquad
 c_2|_U=\begin{pmatrix}0&1\\0&\delta\end{pmatrix}.
\]
On $(122,212,221)$ the original $c_1,c_2$ matrices interchange; applying
$B_3$ and swapping its last two rows again gives $\epsilon\oplus U$.
The weight spaces $111$ and $222$ are each $\epsilon$. Since the determinants
are units, this proves $L^{\otimes3}=\epsilon^4\oplus U^2$ over $R_0$,
not just over its fraction field.

By $P_i=c_i\otimes c_i+(c_i-\delta)\otimes(c_i-\delta)$,
$\epsilon\otimes\epsilon=\epsilon_+$, and each mixed tensor is
$N(\delta)$. On $U\otimes U$ in the order $(11,12,21,22)$,
\[
 P_1=\begin{pmatrix}D&0&0&0\\\delta&0&0&0\\\delta&0&0&0\\
                         2&-\delta&-\delta&D\end{pmatrix},\quad
 P_2=\begin{pmatrix}D&-\delta&-\delta&2\\0&0&0&\delta\\
                         0&0&0&\delta\\0&0&0&D\end{pmatrix}.
\]
Take
\[
 B_4=\begin{pmatrix}
 \delta&-1&-1&\delta\\0&-1&1&0\\
 2-D&-\delta&-\delta&D\\-D&\delta&\delta&D-2
 \end{pmatrix},\qquad \det B_4=8(D-1)^2.
\]
Direct multiplication yields
\[
 B_4P_1B_4^{-1}=\begin{pmatrix}D&0&0&0\\0&0&0&0\\0&0&D&0\\0&0&D-2&0\end{pmatrix},
\quad
 B_4P_2B_4^{-1}=\begin{pmatrix}D&0&0&0\\0&0&0&0\\0&0&0&D-2\\0&0&0&D\end{pmatrix}.
\]
This is $\epsilon_+\oplus\epsilon_-\oplus N(D-2)$, again over $R_0$.
In $(\epsilon^4\oplus U^2)\otimes(\epsilon^4\oplus U^2)$ the unmixed
trivial contribution has multiplicity $16$, the two mixed contributions
have total multiplicity $16$, and $U\otimes U$ occurs four times.
This proves \eqref{eq:r3tensor}, whose rank is $20+4+32+8=64$.

It remains to match the named action to the center model, rather than
merely match dimensions. For $a=\delta,D-2$, put $q_a=a/D$, a unit, and
\[
 S_a=\begin{pmatrix}1&0\\-q_a^2&q_a\end{pmatrix}.
\]
The matrices of $S_a(P_1/D)S_a^{-1}$ and $S_a(P_2/D)S_a^{-1}$ are
respectively
\[
 \begin{pmatrix}1&0\\0&0\end{pmatrix},\qquad
 \begin{pmatrix}1-z_i&1\\z_i(1-z_i)&z_i\end{pmatrix},
 \quad z_i=1-q_a^2.
\]
They are exactly the standard row module for $M_2(C)$ evaluated on branch
$i$. The right-module formula with $\mathcal Hom$ follows. The two branch
evaluations jointly embed $C$ into $R_0^2$, so their modules have joint
annihilator zero; both scalar blocks also occur. At $v=1$ the branch
modules become the same residue module of $k[\varepsilon]/\varepsilon^2$,
and the central nilpotent acts by zero. Its matrix ideal is the radical
of \fid{F19}, proving the last assertion.
\end{proof}
The generic multiplicities themselves are already in IV \S12.3; the
additional assertion here is the displayed regular action isomorphism.
The coherent sheaf $\mathcal F$ is not locally free on the collision
component: its two generic branch ranks are $16$ and $4$. This result
does not identify a simultaneous canonical left quantum-group action
or VIII's based tensor lattice with the monomial lattice.

\begin{finding}[\fid{F22}: constant nearby-cycle purity does not detect collision order]\label{f:collision}
Let $R=k[[t]]$ and, for $m\ge1$, set
\[
 C_m=\{(u,w)\in R^2:u\equiv w\pmod{t^m}\}
     \simeq R[x]/(x(x-t^m)).
\]
For $m\ge n\ge1$, $C_m\subset C_n$ induces a finite universal
homeomorphism $\Spec C_n\to\Spec C_m$ over $\Spec R$, although
\[
 \operatorname{length}_R(R^2/C_m)=m,\qquad
 \operatorname{cond}(C_m,R^2)=t^mR^2.
\]
Their corresponding ordinary \'etale coefficient categories and
cohomology are equivalent. With $k=\F_q$ and unshifted constant
$\Q_\ell$ coefficients, the unique special-point nearby and vanishing
cycles are
\[
 R\Psi\Q_\ell=\Q_\ell^2[0],\qquad
 R\Phi\Q_\ell=\Q_\ell[0],
\]
with trivial inertia and Frobenius. These constant data, including every
finite-extension trace and tensor power, are independent of $m$.
The completed center of \fid{F19} in good characteristic is $C_2$.
\end{finding}
\begin{proof}
Map $x$ to $(0,t^m)$ to obtain the displayed description. For $m\ge n$,
$C_n$ is a finite $C_m$-module since it is already generated by two
elements over the subring $R$. After inverting $t$, the map is the
identity on $k((t))^2$. Each special fiber has one point, with residue
field $k$, and the map induces the identity on that field. It is
therefore integral, surjective and radicial, hence a universal
homeomorphism. The difference map
$(u,w)\mapsto w-u\bmod t^m$ has kernel $C_m$ and gives
$R^2/C_m\simeq R/(t^m)$. The integral closure is $R^2$; a pair multiplies
every element of $R^2$ back into $C_m$ precisely when both coordinates
are divisible by $t^m$. This proves the coherent formulas.

Topological invariance of the \'etale site gives the claimed equivalence,
compatible with the base trait; precise locators are Stacks
Theorem~59.45.2, tag~04DZ, and Proposition~59.45.4, tag~03SI,
the former citing EGA~IV, 18.1.2 \cite{StacksNil}.
This is an equivalence of coefficient-sheaf topoi, not an identification
of coherent structure sheaves as ringed objects.

After strict henselization at the special point, the geometric generic
fiber is two points. Indeed the finite local algebra over the strictly
henselian trait is strictly henselian, and its two branches are already
split over the fraction field. The defining nearby-cycle stalk is
therefore the cohomology of two geometric points: $\Q_\ell^2$ in degree
zero, with trivial inertia and Frobenius. The specialization map from
$i^*\Q_\ell=\Q_\ell[0]$ is diagonal. The defining triangle
\[
 i^*\Q_\ell\longrightarrow R\Psi\Q_\ell
 \longrightarrow R\Phi\Q_\ell\longrightarrow i^*\Q_\ell[1]
\]
thus gives the stated cokernel in degree zero and no other groups.
This also follows by finite base change, Stacks Lemma~59.91.9,
tag~0A4E. Every object and map just described is independent of $m$.
The traces of tensor powers $r$ of $R\Psi$ and $R\Phi$ are respectively
$2^r$ and $1$ over every finite extension.

Finally \eqref{eq:r3collision} says that the actual root difference is
a unit times $(v-1)^2$. Translate $z$ by a root and rescale the translated
coordinate by that unit after completion. The relation becomes
$x(x-t^2)=0$, where $t=v-1$. This identifies the center with $C_2$.
\end{proof}
This is an information limit for the specified constant coefficient
families. A nilpotent algebra action supplied separately, or an
$m$-dependent coefficient construction, is additional information and
is not ruled out. A sign criterion applied only to these common traces
can establish their correct weight-zero purity but cannot determine
$m$. The next finding exhibits coherent information that does retain it.

\begin{finding}[\fid{F25}: the commutant defect, coherent extensions, and a regular basis]\label{f:r3deformation}
Use \fid{F21} over $R=\Q[v,v^{-1}]_{(v-1)}$ and write
$\Delta=z_1-z_2$. Then
\[
 \End_A(M)=M_{20}(R)\times M_4(R)\times M_{16}(R)\times M_4(R).
\]
This is free of rank $688$. Its base-change map to the full special
commutant
\[
 \End_{A_1}(M_1)=M_{20}(k)\times M_4(k)\times M_{20}(k)
\]
is injective with cokernel of dimension $128$; the latter commutant has
dimension $816$. Moreover
\[
 \operatorname{Ext}^1_C(L_1,L_2)
 =\operatorname{Ext}^1_C(L_2,L_1)=R/(\Delta),\qquad
 \operatorname{Ext}^1_C(L_i,L_i)=0,
\]
and $\operatorname{Ext}^1_A(M,M)=(R/(\Delta))^{128}$, of $R$-length $256$.

There is also a regular center basis $1,y$, with
\begin{equation}\label{eq:r3basis}
 y=D^2(z_1-z),\qquad
 y^2=\lambda y,\qquad
 \lambda=(v-v^{-1})^2(v^2+1+v^{-2}).
\end{equation}
Together with the two scalar-block units, $E_{ij},yE_{ij}$ form a basis
of the actual rank-ten algebra over $R_0$, with structure constants
among $0,1,\lambda$. This elementary basis is not asserted to be VIII's
designated canonical basis.
\end{finding}
\begin{proof}
Central blocks separate the two scalar sectors and the matrix sector.
Morita equivalence identifies the last sector's multiplicity module
with $L_1^{16}\oplus L_2^4$. A homomorphism $L_1\to L_2$ is determined
by $r\in R$ satisfying $(z_2-z_1)r=0$. Since $\Delta$ is a nonzero
element of the domain $R$, this Hom group is zero, as is the reverse
one; $\End_C(L_i)=R$. This gives the asserted commutant. At $v=1$ the
two branch modules become the same simple residue module. Their
multiplicity spaces merge into $k^{20}$. The map on these factors is
the block-diagonal embedding $M_{16}(k)\times M_4(k)\hookrightarrow
M_{20}(k)$; the missing off-diagonal blocks each have dimension $16\cdot4$.

For the extension calculation put $f_i=z-z_i$. The following is an
exact free resolution of $L_1$:
\[
 \cdots\xrightarrow{f_1}C\xrightarrow{f_2}C
             \xrightarrow{f_1}C\longrightarrow L_1\longrightarrow0.
\]
To check exactness, lift to the domain $R[z]$ and cancel $f_1$ or $f_2$
from divisibility by $f_1f_2$. Applying $\Hom_C(-,L_2)$ gives
$R\xrightarrow{-\Delta}R\xrightarrow0R\xrightarrow{-\Delta}R\to\cdots$,
so degree one is $R/(\Delta)$. Exchanging indices gives the reverse
group. Applying $\Hom_C(-,L_1)$ instead begins
$R\xrightarrow0R\xrightarrow\Delta R$ and has zero degree one.
Morita equivalence and the multiplicities in \eqref{eq:r3tensor} give
$16\cdot4+4\cdot16=128$ copies; scalar blocks contribute no extensions
of their free $R$-modules. Since $\Delta$ is a unit times $(v-1)^2$,
each crossed quotient has length two. For the general $C_m$ the same
resolution gives crossed extension group $R/(t^m)$, retaining precisely
the collision order omitted by the constant \'etale data.

Finally $D$ is a unit and the affine change $z\mapsto y$ is invertible.
The two evaluations of $y$ are $0$ and $D^2(z_1-z_2)$. Their product
relation yields $y^2=D^2\Delta y$. Since
\[
 D^2\Delta=(D-1)(D-4),\quad
 D-4=(v-v^{-1})^2,\quad D-1=v^2+1+v^{-2},
\]
this is \eqref{eq:r3basis}. Matrix-unit multiplication proves the
structure-constant assertion. At $v=1$, $\lambda=0$ and the four
$yE_{ij}$ span the radical.
\end{proof}
IV \S12.3 already records the numbers $688$ and $816$. Here the free
rank-$688$ commutant family is kept distinct from the full
special-fiber commutant, and their comparison map is explicit.

The scalar in \eqref{eq:r3basis} has a precise ordinary geometric trace:
with $Q=v^2$,
\[
 \lambda=Q^{-2}(Q-1)^2(1+Q+Q^2)
       =Q^{-2}\#(\mathbb G_m^2\times\PP^2)(\F_Q).
\]
It is the compact-support alternating trace with Tate twist $(2)$.
The remaining factor $v^{-2}(1+v^2+v^4)$ is the centered Tate polynomial
of $\PP^2$. This mixed-cohomology trace interpretation of one scalar
does not supply functorial multiplication or associativity constraints,
and it does not assert VIII's consecutive-$v$ unimodality convention.
Indeed $\lambda=v^4-v^2-v^{-2}+v^{-4}$ has negative Laurent coefficients.


\subsection{A positive joint action on the actual rank-three tensor lattice}
The positive right matrices are already visible in \fid{F21}. The
following calculation strengthens that conclusion to the actual
ordinary left operators, tensor-product bar and full regular
commutant. Revised IV already gives a rank-three based construction
in Example~19.2, pp.~132--133; we do not claim priority or identify
our adapted basis with that particular designated basis.

\begin{finding}[\fid{F43}: actual positive actions, bar and graded paths]\label{f:positivejoint}
Put
\[
 \Lambda=\Z[v,v^{-1}],\quad \delta=v+v^{-1},\quad
 D=\delta^2,\quad h=D-1,
 \qquad S=\Lambda[(2h)^{-1}].
\]
The actual tensor lattice $M_\Lambda=(\Lambda^2\otimes\Lambda^2)^{\otimes3}$
has a free full-rank sublattice with basis $\mathcal B$ whose
change-of-basis determinant is $\pm2^{12}h^{72}$.
Thus $\mathcal B$ is a basis of the actual $M_S$ and is regular
at $v=1$ in characteristic different from $2,3$.
On this same basis, the actual right $P_1,P_2$, the ordinary left
$E_V,F_V,E_W,F_W,K_V^{\pm1},K_W^{\pm1}$, and the elementary
matrix units of the full commutant have nonnegative Laurent matrices.
Every vector of $\mathcal B$ is fixed by IV's actual ordinary
tensor-product bar. The commutant is
\[
 \mathcal C_S=\End_{A_S}(M_S)
 \simeq M_{20}(S)\times M_4(S)\times M_{16}(S)\times M_4(S),
\]
and its displayed matrix-unit basis acts by $0/1$ matrices.

For a nonempty right-generator word $w$ of length $L$, let $2s$
be its number of cyclic changes of letter. Then
\begin{equation}\label{eq:allwordpositive}
 \Tr_M(P_w)=20D^L+16D^{L-s}
              +4D^{L-2s}(D-2)^{2s}.
\end{equation}
Every entry of every operator sequence in the listed positive
operators counts paths in one explicit 64-state graded graph.
For explicitly written words this gives polynomial-time coefficient
evaluation and $O(L)$-bit positive witnesses, at fixed tower rank three.
\end{finding}

\paragraph{Left and right conventions.}
Right-generator matrices multiply in the written algebra order.
We also display a left operator by rows recording its coefficients
on basis vectors. For left algebra products these row matrices
reverse composition order: $\rho_{\rm row}(ab)=
\rho_{\rm row}(b)\rho_{\rm row}(a)$.
The path statements concern an explicitly ordered sequence of
applied operators. Reversing the written left-algebra word leaves
all the asserted positivity and complexity properties unchanged.

\begin{proof}
\par\smallskip\noindent\textbf{The ordinary tensor factors.}\quad
Use $\mathbf v_{ijk}$ for tensor monomials. The F21 matrix
\[
 B_3=\begin{pmatrix}
 v^2&v&1\\-1-v^{-2}&v&1\\-v^{-1}&-v^{-2}&\delta
 \end{pmatrix},\qquad \det B_3=h^2,
\]
acts on $(112,121,211)$ and on $(122,212,221)$, with its last
two rows interchanged in the second weight space. Write
\[
 b_0=\mathbf v_{111},\quad
 b_1=v^2\mathbf v_{112}+v\mathbf v_{121}+\mathbf v_{211},
\]
\[
 b_2=v^2\mathbf v_{122}+v\mathbf v_{212}+\mathbf v_{221},
 \qquad b_3=\mathbf v_{222}.
\]
Let $u_{+,1},u_{+,2}$ be the last two rows of $B_3$ in the first
weight space, and $u_{-,1},u_{-,2}$ the corresponding reversed
rows in the second. Set $u_{\pm,j}=g_\pm\otimes e_j$.
Let $W_3$ have basis $b_0,b_1,b_2,b_3$, $W_1$ basis $g_+,g_-$,
and $U$ basis $e_1,e_2$.
Here the subscripts are $\mathfrak{sl}_2$ highest weights;
the polynomial $\GL_2$ labels are $(3)$ and $(2,1)$, so
$W_1$ retains its determinant twist and polynomial degree three.
The actual ordinary Hecke action from \fid{F21} is
\[
 (S^2)^{\otimes3}=(W_3\otimes\epsilon)\oplus(W_1\otimes U),
 \qquad c_i|_\epsilon=\delta,
\]
\[
 c_1|_U=\begin{pmatrix}\delta&0\\1&0\end{pmatrix},\qquad
 c_2|_U=\begin{pmatrix}0&1\\0&\delta\end{pmatrix}.
\]
The determinant of the rank-eight basis change is $\pm h^4$.

The actual ordinary coproduct is
\cite[\S4.1, equation~(29), p.~22]{IV}
\[
 \Delta(E)=E\otimes K^{-1}+1\otimes E,
 \qquad\Delta(F)=F\otimes1+K\otimes F.
\]
Thus on a tensor word, $E$ changes a $2$ to $1$ with coefficient
$v^{-\sum\operatorname{wt}(i)}$, summed over later positions;
$F$ changes a $1$ to $2$ with coefficient
$v^{\sum\operatorname{wt}(i)}$, summed over earlier positions,
where $\operatorname{wt}(1)=1$, $\operatorname{wt}(2)=-1$.
Substituting the eight displayed basis vectors gives
\[
 Fb_0=b_1,\quad Fb_1=\delta b_2,\quad Fb_2=hb_3,\quad Fb_3=0,
\]
\[
 Eb_0=0,\quad Eb_1=hb_0,\quad Eb_2=\delta b_1,\quad Eb_3=b_2,
\]
\[
 Fg_+=g_-,\quad Fg_-=0,\quad Eg_-=g_+,\quad Eg_+=0.
\]
The $K$ eigenvalues are $v^{3-2j}$ on $b_j$ and $v^{\pm1}$
on $g_\pm$. All coefficients are nonnegative Laurent polynomials.

The bar assertion is also explicit. In the first weight space,
IV Figure~1, p.~9, gives the upper canonical vectors
\[
 c_{112}=\mathbf v_{112},\quad
 c_{121}=\mathbf v_{121}-v^{-1}\mathbf v_{112},\quad
 c_{211}=\mathbf v_{211}-v^{-1}\mathbf v_{121}.
\]
They are fixed by the actual tensor bar
\cite[\S6.2, Theorem~6.4, p.~34]{IV}.
The coordinates of $(b_1,u_{+,1},u_{+,2})$ in this basis are
\[
 \begin{pmatrix}h&\delta&1\\0&\delta&1\\0&1&\delta\end{pmatrix}.
\]
Every coefficient is bar invariant. The second weight space has
the same calculation with its last two rows reversed, and the two
endpoint vectors are canonical. Hence all eight adapted vectors
are fixed by the actual bar.

\par\smallskip\noindent\textbf{The full bimodule.}\quad
Regroup the rank-64 tensor module into the two ordinary rank-eight
factors and apply the preceding basis in each. On every Specht
factor $U\otimes U$ apply the F21 matrix
\[
 B_4=\begin{pmatrix}
 \delta&-1&-1&\delta\\0&-1&1&0\\
 2-D&-\delta&-\delta&D\\-D&\delta&\delta&D-2
 \end{pmatrix},\qquad \det B_4=8h^2.
\]
It gives $\epsilon_+\oplus\epsilon_-\oplus N(D-2)$,
while the mixed factors give $N(\delta)$, with
\[
 N(a):\quad
 P_1=\begin{pmatrix}D&0\\a&0\end{pmatrix},\qquad
 P_2=\begin{pmatrix}0&a\\0&D\end{pmatrix}.
\]
Consequently the actual bimodule is
\begin{align*}
 M_S={}&(W_3\otimes W_3\oplus W_1\otimes W_1)\otimes\epsilon_+
       \oplus(W_1\otimes W_1)\otimes\epsilon_-\\
 &\oplus(W_3\otimes W_1\oplus W_1\otimes W_3)\otimes N(\delta)
       \oplus(W_1\otimes W_1)\otimes N(D-2).
\end{align*}
The multiplicity ranks are $20,4,16,4$.
The $B_4$ change acts only on Specht factors, so the actual
ordinary left matrices remain the positive matrices just computed.
The right coefficients are $D=v^2+2+v^{-2}$,
$\delta=v+v^{-1}$ and $D-2=v^2+v^{-2}$.
The tensor bar for the direct-product quantum group is the product
of the two ordinary tensor bars after regrouping: its quasi-$R$
factors $\Theta_V,\Theta_W$ act on separate factors.
All $B_4$ coefficients are bar invariant, proving the full bar claim.

The rank-eight tensor-square basis change has determinant $h^{64}$,
by $\det(T\otimes T)=\det(T)^{16}$.
The four $B_4$ blocks contribute $(8h^2)^4=2^{12}h^8$.
All rows have coefficients in $\Lambda$, so their free span is a
sublattice with determinant $\pm2^{12}h^{72}$, and it becomes the
whole actual lattice over the stated $S$.

For the commutant, different generic simple sectors have zero
cross homomorphisms over the domain $S$.
This can also be seen from $(P_1-P_2)^2$, whose scalar values
on the scalar, $N(\delta)$ and $N(D-2)$ sectors are respectively
$0,D(D-1),4(D-1)$; $P_1$ separates the two scalar sectors.
The endomorphisms of each $N(a)$ are precisely scalars:
commuting a matrix $\left(\begin{smallmatrix}x&y\\z&t\end{smallmatrix}\right)$
with $P_1,P_2$ forces $y=z=0$ and $a(x-t)=0$, hence $x=t$.
The multiplicities give the stated product of matrix rings.
Each elementary left matrix unit changes only a multiplicity basis
vector, retaining its right factor, and therefore has a $0/1$
action on $\mathcal B$. Its usual multiplication constants are
also $0/1$. This is the actual regular commutant; identification
with every integral form of the dual FRT algebra is not asserted.

\par\smallskip\noindent\textbf{All right words.}\quad
In $N(a)$ put $u_1=(D,a)^T$, $u_2=(a,D)^T$, so $P_i=u_i e_i^T$.
If $w=i_1\cdots i_L$ has $t$ ordinary adjacent changes, multiplying
rank-one matrices gives
\[
 P_w=D^{L-1-t}a^t u_{i_1}e_{i_L}^T.
\]
Taking the trace introduces the final cyclic comparison.
The cyclic change count is even, say $2s$, and therefore
$\Tr_{N(a)}P_w=D^{L-2s}a^{2s}$.
Add the scalar-sector contribution $20D^L$, use
$\delta^{2s}=D^s$, and obtain \eqref{eq:allwordpositive}.
The empty word has trace $64$.

\par\smallskip\noindent\textbf{Positive paths and their size.}\quad
Replace an entry $\sum_d n_dv^d$ of a selected operator by $n_d$
directed edges of degree $d$ between its two basis states.
For example $P_1$ resets a rank-two block to state 1, with loop
degrees $-2,0,0,2$ and other-edge degrees $-1,1$ or $-2,2$;
$P_2$ resets it to state 2. Matrix multiplication proves, by
induction, that an operator sequence counts its color-prescribed
paths, with exponents equal to total degrees. Closed paths give
the trace. These are the actual conjugated operators, rather than
a pair chosen only to reproduce their traces.

For an explicitly written length-$L$ sequence, right generators
have at most four outgoing labeled edges, $E/F$ at most three,
and $K^{\pm1}$ or elementary commutant units at most one.
Every step has degree between $-3$ and $3$.
There are at most $4^L$ paths from a fixed start and at most
$64\cdot4^L$ closed paths, so witnesses and counts have $O(L)$
bits. Dynamic programming on 64 states and degrees in $[-3L,3L]$
uses polynomial bit time. This fixes tower rank three and uses
explicit words, not compressed binary powers or general
Kronecker/plethysm inputs.
\end{proof}

\paragraph{An exact graded model, and its summand boundary.}
For each matrix entry take a finite graded vector space with the
indicated homogeneous basis. These kernels define exact functors
on 64-tuples of finite graded vector spaces by direct sums and
tensor products. Composition is path concatenation. Their
Grothendieck-group action is exactly the displayed action on
$\Lambda^{64}$, and a word has the finite graded trace object
given by the sum of its diagonal kernel spaces.
The quadratic relation lifts by the tensor symmetry after a reset.
No coherent lift of all linear nonstandard relations or a
designated GCT categorification is claimed from the $K_0$ equality.

Closure under arbitrary direct summands has a proved limitation.
Keep just one copy of each right simple type, giving a faithful
rank-six module over the parameter ring (its special fiber need not
be faithful). At $v=1$, the two matrices $N(\delta)$
and $N(D-2)$ become identical. Thus every element of the actual
regular algebra acts identically on those two specialized blocks.
But the graded $P_1$ kernel has a direct summand consisting of
one degree-zero loop at state 1 of the $N(\delta)$ block, zero
elsewhere. Its action is a matrix unit supported on only that
branch and fails the preceding specialization condition.
It is outside the actual regular algebra. Hence closing this
semisimple matrix-functor model under all summands adds operations
beyond the nonstandard order, even with multiplicity copies removed.

The formal grading is also distinct from an arithmetic realization.
Closed-word traces in \eqref{eq:allwordpositive} have only even
$v$ powers, whereas individual $N(\delta)$ entries have odd powers.
Thus $v=q^{n/2}$ needs extra weight-one data at the entry level;
it is not automatically a geometric Tate construction.
Choosing $v=q^n$ uses a different, full-Tate normalization.

\paragraph{Two limits of simultaneous positivity.}
The actual factor swap $\tau$ fixes the first, third and fourth
rows of $B_4$ and negates its second row.
On the $\epsilon_-$ multiplicity space it is consequently
$-\mathrm{flip}$ on $W_1\otimes W_1$, of trace $-2$.
No basis of that summand makes $\tau$ nonnegative.
More generally, if $e_-$ is its actual projector, then
$\Tr(e_-\tau)=-2$, so $e_-$ and $\tau$ cannot both have
nonnegative matrices in any basis. The positive commutant cone
contains $e_-$. This agrees with IV's signed extension and is not
an objection to a based-module definition allowing that sign.

For $Q_i=DI-P_i$, negating the second vector of every rank-two
right block makes both $Q_i$ positive, preserves the actual bar
and leaves the ordinary left matrices positive.
But simultaneous positivity of both pairs $P_i,Q_i$ is impossible:
their off-diagonal entries are negatives of each other, so all
would have to vanish. This would make $P_1,P_2$ commute, whereas
their rank-two commutator has nonzero entry $Da$.

\paragraph{A strict positive center basis is also excluded.}
Over $R=\Q[v,v^{-1}]_{(v-1)}$, the actual rank-two center cannot
have an $R$-basis $\{1,b\}$ with
\[
 b^2=a(v)+c(v)b,\qquad a,c\in\N[v,v^{-1}].
\]
Its dual-number fiber forces the discriminant
$c(1)^2+4a(1)$ to vanish. Nonnegative coefficients imply
$a(1)=c(1)=0$, hence $a=c=0$ identically.
The generic algebra would then be $K[b]/b^2$, contradicting its
actual split reduced generic fiber. This excludes exactly the
stated unit-containing rank-two basis. F25's factor
$(v-v^{-1})^2(D-1)$ has negative expanded coefficients but is of
the factored kind allowed in VIII; the argument does not refute
that weaker positivity convention.

\paragraph{The prescribed-lattice gate remains.}
The adapted determinant is not a unit in $\Lambda$, so this basis
does not identify the original global tensor-monomial lattice.
It also does not lie in the ordinary upper lattice at infinity:
$b_1\otimes b_0$ has a $v^2$ tensor-monomial coefficient.
IV's paragraph after Conjecture~19.1 explicitly distinguishes its
desired nonstandard integral form and local lattice from these
ordinary ones. We have not identified our basis with those
prescribed nonstandard data or with Example~19.2.
Indeed, scaling an entire isotypic sector by the bar-invariant
unit $h^r$ preserves all the proved positive matrices and bar
properties over $S$, while changing the global and local lattices.
These properties therefore do not select VIII's prescribed
crystal reduction, duality and minimum-degree globalization.


\subsection{Coherent Frobenius: the collision obstruction and an actual normalization repair}
Here Frobenius means absolute Frobenius on characteristic-$p$ coherent
objects. It is not yet geometric Frobenius on an $\ell$-adic Weil
realization. The distinction between pullback on the singular center
and pullback on the parameter trait gives different, computable answers.

\begin{finding}[\fid{F36}: unit Frobenius depends on the coherent base]\label{f:coherentfrobenius}
Let $k=\F_q$ have characteristic $p>3$, let $R=k[[t]]$, and put
\[
 \Delta=t^m u(t),\quad u(t)\in R^*,\quad m\ge1,\qquad
 C=R[x]/(x(x-\Delta)),\quad L_1=C/(x),\quad L_2=C/(x-\Delta).
\]
For absolute Frobenius $F_C(c)=c^p$,
\[
 F_C^*L_i\simeq R\oplus R/(\Delta^{p-1})
 \quad\text{as }R\text{-modules}.
\]
Thus no unit $C$-linear Frobenius isomorphism $F_C^*L_i\simeq L_i$
exists. Every positive Tor degree of the derived pullback is also
nonzero. In the actual algebra and 64-dimensional module of
\fid{F19}/\fid{F21}, the center has $m=2$ and underived algebra
pullback adds torsion of $R$-length $80(p-1)$.

Nevertheless the actual module has a unit \emph{parameter-base}
semilinear Frobenius structure, compatible with the intrinsic
entrywise Frobenius on $A=R\times R\times M_2(C)$ after transport
through its regular adapted basis. On the normalized center the
actual module has a unit coherent Frobenius model preserving its
original $A$-action under restriction. Pushing this model back through
the collision gives the nonunit surjection just computed.
\end{finding}
\begin{proof}
\par\smallskip\noindent\textbf{The two branches.}\quad
Since $x^2=\Delta x$, characteristic $p$ gives
\[
 x^p=\Delta^{p-1}x,\qquad
 (x-\Delta)^p=\Delta^{p-1}(x-\Delta).
\]
Tensoring the cyclic presentations along $F_C$ gives
\[
 F_C^*L_1=C/(\Delta^{p-1}x),\qquad
 F_C^*L_2=C/(\Delta^{p-1}(x-\Delta)).
\]
The ring $C$ is free over $R$ with basis $1,x$. The first ideal is
exactly $R\Delta^{p-1}x$, because
$x(a+bx)=(a+b\Delta)x$. This proves the stated $R$-module
decomposition for $L_1$; use the basis $1,x-\Delta$ for $L_2$.
The torsion has length $m(p-1)$, while each $L_i$ is $R$-free.
Hence no possible $C$-linear isomorphism can be a unit structure.
The decomposition asserted here is only over $R$, not a split
decomposition of $C$-modules.

There is a natural surjection $\theta_i:F_C^*L_i\to L_i$,
since $(x_i^p)\subset(x_i)$, where $x_1=x$ and $x_2=x-\Delta$.
It linearizes the branch map $r\mapsto r^p$. Its precise kernels are
\[
 \ker\theta_1\simeq L_2/\Delta^{p-1}L_2,
 \qquad \ker\theta_2\simeq L_1/\Delta^{p-1}L_1:
\]
the first is generated by the class of $x$, annihilated by $x-\Delta$,
and the second by the class of $x-\Delta$, annihilated by $x$.
For the $p^e$-power Frobenius all formulas replace $p-1$ by
$p^e-1$; in particular the coefficient-field-linear $q$th-power
version has length $m(q-1)$ per branch.

\par\smallskip\noindent\textbf{Derived pullback on the same object.}\quad
Put $f_1=x$, $f_2=x-\Delta$. Cancellation in the domain $R[X]$
gives $\operatorname{ann}_C(f_1)=(f_2)$ and
$\operatorname{ann}_C(f_2)=(f_1)$, so
\[
 \cdots\xrightarrow{f_1}C\xrightarrow{f_2}C
       \xrightarrow{f_1}C\longrightarrow L_1\longrightarrow0
\]
is an exact free resolution. Frobenius pullback replaces its maps by
$\Delta^{p-1}f_1$ and $\Delta^{p-1}f_2$.
The scalar $\Delta$ is a nonzerodivisor on the $R$-free ring $C$,
so their kernels are still $(f_2)$ and $(f_1)$, while the preceding
images are only $\Delta^{p-1}(f_2)$ and $\Delta^{p-1}(f_1)$.
Writing $C_{F_C}$ for the right module twisted by Frobenius, we obtain
\[
 \operatorname{Tor}^{C}_{2j+1}(C_{F_C},L_1)
     \simeq L_1/\Delta^{p-1}L_1\quad(j\ge0),
\]
\[
 \operatorname{Tor}^{C}_{2j}(C_{F_C},L_1)
     \simeq L_2/\Delta^{p-1}L_2\quad(j\ge1).
\]
Exchange the branches for $L_2$. Every positive degree has
underlying $R$-module $R/(\Delta^{p-1})$.
Thus $\mathbf L F_C^*L_i$ is unbounded below and cannot equal the
degree-zero branch object. Changing its resolution cannot repair
this; replacing the object or using an unbounded category is a
different construction.

\par\smallskip\noindent\textbf{The actual algebra action and multiplicities.}\quad
The completed good-characteristic model from \fid{F19} has
\[
 A=R\times R\times M_2(C).
\]
Define $\phi_A$ by $F_R$ on the scalar blocks and entrywise $F_C$
on the matrix block. It is a ring homomorphism because, in the
commutative coefficient ring $C$,
\[
 \left(\sum_j a_{ij}b_{j\ell}\right)^p
       =\sum_j a_{ij}^p b_{j\ell}^p.
\]
This is not the noncommutative matrix-power map of \fid{F11}.
It fixes the matrix units, so its extension-of-scalars functor
commutes with the standard Morita equivalence.

The actual right module is
\[
 M=R^{20}\oplus R^4\oplus
    (L_1\otimes_C C^{1\times2})^{16}
       \oplus(L_2\otimes_C C^{1\times2})^4,
\]
through the regular action-preserving basis of \fid{F21}. Under
$F_A^*M=M\otimes_{A,\phi_A}A$, each branch multiplicity module
becomes $F_C^*L_i$; the two scalar sectors stay free. Therefore
\begin{equation}\label{eq:frobtorsion}
 F_A^*M\simeq R^{64}\oplus
          (R/(\Delta^{p-1}))^{40}
 \quad\text{as }R\text{-modules}.
\end{equation}
The branch multiplicity sheaf has 20 copies and torsion length
$20m(p-1)$; the actual row module has 40 copies and length
$40m(p-1)$. The natural algebra-linearized surjection discards
exactly this torsion. Derived Morita equivalence transfers the
resolution above as well: every positive Tor degree has 40 copies
of that torsion module, and the scalar sectors contribute none.

For the original center take $t=v-1$, $x=z_1-z$ and
$\Delta=z_1-z_2$. Its formula is
\[
 \Delta=\frac{(D-1)(D-4)}{D^2}
        =\frac34t^2+O(t^3).
\]
In characteristic greater than three its order is exactly two,
giving the claimed length $80(p-1)$.

\par\smallskip\noindent\textbf{Parameter pullback and the regular unit repair.}\quad
Parameter-base pullback of the center gives
\[
 C^{(p)}=R\otimes_{F_R,R}C
       \simeq R[X]/(X(X-\Delta^p)).
\]
Relative Frobenius sends $X$ to $x^p=\Delta^{p-1}x$.
In the two free bases its matrix is
$\operatorname{diag}(1,\Delta^{p-1})$; it is injective with
cokernel of length $m(p-1)$. Its source has branch separation of
order $mp$, whereas the target has order $m$. Extension over this
relative map is the additional operation in center pullback which
causes the torsion.

In the adapted free $R$-basis of $M$, take coefficientwise
Frobenius on every scalar and branch row. Branch evaluation
$C\to R$ commutes with Frobenius, so this semilinear map satisfies
\[
 \tau_{\rm ad}(ma)=\tau_{\rm ad}(m)\phi_A(a).
\]
Its parameter linearization
$R\otimes_{F_R,R}M\to M$ is an isomorphism, sending each
$1\otimes e_j$ to $e_j$. This tensor product is different from
the algebra tensor product in \eqref{eq:frobtorsion}.

The transport to the original tensor basis is explicit. With the
right-row convention of \fid{F21}, write
\[
 \rho_{\rm ad}(a)=B\rho_{\rm old}(a)B^{-1},\qquad
 r_{\rm old}=r_{\rm ad}B,\qquad
 U=(B^{[p]})^{-1}B,
\]
where $B^{[p]}$ means entrywise Frobenius and $B$ is the regular
invertible adapted basis matrix. Then
\begin{equation}\label{eq:frobtransport}
 \tau_{\rm old}(r)=r^{[p]}U,\qquad
 \rho_{\rm old}(\phi_A(a))
       =U^{-1}\rho_{\rm old}(a)^{[p]}U.
\end{equation}
Indeed transform $r$ to $rB^{-1}$, apply entrywise Frobenius,
and multiply back by $B$; applying the same operation to the
action matrix proves the second equality. It also directly checks
$\tau_{\rm old}(r\rho(a))=\tau_{\rm old}(r)\rho(\phi_A(a))$.
The matrix $U$ is invertible, proving the unit parameter structure
on the actual lattice.

Naive entrywise Frobenius in the original tensor basis is not
assumed to preserve the original algebra image. Although its
generator entries have the numerical form $P_i(v)^{[p]}=P_i(v^p)$,
the intrinsic algebra map acts by the conjugated expression in
\eqref{eq:frobtransport}. The regular conjugation cannot be omitted
without an additional argument.

\par\smallskip\noindent\textbf{A unit model on the normalization.}\quad
The normalization is
\[
 C\hookrightarrow\bar C=R\times R,
 \qquad a+bx\longmapsto(a,a+b\Delta).
\]
Its image is the pairs congruent modulo $\Delta$, and its quotient
has length $m$. It is the full normalization: the branch idempotents
are integral, and each factor $R$ is integrally closed.
Extending the algebra to this normalized center gives
\[
 A\hookrightarrow\bar A
       =R\times R\times M_2(R)\times M_2(R),
 \qquad \operatorname{length}_R(\bar A/A)=4m.
\]
Let $f:\bar S\to S$ normalize the three-component scheme in
\fid{F21}; it has four components upstairs. Choose free bundles
\[
 \bar E=(R,R,R^2,R^2)=f^*E,\qquad
 \bar{\mathcal F}=(R^{20},R^4,R^{16},R^4).
\]
They satisfy exactly
\[
 f_*\bar{\mathcal F}=\mathcal F,\qquad
 f_*\mathcal Hom(f^*E,\bar{\mathcal F})
       \simeq\mathcal Hom(E,\mathcal F).
\]
The last is ordinary tensor--Hom adjunction for the finite locally
free $E$. Its right $A$ action is restriction of the $\bar A$
action along the displayed inclusion. It therefore realizes the
same actual module, not just the same dimension or character.

On the normalized free branches, the chosen bases give unit
coherent Frobenius structures on $\bar E$ and
$\bar{\mathcal F}$ and hence the corresponding row module.
Entrywise Frobenius on $\bar A$ commutes with its inclusion of
$A$, since branch evaluations commute with absolute Frobenius.
But the natural comparison
\[
 F_S^*f_*\bar{\mathcal F}
       \longrightarrow f_*F_{\bar S}^*\bar{\mathcal F}
       \simeq f_*\bar{\mathcal F}
\]
is, on each branch, precisely $C/(x_i^p)\twoheadrightarrow C/(x_i)$.
It has torsion length $20m(p-1)$ on the multiplicity sheaf and
$40m(p-1)$ on the actual row module. Thus the unit upstairs
becomes the computed nonunit map downstairs. After inverting $t$
the normalization is an isomorphism and the entire defect vanishes.
\end{proof}


Thus the failed claim is a unit center/algebra pullback on the
specified branch object; the successful claims are the unit
parameter-base structure and the action-compatible normalized model.
Replacing the object, permitting unbounded constructions, or passing
to an \'etale/Weil realization are separate operations. None here
identifies quantum degree with a Weil weight, chooses VIII's basis,
or proves the opposite path-output sign needed by the main criterion.

\subsection{Testing the smooth proper noncommutative Weil framework}
There is a precise finite-field noncommutative Weil conjecture in a
different framework: for a smooth proper dg category, cyclotomic
Frobenius acts on the two finite-dimensional groups $TP_0[1/p]$
and $TP_1[1/p]$, with expected conjugate moduli $1$ and $\sqrt q$.
Tabuada \cite[\S1, pp.~2--3; \S5.1, pp.~13--14; \S6.1,
pp.~15--16]{TabuadaWeil} specifies these objects and proves the
comparison with smooth proper schemes in Theorem~1.5.
This cyclotomic Frobenius is not defined by raising matrices to
their $p$th powers. Its periodicity variable is also not the GCT
quantum parameter. The following test concerns the actual algebra
already constructed, without identifying the two conjectural programs.

\begin{finding}[\fid{F39}: the actual collision fiber fails homological smoothness]\label{f:ncweil}
In characteristic $p>3$, the abstract specialization of the
rank-three algebra in \fid{F19},
\[
 A_1=k\times k\times M_2(B),\qquad B=k[\epsilon]/(\epsilon^2),
\]
is proper but not smooth as a dg algebra concentrated in degree zero.
The same holds for its category of perfect complexes. Therefore this
particular category is not an input to the cited smooth proper
noncommutative Weil conjecture.

In contrast, at a split finite-field parameter point where the two
center roots are distinct and the stated denominators are units,
the algebra is $k^2\times M_2(k)^2$. Its perfect category is Morita
equivalent to four rational points. Its normalized even Frobenius
spectrum consists of four ones and its odd spectrum is empty.
Replacing $A_1$ by its semisimple quotient instead gives three
points and discards the nilpotent collision data.
\end{finding}
\begin{proof}
Properness of the one-object dg algebra means that its underlying
complex over $k$ is perfect; this holds because $A_1$ is finite
dimensional. Homological smoothness would make its diagonal bimodule
perfect over $A_1\otimes_k A_1^{\rm op}$.
We test the dual-number block directly.
Put
\[
 S=B\otimes_k B=k[x,y]/(x^2,y^2),\qquad
 u=x-y,\quad w=x+y.
\]
There is a free $S$-resolution of the diagonal $B$,
\[
 \cdots\xrightarrow{u}S\xrightarrow{w}S
       \xrightarrow{u}S\longrightarrow B\longrightarrow0.
\]
Here the rightmost positive differential is multiplication by $u$.
To check exactness in full, use the basis $1,x,y,xy$ of $S$.
The kernel of the augmentation $x,y\mapsto\epsilon$ is
$\langle x-y,xy\rangle=\operatorname{im}(u)$.
Multiplication by $u$ sends
$a+bx+cy+dxy$ to $a(x-y)+(c-b)xy$, so its kernel is
$\langle x+y,xy\rangle=\operatorname{im}(w)$.
Multiplication by $w$ has kernel
$\langle x-y,xy\rangle=\operatorname{im}(u)$.
Also $uw=0$. These identities prove exactness at every stage.

After tensoring over $S$ with the diagonal $B$, the differentials
alternate between zero and multiplication by $2\epsilon$. Hence
\[
 HH_0(B)=B,\qquad
 HH_{2j+1}(B)=B/(\epsilon)\ (j\ge0),\qquad
 HH_{2j}(B)=k\epsilon\ (j\ge1).
\]
In particular $\operatorname{Tor}^{S}_i(B,B)$ is nonzero for
every $i>0$. A perfect $S$-module would have bounded derived tensor
product with $B$, so $B$ is not homologically smooth.
If $A_1$ were smooth, restricting its diagonal along the central
matrix-block idempotent would make the diagonal of $M_2(B)$
perfect. The matrix Morita equivalence, obtained by the full matrix
unit $E_{11}$ on each side, is exact and preserves finite projective
resolutions; it would then make $B$ perfect over $S$, a contradiction.
This also proves the assertion for the Morita-invariant perfect dg
category, using the smoothness definition in the cited source.

At a split good parameter the distinct-root Chinese remainder
decomposition of \fid{F19} gives $k^2\times M_2(k)^2$.
Applying ordinary matrix Morita equivalence to its two matrix factors
identifies its perfect category with that of $\coprod_1^4\Spec k$.
The latter is smooth proper. The crystalline comparison
\cite[\S7, equations~(7.1)--(7.4), p.~17]{TabuadaWeil}
has only four degree-zero classes, each with Frobenius one; this
proves the stated normalized spectrum. The radical of $A_1$ is
the square-zero matrix ideal $M_2(k\epsilon)$, so its quotient is
$k^2\times M_2(k)$ and the same calculation gives three points.
\end{proof}

The dual-number resolution is the familiar elementary computation;
its use here is a test of the actual nonstandard collision fiber,
not a new claim about dual numbers or priority over noncommutative
motivic theory. A smooth categorical replacement or a relative theory
is not excluded. Such a replacement would need a specified comparison
retaining the coherent specialization and the designated based
operations. The pure spectrum at the good fibers does not supply
that comparison, and semisimplification expressly loses data measured
by \fid{F25}. The smoothness failure also does not refute any
conjecture whose domain already requires smoothness.

\subsection{A smooth category retaining the actual collision corner}
The direct smoothness failure in \fid{F39} has an explicit repair
at the level of a larger category. The following is an application
of the classical Auslander construction to the actual algebra and
module already identified in \fid{F19}/\fid{F21}; no priority claim
for the Auslander algebra is intended.

\begin{finding}[\fid{F40}: a smooth proper corner and an exact module lift]\label{f:auslander}
Let $k=\F_q$ have characteristic greater than three, and put
\[
 B=k[\epsilon]/(\epsilon^2),\qquad
 \Gamma=\End_B(B\oplus k),\qquad
 \Lambda=k\times k\times M_2(\Gamma).
\]
Then $\dim\Gamma=5$, $\dim\Lambda=22$, and both algebras are
smooth proper of global dimension two. For an explicit idempotent
$f\in\Lambda$,
\[
 f\Lambda f=A_1=k\times k\times M_2(B).
\]
There are projective right $\Lambda$-modules
$\mathcal P=f\Lambda$ and $\widetilde M$, with
$\dim_k\widetilde M=104$, such that
\[
 \End_\Lambda(\mathcal P)=A_1,\qquad
 \operatorname{Hom}_\Lambda(\mathcal P,\widetilde M)
   \simeq M_1
\]
as the actual right $A_1$-module of dimension 64 in \fid{F21}.
Thus the full algebra, including its nilpotents, and the specified
specialized tensor action are recovered.

The induced functor
$\operatorname{Perf}(A_1)\to\operatorname{Perf}(\Lambda)$ is fully
faithful but not right-admissible. The noncommutative Chow motive
of $\Lambda$ is a sum of four unit motives. Its normalized
cyclotomic Frobenius therefore has four eigenvalues one on
$TP_0[1/p]$, and $TP_1[1/p]=0$.
No relative deformation or based-operation comparison is asserted.
\end{finding}

\begin{proof}
Let $e_1,e_2$ be the idempotents of the summands $B,k$.
Let $a:k\to B$ send $1$ to $\epsilon$, let $b:B\to k$ be
the quotient, and let $c=ab$. Multiplication is composition,
with the right factor applied first. The five elements
$e_1,e_2,a,b,c$ form a basis and satisfy
\[
 a=e_1ae_2,\qquad b=e_2be_1,\qquad ba=0,\qquad
 c=ab\ne0,\qquad c^2=0.
\]
Indeed the four matrix Hom spaces have dimensions $2,1,1,1$,
and these maps span them. The radical is
$J=\langle a,b,c\rangle$, with $J^2=kc$, $J^3=0$ and
$\Gamma/J=k^2$.

Write
\[
 \mathsf P_1=e_1\Gamma=\langle e_1,a,c\rangle,\qquad
 \mathsf P_2=e_2\Gamma=\langle e_2,b\rangle
\]
and let $S_i$ be their simple tops.
Left multiplication by $a$ sends $e_2,b$ to $a,c$, and left
multiplication by $b$ sends $e_1$ to $b$ while killing $a,c$.
Thus the following are exact projective resolutions:
\[
 0\longrightarrow\mathsf P_2\xrightarrow{a}\mathsf P_1
   \longrightarrow S_1\longrightarrow0,
\]
\begin{equation}\label{eq:auslandersimple}
 0\longrightarrow\mathsf P_2\xrightarrow{a}\mathsf P_1
   \xrightarrow{b}\mathsf P_2\longrightarrow S_2
   \longrightarrow0.
\end{equation}
Every right module has a radical filtration of length at most
three, with semisimple quotients. Direct sums of these resolutions
and the extension bound for projective dimension show that every
module has projective dimension at most two. Applying
$\operatorname{Hom}_\Gamma(-,S_2)$ to
\eqref{eq:auslandersimple} gives
$\operatorname{Ext}^2_\Gamma(S_2,S_2)=k$, since
$\operatorname{Hom}(\mathsf P_1,S_2)=0$ and
$\operatorname{Hom}(\mathsf P_2,S_2)=k$. Hence the global
dimension is exactly two.

We prove smoothness directly. All tensor products in the following
diagonal resolution are over $k$. Set
\[
 Q_0=(\Gamma e_1\otimes e_1\Gamma)
          \oplus(\Gamma e_2\otimes e_2\Gamma),
\]
\[
 Q_1=(\Gamma e_1\otimes e_2\Gamma)_a
          \oplus(\Gamma e_2\otimes e_1\Gamma)_b,\qquad
 Q_2=\Gamma e_2\otimes e_2\Gamma.
\]
These are projective $\Gamma$-bimodules. Let $\mu:Q_0\to\Gamma$
be multiplication and define
\[
 d_1(x\otimes y)_a=xa\otimes y-x\otimes ay,\qquad
 d_1(x\otimes y)_b=xb\otimes y-x\otimes by,
\]
\[
 d_2(x\otimes y)=(xb\otimes y)_a+(x\otimes ay)_b.
\]
The subscripts specify the summands; the terms of $d_1$ lie in
the corresponding source and target vertex summands of $Q_0$.
These are bimodule maps, with $\mu d_1=0$ and
\[
 d_1d_2(x\otimes y)=xba\otimes y-x\otimes bay=0.
\]
The left ideals have bases
$\Gamma e_1=\langle e_1,b,c\rangle$ and
$\Gamma e_2=\langle e_2,a\rangle$. Consequently
\[
 \dim Q_2=4,\quad\dim Q_1=12,\quad\dim Q_0=13,\quad
 \dim\Gamma=5.
\]
The map $\mu$ is surjective. The image of $d_1$ contains the
eight independent vectors
\[
 \begin{gathered}
 a\otimes e_2-e_1\otimes a,\quad
 a\otimes b-e_1\otimes c,\quad
 b\otimes e_1-e_2\otimes b,\quad
 c\otimes e_1-a\otimes b,\\
 b\otimes a,\quad b\otimes c,\quad
 c\otimes a,\quad c\otimes c.
 \end{gathered}
\]
Independence follows from the tensor bases: the first four have
unique coordinates at $a\otimes e_2,e_1\otimes c,e_2\otimes b,
c\otimes e_1$, respectively, and the last four are distinct
remaining coordinates. Thus $\operatorname{im}d_1=\ker\mu$.
The $a$-component of $d_2$ is the tensor product of the injection
$x\mapsto xb:\Gamma e_2\to\Gamma e_1$, which sends $e_2,a$
to $b,c$, with the identity of $e_2\Gamma$. Hence $d_2$ is
injective. Its image has dimension four, equal to
$\dim\ker d_1$, proving exactness of
\begin{equation}\label{eq:auslanderdiagonal}
 0\longrightarrow Q_2\xrightarrow{d_2}Q_1
   \xrightarrow{d_1}Q_0\xrightarrow{\mu}\Gamma
   \longrightarrow0.
\end{equation}
This is a finite projective diagonal resolution. It proves
homological smoothness; finite dimensionality proves properness.
Matrix Morita equivalence and finite products give the same
properties and global dimension for $\Lambda$.

Now put $f=(1,1,I_2\otimes e_1)$. The corner equality follows
from $e_1\Gamma e_1=\langle e_1,c\rangle\simeq B$, with
$\epsilon\mapsto c\ne0$. In particular all named generators in
\fid{F19} are retained: its normalized matrix formulas use
$z=3/4+c$ at the collision. Left multiplication identifies
$\End_\Lambda(f\Lambda)=f\Lambda f=A_1$.

For the module recovery, write the actual specialization from
\fid{F21} as
\[
 M_1=\epsilon_+^{20}\oplus\epsilon_-^4\oplus S^{20},
\]
where $S=k^2$ is the right row module for $M_2(B)$, with
$\epsilon$ acting as zero. The projective right $M_2(\Gamma)$-module
\[
 Q=(E_{11}\otimes e_2)M_2(\Gamma)
     \simeq(\mathsf P_2,\mathsf P_2)
\]
has dimension four. Its corner is
\[
 Q(I_2\otimes e_1)=(\mathsf P_2e_1)^2
       =(kb)^2\simeq S,
\]
as an actual module, since $bc=bab=0$. Therefore
\[
 \widetilde M=\epsilon_+^{20}\oplus\epsilon_-^4\oplus Q^{20}
\]
is projective of dimension $104$, and
$\widetilde M f=M_1$. Evaluation at $f$ identifies
$\operatorname{Hom}_\Lambda(f\Lambda,\widetilde M)$ with this
corner. Precomposition by $\End_\Lambda(f\Lambda)$ is precisely
the right $A_1$-action, proving the claimed action-preserving
comparison.

For the embedding and its limitation, first consider
\[
 I:\operatorname{Perf}(B)\longrightarrow\operatorname{Perf}(\Gamma),
 \qquad X\longmapsto X\otimes_B^{\mathbf L}e_1\Gamma.
\]
It sends the generator $B$ to the projective $\mathsf P_1$.
Its right adjoint on ambient derived categories is corner
restriction $Y\mapsto Ye_1$. The unit is an isomorphism because
\[
 (X\otimes_B^{\mathbf L}e_1\Gamma)e_1
       =X\otimes_B^{\mathbf L}B=X.
\]
Thus $I$ is fully faithful. The matrix and scalar factors give
the assertion for $A_1,\Lambda$.

However $\mathsf P_2e_1=kb\simeq k$ is not perfect over $B$:
the exact resolution
\[
 \cdots\xrightarrow{\epsilon}B\xrightarrow{\epsilon}B
   \xrightarrow{\epsilon}B\longrightarrow k\longrightarrow0
\]
gives $\operatorname{Tor}^B_i(k,k)=k$ in every degree.
If $I$ had a right adjoint valued in $\operatorname{Perf}(B)$,
its value at $\mathsf P_2$ would represent
$X\mapsto\operatorname{Hom}_{D(B)}(X,k)$ on perfect $X$.
Evaluation at $B[n]$ would force the representing object to
have only the cohomology module $k$ in degree zero; naturality
for $\End_B(B)=B$ fixes its $B$-action. It would be isomorphic
to the nonperfect object $k$, a contradiction. This also
obstructs right admissibility of the actual matrix corner.

Finally $\langle S_1,\mathsf P_2\rangle$ is a full exceptional
pair in $\operatorname{Perf}(\Gamma)$.
The projective $\mathsf P_2$ has endomorphism ring $k$.
The first simple resolution gives
$\End(S_1)=k$ and no positive self-Ext, because
$\operatorname{Hom}(\mathsf P_2,S_1)=S_1e_2=0$.
Also $R\operatorname{Hom}(\mathsf P_2,S_1)=0$, so the pair
is semiorthogonal. Its triangle recovers $\mathsf P_1$, hence
both projective generators. Each exceptional subcategory has
both adjoints given by finite derived Hom and its dual, so this
is a semiorthogonal decomposition into two copies of
$\operatorname{Perf}(k)$.

Known semiorthogonal additivity
\cite[Example~12.3, p.~24]{TabuadaWeil} now gives
\[
 U(\Gamma)_{\Q}=U(k)_{\Q}^{\oplus2},\qquad
 U(\Lambda)_{\Q}=U(k)_{\Q}^{\oplus4}.
\]
The Frobenius-enriched TP functor
\cite[\S6.2, equation~(6.9), p.~16]{TabuadaWeil}
preserves these isomorphisms with Frobenius. The point comparison
\cite[\S7, equations~(7.1)--(7.4), p.~17]{TabuadaWeil}
then gives $TP_0(\Lambda)[1/p]=K^4$ with $F_0=1$ and
$TP_1(\Lambda)[1/p]=0$, where $K=\operatorname{Frac}W(k)$.
This is also an instance of the cited connective smooth proper
theorem and its motive computation
\cite[Theorem~2.12, p.~8; proof, equation~(12.6), p.~26]{TabuadaWeil}.
The source hypotheses hold by \eqref{eq:auslanderdiagonal} and
finite dimensionality. The special-case exceptional calculation
above avoids any assumption of nilinvariance for arbitrary
nonsmooth algebras.
\end{proof}

The original $M_1$ is not perfect over $A_1$. Its recovery here
is a corner restriction of a projective object, not the image of
$M_1$ under the perfect-module embedding. Its higher coherent
self-Ext need not agree with that of $\widetilde M$, and the
relative collision order in \fid{F25} has not been recovered
from this single fiber. The pure spectrum of the entire smooth
category does not identify a Frobenius realization of every
specified corner morphism, a positive nonstandard basis, or its
structure constants. The direct smoothness gate is nevertheless
strictly narrowed: this category retains the actual collision
algebra and tensor action as stated, without semisimplifying them.


\subsection{A smooth relative corner retaining the tensor lattice}
The special-fiber construction in \fid{F40} extends through the
actual collision. The smaller generator below suffices; no
classification of Cohen--Macaulay modules over a tacnode is used.

\begin{finding}[\fid{F44}: relative smoothness, the actual action,
and two different Frobenius unit conditions]\label{f:relativeauslander}
Let $R$ be a commutative domain, let $0\ne\Delta\in R$, and set
\[
 C=R[x]/(x(x-\Delta)),\qquad L_1=C/(x),\qquad
 L_2=C/(x-\Delta),\qquad
 \Gamma=\End_C(C\oplus L_1).
\]
Then $\Gamma$ is free of rank five over $R$ and has a projective
relative diagonal resolution of length two. It is smooth and
proper over $R$; if $R$ is a regular DVR, its absolute global
dimension is at most three.

The smooth proper algebra
$\Lambda=R\times R\times M_2(\Gamma)$ has an idempotent $f$ such that
\[
 f\Lambda f=A=R\times R\times M_2(C).
\]
For the actual rank-three family, with $\Delta=z_1-z_2$ and
$x=z_1-z$, this is the specified algebra of \fid{F19}. There is an
$R$-free perfect right $\Lambda$-module $\widetilde M$ of rank
$96$ such that, for $\mathcal P=f\Lambda$,
\[
 \End_\Lambda(\mathcal P)=A,\qquad
 \operatorname{Hom}_\Lambda(\mathcal P,\widetilde M)
       \simeq M
\]
as the actual rank-$64$ right $A$-module of \fid{F21}.
Moreover
\[
 \operatorname{Ext}^1_\Lambda(\widetilde M,\widetilde M)
       \simeq(R/(\Delta))^{64}.
\]
This retains the collision order in a coherent extension inside
the smooth relative category.

In characteristic $p>3$ there is an explicit ring Frobenius lift
on $\Gamma$ inducing ordinary Frobenius on $C$, with unit
derived pullback structures on the displayed perfect objects.
At a nonunit $\Delta$, those structures are not unit after
forgetting to $R$. More generally, no lift in the entire
fixed-vertex family considered below makes $\widetilde M$
both a unit object for the larger-algebra pullback and a unit
object for the underlying $R$-linearization.
This statement does not negate the independent $R$-unit
structure on the original corner module constructed in
\fid{F36}.
\end{finding}

\begin{proof}
\par\smallskip\noindent\textbf{The actual endomorphism algebra.}\quad
Let $e_1,e_2$ be the idempotents for $C,L_1$. Define
$a:L_1\to C$ by $a(1)=x-\Delta$ and let $b:C\to L_1$ be
evaluation at zero. Multiplication is composition, with the
right factor applied first. Put $c=ab$. Then
\begin{equation}\label{eq:relativearrows}
 a=e_1ae_2,\qquad b=e_2be_1,\qquad
 ba=-\Delta e_2,\qquad c=ab,\qquad c^2=-\Delta c.
\end{equation}
These formulas describe the endomorphism algebra itself.
Indeed the four matrix Hom spaces are $C,R,Rb,Ra$, because
\[
 \operatorname{ann}_C(x)=R(x-\Delta).
\]
To check the annihilator, write an element as $u+vx$.
Its product with $x$ is $(u+v\Delta)x$, which vanishes
exactly when $u=-v\Delta$. Thus $e_1,e_2,a,b,c$ is an
$R$-basis of $\Gamma$. The relation $ba=-\Delta e_2$
reduces every longer alternating path to this basis, and
\[
 e_1\Gamma e_1=R[c]/(c(c+\Delta))\simeq C,\qquad
 x=c+\Delta e_1.
\]
For the actual family choose $x=z_1-z$. Then $L_1$ is
the multiplicity-$16$ branch $z=z_1$ and $L_2$ is the
multiplicity-$4$ branch $z=z_2$. In all named matrix formulas
the original coordinate is therefore
\begin{equation}\label{eq:relativez}
 z=z_1-x=z_2-c.
\end{equation}
This specifies the original generators, as well as the abstract
algebra type. At $\Delta=0$ the algebra is the one in
\fid{F40}, with the stated sign change in the corner coordinate.
After inverting $\Delta$, the two branches split and
$\Gamma\simeq M_2(R)\times R$.

\par\smallskip\noindent\textbf{The complete relative diagonal resolution.}\quad
All tensor products in this paragraph are over $R$. Set
\[
 Q_0=(\Gamma e_1\otimes e_1\Gamma)
        \oplus(\Gamma e_2\otimes e_2\Gamma),
\]
\[
 Q_1=(\Gamma e_1\otimes e_2\Gamma)_a
        \oplus(\Gamma e_2\otimes e_1\Gamma)_b,\qquad
 Q_2=\Gamma e_2\otimes e_2\Gamma.
\]
These are finite projective modules over
$\Gamma\otimes_R\Gamma^{\mathrm{op}}$ and have $R$-ranks
$13,12,4$, respectively. Let $\mu:Q_0\to\Gamma$ be
multiplication and define
\[
 d_1(u\otimes v)_a=ua\otimes v-u\otimes av,\qquad
 d_1(u\otimes v)_b=ub\otimes v-u\otimes bv,
\]
\[
 d_2(u\otimes v)=(ub\otimes v)_a+(u\otimes av)_b.
\]
The subscripts identify the two arrow summands. These are
bimodule maps and $\mu d_1=0$. The scalar in the deformed
relation cancels:
\[
 d_1d_2(u\otimes v)=uba\otimes v-u\otimes bav
     =-\Delta u\otimes v+u\otimes\Delta v=0.
\]

We prove exactness over $R$ directly. The left ideal bases are
$\Gamma e_1=\langle e_1,b,c\rangle$ and
$\Gamma e_2=\langle e_2,a\rangle$; the right ideal bases are
$e_1\Gamma=\langle e_1,a,c\rangle$ and
$e_2\Gamma=\langle e_2,b\rangle$.
The image of $d_1$ contains
\[
 \begin{gathered}
 a\otimes e_2-e_1\otimes a,\qquad
 a\otimes b-e_1\otimes c,\qquad
 b\otimes e_1-e_2\otimes b,\qquad
 c\otimes e_1-a\otimes b,\\
 b\otimes a+\Delta e_2\otimes e_2,\qquad
 b\otimes c+\Delta e_2\otimes b,\qquad
 c\otimes a+\Delta a\otimes e_2,\qquad
 c\otimes c+\Delta a\otimes b.
 \end{gathered}
\]
Use the pivot coordinates
\[
 a\otimes e_2,\ e_1\otimes c,\ e_2\otimes b,\ c\otimes e_1,\
 b\otimes a,\ b\otimes c,\ c\otimes a,\ c\otimes c.
\]
Their $8$-by-$8$ minor is triangular in two blocks, with
diagonal entries $1,-1,-1,1,1,1,1,1$. Its determinant is
a unit, irrespective of $\Delta$. Hence their span $V$
is a free direct summand of $Q_0$, and $Q_0/V$ is free of
rank five. Multiplication annihilates $V$ and induces a
surjection $Q_0/V\to\Gamma$ between free modules of rank
five. A surjective square matrix over a commutative ring
has unit determinant, so this map is an isomorphism.
It follows that $V=\ker\mu=\operatorname{im}d_1$.

The $a$-component of $d_2$ uses $u\mapsto ub$, which sends
$e_2,a$ to the two basis vectors $b,c$. Coordinate projection
gives a left inverse. Thus $d_2$ is a split $R$-module
injection with free quotient of rank eight. The induced
surjection $Q_1/\operatorname{im}d_2\to V$ is an
isomorphism by the same square-matrix argument. We have
therefore proved the entire exact sequence
\begin{equation}\label{eq:relativediagonal}
 0\longrightarrow Q_2\xrightarrow{d_2}Q_1
   \xrightarrow{d_1}Q_0\xrightarrow{\mu}\Gamma
   \longrightarrow0.
\end{equation}
This proves relative smoothness. The finite free
$R$-module $\Gamma$ is proper over $R$.
The resolution proof works over any commutative base;
the domain and nonzerodivisor hypotheses enter the later
Ext and Frobenius arguments.

If $R$ is a regular DVR, tensor
\eqref{eq:relativediagonal} on the left with any right
$\Gamma$-module $N$. As a left-module resolution of the
projective module $\Gamma$, the bounded resolution is split
exact, so tensoring is exact. Its terms are finite sums of
$Ne_i\otimes_R e_j\Gamma$. Every $R$-module has projective
dimension at most one; since $e_j\Gamma$ is $R$-free and
projective over $\Gamma$, each such term has projective
dimension at most one over $\Gamma$. A length-two resolution
by these terms gives $\operatorname{pd}_\Gamma N\le3$.
No sharpness assertion for this bound is required.

\par\smallskip\noindent\textbf{The exact lift of the tensor lattice.}\quad
Put
\[
 \Lambda=R\times R\times M_2(\Gamma),\qquad
 f=(1,1,I_2\otimes e_1),\qquad \mathcal P=f\Lambda.
\]
Finite products and matrix Morita equivalence preserve
smoothness and properness, so $\Lambda$ is smooth proper
and $R$-free of rank $22$. Left multiplication gives
\[
 \End_\Lambda(\mathcal P)=f\Lambda f
       =R\times R\times M_2(C)=A.
\]
This identification uses \eqref{eq:relativez}.

Write $\mathsf P_i=e_i\Gamma$. Left multiplication by
$a$ sends the basis $e_2,b$ to the independent vectors
$a,c$. Hence
\begin{equation}\label{eq:relativeS}
 0\longrightarrow\mathsf P_2\xrightarrow{a}\mathsf P_1
      \longrightarrow S\longrightarrow0
\end{equation}
is exact. The quotient $S$ is $R$-free of rank one;
$e_1$ acts by one and $e_2,a,b,c$ by zero.
In particular $S$ is perfect of projective dimension at
most one. Its corner and that of $\mathsf P_2$ are
\[
 \mathsf P_2e_1=Rb\simeq L_1,\qquad Se_1=S\simeq L_2.
\]
Indeed $bx=bc+\Delta b=-\Delta b+\Delta b=0$,
whereas $x$ acts on $S$ by $\Delta$.

For a right $\Gamma$-module $N$, let
$\operatorname{Row}(N)=(N,N)$ be the standard right
$M_2(\Gamma)$-Morita module. Define
\begin{equation}\label{eq:relativelift}
 \widetilde M=\epsilon_+^{20}\oplus\epsilon_-^4
       \oplus\operatorname{Row}(\mathsf P_2)^{16}
       \oplus\operatorname{Row}(S)^4.
\end{equation}
It is $R$-free of rank $24+16\cdot4+4\cdot2=96$
and perfect over $\Lambda$, of projective dimension
at most one. Taking the corner gives exactly
\[
 \widetilde M f=
 \epsilon_+^{20}\oplus\epsilon_-^4\oplus N_1^{16}
       \oplus N_2^4=M.
\]
Evaluation at $f$ identifies
$\operatorname{Hom}_\Lambda(f\Lambda,\widetilde M)$
with $\widetilde M f$. Precomposition by left multiplication
in $f\Lambda f$ is the original right $A$-action.
Thus both the named algebra and the tensor lattice are
recovered over the family, including the collision.
At the closed fiber this rank-$96$ perfect lift differs
from the rank-$104$ projective lift chosen in \fid{F40};
both have the same actual rank-$64$ corner.

\par\smallskip\noindent\textbf{A coherent extension preserving the order of collision.}\quad
Applying $\operatorname{Hom}_\Gamma(-,\mathsf P_2)$ to
\eqref{eq:relativeS} gives
\[
 Rb\xrightarrow{-\Delta}Re_2,
\]
because composition is $ba=-\Delta e_2$.
As $\Delta$ is a nonzerodivisor,
\[
 \operatorname{Hom}_\Gamma(S,\mathsf P_2)=0,\qquad
 \operatorname{Ext}^1_\Gamma(S,\mathsf P_2)=R/(\Delta).
\]
The opposite derived Hom vanishes:
$\mathsf P_2$ is projective and
$\operatorname{Hom}_\Gamma(\mathsf P_2,S)=Se_2=0$.
Both self-Hom rings are $R$, and both positive self-Ext
groups vanish, using \eqref{eq:relativeS} for $S$.
Consequently exactly the $4\cdot16$ oriented cross
pairs in \eqref{eq:relativelift} contribute:
\[
 \operatorname{Ext}^1_\Lambda(\widetilde M,\widetilde M)
       =(R/(\Delta))^{64}.
\]
There is no extra factor from matrix Morita equivalence.
If $\Delta=ut^m$ over a DVR, with $u$ a unit, this
group has length $64m$, hence length $128$ for the actual
order-two collision. The original corner Ext in \fid{F25}
has both orientations and is $(R/(\Delta))^{128}$,
of length $256$ when $m=2$. These two Ext groups are
not identified.

\par\smallskip\noindent\textbf{The relative exceptional pair and the corner limitation.}\quad
The two self-Hom complexes of $S,\mathsf P_2$ are $R$,
the reverse derived Hom vanishes, and the triangle
\eqref{eq:relativeS} recovers $\mathsf P_1$. Thus they
generate. The functors from $\operatorname{Perf}(R)$
defined by these two objects have adjoints using derived
Hom and its $R$-dual. These complexes are perfect over
$R$: all objects involved are perfect over the
$R$-proper algebra $\Gamma$. Therefore
\[
 \operatorname{Perf}(\Gamma)
       =\langle\operatorname{Perf}(R),\operatorname{Perf}(R)\rangle.
\]
On each finite-field fiber the same known additivity and
Frobenius-enriched realization as in \fid{F40} apply:
\[
 U(\Gamma)_{\Q}=U(k)_{\Q}^{\oplus2},\qquad
 U(\Lambda)_{\Q}=U(k)_{\Q}^{\oplus4}.
\]
More explicitly, the invoked primary statements are
\cite[Example~12.3, p.~24; \S6.2, equation~(6.9), p.~16;
\S7, equations~(7.1)--(7.4), p.~17]{TabuadaWeil}.
The normalized TP spectrum on such a $\Lambda$-fiber
is four eigenvalues one in even degree and no odd part.
No fraction field of an arbitrary trait is being assumed
finite.

The corner functor
\[
 \operatorname{Perf}(C)\longrightarrow\operatorname{Perf}(\Gamma),
 \qquad X\longmapsto X\otimes_C^{\mathbf L}e_1\Gamma
\]
is fully faithful: its ambient derived right adjoint is
corner restriction, and the unit is the identity
$X\otimes_C^{\mathbf L}C=X$.
If $\Delta$ is a nonunit, this functor is not
right-admissible. Indeed the alternating resolution
with maps $x$ and $x-\Delta$ of $L_1$ gives
\[
 \operatorname{Tor}_{2j+1}^C(L_1,L_1)=R/(\Delta)
       \quad(j\ge0).
\]
The annihilator identities needed for this resolution
follow directly from the basis $1,x$. Thus $L_1$ is
not perfect. Since $\mathsf P_2e_1=L_1$, a right
adjoint on perfect categories would have to represent
$\operatorname{Hom}_{D(C)}(-,L_1)$ on perfect objects.
Testing $C[n]$ and its endomorphisms forces the
representing object to be $L_1$, a contradiction.
After inverting $\Delta$, the relation
$ba=-\Delta e_2$ makes $e_1$ full, and the corner
is instead a Morita equivalence.

\par\smallskip\noindent\textbf{A genuine noncommutative ring Frobenius lift.}\quad
Assume now $\operatorname{char}R=p>3$. Define the
$F_R$-semilinear ring map $\varphi:\Gamma\to\Gamma$ by
\begin{equation}\label{eq:relativefrobenius}
 \varphi(e_i)=e_i,\qquad
 \varphi(a)=a,\qquad
 \varphi(b)=\Delta^{p-1}b,\qquad F_R(r)=r^p.
\end{equation}
It respects the defining relation because
\[
 \varphi(b)\varphi(a)
       =-\Delta^p e_2=\varphi(-\Delta e_2).
\]
It is a ring map, not the operation taking every
noncommutative element to its $p$th power: for example
it fixes the nonzero square-zero arrow $a$.
On the commutative corner,
\[
 \varphi(c)=\Delta^{p-1}c=c^p,\qquad
 \varphi(x)=\Delta^{p-1}(c+\Delta e_1)=x^p.
\]
Here $c^2=-\Delta c$, $x^2=\Delta x$, and $p$ is odd.
The induced map is ordinary scheme Frobenius, including
on the actual coordinate $z=z_2-c$. Extend $\varphi$
to $\Lambda$ entrywise and on its scalar factors; it
fixes $f$.

For right modules write
$\mathbf L\varphi^*N=N\otimes_{\Gamma,\varphi}^{\mathbf L}\Gamma$.
The fixed vertex idempotents give
$\mathbf L\varphi^*\mathsf P_i\simeq\mathsf P_i$.
The differential $a$ in \eqref{eq:relativeS} remains
$a$ under $\varphi$, so
$\mathbf L\varphi^*S\simeq S$, with no additional
derived cohomology. Thus the displayed perfect objects
$\mathcal P,\widetilde M$ have unit structures for
pullback over their larger algebras.

\par\smallskip\noindent\textbf{The independent unit condition over the base.}\quad
The induced semilinear map on the $R$-basis $e_2,b$
of $\mathsf P_2$ sends these vectors to
$e_2,\Delta^{p-1}b$. Its $R$-linearization is
\begin{equation}\label{eq:relativebaseunit}
 \operatorname{diag}(1,\Delta^{p-1}),
\end{equation}
which is not invertible at a nonunit $\Delta$.
The induced $R$-linearization on $S$ is the identity.
There are $32$ copies of $b$ in the sixteen row copies
of $\mathsf P_2$ in \eqref{eq:relativelift}. Therefore
the determinant on the rank-$96$ lift is
$\Delta^{32(p-1)}$, up to a unit.
The recovered rank-$64$ corner has the same defect.
Changing the unit algebra-pullback isomorphism composes
it with an automorphism of the target, whose underlying
$R$-determinant is a unit; it cannot repair this defect.

We can classify this tradeoff for every
$F_R$-semilinear ring endomorphism fixing $e_1,e_2$.
The two arrow spaces have rank one, so its values must be
\[
 a\longmapsto ua,\qquad b\longmapsto vb
       \quad(u,v\in R).
\]
The relation and the nonzerodivisor $\Delta$ force
\begin{equation}\label{eq:relativetradeoff}
 uv=\Delta^{p-1}.
\end{equation}
Conversely this equality defines such a ring map and
gives ordinary Frobenius on $C$.
The pullback of \eqref{eq:relativeS} has differential
$ua$. Since $u\ne0$, it is injective, and its cokernel
as an $R$-module is
\[
 R\oplus(R/uR)^2.
\]
It is isomorphic to the free rank-one module $S$ only
if $u$ is a unit; a unit $u$ also suffices by rescaling
the resolution. The module $\mathsf P_2$ always has a
unit algebra-pullback structure, but its $R$-linearization
is $\operatorname{diag}(1,v)$, up to a unit change,
and is invertible exactly when $v$ is a unit.
For a nonunit $\Delta$, \eqref{eq:relativetradeoff}
prevents both conditions simultaneously.

For the full $\widetilde M$, a nonunit $u$ contributes
actual $R$-torsion from its $S$-summands, so it cannot
even be isomorphic to its algebra pullback. If $u$ is
a unit, $v$ is a nonunit and its underlying
$R$-linearization has a nonunit determinant. This
proves the assertion for the entire fixed-vertex
family, with the indicated scope.
\end{proof}

The conclusion concerns this chosen smooth lift and its
larger-algebra unit structures. The original corner module
still has the independent $R$-unit semilinear structure for
its intrinsic center Frobenius constructed in \fid{F36}.
The calculation above prevents that base-unit property
from being obtained through the unit pullback structures
of this particular lift; it does not prohibit the
\fid{F36} construction or every possible smooth model.
The ring map $\varphi$, scheme Frobenius on the commutative
corner, algebra-pullback unit structures, and cyclotomic
Frobenius on TP are distinct constructions. No equality
with cyclotomic Frobenius, Weil realization of all
specified morphisms, preservation of both original Ext
orientations, or positive nonstandard basis follows.


\subsection{Projective motives and Lefschetz symmetry for the actual word traces}
The whole-word trace from \fid{F43} has a stronger realization than
coefficientwise positivity. An elementary symmetric-square identity
and a rational algebraic projector supply a smooth-projective
coefficient motive with a specified Lefschetz operator. The motive
depends on the word; no common operator realization is assumed.

\begin{finding}[\fid{F45}: a projective word-trace realization]\label{f:projectiveword}
Let $w$ be a nonempty word in $P_1,P_2$, of length $L$ and
cyclic transition count $2s$. Thus $0\le2s\le L$, and its actual
trace is
\[
 T_w(v)=20D^L+16D^{L-s}+4D^{L-2s}(D-2)^{2s},
 \qquad D=(v+v^{-1})^2.
\]
This is the character of an actual finite-dimensional
$\mathrm{SL}_2$-representation with even weights. Its coefficients
are symmetric and unimodal in $v^2$.

Over every $k=\F_q$, with $\ell\ne\operatorname{char}k$, there is
an explicit rational Chow-motive projector on products of
projective lines whose centered realization $\mathcal K_w$ has
\begin{equation}\label{eq:projectivewordtrace}
 \sum_i(-1)^i\Tr(F^n\mid H^i(\mathcal K_w))
             =T_w(q^{n/2})\qquad(n\ge1).
\end{equation}
Only even degrees occur, and if $T_w(v)=\sum_r c_rv^{2r}$, then
\[
 H^{2r}(\mathcal K_w)=\Q_\ell(-r)^{\oplus c_r}.
\]
Thus $\mathcal K_w$ is pure of weight zero as a complex in BBD's
convention, while its total ungraded Weil space is mixed.
An invariant ample class gives hard Lefschetz on the projector
image and on the centered direct sum.
\end{finding}
\begin{proof}
\par\smallskip\noindent\textbf{An actual character.}\quad
Let $V$ be the standard two-dimensional characteristic-zero
$\mathrm{SL}_2$-representation, with character $\delta=v+v^{-1}$,
and put $W=V^{\otimes2s}$. For eigenvalues $z_i$, the identity
\[
 2\sum_{i\le j}z_iz_j=(\sum_i z_i)^2+\sum_i z_i^2
\]
proves
\[
 \chi_{\Sym^2W}(v)
   =\tfrac12\bigl(D^{2s}+(v^2+v^{-2})^{2s}\bigr)
   =\tfrac12\bigl(D^{2s}+(D-2)^{2s}\bigr).
\]
All tensor exponents in the actual representation
\begin{align*}
 \mathcal R_w={}&(V^{\otimes2L})^{\oplus16}
 \oplus(V^{\otimes2(L-s)})^{\oplus16}\\
 &\oplus\left(V^{\otimes(2L-4s)}
                    \otimes\Sym^2(V^{\otimes2s})\right)^{\oplus8}
\end{align*}
are nonnegative. Its character is
\[
 16D^L+16D^{L-s}
 +8D^{L-2s}\tfrac12\bigl(D^{2s}+(D-2)^{2s}\bigr)=T_w.
\]
All constituents have even highest weights. Their characters are
the centered strings $v^{-2a}+v^{-2a+2}+\cdots+v^{2a}$.
A nonnegative sum of these strings is symmetric and nondecreasing
toward the center in $v^2$. This proves the first assertion.
For $s=0$, the zero-fold tensor power and its symmetric square
are the unit, and the formula gives $40D^L$.

\par\smallskip\noindent\textbf{The projective schemes and projector.}\quad
For $d\ge0$ put
\[
 X_d=(\PP^1_k)^{2d},\qquad
 K_d=R\Gamma(X_{d,\bar k},\Q_\ell)(d)[2d].
\]
For $d=0$ these are the point and unit complex. The scheme $X_d$
is smooth projective of dimension $2d$.
On $X_L$, exchange two ordered blocks of $2s$ factors and leave
the other $2L-4s$ factors fixed. This defines a $k$-rational
involution $\sigma$, with $2s$ transposed pairs; when $s=0$ it
is the identity. Its graph gives the actual rational correspondence
\[
 \pi_s=\tfrac12(\Delta_{X_L}+\Gamma_\sigma)
       \in CH^{2L}(X_L\times X_L)_\Q.
\]
Since $\Gamma_\sigma^2=\Delta_{X_L}$, this is an idempotent.
Its image exists in the idempotent-complete category of rational
Chow motives and realizes the invariant projector $(1+\sigma^*)/2$.
No smoothness of a quotient variety is used. Rational coefficients
also permit characteristic two: the scalar one-half is in the
characteristic-zero coefficient field.

Define
\begin{equation}\label{eq:wordcomplex}
 \mathcal K_w=K_L^{\oplus16}\oplus K_{L-s}^{\oplus16}
                         \oplus(\pi_sK_L)^{\oplus8}.
\end{equation}
Before the centering shifts, this comes from the rational motives
$h(X_L)(L)$, $h(X_{L-s})(L-s)$ and the image of $\pi_s$ on
$h(X_L)(L)$, with multiplicities $16,16,8$.
The first and third realizations are shifted by $2L$, the second
by $2(L-s)$, exactly as displayed. All Tate twists are integers;
the notation $q^{n/2}$ only evaluates an even Laurent polynomial.

\par\smallskip\noindent\textbf{All extension traces.}\quad
The projective-space and projective-bundle calculations
\cite[Example~16.3, p.~109; Theorem~23.2, pp.~139--140]{Milne}
give the product basis
\[
 H^*(X_{d,\bar k})=\Q_\ell[x_1,\ldots,x_{2d}]/(x_i^2),
 \qquad \deg x_i=2.
\]
The monomial $x_I=\prod_{i\in I}x_i$ spans a Tate line with
geometric Frobenius $F^n$ eigenvalue $Q^{|I|}$, where $Q=q^n$.
There are no odd groups. The even shift changes no Euler sign,
and consequently
\[
 \Tr(F^n\mid K_d)=Q^{-d}(1+Q)^{2d}=D_n^d,
 \qquad D_n=Q+2+Q^{-1}.
\]
Here trace of a complex means its finite alternating trace.

The involution permutes the subset basis without Koszul signs,
since all factor classes have even degree. To be fixed, a subset
may choose each fixed coordinate independently and must choose
both or neither coordinate in each transposed pair. Thus
\[
 \Tr(\sigma F^n\mid K_L)
   =Q^{-L}(1+Q)^{2L-4s}(1+Q^2)^{2s}
   =D_n^{L-2s}(D_n-2)^{2s}.
\]
The projector commutes with Frobenius, so its actual image has trace
\[
 \Tr(F^n\mid\pi_sK_L)
   =\tfrac12\left[D_n^L+D_n^{L-2s}(D_n-2)^{2s}\right].
\]
Combining the $16,16,8$ summands gives precisely
$20D_n^L+16D_n^{L-s}+4D_n^{L-2s}(D_n-2)^{2s}$,
proving \eqref{eq:projectivewordtrace} for every $n$.
The denominator two is the actual projector denominator.
In each degree its invariant space has an integral-dimensional
basis of orbit sums of the subset basis; this is not a division
of an unexplained numerical count.

\par\smallskip\noindent\textbf{Centered purity and hard Lefschetz.}\quad
Directly,
\[
 H^{2r}(K_d)=H^{2(r+d)}(X_d)(d)
   =\Q_\ell(-r)^{\oplus\binom{2d}{r+d}},\qquad -d\le r\le d.
\]
The invariant projector acts within each scalar Tate group, so
its image has the same degree-weight relation. This proves the
claimed formula for the cohomology of \eqref{eq:wordcomplex}.
At a point, BBD~5.1.8, pp.~126--127, says that purity of a
complex of weight zero means precisely that its degree-$i$
cohomology is pure of weight $i$ \cite{BBD}.
This is the conclusion here. Its ungraded total has nonzero
weight $-2L$ and $2L$ pieces for $L\ge1$, and is mixed.
Without the centering shifts, the differently twisted dimensions
in \eqref{eq:wordcomplex} need not have the same weight in a
fixed absolute cohomological degree.

Let $\eta_d=\sum_i c_1(\mathcal O_i(1))$.
This is the Chern class of the ample line bundle
$\mathcal O(1,\ldots,1)$, defined over $k$, and $\eta_L$ is
$\sigma$-invariant. Hard Lefschetz applies to the projective
map $X_d\to\Spec k$ and the pure constant intersection complex
\cite[Theorem~5.4.10, p.~144]{BBD}, giving
\begin{equation}\label{eq:wordhardlefschetz}
 \eta_d^{2r}:H^{-2r}(K_d)
          \xrightarrow{\sim} H^{2r}(K_d)(2r),\qquad r\ge0.
\end{equation}
Both this map and its inverse commute with $\sigma$, so it
restricts to an isomorphism on the invariant summand. On
\eqref{eq:wordcomplex}, use the block diagonal operator with
$\eta_L$, $\eta_{L-s}$ and the restriction of $\eta_L$.
All blocks are centered at zero, so the same Lefschetz
isomorphisms hold on their sum.

Below the center, cup product
$H^{2r}\to H^{2r+2}(1)$ is injective: a vector it kills is
also killed by the further power furnishing the Lefschetz
isomorphism to degree $-2r$. Hence $c_r\le c_{r+1}$ for $r<0$,
and \eqref{eq:wordhardlefschetz} gives $c_{-r}=c_r$.
This is the geometric symmetry and unimodality assertion in $v^2$.
In this product example one can also check the Lefschetz action
directly: multiplication by $x$, contraction $x\mapsto1$, and
the degree operator on each $H^*(\PP^1)$ form the standard
$\mathfrak{sl}_2$ representation. Their tensor operators commute
with factor permutations, and the raising operator is cup product
by the displayed invariant ample class.
\end{proof}

\paragraph{Preserving the actual summand is a stronger requirement.}
The trace realization does not extend to every action entry:
$D-2=v^2+v^{-2}$ already has coefficients $(1,0,1)$ in $v^2$.
There is also a basis-independent obstruction. On the actual
rank-two module $N(D-2)$,
\[
 \Tr(P_1P_2)=(D-2)^2=v^4+2+v^{-4}.
\]
This is not an $\mathrm{SL}_2$ character. Any irreducible
constituent supplying weight 4 also supplies weight 2, whereas
the latter coefficient is zero.
Let $e_2$ be the actual commutant projector onto the
multiplicity-four $N(D-2)$ sector of $M$. Then
\[
 \Tr_M(e_2P_1P_2)=4(v^4+2+v^{-4})
\]
is likewise not such a character. The projector belongs to the
regular commutant of \fid{F43}; it is not asserted to belong to
the nonsplit regular right algebra.

Suppose matrix kernels for $P_1,P_2,e_2$ had entries actual
finite-dimensional $\mathrm{SL}_2$ representations, with direct
sum and tensor product as composition, and their character
matrices reproduced these actual operators in some basis.
The sum of the diagonal entries of their composite kernel would
be an actual $\mathrm{SL}_2$ representation. Its character would
equal the last displayed trace, a contradiction. Trace invariance
makes this argument independent of the chosen basis.
Thus that positive kernel framework cannot preserve all three
specified operations simultaneously. This does not exclude the
word-dependent motive just constructed, a model omitting the
projector, virtual kernels, or a category with different operations.

\paragraph{The exact scope of the projective repair.}
The empty word is separately realized by 64 unit objects.
For a nonempty word the total dimension here is
\[
 T_w(1)=20\cdot4^L+20\cdot4^{L-s}.
\]
It varies with the word, consistently with \fid{F42}'s
fixed-common-space obstruction. The present object depends on
$(L,s)$, just as the scalar trace does. We have supplied an actual
projector, polarization and Lefschetz maps, rather than inferred
them from numerical positivity. No action of both nonstandard
generators on these motive objects, compatibility with all word
concatenations, or identification with VIII's designated canonical
basis and structure constants has been proved.


\subsection{An exact pure Weil realization of the actual family}
The actual algebra and module admit a finite realization by a
specified linear Frobenius. This retains the collision nilpotent
and both commuting actions. Its comparison with extension-dependent
quantum grading can then be tested directly.

\begin{finding}[\fid{F46}: actual End/Hom, coefficient geometry and pure Weil data]\label{f:cayleyweil}
Put
\[
 \delta=v+v^{-1},\quad D=\delta^2,\quad h=D-1,\quad
 z_1=\frac hD,\quad z_2=\frac{4h}{D^2},\quad
 \Delta=z_1-z_2=\frac{(D-1)(D-4)}{D^2},
\]
and use the finite coefficient localization
\[
 S_0=\Q(i)[v,v^{-1},\delta^{-1},h^{-1},(D-2)^{-1}],
 \qquad S=S_0[(1-\Delta)^{-1},(1+\Delta^2)^{-1}].
\]
Let $C=S[x]/(x(x-\Delta))$, with $x=z_1-z$ in the actual
presentation, and let $L_1=C/(x)$, $L_2=C/(x-\Delta)$.
Give the regular rank-two module $J=C$ the linear operator
\[
 F_J=\frac{1+ix}{1-ix}=1+s x,\qquad
 s=\frac{2i}{1-i\Delta},\qquad
 \lambda=\frac{1+i\Delta}{1-i\Delta}.
\]
For a scalar $a$, write $S_a$ for a rank-one module with
Frobenius $a$. Give $L_1,L_2$ Frobenius $1,\lambda$ and set
\[
 H=S_{-1}\oplus S_i\oplus J^{\oplus2},\qquad
 K=S_{-1}^{20}\oplus S_i^4\oplus L_1^{16}\oplus L_2^4.
\]
These have ranks $6,44$. Then
\[
 \End_F(H)=A=S\times S\times M_2(C),\qquad
 \operatorname{Hom}_F(H,K)=M,
\]
as the actual rank-three algebra and its rank-$64$ right tensor
module in \fid{F19}/\fid{F21}, with all named generators retained.
These identities commute with every coefficient base change.
Postcomposition also identifies
\[
 \End_F(K)\simeq\End_A(M),
\]
so the actual ordinary left operators from \fid{F43} lift to $K$.
The latter commutant has rank $688$ and special-fiber dimension
$816$ at $v=1$.

There are perfect coherent objects $\mathscr H,\mathscr K$ on
the smooth proper coefficient scheme $\mathbf P^1_S$ with those
same End/Hom groups and
\[
 \operatorname{Ext}^1_{\mathbf P^1_S}(\mathscr K,\mathscr K)
       =S^{688}\oplus(S/(\Delta))^{128}.
\]
Thus both crossed extension directions and the order-two collision
are retained, with additional self-Ext explicitly present.

For each prime power $q>1$, all evaluations $v=q^{\pm1/2}$ are
defined and give finite Weil representations pure of weight zero
under every complex embedding. At $v=1$, $J$ retains a nontrivial
Jordan block with eigenvalue one. These are chosen linear Weil
data; they are not scheme Frobenius on the coefficient geometry.
They do not satisfy extension-wise parameter retuning.
\end{finding}

\begin{proof}
\par\smallskip\noindent\textbf{The finite coefficient domain and cyclic operator.}\quad
The parameters $z_1,z_2$ and $\Delta$ belong to $S_0$. Both $1-i\Delta$ and
$1+i\Delta$ are units in $S$ because their product is
$1+\Delta^2$. For real $v\ne0$,
\[
 D\ge4,\qquad
 0\le\Delta=1-\frac5D+\frac4{D^2}<1.
\]
Thus every displayed denominator is nonzero at every such real
parameter, in particular at $q^{\pm1/2}$ and at $1$. This finite
localization admits the required evaluations; it does not
invert $\Delta$.

The relation $x^2=\Delta x$ gives
\[
 \frac{1+ix}{1-ix}=1+\frac{2i}{1-i\Delta}x,\qquad
 \lambda=1+s\Delta.
\]
The inverse is obtained by replacing $i$ by $-i$.
In the column basis $(1,x)$,
\begin{equation}\label{eq:cayleymatrix}
 F_J=\begin{pmatrix}1&0\\s&\lambda\end{pmatrix}.
\end{equation}
Since $s$ is a unit, multiplication by $x$ is
$s^{-1}(F_J-1)$. Therefore $S[F_J]=C$, and a linear map
commutes with $F_J$ exactly when it is $C$-linear.
The centralizer on the regular module is consequently
\[
 \End_{F_J}(J)=\End_C(C)=C.
\]
The argument works after every coefficient base change.
Equivalently $x=(F_J-1)/(i(F_J+1))$.
The vector $1$ is cyclic, since $(1,F_J1)$ has determinant $s$.
The characteristic polynomial is $(T-1)(T-\lambda)$.
At $v=1$ it remains cyclic and becomes
\[
 F_J=1+2ix,\qquad x^2=0,\quad x\ne0.
\]
Thus its repeated eigenvalue is a genuine Jordan block.

\par\smallskip\noindent\textbf{Exact algebra, module and both actions.}\quad
The scalar sectors are separated by units:
\[
 F_J+1=\frac2{1-ix},\qquad
 F_J-i=\frac{(1-i)(1-x)}{1-ix}.
\]
The second is invertible because the norm of $1-x$ is
$1-\Delta$, a unit of $S$. Also $-1-i$ is a unit.
Hence no intertwiners occur between different scalar sectors
or between those sectors and $J$. The same identities hold
on the branch quotients. The matrix sector of $\End_F(H)$ is
therefore $M_2(C)$, giving the asserted algebra.

An intertwiner from $J$ to a branch commutes with
$x=s^{-1}(F-1)$ and is $C$-linear. Hence
\[
 \operatorname{Hom}_F(J^{\oplus2},L_j)
    =\operatorname{Hom}_C(C^{\oplus2},L_j)
    =\operatorname{Row}(L_j).
\]
The scalar Hom groups contribute ranks $20,4$ and the branch
groups ranks $32,8$. In the column convention for $H$,
precomposition is right multiplication on these rows.
Thus this is precisely
\[
 M=\epsilon_+^{20}\oplus\epsilon_-^4
       \oplus\operatorname{Row}(L_1)^{16}
       \oplus\operatorname{Row}(L_2)^4.
\]
The chosen coordinate makes $L_1$ the multiplicity-$16$
branch $z=z_1$ and $L_2$ the multiplicity-$4$ branch $z=z_2$.
In particular the named normalized generator matrices are
\[
 \begin{pmatrix}1&0\\0&0\end{pmatrix},\qquad
 \begin{pmatrix}1-z&1\\z(1-z)&z\end{pmatrix},
 \qquad z=z_1-x,
\]
with the original generators obtained by multiplying by $D$,
and their scalar values $D,0$ as in \fid{F19}.
The certified regular lattice comparison of \fid{F21}
is defined over $S_0$, so this retains the actual action.
All arguments used unit identities and regular $C$-modules,
and therefore commute with arbitrary coefficient base change.
No diagonalization using $\Delta^{-1}$ was used.

Over the domain $S$, a map between the branch lines is
annihilated by $\lambda-1=s\Delta$ and hence is zero.
It follows that
\begin{equation}\label{eq:cayleycommutant}
 \End_F(K)=M_{20}(S)\times M_4(S)\times M_{16}(S)\times M_4(S).
\end{equation}
At $v=1$ the branch eigenvalues coincide and their multiplicities
combine into $20$, so this dimension increases from $688$ to
$400+16+400=816$. The base-change map is block diagonal
on those $20$ coordinates and has cokernel dimension $128$.
There is no such jump in $\End_F(H)$ or $\operatorname{Hom}_F(H,K)$,
because their source $J$ remains cyclic at collision.

Postcomposition gives an isomorphism
\[
 \End_F(K)\longrightarrow\End_A(\operatorname{Hom}_F(H,K)).
\]
For the scalar sectors this is immediate. On the matrix sector
put $K_C=L_1^{16}\oplus L_2^4$. The two diagonal matrix units
force an $M_2(C)$-commuting map on the row module $K_C^2$
to preserve its two coordinates; the two off-diagonal units
force the same $C$-linear map on each coordinate.
This is exactly postcomposition by a unique endomorphism of
$K_C$. The proof works over every base ring.
In the multiplicity bases of \fid{F43}, this is the identical
product-matrix action, so each actual ordinary left operator
lifts to $K$ and commutes with its chosen $F$.
Both commutants have dimension $816$ at collision, while
the additional $128$ maps need not lift over $S$.

\par\smallskip\noindent\textbf{Smooth proper coefficient geometry and extensions.}\quad
On $\mathbf P^1_S$ let $T$ denote the affine coordinate and put
\[
 Z=V((T-1)(T-\lambda)).
\]
It is a finite flat degree-two divisor, avoiding both infinity
and zero. The homomorphism
\[
 S[T,T^{-1}]/((T-1)(T-\lambda))
       \longrightarrow C,\qquad T\longmapsto1+s x
\]
is an isomorphism with inverse $x=s^{-1}(T-1)$, because its
relation becomes $s^2x(x-\Delta)=0$.
Write $O_a$ for the section sheaf at $T=a$, and define
\[
 \mathscr H=O_{-1}\oplus O_i\oplus O_Z^{\oplus2},\qquad
 \mathscr K=O_{-1}^{20}\oplus O_i^4
       \oplus O_1^{16}\oplus O_\lambda^4.
\]
Each section has its Cartier resolution by $O(-1),O$,
and $Z$ has the corresponding resolution by $O(-2),O$.
Thus these objects are perfect on the smooth proper
relative projective line. Their supports are finite over
affine $\Spec S$, so derived global sections are their
finite free modules in degree zero, namely $H,K$.
Computing sheaf Hom on the affine chart containing their
supports gives
\[
 \End_{\mathbf P^1_S}(\mathscr H)=A,\qquad
 \operatorname{Hom}_{\mathbf P^1_S}(\mathscr H,\mathscr K)=M.
\]
Multiplication by $T$ supplies the extra linear operators.

For the coherent Ext computation use the section resolution
\[
 0\longrightarrow O(-1)\xrightarrow{T-a}O
       \longrightarrow O_a\longrightarrow0.
\]
Hom into $O_b$ gives the two-term complex
$S\xrightarrow{b-a}S$. Its terms have no higher cohomology,
being supported on an affine section. Consequently
\[
 \operatorname{Hom}(O_a,O_b)=\operatorname{ann}_S(b-a),
 \qquad
 \operatorname{Ext}^1(O_a,O_b)=S/(b-a),
\]
and there are no higher Ext groups. For $a=b$ the two
groups are $S$. All distinct scalar-sector differences
are units; the sole nonunit cross difference is
$\lambda-1=s\Delta$. Both branch directions therefore give
$S/(\Delta)$, and
\begin{equation}\label{eq:cayleyext}
 \operatorname{Ext}^1_{\mathbf P^1_S}(\mathscr K,\mathscr K)
      =S^{688}\oplus(S/(\Delta))^{128}.
\end{equation}
The free rank is $20^2+4^2+16^2+4^2$; the torsion multiplicity
is $2\cdot16\cdot4$. At the actual order-two collision its
torsion length is $256$. This retains both crossed groups,
but the additional free self-Ext is not the self-Ext of the
original $C$-branch modules. No full derived equivalence
or extra Frobenius-equivariant derived Ext comparison is
asserted.

\par\smallskip\noindent\textbf{All-conjugate purity and continuous realizations.}\quad
Fix a prime power $q>1$ and evaluate $v=q^{1/2}$ or $q^{-1/2}$
in $\Q(i,\sqrt q)$. Then $D=q+2+q^{-1}$ and $\Delta$ is
rational with $0<\Delta<1$. The eigenvalues of $H,K$ lie in
\[
 \{-1,i,1,\lambda\},\qquad
 \lambda=\frac{1+i\Delta}{1-i\Delta}.
\]
Every complex embedding fixes $\Delta$ and sends $i$ to
$i$ or $-i$. Thus it sends $\lambda$ to $\lambda$ or
$\lambda^{-1}$, both of modulus one. The possible sign
change of $\sqrt q$ does not change $D$ or $\Delta$.
After an embedding into a finite extension of $\Q_\ell$,
these are finite Weil representations pure of weight zero
in the eigenvalue sense. No integrality assertion is
part of this purity conclusion.
For a general real algebraic parameter, total reality is a
sufficient additional hypothesis for the analogous all-conjugate
assertion; it is not asserted to be necessary.

At $v=1$ the same construction retains the nontrivial Jordan
block \eqref{eq:cayleymatrix} with eigenvalue one. It can be
viewed as a pure Weil object over any chosen finite field;
this does not posit a field of cardinality one.
Geometric restriction to a point forgets that arithmetic
Jordan data, so there is no contradiction with geometric
semisimplicity.

The ordinary internal Hom has dimension $6\cdot44=264$ and
Frobenius
\[
 \Phi(u)=F_KuF_H^{-1}.
\]
Its eigenvalues are ratios of the preceding unit-modulus
eigenvalues, so it too is pure of weight zero, with Jordan
blocks allowed. Its fixed subspace is exactly $M$, on which
$\Phi$ is the identity. Likewise $A$ is the fixed subalgebra
of the $36$-dimensional internal End object.

A finite Weil representation on a point is an invertible
linear operator, since the point Weil group is discrete
infinite cyclic. Extending it to a continuous representation
of the profinite Galois group requires its eigenvalues to be
$\ell$-adic units. This condition is necessary: compactness
of the image bounds all positive and negative powers.
It is also sufficient. After a finite coefficient extension,
choose Jordan form with unit eigenvalues. The matrix and
its inverse preserve a suitable integral lattice, so their
powers lie in a compact matrix group. Modulo each lattice
congruence subgroup the cyclic image is finite; these maps
extend compatibly from $\Z$ to $\widehat{\Z}$.
For each fixed $q$, the algebraic number $\lambda$ is a unit
outside finitely many finite places. Choosing
$\ell\nmid q$ outside that set therefore gives continuous
$\ell$-adic realizations retaining every displayed comparison.
No uniform choice for all $q$, or validity at every $\ell$,
is asserted.

\par\smallskip\noindent\textbf{A proved extension mismatch and an integrality boundary.}\quad
Write $\lambda(q)$ for the branch value at $v=q^{1/2}$.
For $q>1$,
\[
 \lambda(q)\ne1,\qquad
 \Delta(q^n)\longrightarrow1,\qquad
 \lambda(q^n)\longrightarrow i.
\]
If $\lambda(q)^n=\lambda(q^n)$ for every sufficiently
large $n$, the ratio of consecutive terms on the right
would tend to one, because their limit is nonzero.
The same ratio on the left is $\lambda(q)$, a contradiction.
Thus raising the fixed Frobenius to its $n$th power is
not retuning this named branch at $q^n$.

For a concrete arithmetic boundary, at $q=2$ one has
\[
 D=\frac92,\quad \Delta=\frac7{81},\quad
 \lambda=\frac{81+7i}{81-7i}
       =\frac{3256+567i}{3305},\quad
 \lambda+\lambda^{-1}=\frac{6512}{3305}\notin\Z[1/2].
\]
If $\lambda$ were integral over $\Z[1/2]$, its conjugate
and their rational sum would be integral, and that sum
would lie in the integrally closed ring $\Z[1/2]$.
This contradiction is unchanged by multiplying $\lambda$
by any integer power of $2$.
In particular it cannot be an eigenvalue in a Tate twist
of constant-coefficient cohomology of a smooth projective
variety over $\F_2$: the monic integral characteristic
polynomial in Deligne's theorem
\cite[Theorem~(1.6), p.~276; argument on pp.~276--277]{WeilI}
makes the untwisted eigenvalues algebraic integers, and
Tate twists introduce only powers of $2$ into denominators.
This excludes that extra comparison, not the constructed
pure Weil objects. For example $\lambda$ is a unit at
$\ell=3$ and gives a continuous realization there.
At $\ell=5$, the residues $i=\pm2$ make exactly one of
$81+7i$ and $81-7i$ divisible by $5$, so it is not a unit.
\end{proof}

The projective line above is over the characteristic-zero
parameter/coefficient base. It is not thereby a variety
over $\F_q$ with geometric Frobenius equal to multiplication
by $T$. The evaluation $v=\sqrt q$ occurs in a characteristic-zero
coefficient field and cannot simultaneously be that same
characteristic-$p$ base evaluation, where $q$ becomes zero
while $v$ must remain invertible. The coherent geometry and
the chosen linear Weil operator are separate data.

This realization closes the existence question for exact
finite pure coefficient objects and ordinary smooth proper
coefficient geometry retaining this rank-three action.
It does not identify the chosen operator with absolute
Frobenius on the center, the lift of \fid{F44}, or cyclotomic
Frobenius on TP. The unit-modulus spectrum does not supply
a positive self-adjoint structure, the designated VIII
globalization, or the \fid{F24} path-sign comparison.
The extension mismatch and integrality calculation above
are explicit limits of those additional interpretations.


\subsection{Positive duality and the radical at the collision}
\begin{finding}[\fid{F47}: a positive adjoint realization must lose the radical]\label{f:positivecollision}
The actual algebra
\[
 A_1=\C\times\C\times M_2(\C[\varepsilon]/\varepsilon^2)
\]
cannot be the complexified Grothendieck ring of a multifusion
category. The same conclusion applies to a graded categorification
whose specialization at $v=1$ is explicitly required to have that form.
It is not asserted that forgetting any grading automatically produces
such a category.

More directly, if $\rho:A_1\to\End(H)$ is a representation on a
finite-dimensional positive-definite complex Hermitian space and
both named generators $\rho(P_1),\rho(P_2)$ are self-adjoint, then
\[
 0\times0\times M_2(\varepsilon\C)\subseteq\ker\rho.
\]
Thus this positive adjoint realization factors through the actual
rank-six classical quotient and cannot retain the full rank-ten
collision algebra faithfully.

On the regular rank-two center module, the natural trace pairing
is positive away from the collision and has determinant $\Delta^2$,
of order four at $v=1$. This is the smallest possible determinant
order of a regular symmetric form making the center coordinate
self-adjoint and positive at nearby real noncollision parameters.
The actual rank-64 tensor module, which kills the radical after
specialization, does have a positive form regular through the collision.
\end{finding}

\begin{proof}
The categorical input is the known finite-based-ring theorem:
the Grothendieck ring of a multifusion category is based, and its
complexification is semisimple
\cite[Definition~1.42.2 and Proposition~1.42.4, pp.~76--77;
Corollary~1.43.5, p.~80]{EGNO2009}.
The nonzero square-zero ideal
$J=0\times0\times M_2(\varepsilon\C)$ is the Jacobson radical
of $A_1$, by \fid{F19}. This contradicts the proposed semisimplicity
and proves the first assertion. The known theorem concerns finite
based rings, not every ring with nonnegative structure constants.

For the operator assertion put $e=P_1/4$, $f=P_2/4$ and
$z=(e-f)^2$. Finding~\fid{F19} shows that $z$ is central, zero
on the two scalar factors, and $3/4+\eta$ on the matrix factor,
where $\eta$ is a nonzero central square-zero element generating
$J$ as an ideal. The real polynomial
\[
 p(T)=1-\frac{(T-3/4)^2}{(3/4)^2}
\]
makes $p(z)$ the central projector onto the matrix factor.
Consequently $n=p(z)(z-3/4)$ is zero on the scalar factors,
equals $\eta$ on the matrix factor, satisfies $n^2=0$, and
generates $J$. Since $\rho(e)$ and $\rho(f)$ are self-adjoint,
$\rho(z)$ and then $\rho(n)$ are self-adjoint. For every $u\in H$,
\[
 \|\rho(n)u\|^2=\langle u,\rho(n)^2u\rangle=0.
\]
Thus $\rho(n)=0$ and $J\subseteq\ker\rho$. The quotient is
$\C\times\C\times M_2(\C)\simeq\C[S_3]$, with the actual
specialized generator map already identified in \fid{F19}.

For the pairing assertion work over
$R=\R[v,v^{-1}]_{(v-1)}$, set $t=v-1$, and write
\[
 C=R[x]/(x(x-\Delta)),\quad x=z_1-z,\qquad
 \Delta=\frac{D-1}{D^2}(v-v^{-1})^2=t^2u(t),\quad u(0)=3/4.
\]
The regular algebra trace gives, in the basis $(1,x)$,
\[
 G_{\rm tr}=\begin{pmatrix}2&\Delta\\\Delta&\Delta^2\end{pmatrix},
 \qquad \det G_{\rm tr}=\Delta^2.
\]
Indeed the traces of multiplication by $1,x,x^2$ are
$2,\Delta,\Delta^2$. The leading principal minor is 2 and the
determinant is positive at each nearby real point with
$\Delta\ne0$, proving positive definiteness and order four.

To prove optimality, right multiplication by $x$ has row matrix
\[
 X=\begin{pmatrix}0&1\\0&\Delta\end{pmatrix}.
\]
For a symmetric regular form
$G=\left(\begin{smallmatrix}a&b\\b&c\end{smallmatrix}\right)$,
the relation $XG=GX^{\mathsf T}$ is precisely $c=\Delta b$.
Positive definiteness for small real $t>0$ gives
\[
 a(t)>0,\qquad 0<b(t)<a(t)\Delta(t),\qquad
 \det G=b(a\Delta-b)>0.
\]
Let $s=\operatorname{ord}_t a$ and $r=\operatorname{ord}_t b$;
both are finite nonnegative integers, since generic positivity
excludes their identically vanishing. The inequality forces
$r\ge s+2$: otherwise $b/(a\Delta)$ tends to positive infinity
as $t\to0^+$, contradicting its being less than one. Therefore
\[
 \operatorname{ord}_t\det G
 \ge r+(s+2)\ge2s+4\ge4.
\]
The trace form attains the lower bound. The identical proof for a
positive $\Delta=t^m u(t)$ on a punctured real half-neighborhood
gives the optimal bound $2m$.

Finally, on either actual right block $N(a)$ of \fid{F21}, where
$a=\delta$ or $a=D-2$, put
\[
 G_a=\begin{pmatrix}1&a/D\\a/D&1\end{pmatrix}.
\]
Direct multiplication with
$P_1=\left(\begin{smallmatrix}D&0\\a&0\end{smallmatrix}\right)$
and $P_2=\left(\begin{smallmatrix}0&a\\0&D\end{smallmatrix}\right)$
gives $P_iG_a=G_aP_i^{\mathsf T}$. For every real $v>0$,
$\delta\ge2$, $D\ge4$, and
\[
 0<\delta/D=1/\delta<1,\qquad 0<(D-2)/D<1.
\]
Thus both forms are positive and specialize to
$\left(\begin{smallmatrix}1&1/2\\1/2&1\end{smallmatrix}\right)$.
Give the scalar sectors positive standard forms and take the
direct sum with multiplicities $20,4,16,4$. This is the asserted
positive form on the actual tensor module. The change of basis
of \fid{F21} is regular and invertible at $v=1$, so the form
transports to its original tensor basis without a pole there.
Its special action kills $J$, exactly as the operator theorem requires.
\end{proof}

An indefinite form can have nonzero self-adjoint nilpotents,
and purity alone imposes no positive-adjoint condition on the
named generators. A smooth proper dg category may also have
this algebra as the endomorphism algebra of one object, as
\fid{F40} and \fid{F44} show: this is a different role from the
whole based Grothendieck algebra. The result excludes the
specified finite rigid and faithful positive-adjoint routes,
not those other realizations or VIII's weaker signed-factor
positivity convention.


\subsection{Tensor realization and the exact arithmetic comparison gate}
The fixed-parameter action extends through all tensor degrees.
The stronger requirement that finite Weil centralizers recover the
entire nonstandard endomorphism system has different consequences.
Rational dependence on the quantum parameter makes the extension
comparison especially restrictive.

\begin{finding}[\fid{F48}: faithful tensor coefficients and full-centralizer obstructions]\label{f:tensorweilgate}
In the actual two-row Kronecker setting let $X=V\otimes W$,
with $\dim V=\dim W=2$, and fix a coefficient parameter at which
the indicated source maps are defined.
The finite-dimensional nonstandard polynomial-comodule category
has an exact faithful strong monoidal realization by finite
Weil coefficient objects, all pure of weight zero after an
arithmetic coefficient specialization. In IV's actual Hopf
setting this realization preserves the existing rigid duals.
It retains every named tensor action and every specified
source map. Nontrivial finite-order tensor Frobenii also exist.

The realization cannot be fully faithful into the category of
finite-dimensional Weil representations on a point.
The actual polynomial endomorphism algebras already satisfy
\[
 \End(X)=E,\qquad
 \End(X^{\otimes2})=E\,1\oplus E\,P,
\]
where $P$ is the rank-ten nonstandard symmetric projector.

There are additional arithmetic restrictions independent of
that categorical fullness condition. Let $b>1$ and let
$F(v)\in GL_N(\C(v))$ be regular and invertible at $b$ and
at all sufficiently large $b^n$. If
\begin{equation}\label{eq:tensorweilspectral}
 \operatorname{Spec}F(b^n)
   =\operatorname{Spec}(F(b)^n)
\end{equation}
as eigenvalue multisets for every sufficiently large integer $n$,
then, for finitely many nonnegative integers $r_m$,
\[
 \operatorname{Tr}F(v)=\sum_m r_m v^m,\qquad
 \det(TI-F(v))=\prod_m(T-v^m)^{r_m}.
\]
The $r_m$ are the actual positive multiplicities of the
eigenvalues that occur.

Consequently no such rational family of one fixed integral
pure weight can recover the actual nonsimple semisimple
rank-three algebra at $b>1$ as its full Frobenius centralizer,
at any auxiliary rank. If $F(v)$ is also regular at $v=1$,
it cannot recover the full actual collision algebra there,
even with arbitrary mixed weights.

Under the stronger equality of matrices
$F(b^n)=F(b)^n$ in one fixed rational trivialization,
there is a constant invertible matrix $S$ and integers $m_j$
such that
\[
 F(v)=S\operatorname{diag}(v^{m_1},\ldots,v^{m_N})S^{-1}.
\]
These restrictions do not exclude the faithful actions as
specified subalgebras supplied by the first part.
\end{finding}

\begin{proof}
\par\smallskip\noindent\textbf{Actual source data and conventions.}\quad
Write $B$ for IV's nonstandard matrix coordinate bialgebra.
Its defining relation is
\[
 P(u\otimes u)=(u\otimes u)P;
\]
see IV, Definition~8.1, equation~(61), \S8.1, p.~39.
In this Kronecker setting, Proposition~10.3, \S10.1, p.~56,
constructs the Hopf algebra after adjoining the inverse
nonstandard determinant. Proposition~10.5, \S10.2, p.~57,
gives its Hopf map to the product of the ordinary quantum
$GL(V)$ and $GL(W)$ coordinate algebras; the determinant
maps to $D_v(V)^2D_v(W)^2$.
Theorem~12.1, p.~68, identifies the degree-$r$ coordinate
dual with the commutant of the adjacent projectors
\cite{IV}.

The remark immediately after that theorem corrects its
displayed $X$ to $X^*$ relative to the conventions in
\S\S8--10, and its proof uses the dual consistently.
Here one matrix corepresentation and its induced tensor
action are fixed throughout. Passing to its dual or
transposed presentation preserves the endomorphism
dimensions and all obstruction arguments below; the
source projector is symmetric in its chosen coordinates.

The Hopf assertion just cited is specific to IV.
VII, \S2, explicitly allows the general plethysm determinant
to vanish. Its bialgebra category still has tensor products;
no general VII antipode or rigid duals are inferred here.
At an algebraic specialization we use only the coefficient
maps actually defined there, without importing a generic
compact-real-form theorem into an unverified specialization.

Put
\[
 \delta=v+v^{-1},\qquad D=\delta^2,\qquad
 P=P_+^V\otimes P_+^W+P_-^V\otimes P_-^W.
\]
The ordinary symmetric and alternating projectors have ranks
$3$ and $1$. Thus $P$ has rank $3\cdot3+1\cdot1=10$,
and $1-P$ has rank $6$. On $X^{\otimes r}$ the actual
right generators are $DP_i$ on consecutive tensor factors.
These assertions hold in particular at positive real
parameters and at $v=1$.

\par\smallskip\noindent\textbf{The faithful tensor realization.}\quad
Let $\omega$ be the ordinary forgetful functor on finite
$B$-comodules. Send
\[
 N\longmapsto(\omega N,1)
\]
to finite-dimensional vector spaces with an invertible
linear operator. Kernels and cokernels of comodule maps
are computed on underlying vector spaces, so this functor
is exact and faithful. The bialgebra multiplication defines
the source tensor product; its underlying tensor product
is the ordinary one. Since $1\otimes1=1$, the ordinary
associativity and unit comparisons are Weil isomorphisms
with the required coherence.
In a rigid Hopf category the ordinary underlying dual,
evaluation and coevaluation are unchanged, and the dual
of the identity is the identity. Thus this also retains
the specified IV dual maps and determinant inversion.
After an arithmetic coefficient embedding all eigenvalues
are one, giving purity of weight zero and a continuous
point-Galois representation.

The right generators on each image remain exactly the
original matrices $DP_i$. Tensoring two consecutive blocks
gives their actual block inclusion, and all original
relations remain matrix equalities.
Every source comodule map remains the same linear map.
In addition, ordinary left operators that are not comodule
intertwiners, including the named operators of \fid{F43},
are still Weil endomorphisms of these coefficient objects,
because they commute with identity Frobenius.
This is retention of the actual operations, not a claim
that arbitrary linear maps become source comodule maps.

One may assign a fixed pure unit scalar $c$ to $X$ and
$c^r$ to homogeneous degree $r$, with inverse on dual
degree, when that grading of all maps is part of the
data. The identity construction requires no additional
grading qualification and is the unconditional version.

\par\smallskip\noindent\textbf{Nontrivial pure tensor characters.}\quad
Choose roots of unity $\zeta,\eta$ in the coefficient field.
The diagonal evaluations
\[
 A=\operatorname{diag}(\zeta,\zeta^{-1}),\qquad
 B'=\operatorname{diag}(\eta,\eta^{-1})
\]
are characters of the ordinary quantum $GL_2$ coordinate
algebras: off-diagonal generators go to zero, the remaining
relations hold for commuting diagonal scalars, and the
quantum determinants go to one. Compose with IV's actual
Hopf map to obtain a character $\chi$ of the nonstandard
Hopf algebra. On $X$ its matrix is $A\otimes B'$.
Directly in the FRT bialgebra the same character exists
because $(A\otimes B')^{\otimes2}$ commutes with $P$.

For a right comodule with coaction $\rho_N$, define
\[
 F_N=(1\otimes\chi)\rho_N.
\]
It is natural for all comodule maps. Multiplicativity of
$\chi$ gives
\[
 F_{N\otimes N'}=F_N\otimes F_{N'}.
\]
In the Hopf setting the inverse uses $\chi$ composed with
the antipode; on a dual the operator is the inverse
transpose. Choose a common order $L$ of $\zeta,\eta$.
Then $F_X^L=1$, and the same holds on every tensor product,
dual, stable subspace and quotient in the tensor category
they generate. In characteristic zero these operators
are semisimple with root-of-unity eigenvalues.
They are pure of weight zero under every embedding and
give continuous representations at every coefficient prime.

All actual right generators and source comodule maps
commute with these Frobenii. The ordinary left raising
and lowering operators have the precise covariance
\[
\begin{aligned}
 F_XE_VF_X^{-1}&=\zeta^2E_V,&
 F_XF_VF_X^{-1}&=\zeta^{-2}F_V,\\
 F_XE_WF_X^{-1}&=\eta^2E_W,&
 F_XF_WF_X^{-1}&=\eta^{-2}F_W.
\end{aligned}
\]
This follows first for the elementary upper and lower
matrices on each two-dimensional factor; the source
coproduct preserves the same equations on tensor powers.
Thus, for example, $E_V$ is a Weil map
$L_{\zeta^2}\otimes X\to X$, where $L_a$ is the rank-one
coefficient object with Frobenius $a$.
With identity Frobenius all these named operators are
ordinary invariant endomorphisms instead.
The character construction fixes the parameter; it does
not impose the extension-retuning rule.

\par\smallskip\noindent\textbf{Actual endomorphisms and failure of fullness.}\quad
In coordinate degree one there are no FRT relations.
The degree-one dual therefore acts as $\End_E(X)$,
whose commutant is $E$. In degree two the sole relations
are the entries of $PU-UP$, where $U$ is the matrix
of the universal tensor-square corepresentation.
Their annihilator identifies the acting coordinate dual
with
\[
 \{a\in\End_E(X^{\otimes2}):aP=Pa\}
   =\End(\operatorname{im}P)\oplus\End(\ker P).
\]
Taking the commutant of these two full matrix blocks gives
\[
 \End_B(X)=E,\qquad
 \End_B(X^{\otimes2})=E\,1\oplus E\,P.
\]
The blocks are nonzero of ranks ten and six.
This is also the degree-two instance of the source
proof of Theorem~12.1.

For the IV Hopf category, the facts needed for the same
obstruction are $\End_H(X)=E$ and the independent
endomorphisms $1,P$ of $X^{\otimes2}$.
The first follows by ordinary quantum-product restriction:
the elementary raising and lowering matrices on the two
fundamental factors generate all four-by-four matrix
units on $X$, so a commuting endomorphism is scalar.
The second is the defining FRT relation after determinant
inversion. No assertion that arbitrary localization
preserves every polynomial Hom space is needed.

Suppose a fully faithful $E$-linear strong monoidal
functor sent $X$ to a nonzero Weil object $(Y,F)$.
Fullness in degree one gives $\End_F(Y)=E$.
Since $F$ itself belongs to this algebra, $F=cI$.
Its full centralizer is then $\End_E(Y)$, so $\dim Y=1$.
Strong monoidality makes the image of $X^{\otimes2}$
one-dimensional as well, contradicting its independent
endomorphisms $1,P$. This proves the fullness obstruction.

Equivalently, auxiliary objects $H_r$ with
$H_{r+s}\simeq H_r\otimes H_s$ cannot identify the
degree-one and degree-two nonstandard algebras with
their entire Frobenius centralizers.
The rank-six $H$ of \fid{F46} is an auxiliary degree-three
presentation object and has not been asserted to be
such a tensor power.
This obstruction already occurs in the generic
semisimple setting; it is distinct from the collision
radical obstruction in \fid{F47}.

\par\smallskip\noindent\textbf{A scalar extension-law classification.}\quad
For later comparison, let $f\in\C(v)^\times$ be regular
and nonzero at $b>1$ and suppose
$f(b^n)=f(b)^n$ for all sufficiently large $n$.
Write its leading Laurent term at infinity as
\[
 f(v)=a v^m(1+O(v^{-1})),\qquad a\ne0,\quad m\in\Z.
\]
The ratio of consecutive evaluations tends to $b^m$,
but is identically $f(b)$ by the extension law.
Thus $f(b)=b^m$. The rational function $f(v)-v^m$
vanishes at infinitely many $b^n$ and is zero.
Conversely every $v^m$ satisfies the law.
In particular a nontrivial constant phase is not allowed
by this all-extension equation.
For $b=Q^{1/2}$ its weight is $m$, so weight zero forces
$f=1$. The inverse parameter convention is obtained by
replacing $v$ by $v^{-1}$.
The relation $f(v^{-1})=f(v)^{-1}$ already holds for
these monomials. This statement concerns a rational
labeled branch, not an unspecified multivalued
algebraic continuation.

\par\smallskip\noindent\textbf{Unlabeled rational spectra.}\quad
We now prove the stronger matrix-spectral statement
without labeling branches or assuming diagonalizability.
Put $f(v)=\operatorname{Tr}F(v)$.
Under \eqref{eq:tensorweilspectral},
\[
 f(b^n)=\sum_\alpha m_\alpha\alpha^n,
\]
where the distinct nonzero $\alpha$ are eigenvalues of
$F(b)$ and $m_\alpha>0$ are their actual multiplicities.
This finite power sum satisfies a constant-coefficient
recurrence
\[
 \sum_{j=0}^s c_j f(b^{n+j})=0,\qquad c_0\ne0.
\]
The constant term is nonzero because all roots are nonzero.
The rational function
\[
 G(v)=\sum_{j=0}^s c_j f(b^jv)
\]
vanishes at infinitely many $b^n$, so it is identically
zero. If $f$ had a nonzero finite pole $\rho$ of maximal
modulus, a pole of $f(b^jv)$ for $j\ge1$ would have modulus
at most $|\rho|/b^j<|\rho|$.
Only $c_0f(v)$ could have a pole at $\rho$, a contradiction.
Thus $f$ has poles only at zero and infinity and is a
Laurent polynomial. This is the rational dilation
argument of \fid{F26}, reproved here in the needed form.

Write $f(v)=\sum_m a_m v^m$. Evaluating gives
\[
 \sum_m a_m(b^m)^n=\sum_\alpha m_\alpha\alpha^n.
\]
Distinct nonzero finite power roots are linearly
independent: any zero relation at sufficiently many
consecutive integers gives an invertible Vandermonde
matrix, times a nonzero diagonal matrix of root powers.
Therefore every $\alpha$ is $b^m$, and $a_m$ is its
positive integer multiplicity $r_m$.
This proves the Laurent trace assertion.

At $v=b^n$ the characteristic polynomial is now
$\prod_m(T-b^{mn})^{r_m}$.
Each rational coefficient agrees at infinitely many
points with the corresponding coefficient of
$\prod_m(T-v^m)^{r_m}$, proving the family identity.
Virtual alternating traces with cancelled roots would
not justify this inference about invisible eigenvalues;
the hypothesis here is the full matrix spectrum.

\par\smallskip\noindent\textbf{One pure weight and the generic actual algebra.}\quad
Suppose $b=Q^{1/2}$ and $F(b)$ is pure of one fixed
integral weight $w$. Normalize
\[
 F_0(v)=v^{-w}F(v).
\]
It is rational, has the same full centralizer, and still
satisfies \eqref{eq:tensorweilspectral}, since the
normalizing factor at $b^n$ is $(b^{-w})^n$.
Every eigenvalue of $F_0(b)$ has modulus one.
The spectral classification forces them all to be one,
so $N_0=F_0(b)-I$ is nilpotent.

At $b>1$, the actual rank-three algebra of \fid{F19} is
\[
 A_b\simeq E\times E\times M_2(E)\times M_2(E)
\]
over its splitting coefficient field.
Indeed $D=(b+b^{-1})^2>4$ and
$\Delta=(D-1)(D-4)/D^2\ne0$, so its center branches split.
This algebra is semisimple and not simple.
If it were the full centralizer of $F_0(b)$, the
nilpotent $N_0$ would be central in that semisimple algebra,
because every centralizer element commutes with it.
A semisimple finite algebra has reduced center, so
$N_0=0$. But the centralizer of $I$ is $M_N(E)$,
which is simple, a contradiction.
The argument permits every auxiliary rank and every
Jordan form. For $v=Q^{-1/2}$ the normalization is
$v^wF(v)$, or equivalently the preceding proof after
inverting the parameter.

\par\smallskip\noindent\textbf{Mixed weights and the retained collision.}\quad
We next prove the center fact needed at $v=1$.
For a nilpotent operator $N_0$ of nilpotency index $s$
on a nonzero finite-dimensional $k$-space,
\begin{equation}\label{eq:nilpotentcentralizercenter}
 Z(\End_{N_0}(Y))=k[N_0]\simeq k[t]/(t^s).
\end{equation}
Regard $Y$ as the $k[t]$-module
$\bigoplus_i k[t]/(t^{r_i})$, with $\max r_i=s$.
A central endomorphism commutes with the projections
onto the cyclic summands, so is block diagonal.
Its restriction to each summand is multiplication by
a polynomial modulo $t^{r_i}$.
Choose a largest summand $k[t]/(t^s)$ and let $f$
be that restriction. The quotient map from it to each
smaller summand, sending $1$ to $1$, is an intertwiner.
Commuting with these maps forces every restriction to
be $f$ modulo the appropriate power of $t$.
This also handles multiple largest summands.
Conversely every polynomial in $N_0$ is central.
The largest summand makes the kernel exactly $(t^s)$,
proving \eqref{eq:nilpotentcentralizercenter}.
Its center is local, over any field.

If the rational family $F(v)$ is regular at $v=1$,
the proved characteristic-polynomial identity gives
\[
 \det(TI-F(1))=(T-1)^N.
\]
Thus $F(1)-I$ is nilpotent and its full centralizer has
local center. The actual collision algebra, however, is
\[
 A_1=k\times k\times M_2(k[\epsilon]/(\epsilon^2)),
\qquad
 Z(A_1)=k\times k\times k[\epsilon]/(\epsilon^2).
\]
Its center has three nonzero primitive central idempotents
and is not local. Hence it cannot be the full centralizer
of $F(1)$.
No purity assumption was used in this collision argument.
A rational change of basis giving a regular matrix model
does not change its characteristic-polynomial identity
or its local-center conclusion.
The assumptions are a finite fixed-rank family, regularity
at the collision, and recovery of the full special
centralizer; a nonflat coherent object of different
special rank or a specified proper subalgebra is outside
this assertion.

\par\smallskip\noindent\textbf{The stronger equality of matrices.}\quad
Finally suppose $F(b^n)=F(b)^n$ in a fixed rational
trivialization. Choose a constant complex Jordan basis
$S$ for $F(b)$ and set $G(v)=S^{-1}F(v)S$.
Every diagonal entry satisfies
$g_j(b^n)=\alpha_j^n$, so the scalar argument above gives
$g_j(v)=v^{m_j}$ and $\alpha_j=b^{m_j}$.

If a Jordan block had length at least two, its first
superdiagonal entry $h(v)$ would satisfy
\[
 h(b^n)=n\alpha^{n-1}.
\]
The ratio of consecutive values tends to
$\alpha=b^m$, forcing the leading Laurent exponent
of $h$ to be $m$.
But $h(b^n)/b^{mn}$ tends to its finite leading
coefficient, whereas the displayed value is
$n/\alpha$, unbounded. This is impossible.
All Jordan blocks have length one.
Every off-diagonal entry now vanishes at infinitely
many $b^n$ and is zero rationally.
Consequently
\[
 F(v)=S\operatorname{diag}(v^{m_1},\ldots,v^{m_N})S^{-1}.
\]
Its converse is immediate. Pure weight zero forces
all $m_j=0$ and hence $F(v)=I$.
This final classification requires equality in a
fixed rational trivialization; spectral equality,
parameter-dependent conjugacy or independently chosen
isomorphisms do not imply that stronger hypothesis.
\end{proof}

\paragraph{What has been closed, and what still requires a comparison.}
There are exact finite pure coefficient tensor families
carrying the actual fixed-parameter actions and the IV
rigid dual maps. Thus absence of any pure tensor realization
is not the remaining issue.
Full recovery from finite point-Weil centralizers is a
stronger requirement and fails as above.
Rational all-extension spectral retuning with one fixed
pure weight already conflicts with the actual generic
full centralizer; allowing mixed weights still cannot
retain the full collision centralizer in a regular
fixed-rank family.
These statements allow subalgebras of larger centralizers,
additional comodule or geometric structure, nonrational
comparison data, and separately chosen word objects.
They do not supply the \fid{F24} opposite-sign output,
the required grading comparison, or VIII's prescribed
global integral and infinity lattices and basis selection.
The results concern finite matrix spectra with genuine
multiplicities, and do not bound eigenvalues invisible
in a virtual trace.


\subsection{A projector that cannot descend to the designated global form}
The determinant of one adapted basis does not by itself rule out
an alternative basis of the original tensor lattice. An actual
commutant projector gives a stronger test.

\begin{finding}[\fid{F49}: an intrinsic global-form obstruction]\label{f:globallattice}
Put
\[
 R=\Q[v,v^{-1}],\quad F=\Q(v),\quad
 \delta=v+v^{-1},\quad D=\delta^2,\quad h=D-1,
\]
and let
\[
 M_R=(R^2\otimes_R R^2)^{\otimes3}
\]
be the original tensor lattice with its actual nonstandard action.
There is a projector $p\in\End_F(M_F)$ which commutes with both
named right generators and all the named ordinary left operators,
is a sum of the diagonal matrix units in \fid{F43}, and is induced
by a literal coherent direct-summand projector in \fid{F46}, but
\[
 p(M_R)\not\subseteq M_R.
\]
Consequently no basis of this same lattice can make all the specified
\fid{F43} commutant matrix units integral. Algebraic coefficient
extensions with the prescribed integral closure do not remove this
obstruction.

Under the generic nonstandard Schur--Weyl identification, $p$ lies
in the Schur algebra but not in VIII's specified dual global form.
It is nevertheless compatible with the actual ordinary tensor bar
and preserves the two explicit monomial-coordinate lattices at
$v=0$ and $v=\infty$. No identification of those two lattices with
VIII's conjectural dual local crystal lattice is asserted.
\end{finding}

\begin{proof}
\par\smallskip\noindent\textbf{The actual ordinary projector.}\quad
Regroup $M_F$ as $V_F^{\otimes3}\otimes W_F^{\otimes3}$.
On the $V$-weight space with ordered monomials
$(m_{112},m_{121},m_{211})$, the ordinary upper canonical basis is
\[
 c_{112}=m_{112},\qquad
 c_{121}=m_{121}-v^{-1}m_{112},\qquad
 c_{211}=m_{211}-v^{-1}m_{121};
\]
see IV, Figure~1, p.~9, and Theorem~6.4, p.~34 \cite{IV}.
This change is unimodular over $R$.
The actual decomposition of \fid{F43} uses
\[
 b_1=v^2m_{112}+v m_{121}+m_{211},
\]
\[
 u_1=-(1+v^{-2})m_{112}+v m_{121}+m_{211},\qquad
 u_2=-v^{-1}m_{112}-v^{-2}m_{121}+\delta m_{211}.
\]
Their right-row coordinate matrix in the canonical basis is
\[
 \begin{pmatrix}
 h&\delta&1\\
 0&\delta&1\\
 0&1&\delta
 \end{pmatrix}.
\]
In particular $b_1=h c_{112}+u_1$.
The other three-dimensional weight space has the same rows,
with the last two interchanged; its first vector is
$b_2=v^2m_{122}+v m_{212}+m_{221}$.
Together with $b_0=m_{111}$ and $b_3=m_{222}$ these give
\[
 V_F^{\otimes3}
 =(W_3\otimes\epsilon)\oplus(W_1\otimes U).
\]
The actual coproduct calculation gives
\[
 Fb_0=b_1,\quad Fb_1=\delta b_2,\quad Fb_2=h b_3,\quad Fb_3=0,
\]
\[
 Eb_0=0,\quad Eb_1=h b_0,\quad Eb_2=\delta b_1,\quad Eb_3=b_2.
\]
The two $W_1$ copies have $Fg_+=g_-$ and $Eg_-=g_+$,
with the other actions zero. The ordinary right Hecke action
is scalar $\delta$ on $\epsilon$, and on $U$ is
\[
 c_1=\begin{pmatrix}\delta&0\\1&0\end{pmatrix},
 \qquad
 c_2=\begin{pmatrix}0&1\\0&\delta\end{pmatrix}.
\]
Thus these are actual invariant summands.

Let $p_V$ project onto $W_3\otimes\epsilon$ along $W_1\otimes U$.
Its two weight-space matrices, in canonical and monomial
coordinates respectively, are
\[
 [p_V]_{\mathrm{can}}=
 \begin{pmatrix}
 1&\delta/h&1/h\\
 0&0&0\\
 0&0&0
 \end{pmatrix},
 \qquad
 [p_V]_{\mathrm{mon}}=
 \frac1h
 \begin{pmatrix}
 v^2&v&1\\
 v&1&v^{-1}\\
 1&v^{-1}&v^{-2}
 \end{pmatrix}.
\]
The first follows from the preceding coordinate matrix.
For the second put $u=(v,1,v^{-1})$.
Its numerator is $u^{\mathsf t}u$, and $u u^{\mathsf t}=h$,
so the matrix is idempotent. It fixes $b_1=v u$ and kills
$u_1,u_2$. The lower weight space has the same matrix;
on each endpoint weight space $p_V$ is one.

In particular
\[
 p_V(m_{112})
 =c_{112}+\frac{\delta}{h}c_{121}+\frac1h c_{211}
 =\frac{v^2m_{112}+v m_{121}+m_{211}}h.
\]
Since $h=v^2+1+v^{-2}$ is a nonunit of $R$, the coefficient
$1/h$ is not in $R$. Hence $p_V$ does not preserve the
original lattice.

Now put $p=p_V\otimes1$ on $M_F$. It commutes with
\[
 P_i=c_i^V\otimes c_i^W+
       (c_i^V-\delta)\otimes(c_i^W-\delta)
\]
because $p_V$ commutes with $c_i^V$.
It also commutes with the ordinary left $E,F,K^{\pm1}$
operators on both factors. The coefficient of
$m_{211}\otimes m_{111}$ in
$p(m_{112}\otimes m_{111})$ is again $1/h$, proving
$p\notin\End_R(M_R)$.
If any $R$-basis made its matrix integral, it would preserve
$M_R$, a contradiction. This is independent of the determinant
of any displayed change of basis.

If $1/h$ were integral over $R$, multiply a monic integral
equation by the appropriate power of $h$ to obtain $1\in hR$.
This is impossible. Therefore $1/h$ does not lie in the
integral closure of $R$ in any algebraic extension of $F$.
The proof is unchanged after adjoining $i$.
If a parameter convention first adjoins $v$ as a square root,
integrality over the smaller Laurent ring would imply
integrality over $R$, so it cannot repair this coefficient.

\par\smallskip\noindent\textbf{The specified commutant and coherent projector.}\quad
The multiplicity decomposition in \fid{F43} is
\[
 \begin{aligned}
 M={}&(W_3\otimes W_3\oplus W_1\otimes W_1)\otimes\epsilon_+\\
 &\oplus(W_1\otimes W_1)\otimes\epsilon_-\\
 &\oplus(W_3\otimes W_1\oplus W_1\otimes W_3)\otimes N(\delta)\\
 &\oplus(W_1\otimes W_1)\otimes N(D-2).
 \end{aligned}
\]
The multiplicity ranks are $20,4,16,4$.
The operator $p$ selects the first sixteen multiplicities in
the $\epsilon_+$ block and the first eight in the $N(\delta)$
block, and kills the others. It is therefore a sum of diagonal
matrix units of the full commutant
$M_{20}\times M_4\times M_{16}\times M_4$.
If all those units preserved $M_R$, so would $p$.

On \fid{F46}'s coefficient base $S$, retain its framed object
\[
 K=S_{-1}^{20}\oplus S_i^4\oplus L_1^{16}\oplus L_2^4.
\]
Its direct-summand projector
\[
 g=\operatorname{diag}(I_{16},0_4)\oplus0_4
   \oplus\operatorname{diag}(I_8,0_8)\oplus0_4
\]
is an actual endomorphism of both the coefficient module and
its corresponding finite-support coherent sheaf on $\mathbf P^1_S$.
It commutes with the coordinate $T$ and the chosen Frobenius.
Under the proved postcomposition isomorphism
\[
 \End_F(K)\simeq
 \End_A\operatorname{Hom}_F(H,K)
\]
it induces $p$, with the multiplicity framing above.
The branch comparison changes the two-dimensional right factor
only, and hence leaves these multiplicity selectors unchanged.
The obstruction concerns descent of this particular actual
projector; the coherent realization itself has the projector.

\par\smallskip\noindent\textbf{Bar and the two boundary tests.}\quad
The canonical vectors above are fixed by the actual tensor bar.
All coefficients in their displayed adapted coordinate matrix
are bar invariant. The endpoints and lower weight analogue
have the same property. Therefore both summands defining $p_V$
are bar stable and $p_V$, hence $p$, commutes with the actual bar.
This does not require the tensor monomials to be bar fixed.

For the framed coherent comparison one may send $v$ to $v^{-1}$
and fix $i,x,T$ in \fid{F46}'s coordinate presentations.
The coefficients $\delta,D,h,\Delta,s,\lambda$ are invariant.
The comparison of a right block $N(a)$, $a=\delta,D-2$, with
the corresponding branch presentation has invariant matrix
\[
 S_a=\begin{pmatrix}1&0\\-(a/D)^2&a/D\end{pmatrix},
\]
which is invertible on that coefficient base.
With the selected multiplicity frames this gives the compatible
bar on the Hom realization, and $g$ is fixed.
It is a framed comparison, not a unique global-basis selection.

Let $A_0,A_\infty$ be the rings of rational functions regular
at zero and infinity. The nonzero entries of the monomial
projector are $v^j/h$ with $-2\leq j\leq2$.
Since
\[
 \frac{v^j}{h}=\frac{v^{j+2}}{v^4+v^2+1},
\]
their orders at zero and infinity are respectively $j+2$
and $2-j$, both nonnegative. Thus $p$ preserves both explicit
monomial-coordinate lattices $M_{A_0}$ and $M_{A_\infty}$.
The canonical-coordinate entries $1,\delta/h,1/h$ are also
regular at both boundaries.
This does not assert membership in VIII's conjectural dual
local crystal lattice. It proves that these two concrete
boundary tests and actual bar compatibility do not imply
the global Laurent condition: the latter fails at finite
nonzero zeros of $h$.

\par\smallskip\noindent\textbf{VIII's precise global form.}\quad
Write $\mathcal S$ for the nonstandard coordinate algebra.
VIII defines $\mathcal S_{\Q}$ as the $K_{\Q}$-subring generated
by the generic matrix entries $u^i_j$, with $K_{\Q}$ the
integral closure of the Laurent parameter ring.
Its degree-three Schur algebra is
$\mathcal A=(\mathcal S_3)^*$ and its specified dual global form is
\[
 \mathcal A_{\Q}
 =\{a\in\mathcal A:
       \langle a,\mathcal S_{3,\Q}\rangle\subseteq K_{\Q}\}.
\]
These are VIII, \S2.2, pp.~7, 9--10, and \S2.3, p.~11
\cite{VIII}.
The global-form part of its Proposition~2.4, pp.~12--13,
has the following direct proof in this situation.
Write the tensor coaction in the original basis as
\[
 e_J\longmapsto\sum_I e_I\otimes u_{I,J}.
\]
Each $u_{I,J}$ is a degree-three product of generating
coordinate entries and belongs to $\mathcal S_{3,\Q}$.
The matrix entries of the induced operator $a$ are
$\langle a,u_{I,J}\rangle$, hence belong to $K_{\Q}$ whenever
$a\in\mathcal A_{\Q}$. Consequently
\[
 \eta(\mathcal A_{\Q})
 \subseteq\End_{K_{\Q}}(X_{K_{\Q}}^{\otimes3}).
\]
For $X=V\otimes W$ under the direct-product ordinary quantum
group, its chosen upper canonical basis is the four elementary
products. This is the same original tensor form after scalar
extension.

Generic nonstandard Schur--Weyl duality identifies the full
commutant with $\eta(\mathcal A)$, so $p$ belongs to that algebra.
The remark immediately after IV Theorem~12.1 corrects its
tensor-dual convention relative to \S\S8--10.
Transposing that identification preserves the global entry
criterion; our monomial projector matrices are symmetric,
so the present $p$ is unchanged by that transpose.
Its coefficient $1/h$ and the preceding integral-closure
argument prove
\[
 p\notin\eta(\mathcal A_{\Q}).
\]
The localized matrix-unit basis therefore cannot, under this
actual action, be identified with a basis of the designated
global form; not all its elements even belong to that form.
The global inclusion just used has been proved directly.
The printed proof of VIII Proposition~2.4 is a placeholder;
its additional local-lattice inclusion is conditional on
Conjecture~2.2 and has not been assumed here.
\end{proof}

\paragraph{The exact boundary of this obstruction.}
IV already constructs a based rank-three example in
Example~19.2, pp.~132--133. Its paragraph after Conjecture~19.1,
p.~131, explicitly distinguishes the desired integral form and
local lattice from the original tensor forms \cite{IV}.
Our obstruction is compatible with that construction.
It concerns simultaneous retention of the original global tensor
lattice and the full specified \fid{F43}/\fid{F46} commutant
projectors. It does not obstruct an alternate basis with only
the named nonstandard generators, a changed integral form, a
smaller integral commutant, or VIII's conjectured minimum-degree
globalization within its designated form.


\subsection{An exact Gaussian construction for the path-sign comparison}
The graded coefficient objects of \fid{F43} and \fid{F45} admit an
explicit exponential-sum realization with the required opposite sign.
The same construction embeds point-Weil End/Hom data fully faithfully
in perverse sheaves. Its limitation is also exact: for this construction,
the pointwise estimate retains the entire coefficient-weight question.

\paragraph{Classical input and conventions.}
Let $k_0=\F_q$ have odd characteristic $p$, let $\ell\ne p$, and use a
characteristic-zero coefficient field $E$ containing the character
values. Let $\psi$ be a nontrivial additive character of $k_0$ and
$\psi_n=\psi\circ\Tr_{\F_{q^n}/k_0}$. Frobenius is geometric.
The classical inputs are the Artin--Schreier/Kummer trace construction
and its all-extension formula \cite[\emph{Sommes trig.},
1.7--1.9, especially (1.9.2)]{SGA45}, the rank-one Gaussian
cohomology in degree one \cite[Propositions~4.2--4.3,
pp.~196--197]{SGA45}, and its perfect pairing into $E(-1)$
\cite[Remark~4.4, p.~197]{SGA45}.
The original printed pages, not merely a modern transcription, were
inspected. We use the correct pairing/product identity in Remark~4.4,
not its misprinted absolute-value shorthand. No Gaussian size estimate
or Weil weight theorem is an input to the sign calculation below.

\begin{finding}[\fid{F50}: Gaussian path signs and the exact weight gate]
\label{f:gaussiansign}
There is a one-dimensional Gaussian Weil space $G_\psi$ such that
\[
 R\Gamma_c(\A^1_{\bar k_0},\mathcal L_\psi(x^2))=G_\psi[-1].
\]
For every effective finite-dimensional Weil space $C$ with invertible
Frobenius and algebraic eigenvalues, put
\[
 J_\psi(C)=\mathcal L_\psi(x^2)\otimes
                p^*(C\otimes G_\psi^\vee)[1],
 \qquad p:\A^1\longrightarrow\Spec k_0.
\]
Then $R\Gamma_c(\A^1,J_\psi(C))\simeq C$ naturally. Its trace
function reproduces $\Tr(F^n\mid C)$ exactly after summing over
$\A^1(\F_{q^n})$.

For an unbounded set of even path lengths $k$, its path
hypersurface complex is concentrated in degree $-1$. Its twisted
path sum is the negative trace of an explicit effective finite Weil
space, with zero error in \fid{F04}. Before the normalization in $J$,
the actual Gaussian path cohomology is concentrated in degree $k-1$.

The functor $J_\psi$ is exact and fully faithful into perverse
Weil sheaves in degree zero:
\[
 \Hom(J_\psi(C),J_\psi(C'))=\Hom_F(C,C').
\]
If $C$ is pure of weight $w$, then $J_\psi(C)$ is perverse pure
of weight $w$. Arithmetic Jordan blocks are allowed.

For any specified integer $w$, the pointwise bounds needed to
apply \fid{F03}/\fid{F17} to $J_\psi(C)$ and the reciprocal system
$J_{\psi^{-1}}(C^\vee(-1))$ hold under every embedding if and only if
every conjugate eigenvalue of $C$ has modulus $q^{w/2}$.
Thus the output-sign requirement is proved for this construction;
the two pointwise requirements retain precisely the original
coefficient-purity assertion.
\end{finding}

\begin{proof}
\par\smallskip\noindent\textbf{The Gaussian line without a weight estimate.}\quad
Write $\mathscr L_\psi=\mathcal L_\psi(x^2)$ and
$G_\psi=H_c^1(\A^1_{\bar k_0},\mathscr L_\psi)$.
The cited classical propositions give
\begin{equation}\label{eq:gaussline}
 R\Gamma_c(\A^1,\mathscr L_\psi)=G_\psi[-1],
 \qquad
 g_n:=\sum_{x\in\F_{q^n}}\psi_n(x^2)=-\gamma_\psi^n,
\end{equation}
where $\gamma_\psi$ is the Frobenius eigenvalue on the line.
They also give a perfect pairing
\begin{equation}\label{eq:gausspair}
 G_\psi\otimes G_{\psi^{-1}}\longrightarrow E(-1),
 \qquad \gamma_\psi\gamma_{\psi^{-1}}=q.
\end{equation}

The precise passage from the cited $\mathbb G_m$ calculation to this
$\A^1$ version is elementary. For the finite square map
$\pi:\A^1\to\A^1$, with $j:\mathbb G_m\hookrightarrow\A^1$,
\[
 \pi_*E=E\oplus j_!\mathcal L_\chi,
\]
where $\chi$ is the nontrivial quadratic character. The
anti-invariant summand has zero stalk at zero. The compact
cohomology of the nontrivial linear Artin--Schreier sheaf
$\mathcal L_\psi(x)$ is zero: it is a nontrivial character
summand of the Lang cover $\A^1\to\A^1$, whose source has
compact cohomology $E(-1)[-2]$ with trivial translation action.
Projection formula therefore identifies
$R\Gamma_c(\A^1,\mathscr L_\psi)$ with
$R\Gamma_c(\mathbb G_m,\mathcal L_\psi(x)\otimes\mathcal L_\chi)$.
Proposition~4.2 and Remark~4.4 apply.

The involution $x\mapsto-x$ acts by $-1$ on $G_\psi$,
since this is precisely the anti-invariant summand. For a
varying nonzero scalar $a$, the Gaussian line of
$\mathcal L_\psi(ax^2)$ is consequently $G_\psi$ tensored with
the quadratic Kummer line of $a$: use the etale change of
variable $x\mapsto\sqrt a\,x$ and its descent involution.
Completing the square also gives the exact Fourier identity
\[
 \sum_x\psi_n(x^2+xy)=g_n\psi_n(-y^2/4).
\]
This is a character-sum calculation, not positivity of the
Fourier operator.

\par\smallskip\noindent\textbf{Good even path lengths.}\quad
Let $A_k(h)$ have diagonal entries $1$, adjacent off-diagonal
entries $h/2$, and other entries zero. Then
\[
 x^{\mathsf T}A_k(h)x=\sum_i x_i^2+hP_k(x),\qquad
 P_k(x)=\sum_{i=1}^{k-1}x_ix_{i+1}.
\]
Set $d_k(h)=\det A_k(h)$, with $d_0=d_1=1$. Expansion gives
\begin{equation}\label{eq:gaussdetrec}
 d_k=d_{k-1}-\frac{h^2}{4}d_{k-2}.
\end{equation}
A root $h_0$ of $d_k$ is nonzero because $d_k(0)=1$.
At such a root $d_{k-1}(h_0)\ne0$: otherwise the recurrence
would make every preceding $d_j(h_0)$ zero, contradicting
$d_0=1$. Thus the singular matrix has rank exactly $k-1$
and its leading $(k-1)$-dimensional block is nonsingular.

For even $k=2m$, $d_k$ has degree $k$ and leading coefficient
\[
 \ell_k=(-1)^m/2^k.
\]
It is squarefree whenever $p\nmid k+1$.
Here is a proof retaining the characteristic hypothesis.
Define $U_0=1$, $U_1=2T$ and
$U_k=2TU_{k-1}-U_{k-2}$. The recurrence proves
\[
 d_k(h)=(h/2)^kU_k(1/h),\qquad
 U_k((z+z^{-1})/2)
   =\frac{z^{k+1}-z^{-(k+1)}}{z-z^{-1}}.
\]
If $p\nmid k+1$, the $2(k+1)$-st roots of unity are distinct.
Excluding $z=\pm1$ and identifying inverse pairs gives $k$
distinct roots of $U_k$; two pairs have the same sum
$z+z^{-1}$ only when they solve the same quadratic.
For even $k$ none has sum zero, because this would give
a fourth root of unity of order four while $k+1$ is odd.
Inverting the resulting nonzero roots proves squarefreeness.
Finally $U_{2m}(0)=(-1)^m$ gives the leading coefficient.

Call an even $k$ \emph{good} when $d_k$ is squarefree.
The sufficient condition supplies the unbounded good set
$k=2pr$, $r\ge1$. The condition is sufficient, not necessary:
$k=2$ is good in every odd characteristic.

\par\smallskip\noindent\textbf{Actual middle cohomology.}\quad
For a good even $k$, put $H_k=\{P_k=0\}\subset\A^k$.
We claim
\begin{equation}\label{eq:gaussmiddle}
 H_c^i(H_{k,\bar k_0},
       \mathcal L_\psi(\textstyle\sum_i x_i^2)|_{H_k})=0
       \quad\text{unless }i=k-1.
\end{equation}
This includes the singular vertex; it is not deduced from
the dimension of the hypersurface.

Consider $f:\A^1_h\times\A^k_x\to\A^1_h$ and
\[
 \mathscr Q=\mathcal L_\psi(\textstyle\sum_i x_i^2+hP_k(x)).
\]
Let $Z_k=\{d_k=0\}$, $U_k=\A^1\setminus Z_k$, and
write $j:U_k\hookrightarrow\A^1$, $i:Z_k\hookrightarrow\A^1$.
The finite scheme $Z_k$ is etale. Over $U_k$, etale-local
diagonalization and square roots of the unit coefficients,
together with \eqref{eq:gaussline} and Kunneth, show that the
only fiber compact cohomology is rank one in degree $k$.
Thus the generic system $T_k$ is lisse of rank one.

Near a root $h_0$, Schur complement splits off the nonsingular
leading block, leaving the scalar $d_k(h)/d_{k-1}(h)$.
At the root its free affine-line direction contributes
$E(-1)[-2]$. Its only fiber cohomology is therefore rank
one in degree $(k-1)+2=k+1$. Compact-support base change gives
\[
 R^kf_!\mathscr Q=j_!T_k,\qquad
 R^{k+1}f_!\mathscr Q=i_*V_k,
\]
with all other direct-image sheaves zero and $V_k$ rank one
on $Z_k$. The first equality follows from its zero stalks
at $Z_k$. The compact-support formalism used here follows
from proper base change and compactification
\cite[Theorem~17.10, p.~116; Proposition~18.4,
pp.~121--122]{Milne}; Kunneth and projection formula are
\cite[Theorem~22.1, Remarks~22.3(b),(d), and
Proposition~22.5, pp.~135--137]{Milne}.

At every simple root, the Schur-complement scalar has order
one. Its Gaussian factor has quadratic Kummer monodromy
$-1$, while the nonsingular block is locally unramified.
Hence $T_k$ is geometrically nonconstant. On the connected
nonproper curve $U_k$,
\[
 H_c^0(U_k,T_k)=H_c^2(U_k,T_k)=0.
\]
The first assertion concerns proper-support sections of a
lisse sheaf; the second follows by duality from the absence
of geometric invariants of the dual line
\cite[\emph{Sommes trig.}, 1.18(b)--(d),
pp.~178--179]{SGA45}. No weight bound enters.
The compact-support Leray spectral sequence has only the
terms $(p,q)=(1,k)$ and $(0,k+1)$, both in total degree
$k+1$. Thus $R\Gamma_c(\A^1\times\A^k,\mathscr Q)$ is
concentrated in degree $k+1$.

Integrating $h$ first gives zero when $P_k(x)\ne0$, by the
linear Artin--Schreier calculation, and $E(-1)[-2]$ when
$P_k(x)=0$. Base change and projection formula give
\[
 R\Gamma_c(\A^1\times\A^k,\mathscr Q)
 \simeq
 R\Gamma_c(H_k,\mathcal L_\psi(\textstyle\sum_i x_i^2))(-1)[-2].
\]
This proves \eqref{eq:gaussmiddle}.

\par\smallskip\noindent\textbf{An explicit effective output spectrum.}\quad
Let $C_k$ be the smooth projective completion of
$y^2=d_k(h)$. It is geometrically connected.
The double cover has $k$ simple branch points and no
branch at infinity, so the tame Riemann--Hurwitz formula
\cite[Lemmas~53.12.2 and~53.12.4, tags~0C1D and~0C1F]{StacksRH}
gives genus $(k-2)/2$. The dimension formula
$\dim H^1(C_k,E)=2g(C_k)$ is
\cite[Proposition~14.2, p.~90]{Milne}. Let $\eta_k$ be the quadratic Weil line whose
extension trace is $\chi_n(\ell_k)$; it is the anti-invariant
line of the two points at infinity. Put
\[
 W_k=H^0(Z_{k,\bar k_0},
             \mathcal L_\chi(d_{k-1}|_{Z_k})).
\]
This is well defined by \eqref{eq:gaussdetrec} and has
dimension $k$. It retains the actual Frobenius permutation
of the roots, which need not be rational over $k_0$.
Define the effective Weil space
\begin{align}
 B_{k,\psi}={}&G_\psi^{\otimes k}\otimes
                 (H^1(C_k,E)\oplus\eta_k)\notag\\
 &\oplus E(-1)\otimes G_\psi^{\otimes(k-1)}\otimes W_k .
 \label{eq:gaussoutput}
\end{align}
Its dimension is $2k-1$ and its eigenvalues have genuine
positive multiplicities.

For $Q=q^n$ and
$N_{k,n}^\psi=\sum_{P_k=0}\psi_n(\sum_i x_i^2)$,
orthogonality and diagonalization give
\begin{align}
 QN_{k,n}^\psi
 ={}&g_n^k\sum_{h\in\F_Q}\chi_n(d_k(h))\notag\\
 &+Qg_n^{k-1}
       \sum_{\substack{h\in\F_Q\\d_k(h)=0}}
                         \chi_n(d_{k-1}(h)).
 \label{eq:gausspenciltrace}
\end{align}
Indeed a nonsingular quadratic form contributes its
determinant character times $g_n^k$; at a root Schur
complement gives one free affine-line factor and the
displayed nonsingular minor. The projective curve trace
formula gives, without a curve weight estimate,
\[
 \sum_h\chi_n(d_k(h))
 =-\Tr(F^n\mid H^1(C_k))-\chi_n(\ell_k).
\]
Its affine point count is $Q$ plus the sum on the left;
its infinity fiber has $1+\chi_n(\ell_k)$ points, while
$H^0=E$ and $H^2=E(-1)$.
Since $k$ is even, $g_n^k=\gamma_\psi^{kn}$ and
$g_n^{k-1}=-\gamma_\psi^{(k-1)n}$. Therefore
\begin{equation}\label{eq:gaussnegativesign}
 QN_{k,n}^\psi=-\Tr(F^n\mid B_{k,\psi})
                      \qquad(n\ge1).
\end{equation}
This identifies the character of the twisted middle
cohomology with \eqref{eq:gaussoutput}; no canonical
splitting or arithmetic semisimplicity is asserted.

The algebraicity and embedding requirements are explicit.
The traces of $H^1(C_k)$ are the rational integers
$1+q^n-\#C_k(\F_{q^n})$. Newton identities in its finite
dimension give a rational characteristic polynomial, hence
an algebraic embedding-stable spectrum. Frobenius on $W_k$
is a signed permutation: each cycle has characteristic polynomial
$T^d-\epsilon$, with $\epsilon=\pm1$. It therefore has a rational
characteristic polynomial and finite order. Also $\eta_k$ is quadratic, and
$\gamma_\psi=-g_1$ is cyclotomic. Conjugating
\eqref{eq:gaussoutput} therefore gives exactly the
construction for the conjugated character $\psi$.
No Weil integrality or weight theorem is needed for this step.

At $k=2$, the union of the two axes gives
$N_{2,n}^\psi=2g_n-1$, so
\[
 QN_{2,n}^\psi=-2Q\gamma_\psi^n-Q.
\]
This also checks every sign in \eqref{eq:gaussoutput};
the relation $\gamma_\psi^2\chi(-1)=q$ follows from
\eqref{eq:gausspair} and quadratic scaling.
The good-length condition matters. In characteristic five,
$k=4$ has $d_4=(h^2-1)^2$ and $d_3(\pm1)=3$.
For $q=5$, \eqref{eq:gausspenciltrace} gives
\[
 QN_{4,n}^\psi=5^{3n}-2\cdot5^{2n}
                         -2(-5\gamma_\psi^3)^n.
\]
Here $\gamma_\psi^2=5$. The root $125$ has positive
coefficient and strictly largest modulus, so this bad
length cannot have an exact negative effective spectrum.
The proof supplies an unbounded good even set, not all
even lengths in every characteristic.

\par\smallskip\noindent\textbf{The size follows from the sign theorem.}\quad
For $t_n(x)=\psi_n(x^2)$ the pointwise bound is $1$.
Thus $a=0,d=1$ in \fid{F03}.
Equation~\eqref{eq:gaussline} is its negative rank-one
input spectrum, and \eqref{eq:gaussnegativesign} supplies
the opposite output on the unbounded good set, with zero
error. Finding~\fid{F03} gives $|\gamma_\psi|\le q^{1/2}$.
Apply the same argument to $\psi^{-1}$ and then
\eqref{eq:gausspair}: both absolute values equal $q^{1/2}$.
Every embedding has its corresponding nontrivial character
and output spectrum, so this holds for every complex
absolute value. This recovers the classical Gaussian
result as an application test of the sign criterion.

\par\smallskip\noindent\textbf{Normalization and the exact target trace.}\quad
For an effective finite Weil space $C$, projection formula
and evaluation $G_\psi\otimes G_\psi^\vee\to E$ give
\[
 R\Gamma_c(\A^1,J_\psi(C))\simeq C.
\]
The shifts $[-1]$ and $[1]$ cancel. Its trace function is
\begin{equation}\label{eq:gaussnormalizedtrace}
 t_{C,n}(x)=-\gamma_\psi^{-n}
            \Tr(F^n\mid C)\psi_n(x^2),
 \qquad M_{C,n}=\Tr(F^n\mid C).
\end{equation}
For a good even $k$,
\begin{align*}
 R\Gamma_c(H_k,J_\psi(C)^{\boxtimes k}|_{H_k})
 \simeq{}&C^{\otimes k}\otimes G_\psi^{\vee\otimes k}\\
 &\otimes
 R\Gamma_c(H_k,\mathcal L_\psi(\textstyle\sum_i x_i^2))[k].
\end{align*}
Its only degree is $(k-1)-k=-1$. Thus its effective
twisted output is
\[
 B_{k,C}=C^{\otimes k}\otimes
                    G_\psi^{\vee\otimes k}\otimes B_{k,\psi},
 \qquad
 QN_{k,C,n}=-\Tr(F^n\mid B_{k,C}).
\]
This proves the exact opposite sign with no boundary error.
It concerns effective coefficient spaces, including the
even-degree total realization of \fid{F45}; it does not
declare tensor powers of arbitrary signed virtual complexes
to be effective.

\par\smallskip\noindent\textbf{The pointwise and dual gates.}\quad
For nonzero $C$ and a fixed embedding put
\[
 b_{\max}=2\log_q\max_\alpha|\alpha|,\qquad
 b_{\min}=2\log_q\min_\alpha|\alpha|,
 \qquad \alpha\in\operatorname{Spec}(F\mid C).
\]
These are real thresholds, not yet asserted integer weights.
The Gaussian result just proved gives
\[
 \sup_x|t_{C,n}(x)|
       =q^{-n/2}|\Tr(F^n\mid C)|.
\]
There is a constant $B$ independent of $n$ with
$|t_{C,n}(x)|\le Bq^{na/2}$ if and only if
\begin{equation}\label{eq:gausspointwisegate}
 \max_\alpha|\alpha|\le q^{(a+1)/2}.
\end{equation}
Sufficiency uses $B=\dim C$; necessity is \fid{F02}.
Jordan blocks do not introduce polynomial factors into
traces of powers. The sharp real exponent is therefore
$b_{\max}-1$.

The reciprocal system is
$J_{\psi^{-1}}(C^\vee(-1))$. Its total spectrum is
$\{q/\alpha\}$, with the required perfect pairing into
$E(-1)$, and its sharp real pointwise exponent is
$1-b_{\min}$. Relative to the formally matching exponent
$1-b_{\max}$, the exact gap is $b_{\max}-b_{\min}$.

To invoke \fid{F03}/\fid{F17} as actually stated, fix one
integer target weight $w$ and set $a=w-1$.
The two pointwise bounds at exponents $w-1$ and $1-w$
hold under every embedding precisely when
\begin{equation}\label{eq:gausscommonweight}
 b_{\max,\sigma}\le w\le b_{\min,\sigma}
                    \quad\text{for every }\sigma.
\end{equation}
Equivalently all conjugate eigenvalues have modulus
$q^{w/2}$. Width zero separately in every embedding is
insufficient if the radii differ between embeddings:
the same integer $w$ must work throughout.

\par\smallskip\noindent\textbf{Hom and conditional perverse purity.}\quad
For $C,C'$, the equal shifts and the same invertible
rank-one sheaf and Gaussian line cancel from
$\Hom(J_\psi(C),J_\psi(C'))$. A map between constant
sheaves on connected geometric $\A^1$ is exactly a
linear map of their fibers. Weil compatibility is exactly
Frobenius commutation. Therefore, naturally and
compatibly with composition,
\[
 \Hom(J_\psi(C),J_\psi(C'))=\Hom_F(C,C').
\]
This is both perverse-heart Hom and degree-zero derived
Weil Hom; no higher-Ext comparison is claimed.
The construction is exact. If $C$ is pure of weight $w$,
its unshifted lisse coefficient has weight $w-1$ and its
shift by $[1]$ has complex weight $w$ in BBD's convention
\cite[5.1.8, pp.~126--127]{BBD}.
It is perverse because the base is a smooth curve.
Arithmetic Jordan blocks remain allowed; geometrically
the constituents are copies of one rank-one local system.
\end{proof}

\paragraph{Two precise modes of retaining the actual actions.}
At a fixed arithmetic specialization of \fid{F46}, including
the collision parameter, its actual $H,K$ give
\[
 \End(JH)=A,\qquad \Hom(JH,JK)=M,\qquad
 \End(JK)=\End_A(M).
\]
Precomposition and postcomposition retain the specified
right generators and ordinary left operators, including
the arithmetic collision nilpotent. With \fid{F46}'s
chosen pure weight-zero Weil data, $JH$ and $JK$ are
actual pure weight-zero perverse Weil sheaves.
Their Frobenius powers are those of the fixed operators;
they are not extension-by-extension quantum-parameter
substitutions.

The sign statement applies to identity traces of effective
objects. An arbitrary inserted trace $\Tr(TF^n\mid C)$ is
different: its Frobenius coefficients need not be effective
multiplicities, and its tensor-inserted output need not have
the required sign. Retaining $T$ as a morphism does not
establish that extra assertion.

For \fid{F43}'s graded model, take a point-Weil line $U$
with eigenvalue $q^{1/2}$ and replace every degree-$e$
edge by $U^e$. Each named kernel then has trace matrix
equal to its actual Laurent matrix evaluated at
$v=q^{n/2}$. Odd degrees give arithmetic Weil lines of
odd integer weight; they are not asserted to be integral Tate
motives. Even degrees use the usual Tate lines.
For any constant coefficient space $V$,
\[
 J_\psi(C\otimes V)\simeq p^*V\otimes J_\psi(C).
\]
Thus the kernel functors and their explicitly ordered
word compositions transport through $J$. All the actual
graded $K_0$ matrices of $P_1,P_2$ and the stated left
operators remain those of \fid{F43}, with its opposite
row-composition convention for the left action.
This does not choose a coherent lift of every algebra
relation or commutation beyond that established there.
The word coefficient objects change with the word.

For the ordinary tensor product there is no strong monoidal claim:
$J(E)$ is not the tensor unit, and
\[
 J(C)\otimes J(C')
 =\mathcal L_{\psi^2}(x^2)\otimes
     p^*(C\otimes C'\otimes G_\psi^{\vee\otimes2})[2],
\]
whereas $J(C\otimes C')$ has one Gaussian factor and
shift $[1]$. Nor have distinct geometric simple labels
or a designated canonical basis been constructed.
These exact scope restrictions agree with \fid{F42}
and \fid{F48}.

\paragraph{The actual word objects and their dual gap.}
For a nonempty right word of length $L$ and cyclic
transition count $2s$, write
\[
 T_w(v)=20D^L+16D^{L-s}
          +4D^{L-2s}(D-2)^{2s}
       =\sum_{r=-L}^{L}c_rv^{2r},\qquad c_r\ge0.
\]
The total effective Weil space of \fid{F45} is
$C_w=\bigoplus_r E(-r)^{c_r}$.
Equation~\eqref{eq:gaussnormalizedtrace} reproduces
$T_w(q^{n/2})$ exactly, and its path output has the
proved negative effective spectrum with zero error.
Both extreme coefficients are nonzero, already from
$20D^L$. Hence $b_{\max}=2L$, $b_{\min}=-2L$.
The sharp pointwise exponents are $2L-1$ for the original
system and $2L+1$ for its reciprocal system. At target
weight $2L$, \fid{F17} would need the latter exponent
$1-2L$: the exact gap is $4L$.

For each homogeneous piece $E(-r)^{c_r}$ the two exponents
are instead $2r-1$ and $1-2r$. The criterion recovers
exactly its known pure weight $2r$. The analogous statement
holds for every homogeneous path coefficient of \fid{F43},
using $U^e$ when its degree is odd. Its actual finite path
labels and explicit-word bounds are available from that
finding. No new all-rank positive rule, canonical-basis
identification or binary-input PH1 consequence follows.

\paragraph{The direct positive graph-count route has a different sign.}
Putting the effective coefficient trace $S_n=\Tr(F^n\mid C_w)$
at the origin and zero elsewhere on $\A^1$ gives
\[
 M_n=S_n,\qquad QN_{k,n}=QS_n^k.
\]
This nonzero positive-multiplicity Tate spectrum cannot be
the negative of an effective spectrum for every extension:
their sum would contradict \fid{F02}. This excludes the
exact sign for that point-supported construction. Its
small support still permits loose error estimates with
empty output; no broader \fid{F04} obstruction is inferred
from this example alone.

Affine padding gives a sharper quantitative failure.
For a base $\A^d$ with $d\ge1$, take the constant trace
$t_n(x)=Q^{-d}S_n$, so $M_n=S_n$ and the sharp pointwise
exponent is $a=2L-2d$. For even $k=2m$, separating the
odd and even coordinates makes $P_k$ a nondegenerate
bilinear pairing between two $md$-dimensional spaces:
the block coefficient matrix is triangular with invertible
diagonal. The zero first vector contributes $Q^{md}$
choices, and every nonzero first vector imposes one
nonzero linear equation. Hence
\[
 \#H_k(\F_Q)=Q^{kd-1}+Q^{kd/2}-Q^{kd/2-1},
\]
and
\[
 QN_{k,n}=S_n^k
             \bigl(1+Q^{1-kd/2}-Q^{-kd/2}\bigr).
\]
For $kd>2$ its highest-radius root is $q^{kL}$ with
positive coefficient $c_L^k$; for $kd=2$ that coefficient
is twice as large. An added effective spectrum cannot
cancel it. If this were minus an effective spectrum plus
an \fid{F04} error, \fid{F02} would therefore force
\[
 q^{kL}\le q^{k(a+d)/2+\eta_k}
              =q^{k(L-d/2)+\eta_k},
 \qquad \eta_k\ge kd/2.
\]
The cost is linear in $k$, incompatible with the required
sublinear overhead. Full extension-degree progressions
do not remove the positive top Tate root.
Increasing $a$ by $d$ can absorb the cost, but weakens
the resulting weight estimate by $d$. This is the exact
boundary for the specified normalized affine-padding route.


\subsection{Algebraic coefficients and a trace-only comparison gate}
VIII allows an algebraic extension of the quantum parameter field.
The single-weight obstruction persists there and requires only
ordinary trace compatibility across extensions.

\begin{finding}[\fid{F51}: the algebraic single-weight trace gate]\label{f:algebraictrace}
Fix $b>1$ and let $\mathcal F$ be an invertible finite matrix
whose coefficients lie in an algebraic extension of $\C(v)$.
Suppose it has a regular invertible value $F_b$ at $v=b$,
and for every sufficiently large integer $n$ choose a regular
value $F_{b^n}$ of the same algebraic family.
Assume
\begin{equation}\label{eq:algebraictraceretune}
 \operatorname{Tr}(F_{b^n})=\operatorname{Tr}(F_b^n).
\end{equation}
If $F_b$ is pure of one fixed integral weight $w$ relative
to $Q=b^2$, then every eigenvalue of $b^{-w}F_b$ is one.
If its full Frobenius centralizer is semisimple, that
centralizer is therefore a single full matrix algebra.
In particular it cannot be the actual generic rank-three
algebra
\[
 A_b\simeq\C\times\C\times M_2(\C)\times M_2(\C).
\]
The conclusion allows different branches of the same
algebraic relation at different extension parameters;
it assumes neither rational coefficients nor full
spectral retuning.

More generally, suppose the base eigenvalues have several
integral weights $w_j$, with genuine multiplicities.
The same trace comparison forces the family trace to be
the corresponding positive Laurent polynomial and every
base eigenvalue to be $b^{w_j}$, without a nontrivial phase.
If full spectral retuning is imposed as well, then
\[
 \det(TI-\mathcal F)=\prod_w(T-v^w)^{m_w}.
\]
A regular fixed-rank model at $v=1$ therefore cannot
have the actual collision algebra as its full special
Frobenius centralizer.
\end{finding}

\begin{proof}
We first prove the precise algebraic-function statement.
Let
\[
 u_n=\sum_{j=1}^s c_j\alpha_j^n,\qquad |\alpha_j|=1,
\]
with finitely many roots and arbitrary complex coefficients.
Let $P(v,Y)\in\C[v,Y]$ be irreducible of positive
$Y$-degree, and suppose
\[
 P(b^n,u_n)=0
\]
for every sufficiently large integer $n$.
We claim that
\begin{equation}\label{eq:radialconstant}
 P(v,Y)=a(Y-u_0),\qquad a\ne0,
\end{equation}
and $u_n=u_0$ for every integer $n$.

Write
\[
 P(v,Y)=\sum_{i=0}^d v^iP_i(Y).
\]
Each polynomial $P_i(u_n)$ is a finite power sum on
unit-circle roots: after expanding its powers and grouping
equal products,
\[
 P_i(u_n)=\sum_{\beta\in B_i} a_{i,\beta}\beta^n,
 \qquad |\beta|=1.
\]
Its constant term is represented by the root $\beta=1$.
The assumed relation becomes
\[
 \sum_i\sum_{\beta\in B_i}
       a_{i,\beta}(b^i\beta)^n=0.
\]
Different $i$ give different root moduli $b^i$, so no
roots from those groups coincide.
Within a group equal roots have already been collected.
Evaluating a relation on distinct nonzero power roots
at sufficiently many consecutive integers gives an
invertible Vandermonde matrix times an invertible
diagonal matrix. Thus every $a_{i,\beta}=0$.

Consequently every $P_i(u_n)$ vanishes as an exponential
sequence for all integers, including $n=0$.
Therefore $P_i(u_0)=0$ for every $i$, or
$P(v,u_0)=0$ identically in $v$.
The factor $Y-u_0$ divides $P$ in $\C[v,Y]$.
Irreducibility proves \eqref{eq:radialconstant}.
The eventual equality $u_n=u_0$ then extends to all
integers by the same finite-power-sum uniqueness.
The auxiliary value $n=0$ is an algebraic sequence value,
not an assertion of an extension of degree zero or a
field of cardinality one.

For a genuine $N$-dimensional matrix spectrum,
$u_n=\operatorname{Tr}(G^n)$ and $u_0=N$.
If all eigenvalues of $G$ have modulus one, the preceding
argument gives $u_n=N$ eventually.
A sum of $N$ unit-modulus numbers equals the real number
$N$ only when each summand is one.
Two consecutive sufficiently large $n$ then imply
$\alpha_j^n=\alpha_j^{n+1}=1$, hence $\alpha_j=1$
for every eigenvalue.
This is where genuine positive multiplicities are used.
For a virtual trace the same argument concerns only
its visible grouped roots.

Return to the family and normalize
\[
 \mathcal G=v^{-w}\mathcal F,\qquad G_b=b^{-w}F_b.
\]
The trace $g=\operatorname{Tr}\mathcal G$ is algebraic
over $\C(v)$. Its irreducible minimal polynomial,
after clearing denominators and taking its primitive
part, is an irreducible $P(v,Y)\in\C[v,Y]$.
Every regular evaluation of this same algebraic element
satisfies that relation.
By \eqref{eq:algebraictraceretune},
\[
 g_{b^n}
 =b^{-nw}\operatorname{Tr}(F_b^n)
 =\operatorname{Tr}(G_b^n).
\]
The eigenvalues of $G_b$ all have modulus one.
The preceding argument proves $g(v)=N$ and that every
eigenvalue of $G_b$ is one. Thus $G_b-I$ is nilpotent.

Scalar normalization does not change the full centralizer.
The nilpotent $G_b-I$ is central in that algebra.
If the algebra is semisimple, its center is reduced,
so $G_b-I=0$.
The full centralizer is then $M_N(\C)$.
The actual algebra of \fid{F19} at $b>1$ has its two
center branches split and is the nonsimple semisimple
product displayed in the statement. It cannot be
such a full centralizer.

\par\smallskip\noindent\textbf{Several integral spectral weights.}\quad
We prove the more general algebraic-function assertion
with arbitrary finite complex coefficients first.
Write
\[
 u_n=\sum_j c_j b^{w_jn}\zeta_j^n,\qquad
 w_j\in\Z,\quad |\zeta_j|=1,
 \qquad p(v)=\sum_j c_jv^{w_j}.
\]
Suppose $P(b^n,u_n)=0$ eventually for the irreducible
minimal relation of one algebraic function $f$.
Expand
\[
 P(v,Y)=\sum_{i,k}a_{i,k}v^iY^k.
\]
Every root in its evaluated power-sum expansion has form
\[
 b^{\,i+w_{j_1}+\cdots+w_{j_k}}
       \zeta_{j_1}\cdots\zeta_{j_k}.
\]
Group these by the integer
$r=i+w_{j_1}+\cdots+w_{j_k}$.
For fixed $r$ the coefficient is a finite unit-circle
power sum $S_r(n)$.
Distinct $r$ give distinct root moduli $b^r$, so
Vandermonde uniqueness makes each $S_r(n)$ identically
zero. At $n=0$ all phases disappear.
The value $S_r(0)$ is exactly the coefficient of $v^r$
in $P(v,p(v))$. Hence
\[
 P(v,p(v))=0\quad\hbox{in }\C[v,v^{-1}].
\]
Irreducibility over $\C(v)$ and the rational root $p(v)$
force $f=p$.
This argument is taken over $\C(v)$ when weights are
negative; no polynomial factorization with an omitted
denominator is being assumed.

Thus $u_n=p(b^n)$ eventually and, by finite-power
uniqueness, identically as a sequence.
For genuine positive eigenvalue multiplicities, any
nontrivial phase would give a positive coefficient on
a root absent from the Laurent side. It cannot cancel.
Every base eigenvalue is therefore $b^w$, with its
stated weight multiplicity.

\par\smallskip\noindent\textbf{Full spectra and the collision.}\quad
Now additionally assume full spectral retuning.
Let $a(v)$ be an algebraic characteristic-polynomial
coefficient, with irreducible minimal relation $P_a$.
Its evaluated values agree with the corresponding
Laurent coefficient $p_a(v)$ of
$\prod_w(T-v^w)^{m_w}$.
Thus $P_a(b^n,p_a(b^n))=0$ at infinitely many points.
The rational function $P_a(v,p_a(v))$ is identically
zero, and its rational root forces $a=p_a$.
This proves the stated characteristic-polynomial identity,
even when different allowed sheets are chosen at
different extension parameters.

If the family is regular at a point above $v=1$,
its characteristic polynomial there is $(T-1)^N$.
The center of its full centralizer is therefore local,
by the complete cyclic-module proof in
\eqref{eq:nilpotentcentralizercenter}.
The actual center
$k\times k\times k[\epsilon]/(\epsilon^2)$
has three components, giving the collision contradiction.
\end{proof}

\paragraph{Algebraic sheets and the exact scope.}
The finitely many matrix coefficients generate a finite
extension even if the initial coefficient extension is
not finite. Their trace is one algebraic element.
Allowed evaluations may choose different sheets at
different $b^n$, but must satisfy its same irreducible
relation. Switching to an unrelated algebraic function
or a different constant-field component is not part
of that assumption.
No continuation at infinity, common eigenvalue labeling,
or diagonalization of the family was used.
The trace and single-weight conclusions require no
regularity at $v=1$.
For the inverse quantum-parameter convention, replace
$v$ by $v^{-1}$; the normalization becomes $v^w$
in the original coordinate.

From trace comparison alone, only the base operator
is proved unipotent after single-weight normalization.
Determining the full characteristic polynomial at
other parameters above used the separately stated
full spectral comparison.
The obstruction permits any auxiliary rank and Jordan
blocks before the comparison is imposed.
It addresses one fixed integral pure weight and a
finite genuine spectrum, with the ordinary trace
rule \eqref{eq:algebraictraceretune}.
The several-weight extension requires integral spectral
weights. The collision conclusion additionally requires
full spectral retuning and a regular fixed-rank model.
This extends the rational collision theorem of
\fid{F48} to algebraic entries in that integral-weight
case. Nonintegral spectral weights are not classified
by this argument.
Infinite, non-atomic, regularized and virtual traces
require their own convergence and visibility checks.

The fixed-parameter models of \fid{F46} and the
faithful tensor actions of \fid{F48} remain valid.
What is incompatible is their proposed strengthening
to a full nonsimple semisimple centralizer, one pure
weight, an algebraic parameter family and this
all-extension trace comparison.
This directly covers the algebraic coefficient-field
format allowed in VIII, \S2.2, p.~7 \cite{VIII},
when that additional comparison is required.
It is not a general nonstandard Riemann hypothesis.

\subsection{Algebraic spectral retuning detects the collision order}
The integral-weight hypothesis can be removed if the comparison
retains the whole relative centralizer order. The actual
order-two contact then gives an obstruction that an abstract
special-fiber algebra alone does not detect.

\begin{finding}[\fid{F52}: the algebraic contact-order gate]\label{f:algebraiccontact}
Let $K/\C(v)$ be finite and let $F\in GL_N(K)$.
Fix $b>1$ and regular evaluations of this same algebraic
family such that, for every sufficiently large integer $n$,
\begin{equation}\label{eq:contactretuning}
 \operatorname{Spec}F_{b^n}
     =\operatorname{Spec}(F_b^n)
\end{equation}
as multisets. The evaluations may use different places.
No purity or integral spectral-weight hypothesis is imposed.

Let $R=\C[v,v^{-1}]_{(v-1)}$, let $R'$ be a discrete
valuation ring of $K$ above $R$, and suppose $F$ has
a regular invertible model on a finite free $R'$-module.
With
\[
 \delta=v+v^{-1},\quad D=\delta^2,\quad
 \Delta=\frac{(D-1)(D-4)}{D^2},
\]
its full relative centralizer cannot be isomorphic,
as an $R'$-algebra, to the actual order
\[
 A_{R'}=R'\times R'\times
 M_2\bigl(R'[x]/(x(x-\Delta))\bigr).
\]

There is a first-order algebraic model satisfying
\eqref{eq:contactretuning} whose abstract special
centralizer is nevertheless the same algebra as $A_1$.
It has nonintegral spectral weights and a different
relative order. The actual order embeds in its centralizer
and retains the actual rank-$64$ tensor module by restriction.
On that unchanged auxiliary module, adding ordinary linear
operators cannot make the embedded nonsaturated suborder
their simultaneous centralizer.
\end{finding}

\begin{proof}
\par\smallskip\noindent\textbf{Integral powers of the base eigenvalues.}\quad
Regard $F$ as a $\C(v)$-linear endomorphism of $K^N$.
Its characteristic polynomial over $\C(v)$ annihilates
the original $K$-linear matrix. Clear denominators to
obtain a nonzero $Q(v,Y)\in\C[v,Y]$ with $Q(v,F)=0$.
Outside finitely many parameter values its coefficients
are regular. For every base eigenvalue $\alpha$,
\eqref{eq:contactretuning} gives
\[
 Q(b^n,\alpha^n)=0
\]
eventually.

Expand $Q$ into distinct nonzero monomials
$q_{a,c}v^aY^c$. This is a finite power-sum relation
on roots $b^a\alpha^c$.
Vandermonde uniqueness forces every grouped coefficient
to vanish, so two distinct monomials have equal roots:
\[
 b^a\alpha^c=b^{a'}\alpha^{c'}.
\]
They cannot have $c=c'$, since then $b>1$ would force
$a=a'$ as well. Thus, after reversing signs if needed,
\[
 \alpha^L=b^a\quad\hbox{for some }L>0,\ a\in\Z.
\]
Choose one common $L$ for all base eigenvalues, including
their multiplicities, and write
\begin{equation}\label{eq:contactbasecharacters}
 \alpha_j^L=b^{a_j},\qquad a_j\in\Z.
\end{equation}
Their magnitudes are rational powers of $b$ and their
phases are roots of unity; this conclusion was proved
rather than assumed as a weight convention.

\par\smallskip\noindent\textbf{A monomial characteristic polynomial on a finite cover.}\quad
Adjoin $t$ with $v=t^L$.
The finite algebra
\[
 K\otimes_{\C(v)}\C(t)
\]
is a product of fields: in characteristic zero a
primitive-element polynomial for $K$ is separable,
and its distinct irreducible factors over $\C(t)$
give this product decomposition.
On at least one component $K''$ we claim
\begin{equation}\label{eq:contactcharpoly}
 \det(TI-F)=\prod_{j=1}^N(T-t^{a_j}).
\end{equation}

Indeed restrict \eqref{eq:contactretuning} to $n=Lr$,
and at $v=b^{Lr}$ prescribe the positive value $t=b^r$.
The evaluated spectrum is
\[
 \{\alpha_j^{Lr}\}_j=\{b^{a_jr}\}_j
                  =\{t^{a_j}\}_j.
\]
After discarding finitely many bad parameters, these
evaluations select factors of a finite model of the
tensor product. One component occurs for infinitely
many distinct $r$.
On that component each difference between the family
characteristic-polynomial coefficients and those
in \eqref{eq:contactcharpoly} has infinitely many zeros
over distinct $t=b^r$.

A nonzero algebraic function cannot have those zeros.
Its monic minimal polynomial over $\C(t)$ has nonzero
constant term; outside the finitely many coefficient
poles, a zero of the function forces a zero of that
nonzero rational constant term.
Thus every coefficient difference is zero, proving
\eqref{eq:contactcharpoly}.

The finite-model assertion needs no continuous sheet.
One may invert finitely many polynomials so that the
primitive-element integral model, discriminant, matrix
coefficients and the idempotents selecting the field
components are regular. The chosen evaluation on $K$
and $t=b^r$ then give a map from their tensor product
to $\C$, which selects one factor. Infinitely many
parameters survive every such finite deletion.
Other embeddings may multiply the resulting characters
by roots of unity; the local argument is unaffected.

\par\smallskip\noindent\textbf{The relative-center contradiction.}\quad
Extend $R'$ to a valuation ring $R''$ in that component
$K''$. The extension is a DVR and is torsion free,
hence flat, over $R'$.
The full centralizer is the kernel of
\[
 \operatorname{Mat}_N(R')\longrightarrow
 \operatorname{Mat}_N(R'),\qquad U\longmapsto UF-FU.
\]
Flat base change preserves this kernel.
Thus an assumed isomorphism with $A_{R'}$ would give
the same actual order over $R''$.
Put $e=\operatorname{ord}_{R''}(v-1)>0$.

The operator $F$ is central in its full centralizer.
Its component in the matrix-center factor must
therefore have the form $c_0+c_1x$, with
$c_0,c_1\in R''$.
On the two generic center branches its values are
\[
 \beta_1=c_0,\qquad
 \beta_2=c_0+c_1\Delta.
\]
Hence
\begin{equation}\label{eq:contactdivisibility}
 \beta_2-\beta_1\in\Delta R''.
\end{equation}
But
\[
 D-4=(v-v^{-1})^2
     =(v-1)^2(v+1)^2/v^2,
\]
and the remaining factor in $\Delta$ is a unit at
$v=1$. Therefore
\[
 \operatorname{ord}_{R''}(\Delta)=2e.
\]

The generic actual order is the split semisimple
product of two scalar and two two-by-two matrix
algebras. A central element in it is a tuple of
scalars, so its minimal polynomial is square free
after repeated values are removed. Thus $F$ is
diagonalizable generically.
Its values on the four primitive central blocks
must be distinct. Otherwise every linear map between
two nonzero eigenspaces with the same value would
commute with $F$, contradicting their status as
separate central blocks of the full centralizer.
In particular $\beta_1,\beta_2$ are distinct functions.

By \eqref{eq:contactcharpoly} they are two distinct
monomial characters $t^a,t^c$.
At the chosen point over $v=1$, $t=t_0$ with
$t_0^L=1$. The zero of $t^L-1$ at $t_0$ is simple,
so $\operatorname{ord}_{R''}(t-t_0)=e$.
The divisibility \eqref{eq:contactdivisibility}
forces $t_0^a=t_0^c$.
Distinctness gives $a\ne c$, and
\[
 \left.\frac{d}{dt}(t^a-t^c)\right|_{t=t_0}
       =(a-c)t_0^{a-1}\ne0.
\]
This holds also for negative exponents since
$t_0\ne0$. Thus
\[
 \operatorname{ord}_{R''}(\beta_2-\beta_1)=e,
\]
contradicting the required order at least $2e$.
This proves the full relative-order obstruction.

\par\smallskip\noindent\textbf{A sharp first-order model.}\quad
Let $v=t^3$ and choose a primitive cube root $\omega$
as the point $t=\omega$ over $v=1$. Put
\[
 \gamma_1=t^2,\qquad \gamma_2=t^5,\qquad
 \eta=\gamma_2-\gamma_1=t^2(v-1).
\]
Use the finite localization $T_0$ of $\C[t,t^{-1}]$
which inverts
\[
 \delta(v),\ h(v)=D(v)-1,\ D(v)-2,\quad
 t-1,\ t^2-1,\ t^5-1,\ t^2-t,\ t^5-t.
\]
All these factors are nonzero at $t=\omega$ and
at every positive real $t>1$.
This is an explicit finite coefficient base admitting
both the required local point and the arithmetic
evaluations.

Set
\[
 C_\eta=T_0[x]/(x(x-\eta)),\qquad
 J=C_\eta,\qquad F_J=\gamma_1+x.
\]
In the column basis $(1,x)$,
\[
 F_J=\begin{pmatrix}\gamma_1&0\\1&\gamma_2\end{pmatrix}.
\]
Its determinant is $t^7$, a unit.
Since $x=F_J-\gamma_1$, its centralizer on $J$ is
$C_\eta$, including at collision.
With $(T_0)_a$ denoting a scalar line of Frobenius $a$,
put
\[
 H=(T_0)_1\oplus(T_0)_t\oplus J^{\oplus2}.
\]
The sector differences are units:
the relevant determinants are
$(\gamma_1-1)(\gamma_2-1)$ and
$(\gamma_1-t)(\gamma_2-t)$.
Therefore
\begin{equation}\label{eq:firstordercentralizer}
 \End_F(H)=T_0\times T_0\times M_2(C_\eta).
\end{equation}

At $t=\omega$ the three eigenvalue sectors are
$1,\omega,\omega^2$, and $J$ remains a nontrivial
Jordan block in the last one.
The special algebra is abstractly
\[
 \C\times\C\times M_2(\C[\epsilon]/(\epsilon^2)),
\]
the same as $A_1$.
At $v=b>1$ choose $t_b=b^{1/3}$, and at $v=b^n$
choose $t_n=t_b^n$.
The spectrum consists of characters $1,t,t^2,t^5$
with multiplicities $1,1,2,2$, so it obeys full
spectral retuning.
The lower-left entry need not obey a fixed-basis
matrix-power identity; no such identity is claimed.

Under $v=Q^{1/2}$ these characters have spectral
weights $0,1/3,2/3,5/3$.
This is explicitly a nonintegral-weight model,
not a mixed model with integral Deligne weights.
The point $t=\omega$ is not the positive-real
branch limit $t=1$.
It does not contradict \fid{F51}.
The relative center in
\eqref{eq:firstordercentralizer} has first-order
contact, which explains why it also does not
contradict the present theorem.

\par\smallskip\noindent\textbf{The actual order as a specified subalgebra.}\quad
The element
\[
 c=\frac{\Delta}{\eta}
   =\frac{(D-1)(v-1)(v+1)^2}{D^2t^2v^2}
\]
belongs to $T_0$ and has order one at $t=\omega$.
There is an injection
\begin{equation}\label{eq:actualsuborder}
 T_0[x_{\rm act}]/(x_{\rm act}(x_{\rm act}-\Delta))
   \longrightarrow C_\eta,\qquad x_{\rm act}\longmapsto cx.
\end{equation}
Indeed $(cx)^2=\Delta(cx)$.
Its free-module matrix in the bases $(1,x_{\rm act})$
and $(1,x)$ is $\operatorname{diag}(1,c)$, proving
injectivity over the domain.
Keeping the two scalar sectors and taking two-by-two
matrices embeds the actual order into
\eqref{eq:firstordercentralizer}.

Locally at $t=\omega$, the center cokernel is
$R'/(c)$, of length one, and the matrix-sector
cokernel has length four.
The generic embedding is an isomorphism.
In the special fiber $x_{\rm act}$ maps to zero,
killing the actual four-dimensional matrix radical.
Thus the abstract special centralizer has the right
isomorphism type, but the actual relative embedding
is not that special isomorphism.

Let $L_0=C_\eta/(x)$ and $L_\eta=C_\eta/(x-\eta)$,
with Frobenii $\gamma_1,\gamma_2$, and define
\[
 K=(T_0)_1^{20}\oplus(T_0)_t^4
       \oplus L_0^{16}\oplus L_\eta^4.
\]
The cyclic source and unit sector differences give
\[
 \operatorname{Hom}_F(H,K)
  =\epsilon_+^{20}\oplus\epsilon_-^4
       \oplus\operatorname{Row}(L_0)^{16}
       \oplus\operatorname{Row}(L_\eta)^4,
\]
of rank $64$.
Under \eqref{eq:actualsuborder}, $x_{\rm act}$ acts
on these branches by $0,\Delta$.
This is the actual \fid{F21} tensor lattice after
the indicated coefficient extension, with its named
generators. The regular matrix comparison is defined
because $\delta,h,D-2$ are units on the chosen base.

Postcomposition on the multiplicity factors retains
the actual ordinary left operators as in \fid{F46}.
Over the domain, cross-branch intertwiners vanish
because $\eta$ is a nonzero divisor.
The multiplicity commutant is
$M_{20}(T_0)\times M_4(T_0)\times
 M_{16}(T_0)\times M_4(T_0)$, and its matrix units
lift to $K$.
At collision the branch multiplicities combine as
before. The spectra of $H$ and $K$ retune, with the
stated nonintegral weights, and their invariant Hom
has identity Frobenius.
This does not give the all-extension quantum-word
trace on that one fixed Hom space or the
\fid{F24} opposite-sign output.

\par\smallskip\noindent\textbf{Why extra linear commutators do not repair this embedding.}\quad
Localize the explicit model at $t=\omega$ and retain
this same free $H$. Let $B=\End_F(H)$ and let $A$
be the actual embedded suborder.
Any simultaneous centralizer of ordinary $R'$-linear
operators in $\End_{R'}(H)$ is saturated.
Indeed, if $d\ne0$ and $dU$ commutes with every
specified operator $G$, then
$d(UG-GU)=0$.
The matrix module is torsion free, so $UG=GU$ for
each $G$ and $U$ lies in that centralizer.

The embedded $A$ is not saturated.
Take $U=xE_{11}$ in the $M_2(C_\eta)$ sector.
It is a regular endomorphism in $B$ and is not in
$A$, since its $x$ coefficient is not divisible
by $c$. But $cU\in A$.
Therefore no additional ordinary linear operators
on this unchanged $H$ have simultaneous centralizer
equal to the displayed $A$.
Changing the module lattice or object, retaining
$A$ as explicit algebra data, or making a different
categorical comparison is outside this assertion.
\end{proof}

\paragraph{The exact remaining alternatives.}
The obstruction requires full spectral retuning of
one finite algebraic family and preservation of the
whole regular relative centralizer order.
It allows arbitrary evaluation sheets and auxiliary
rank and needs no purity assumption.
An abstract special-fiber isomorphism, a proper
specified subalgebra, a nonregular or nonflat model,
or a trace-only scalar comparison does not supply
those hypotheses.
The explicit model demonstrates this distinction
while retaining the actual tensor action by restriction.
Its saturation obstruction fixes that model and
embedding; it is not a prohibition on all extra
algebraic structures.
A regular-module wrapper with identity Frobenius
can tautologically retain algebra and module actions,
but does not produce the quantum-grading, canonical-basis
or path-sign comparison.

\subsection{A coherent Gaussian convolution carrying the actual tensor maps}
Finding~\fid{F50} supplied exact Gaussian path signs and a faithful
realization of point-Weil Hom spaces. Its failure of ordinary tensor
monoidality has an explicit geometric repair. The repair uses a
specified convolution and leaves the pointwise weight question unchanged.

\begin{finding}[\fid{F53}: geometric convolution and actual tensor operations]
\label{f:gaussianconvolution}
Retain the odd-characteristic Gaussian data of \fid{F50}:
\[
 L=\mathcal L_\psi(x^2),\qquad
 G=H_c^1(\A^1_{\bar k},L),\qquad
 J(C)=L\otimes\underline{C\otimes G^\vee}[1].
\]
An underline denotes constant pullback from the finite-field point.
There is an actual geometric product
\begin{equation}\label{eq:gaussianstar}
 A\star_\psi B=
 \underline{G(1)}[1]\otimes
 Rm_!\left(A\boxtimes B\otimes\mathcal L_\psi(2xy)\right),
 \qquad m(x,y)=x+y,
\end{equation}
with the associator and symmetry specified below, for which
\[
 J(C)\star_\psi J(C')\simeq J(C\otimes C').
\]
On the full perverse image of $J$, its unit is $J(E)$ and $J$ is
an exact fully faithful symmetric strong monoidal equivalence from
finite Weil spaces. Its image unit differs from its ambient unit,
and its transported symmetry differs by a sign from the naive
kernel swap on two Gaussian image objects.

Composing $J$ with \fid{F48}'s actual fixed-parameter tensor
functor gives an exact faithful strong monoidal realization of the
nonstandard family through all degrees by pure weight-zero
perverse objects. It preserves the actual named maps and their
existing relations. The original external-product path construction
of \fid{F50} retains its exact negative effective spectrum and is
natural in those coefficient maps.

This supplies coherent operations on these chosen objects.
It does not make the source functor full, identify its quantum
symmetry with the new symmetry, select the designated basis, or
improve the exact pointwise/dual estimate of \fid{F50}.
\end{finding}

\begin{proof}
\par\smallskip\noindent\textbf{Normalized constants.}\quad
Let $*$ be ordinary additive convolution
$A*B=Rm_!(A\boxtimes B)$ on constructible Weil complexes.
It has the usual Fubini associator, the derived symmetry
induced by variable interchange, and ambient unit $\delta_0$.
Put
\[
 \Theta=\underline E(1)[2],\qquad
 K(C)=\underline C\otimes\Theta.
\]
The coordinates $(s=x+y,y)$ make $m$ a trivial affine-line bundle.
Consequently
\[
 Rm_!\underline E=\underline E(-1)[-2].
\]
Canonical compact trace gives natural isomorphisms
\begin{equation}\label{eq:normalizedconstantconv}
 K(C)*K(C')\simeq K(C\otimes C').
\end{equation}
Explicitly, its left side is
$\underline{C\otimes C'}\otimes Rm_!\underline E(2)[4]$.

The primary formalism is affine-space cohomology
\cite[Example~16.3, p.~109]{Milne},
base change and compactification
\cite[Theorem~17.10, p.~116; Proposition~18.4,
pp.~121--122]{Milne}, and compact Kunneth and projection
formula \cite[Theorem~22.1, Remarks~22.3(b),(d),
Proposition~22.5, pp.~135--137]{Milne}.
The canonical trace
\[
 R\Gamma_c(\A^1,\underline E(1)[2])\simeq E
\]
is normalized by sending the class of each geometric point
to one \cite[Theorem~24.1(a),(b), p.~144]{Milne}.

We check coherence directly. For the addition map $m_r$
on $r$ coordinates, the fibers are affine $(r-1)$-spaces.
Their canonical relative compact trace gives
\begin{equation}\label{eq:normalizedconstantmulti}
 R(m_r)_!\bigl(K(C_1)\boxtimes\cdots\boxtimes K(C_r)\bigr)
 \simeq K(C_1\otimes\cdots\otimes C_r).
\end{equation}
Every iterated binary comparison is this map.
Indeed the fiber has compact cohomology only in top degree
$2(r-1)$, successive normalized traces send its point
class to one, and Kunneth sends a product of point classes
to the point class in the product.
Changes between fiber coordinates for different parenthesizations
are affine automorphisms. They preserve the canonical trace,
since they send point classes to point classes.
Thus every parenthesization gives
\eqref{eq:normalizedconstantmulti}; this proves the pentagon
and all higher tensor compatibilities.

Taking either coefficient to be $E$ in
\eqref{eq:normalizedconstantconv} gives left and right unit
constraints with unit $\Theta$. The three-factor comparison
proves the triangle identity.
There is also a concrete ambient unit map.
For $i:\{0\}\hookrightarrow\A^1$, cohomological purity and
the $i_*$--$Ri^!$ adjunction give
\[
 Ri^!E=E(-1)[-2],\qquad \delta_0\longrightarrow\Theta.
\]
This is the origin Gysin class
\cite[Theorem~16.1, p.~108; Theorem~16.7 and its
preceding adjunction, p.~111]{Milne}.
Convolution with this map on a normalized constant object
is inverse to its unit comparison: the composite integrates
a point class and is therefore the identity.

Permutations of the inputs preserve the affine-fiber trace.
The shifts $[2]$ are even, so the resulting coefficient
symmetry is the ordinary flip with no additional sign.
The same $r$-factor comparison proves its compatibility
with the associator.
Equal shifts and twists cancel in degree-zero Hom;
connectedness of geometric $\A^1$ identifies constant-sheaf
maps with linear maps, and Weil compatibility is Frobenius
commutation.
Hence $K$ is an exact fully faithful symmetric strong
monoidal equivalence onto the full normalized-constant
subcategory, with its unit $\Theta$.
This unit has full support and is not the ambient unit
$\delta_0$.

\par\smallskip\noindent\textbf{The Gaussian gauge and both units.}\quad
Define an autoequivalence
\[
 \Phi(A)=\underline{G(1)}[1]\otimes L^\vee\otimes A,
 \qquad
 \Phi^{-1}(B)=L\otimes\underline{G^\vee(-1)}[-1]\otimes B.
\]
Evaluation of $G$ against its dual gives
\[
 \Phi J(C)\simeq K(C).
\]
Define the product, associator and symmetry by transport:
\begin{equation}\label{eq:gaussianstartransport}
 A\star_\psi B=\Phi^{-1}(\Phi(A)*\Phi(B)).
\end{equation}
Pulling the constant factors through $m_!$ and applying
projection formula leaves one $\underline{G(1)}[1]$.
The remaining sheaf is
\[
 m^*L\otimes\operatorname{pr}_1^*L^\vee
          \otimes\operatorname{pr}_2^*L^\vee
 =\mathcal L_\psi((x+y)^2-x^2-y^2)
 =\mathcal L_\psi(2xy).
\]
This proves \eqref{eq:gaussianstar}.
The Artin--Schreier tensor identities use the fixed
additive Lang character; the scalar phase also satisfies
\[
 2xy+2(x+y)z=2yz+2x(y+z).
\]
Transport specifies the constant lines, shifts and symmetry
in addition to that polynomial cocycle.

The ambient unit of \eqref{eq:gaussianstartransport} is
\[
 u_\psi=\delta_0\otimes\underline{G^\vee(-1)}[-1].
\]
On the full perverse image of $J$, the normalized-constant
comparisons instead give
\[
 J(C)\star_\psi J(C')\simeq J(C\otimes C'),
 \qquad e_\psi=J(E).
\]
This image unit has full support. The transported Gysin
map gives $u_\psi\to e_\psi$; convolution with an image
object on either side is an isomorphism and supplies
its image unit constraint.
Thus the two units have a precise relationship and are
not identified with each other.

The product is exact and perverse on this image, because
it is the coefficient tensor product under these
comparisons. No global t-exactness is asserted.
Full faithfulness and exactness of $J$ are those of
\fid{F50}; the new coherent comparisons make it a
symmetric strong monoidal equivalence onto its image.

\par\smallskip\noindent\textbf{The symmetry correction.}\quad
Transported symmetry is $J$ of the usual coefficient flip.
In particular it is the identity on $e_\psi\star_\psi e_\psi$
after its unit identification.

The naive kernel swap gives a different map.
Each $J$-input is lisse shifted by $[1]$, so interchanging
two such inputs gives the Koszul sign $-1$.
After the phases combine to $L(x+y)$, the fiber involution
is $y\mapsto s-y$. It preserves normalized top compact
trace, since it is an affine automorphism.
Consequently, on the image,
\[
 c_{\mathrm{naive}}\bigm|_{J(C),J(C')}
       =-J(\operatorname{flip}_{C,C'}).
\]
For $C=C'=E$ it acts by $-1$ on the unit and cannot be its
monoidal symmetry. Transport through $\Phi$ supplies
the correction. Equivalently, $\Phi J(C)=K(C)$ has
even shift $[2]$, where ordinary convolution symmetry
has no extra sign.
Thus the kernel expression alone is not being used
to choose a commutor.

\par\smallskip\noindent\textbf{Weights and duals.}\quad
Finding~\fid{F50} establishes that $G$ has weight one.
If $C$ is pure of integer weight $w$, the lisse
coefficient of $J(C)$ has weight $w-1$ and its shift
is $[1]$, so it is perverse pure of weight $w$.
Likewise $K(C)$ has coefficient weight $w-2$ and
shift $[2]$, hence complex weight $w$
\cite[5.1.8, pp.~126--127]{BBD}.

The internal monoidal dual is $J(C^\vee)$.
Its evaluation and coevaluation are $J$ of the
coefficient maps, with target or source the image unit.
Their triangle identities follow from the proved
tensor coherence.
Literal Verdier duality on the line instead gives
\[
 \mathbb D J_\psi(C)
   =L^\vee\otimes\underline{C^\vee\otimes G}(1)[1].
\]
The Gaussian pairing
$G\otimes G_{\psi^{-1}}\simeq E(-1)$ identifies
$G_{\psi^{-1}}^\vee$ with $G(1)$, so
\begin{equation}\label{eq:gaussianthreeduals}
 J(C)^{\vee_\star}=J_\psi(C^\vee),\qquad
 \mathbb D J_\psi(C)\simeq J_{\psi^{-1}}(C^\vee),\qquad
 J_{\psi^{-1}}(C^\vee(-1))
       \simeq\mathbb D J_\psi(C)(-1).
\end{equation}
The last object is the separately specified reciprocal
system of \fid{F17}, with total eigenvalues $q/\alpha$.
Internal dual, Verdier dual and that reciprocal system
are distinct constructions. A categorical dual does not
prove the second pointwise bound.
Arithmetic Jordan blocks remain allowed.

On the image, compact cohomology is inverse to $J$.
The inverse tensor comparisons make it strong monoidal
there, compatibly with the natural Gaussian evaluation.
No global monoidality for all twisted-convolution
complexes is asserted.

\par\smallskip\noindent\textbf{The actual tensor family.}\quad
Fix a regular algebraic quantum parameter and the actual
IV two-row Kronecker category, $X=V\otimes W$ with both
factors two-dimensional.
Use IV, Definition~8.1 and equation~(61), p.~39;
Propositions~10.3 and~10.5, pp.~56--57; and
Theorem~12.1, p.~68, with its immediately following
dual-corepresentation convention \cite{IV}.
These are precisely the source hypotheses of \fid{F48}.

For the actual comodule forgetful functor $\omega$, put
\[
 \mathcal T_0(N)=J((\omega N,1)).
\]
Forgetting is exact and faithful, and the underlying
tensor product is the ordinary coefficient tensor
product. The Gaussian comparisons therefore make
$\mathcal T_0$ exact, faithful and strong monoidal,
with perverse pure weight-zero images.
For every $r$ there is a coherent identification
\[
 \mathcal T_0(X^{\otimes r})
       \simeq\mathcal T_0(X)^{\star_\psi r}.
\]
The named adjacent generators remain the original
coefficient matrices
\[
 P_i^{\mathrm{ns}}=D\Pi_i,\qquad
 D=(v+v^{-1})^2,
\]
where $\Pi_i$ is the source symmetric projector.
The generator convention of the main paper is unchanged.

Every source map is sent to $J$ of its unchanged
coefficient map. Functoriality and coherence preserve
compositions, tensor products, block inclusions,
multiplication maps and all their existing relations
as actual equalities.
Identity Frobenius also makes every specified ordinary
left operator a Weil morphism, even when it is not
a comodule intertwiner.
Thus the named left and right operators of \fid{F43}
retain their matrices, commutations and original
row/opposite-composition convention as actual perverse
morphisms. These extra maps remain decorations;
they have not been declared source comodule maps.

In IV's Hopf case, existing rigid dual maps and
determinant inversion are retained using the internal
dual above. In VII's general bialgebra case, no antipode
is supplied when the determinant can vanish.
Specialization retains the actual defined action,
including any kernel it already has; faithfulness
on source morphisms does not repair a specialized
tensor action that kills the radical.
The source need not be symmetric.
A quantum Hecke braid or designated projector remains
its actual named morphism, without being identified
with the target's symmetry or exterior-power operation.

The finite-order character functor $\mathcal R_\chi$
of \fid{F48} similarly gives
$\mathcal T_\chi=J\circ\mathcal R_\chi$, exact,
faithful, strong monoidal and perverse pure of weight zero.
For example its covariance
$FE_VF^{-1}=\zeta^2E_V$ becomes the actual map
\[
 J(L_{\zeta^2})\star_\psi\mathcal T_\chi(X)
       \longrightarrow\mathcal T_\chi(X).
\]
The other raising and lowering operators have the
specified inverse phases.
Both choices keep the quantum parameter fixed;
neither imposes extension-wise parameter retuning.

\par\smallskip\noindent\textbf{The original path sign.}\quad
Nothing changes the objects $J(C)$ or their ordinary
external products. If $\gamma$ is Frobenius on $G$,
\[
 t_{C,n}(x)=-\gamma^{-n}\Tr(F^n\mid C)\psi_n(x^2),
 \qquad \sum_x t_{C,n}(x)=\Tr(F^n\mid C).
\]
For every good even path length of \fid{F50},
\[
\begin{aligned}
 R\Gamma_c(H_k,J(C)^{\boxtimes k}|_{H_k})
 \simeq{}& C^{\otimes k}\otimes G^{\vee\otimes k}\\
 &\otimes
 R\Gamma_c(H_k,\mathcal L_\psi(\textstyle\sum_i x_i^2))[k].
\end{aligned}
\]
The unshifted geometric factor has only degree $k-1$,
so the result has only degree $-1$.
Its twisted trace is exactly
\[
 q^nN_{k,C,n}
   =-\Tr(F^n\mid
       C^{\otimes k}\otimes G^{\vee\otimes k}\otimes B_{k,\psi}),
\]
with the effective space of \fid{F50}, zero error,
and an unbounded good even set.
A coefficient map $f$ induces $f^{\otimes k}$ on
the displayed coefficient factor. This proves the
claimed naturality in actual named maps and equations.
The coherent convolution comparison involves the
same coefficient tensor power.
The path itself is still the ordinary external-product
restriction required in \fid{F01}--\fid{F04};
it has not been replaced by a convolution path.
\end{proof}

\paragraph{What the operation does not change.}
The sign statement is about identity traces of effective
objects. An arbitrary inserted trace $\Tr(TF^n)$ need
not have effective spectral multiplicities, and its
tensor-inserted path output need not have the required
sign. Retaining $T$ and its tensor powers as morphisms
does not remove that distinction.

For each nonzero coefficient object, the exact pointwise
thresholds are still $b_{\max}-1$ and $1-b_{\min}$.
For one integer target $w$, the required two bounds
hold under every embedding precisely when
\[
 b_{\max,\sigma}\le w\le b_{\min,\sigma}
                    \quad\text{for every }\sigma.
\]
Thus they still assert the original pure weight.
For \fid{F45}'s total length-$L$ word object the
dual gap remains exactly $4L$.
The identity and finite-order families above already
have weight zero; their successful geometric operations
do not supply an unknown coefficient estimate or the
designated quantum-grading comparison.

The source functor is not full. For identity Frobenius,
\[
 \End(\mathcal T_0(X))\simeq M_4(E),\qquad
 \End_{\mathrm{source}}(X)=E.
\]
Its exact Grothendieck-ring character is
\[
 [\mathcal T_0(N)]=\dim_E(\omega N)[e_\psi],
 \qquad [e_\psi]\star[e_\psi]=[e_\psi].
\]
It is just the dimension character.
Geometrically every $J(C)$ is a sum of copies of the
single Gaussian lisse $[1]$ constituent.
A source simple $S$ therefore contributes $\dim S$
copies of that one constituent, rather than a new
distinguished geometric simple label.
Tensor faithfulness does not identify the source
simple basis with a KL-style perverse basis.

For \fid{F43}'s graded coefficient kernels, entrywise
$J$ and $\star_\psi$ transport direct sums, tensor
factors, and every already specified coefficient
map and word composition.
Odd degrees have odd integer Weil weights; integral
Tate motives are not asserted for them.
This does not turn a Laurent-matrix or graded-$K_0$
equality into a uniquely chosen natural isomorphism,
or provide coherence for relations established only
at $K_0$ level. The actual \fid{F48} matrix relations
are existing equalities of linear maps and are
preserved exactly.
No prescribed integral lattice, bar globalization,
cell or canonical-basis identification follows.

\subsection{Identity traces over algebraic coefficient fields}
Full spectral retuning is unnecessary for classifying the
base roots of a genuine finite trace. Positivity of their
multiplicities supplies a group-algebra argument. For a
single common spectral radius, identity traces already
give the generic full-centralizer obstruction.

\begin{finding}[\fid{F54}, Part 1: genuine algebraic traces and one radius]
\label{f:algebraicgenuinetrace}
Fix $b>1$ and
\[
 u_n=\sum_{i=1}^s m_i\alpha_i^n,\qquad
 m_i\in\Z_{>0},\quad \alpha_i\in\C^\times,\qquad
 N=\sum_i m_i>0.
\]
Suppose a nonzero polynomial $P(v,Y)\in\C[v,Y]$
of positive $Y$-degree satisfies
\begin{equation}\label{eq:genuinetracealgebraic}
 P(b^n,u_n)=0
\end{equation}
for every sufficiently large integer $n$.
Then there are a common integer $L>0$ and integers
$a_i$ such that
\begin{equation}\label{eq:genuinetracepowers}
 \alpha_i^L=b^{a_i}\qquad\text{for every }i.
\end{equation}
Thus the radii are rational powers of $b$ and the phases
are roots of unity. Neither purity nor an integral
spectral-weight hypothesis is assumed.

Now let $F$ be an invertible finite matrix over an algebraic
extension of $\C(v)$, with a regular invertible value $F_b$.
For every sufficiently large $n$, choose a regular value
of the same family above $b^n$, and suppose only
\begin{equation}\label{eq:genuineidentityretuning}
 \Tr(F_{b^n})=\Tr(F_b^n).
\end{equation}
If all base eigenvalues have one common absolute value
$R>0$, then they are all equal to one number $\lambda$.
If the full centralizer of $F_b$ is semisimple, then
\[
 F_b=\lambda I,\qquad
 \End_{F_b}(\C^N)=M_N(\C).
\]
It therefore cannot be the actual generic algebra
$\C\times\C\times M_2(\C)\times M_2(\C)$.
This conclusion requires no regularity at $v=1$ and
no full characteristic-polynomial comparison.
\end{finding}

\begin{proof}
\par\smallskip\noindent\textbf{A genuine trace in a group algebra.}\quad
Let $\Gamma\subset\C^\times$ be the multiplicative group
generated by $b$ and the finitely many $\alpha_i$.
Its torsion subgroup is finite. Choose a positive integer
$m$ killing that subgroup and put
\[
 \Lambda=\Gamma^m=\{g^m:g\in\Gamma\}.
\]
This group is torsion free: if $g^m$ has finite order,
then $g$ has finite order and hence $g^m=1$.
It is therefore free abelian of finite rank $r$, so
\[
 \C[\Lambda]\simeq
 \C[z_1^{\pm1},\ldots,z_r^{\pm1}]
\]
is an integral domain.
Write $[\lambda]$ for its group-basis monomial and set
\[
 B=[b^m],\qquad U=\sum_i m_i[\alpha_i^m].
\]
The element $B$ is transcendental over $\C$: its powers
are distinct group-basis elements because $b^m$ has
infinite order.

A group-algebra element $W=\sum_\lambda c_\lambda[\lambda]$
gives the finite power sequence
\[
 W(r)=\sum_\lambda c_\lambda\lambda^r.
\]
This correspondence respects sums and products and is
injective even when sequences are compared only for
sufficiently large integers.
Indeed, on distinct nonzero roots
$\lambda_1,\ldots,\lambda_d$, vanishing at $d$ consecutive
indices starting at $r_0$ gives a Vandermonde matrix
times the invertible diagonal matrix
$\operatorname{diag}(\lambda_j^{r_0})$.
Every coefficient must vanish.

Applying this uniqueness to
\eqref{eq:genuinetracealgebraic} on $n=mr$ gives
\begin{equation}\label{eq:genuinegroupring}
 P(B,U)=0\quad\text{in }\C[\Lambda].
\end{equation}
Since $B$ is transcendental, $P(B,Y)$ is nonzero.
Thus $U$ is algebraic over $\C(B)$ in the fraction
field of the Laurent polynomial ring.

\par\smallskip\noindent\textbf{Euler derivations and positive support.}\quad
Choose a lattice basis so that
\[
 B=z_1^d\qquad(d>0).
\]
This is integer linear algebra: the positive gcd of
the coordinates of the nonzero vector for $b^m$ is
$d$, and successive extended-Euclidean coordinate
operations send it to $(d,0,\ldots,0)$ by a
unimodular change of basis.

For $j\ge2$, the Euler derivation
\[
 D_j=z_j\frac{\partial}{\partial z_j}
\]
extends to the fraction field and kills $\C(B)$.
Let $Q(Y)$ be the monic minimal polynomial of $U$
over $\C(B)$. It is separable in characteristic zero,
so $Q'(U)\ne0$.
Applying the derivation to its relation gives
\[
 0=D_jQ(U)=Q'(U)D_jU,\qquad D_jU=0.
\]
After grouping equal Laurent monomials, write
\[
 U=\sum_{a\in S}M_a z_1^{a_1}\cdots z_r^{a_r},
 \qquad M_a\in\Z_{>0}.
\]
The coefficient of this monomial in $D_jU$ is
$M_a a_j$.
Monomial independence and characteristic zero give
$a_j=0$ for every support vector and every $j\ge2$.
All support lies on the first lattice axis.

No original summand disappears during the torsion-killing
or grouping step, because all its multiplicities are
positive. Thus for every $i$ there is an integer $c_i$
with
\[
 [\alpha_i^m]=z_1^{c_i}.
\]
Raising this equality to the $d$th power gives
\[
 [\alpha_i^{md}]=[b^{mc_i}].
\]
Equality of group-basis elements is equality in the
original subgroup of $\C^\times$.
Hence \eqref{eq:genuinetracepowers} follows with
$L=md$ and $a_i=mc_i$.
If the lattice has rank one, the same conclusion holds
without any additional derivation.

For completeness, the relative-algebraic-closure fact
underlying this argument has the same elementary proof.
If $h\in\C(z_1,\ldots,z_r)$ is algebraic over $\C(z_1)$,
separability gives $D_jh=0$ for all $j\ge2$.
The kernel of differentiation in a characteristic-zero
rational function field $M(z)$ is $M$:
for $h=0$ the claim is immediate; otherwise write
$h=p/q$ with coprime polynomials.
The equation $p'q-pq'=0$ implies $p\mid p'$ unless
$p$ is constant; comparison of degrees gives $p'=0$,
and then $q'=0$ as well.
Thus both are constant.
Iterating proves $h\in\C(z_1)$.
Its intersection with the Laurent polynomial ring is
$\C[z_1^{\pm1}]$, also directly seen from monomial
derivatives. The displayed support proof already gives
the precise conclusion needed here.

\par\smallskip\noindent\textbf{The same algebraic trace at one common radius.}\quad
The finitely many entries of the matrix family generate
a finite algebraic extension.
Its trace $f=\Tr F$ is therefore one algebraic element.
Every regular evaluation of that element satisfies
its minimal relation, even if different allowed sheets
are selected at different parameters.
Equation \eqref{eq:genuineidentityretuning} permits the
preceding theorem to be applied to the genuine base
eigenvalues with their algebraic multiplicities.

Choose the common $L$ in \eqref{eq:genuinetracepowers}.
If every base eigenvalue has modulus $R$, then all
the integers $a_i$ are one integer $a$, because
\[
 R^L=b^{a_i}.
\]
Consequently at degrees $n=Lr$,
\[
 \Tr(F_b^{Lr})=N b^{ar}.
\]
The algebraic function
\[
 g(v)=f(v)^L-N^L v^a
\]
therefore vanishes at the selected regular places
above infinitely many distinct values $v=b^{Lr}$.
It must vanish identically.
Indeed, a nonzero algebraic function has a monic
minimal polynomial over $\C(v)$ with nonzero rational
constant term. Outside finitely many coefficient
poles, a zero of the function forces a zero of that
nonzero rational function, which has only finitely
many zeros.
Negative $a$ causes no difficulty in $\C(v)$ or
at the nonzero parameter $b$.
Thus
\begin{equation}\label{eq:commonradiustraceidentity}
 f(v)^L=N^L v^a.
\end{equation}

Evaluate this identity at the original regular place
above $b$, where $f_b=\Tr F_b$.
It gives
\[
 |\Tr F_b|=N b^{a/L}=NR.
\]
This is equality in the triangle inequality for $N$
actual eigenvalues of modulus $R$.
Their genuine positive multiplicities force a common
phase, hence every eigenvalue equals a single $\lambda$.

No diagonalizability has been assumed.
The operator $F_b-\lambda I$ is nilpotent and is
central in the full centralizer.
If that algebra is semisimple, its center is reduced,
so the nilpotent is zero.
The full centralizer is then $M_N(\C)$.
The actual generic algebra of \fid{F19} is semisimple
and nonsimple, giving the stated contradiction.
\end{proof}

\paragraph{The common phase and exact trace-only scope.}
In the unrestricted-radius case the common phase need
not be one.
For example $F(v)=\sqrt v\,I$ may be evaluated on the
negative sheet at $b$, and on the sheets
$(-1)^n b^{n/2}I$ at $b^n$.
The same irreducible algebraic family then obeys
\eqref{eq:genuineidentityretuning} with base eigenvalue
$-\sqrt b$.
Its spectral weight relative to $Q=b^2$ is $1/2$,
so it is outside the integral-weight special case
of \fid{F51}.
The new centralizer conclusion only needs the common
eigenvalue, not removal of this phase.

The group-algebra theorem classifies the base roots
from identity traces alone.
It does not determine the full characteristic
polynomial or eigenvalue functions at other parameters,
nor their collision contact order.
In particular it must not be used to exclude a
multiple-radius model preserving only selected
identity traces.
The stronger contact conclusion of \fid{F52} used
full spectral retuning.

Positive coefficients prevent disappearance of roots
after powering and grouping.
The argument does not recover virtual eigenvalues
that cancel from the visible sequence or a powered
progression.
Arbitrary inserted traces do not automatically
satisfy this positivity hypothesis.
The theorem is a necessary classification, not a
claim that every phase assignment satisfying
\eqref{eq:genuinetracepowers} is realized by one
irreducible algebraic trace function.

The fixed-parameter constructions of \fid{F46},
the faithful tensor families of \fid{F48}, and the
geometric convolution of \fid{F53} remain valid.
Their specified subalgebra actions are different
from equality with a full centralizer under the
additional arithmetic comparison.
No regularity at the collision, integral spectral
weight, or all-conjugate estimate was used for the
common-radius theorem. A pure spectrum at any chosen
embedding supplies its radius hypothesis.

\subsection{Two identity traces retaining the actual collision order}
The preceding common-radius result has a genuine multiple-radius
boundary. Both identity traces in the actual End/Hom model can
retune while the full relative order and its tensor action are
preserved.

\begin{finding}[\fid{F54}, Part 2: a genuine trace-only full-order model]
\label{f:traceonlyactualmodel}
There is a finite localization $S$ of $\mathbb Q[v,v^{-1}]$,
admitting $v=1$ and every $v=2^n$, and invertible regular
operators on free modules $H,K$ of ranks $6,44$, such that
\[
 \End_F(H)=A_S,\qquad \operatorname{Hom}_F(H,K)=M_S.
\]
Here $A_S$ is the actual order-two rank-three nonstandard
algebra of \fid{F19}, and $M_S$ is the actual rank-$64$
tensor lattice of \fid{F21}, with its named right generators.
Postcomposition also retains its ordinary left operators.
For every $n\ge1$,
\begin{equation}\label{eq:twogenuinetraces}
 \Tr F_H(2^n)=\Tr(F_H(2)^n),\qquad
 \Tr F_K(2^n)=\Tr(F_K(2)^n).
\end{equation}
The base spectra have integral weights $0,1,2,3$ under
$v=Q^{1/2}$, but full spectral retuning fails.
Thus the two identity traces do not determine the
parameter-dependent spectra or their collision data.
\end{finding}

\begin{proof}
Use
\[
 \delta=v+v^{-1},\quad D=\delta^2,\quad
 \Delta=\frac{(D-1)(D-4)}{D^2},\qquad
 S_0=\mathbb Q[v,v^{-1},\delta^{-1},(D-1)^{-1},(D-2)^{-1}].
\]
The previously proved actual algebra and tensor module are
\[
 A_{S_0}=S_0\times S_0\times M_2(C),\qquad
 C=S_0[x]/(x(x-\Delta)),
\]
\[
 M_{S_0}=\epsilon_+^{20}\oplus\epsilon_-^4
       \oplus N(\delta)^{16}\oplus N(D-2)^4.
\]
The two branch values of $x$ are $0,\Delta$ in that order.
All stated F21 changes of basis are regular over $S_0$.

\par\smallskip\noindent\textbf{Explicit scalar functions.}\quad
Set
\[
 k=\frac{2500}{189},\qquad a=3-v,
\]
\[
 c=-v^3+2v^2+4v-8+k\Delta,\qquad d=c+k\Delta,
\]
\[
 z=6v^3-6v^2-14v+30-6k\Delta.
\]
Collecting terms gives
\begin{align}
 a+z+2c+2d&=1+v+2v^2+2v^3=:S_H(v),
       \label{eq:traceescapeH}\\
 20a+4z+16c+4d&=20+4v+16v^2+4v^3=:S_K(v).
       \label{eq:traceescapeK}
\end{align}
Furthermore,
\[
 \Delta(2)=\frac{189}{625},\quad k\Delta(2)=4,
\]
\begin{equation}\label{eq:traceescapevalues}
 (a,z,c,d)|_{v=2}=(1,2,4,8),\qquad
 (a,z,c,d)|_{v=1}=(2,16,-3,-3).
\end{equation}

\par\smallskip\noindent\textbf{One finite base for all evaluations.}\quad
From $S_0$, invert $a,z,c,d$ and
\[
 a-z,\quad a-c,\quad a-d,\quad z-c,\quad z-d
\]
to obtain $S$.
The difference $d-c=k\Delta$ is not inverted.
Every displayed factor is nonzero at $v=1$ and $v=2$
by \eqref{eq:traceescapevalues}.
For $v>1$, $D>4$, so
\[
 0<\Delta=(1-D^{-1})(1-4D^{-1})<1.
\]
At $v=4$, exact values are
\[
 (a,z,c,d)=
 \left(-1,\frac{17007502}{83521},
       -\frac{1192004}{83521},-\frac{379504}{83521}\right),
\]
giving $c<d<a<0<z$.

For $v\ge8$,
\[
 d-a<-v^3+2v^2+5v-11+2k.
\]
The right side is $-62095/189$ at $8$ and has derivative
$-3v^2+4v+5<0$ on this range.
Thus $c<d<a<0$ there.
For $v\ge4$,
\[
 z>6v^3-6v^2-14v+30-6k.
\]
This lower bound is $11506/63$ at $4$ and has positive
derivative $18v^2-12v-14$ on the range.
Hence $z>0$.
Every factor inverted in $S$ is consequently nonzero
at every $2^n$.
The factors from $S_0$ are also nonzero at $1$ and
every real $v>1$. This proves the finite-base claim.

\par\smallskip\noindent\textbf{The exact algebra and module.}\quad
Let $J=C=S[x]/(x(x-\Delta))$ and give it the operator
\[
 F_J=c+kx
     =\begin{pmatrix}c&0\\ k&d\end{pmatrix}
\]
in the column basis $(1,x)$.
Its determinant $cd$ is a unit, and
$x=(F_J-c)/k$.
Therefore commuting with $F_J$ is the same as being
$C$-linear, and
\[
 \End_{F_J}(J)=C.
\]
This identity persists after every coefficient base change.
At $v=1$, $F_J$ has a genuine nonzero Jordan off-diagonal
entry $k$ and eigenvalue $-3$.

Write $S_r$ for a line with scalar operator $r$, and set
\[
 H=S_a\oplus S_z\oplus J^{\oplus2}.
\]
The scalar difference $a-z$ and the determinants
$(a-c)(a-d)$ and $(z-c)(z-d)$ are units.
They exclude all cross-sector intertwiners.
Consequently
\begin{equation}\label{eq:traceescapefullcentralizer}
 \End_F(H)=S\times S\times M_2(C)=A_S
\end{equation}
over $S$ and after any coefficient base change.
The center is the actual order-two center:
\[
 D-4=(v-1)^2(v+1)^2/v^2
\]
and the remaining factor in $\Delta$ is a unit at $1$.

Let $L_0=C/(x)$ and $L_\Delta=C/(x-\Delta)$, with
operators $c,d$, and put
\[
 K=S_a^{20}\oplus S_z^4\oplus L_0^{16}\oplus L_\Delta^4.
\]
The cyclic source gives
\[
 \operatorname{Hom}_F(J^2,L_i)
       =\operatorname{Hom}_C(C^2,L_i)
       =\operatorname{Row}(L_i).
\]
Precomposition supplies the right matrix action.
Thus
\[
 \operatorname{Hom}_F(H,K)=
 \epsilon_+^{20}\oplus\epsilon_-^4
 \oplus\operatorname{Row}(L_0)^{16}
 \oplus\operatorname{Row}(L_\Delta)^4.
\]
The regular \fid{F21} branch identification makes this
the actual tensor lattice with the same named $P_i$.
Its rank is $20+4+32+8=64$.
The named-action identification uses the explicit regular
comparison, not only the dimension.

Over the domain $S$, maps between the distinct branch
targets vanish because $k\Delta$ is a nonzero divisor.
Their multiplicity matrix units give
\[
 \End_F(K)=M_{20}(S)\times M_4(S)\times
          M_{16}(S)\times M_4(S)
          \cong\End_{A_S}(M_S)
\]
by postcomposition.
This retains the actual ordinary left operators.
At collision, the two target branch sectors merge,
and the two endomorphism algebras acquire the same
cross-branch maps. Their dimension becomes
\[
 20^2+4^2+20^2=816.
\]
The invariant Hom has identity conjugation Frobenius,
as it does in every such invariant-Hom construction.

\par\smallskip\noindent\textbf{The two genuine traces.}\quad
At the base $2$, the eigenvalues $1,2,4,8$ have
multiplicities $1,1,2,2$ on $H$ and $20,4,16,4$ on $K$.
Equations \eqref{eq:traceescapeH}--\eqref{eq:traceescapevalues}
give \eqref{eq:twogenuinetraces} for every $n$.
All multiplicities are genuine and positive.

The spectra themselves fail to retune already at $n=2$:
$a(4)=-1$, whereas every eigenvalue of $F_H(2)^2$ is
positive.
At $Q=4$ the base weights are $0,1,2,3$.
This does not impose those weights on the entire family.
For example $a(8)=-5$ has modulus $5$, not $8^w$ for
any integer $w$.
The model therefore violates neither the single-radius
result of Part~1 nor the integral-weight spectral
comparison of \fid{F51}.
\end{proof}

\paragraph{What the example settles.}
The whole actual relative order, its named tensor action,
the ordinary left operators, and two identity-trace systems
can coexist without full spectral retuning.
The example uses multiple spectral radii and does not
supply a pure family with the arithmetic parameter law.
It does not impose marked-operator traces, a path-sign
comparison, or the canonical integral forms of VIII.
The next part tests those additional trace requirements.

\subsection{Marked traces and the actual generator relations}
The trace-only construction leaves open the arithmetic
comparison for specified operators. A regular central
insertion already detects the collision in pure or
integrally mixed coefficient models. Requiring composition
for just two generator words gives an obstruction without
any purity or algebraicity assumption.

\begin{finding}[\fid{F54}, Part 3: marked-trace comparison gates]
\label{f:markedtracegates}
Let $b>1$, and use the actual rank-$64$ tensor module $M$
and named generators $P_1,P_2$.
Put
\[
 Z=\frac{P_1+P_2}{D}
       -\frac{P_1P_2+P_2P_1}{D^2},\qquad
 z_1=\frac{D-1}{D},\quad z_2=\frac{4(D-1)}{D^2},
\]
\begin{equation}\label{eq:actualmarkedX}
 X=\frac{Z(z_1-Z)}{z_2}.
\end{equation}
This is a regular central operator at $v=1$.
It acts by $\Delta$ on $N(D-2)^4$, a rank-$8$ sector,
and by zero on the other sectors; in particular $X(1)=0$.

Suppose $F$ is an invertible algebraic matrix family on $M$,
regular at a place above $v=1$, commuting with the actual
action, and with regular evaluations at $b$ and all
sufficiently large $b^n$.
The marked trace identity
\begin{equation}\label{eq:markedXretuning}
 \Tr(X_{b^n}F_{b^n})=\Tr(X_bF_b^n)
\end{equation}
is impossible if the eight base eigenvalues on the
$\Delta$-sector either have one common absolute value
or have several integral weights
$|\alpha_j|=b^{w_j}$, $w_j\in\Z$.
No full-centralizer comparison is needed for this assertion.

Alternatively, suppose $F$ acts on a regular auxiliary
family $H$ whose full relative centralizer is the actual
order, with the insertion $X$ corresponding to its
specified central coordinate, and with the same algebraicity,
invertibility and regular-evaluation hypotheses.
Then \eqref{eq:markedXretuning} is impossible without
any weight restriction.

Finally, let $M_n$ be the actual module at $v=b^n$, let
$P_n$ be its named $P_1$, and let $F_n$ be arbitrary
invertible operators, with $F_1P_1=P_1F_1$ at the base.
The two identities
\begin{equation}\label{eq:twowordretuning}
 \Tr(P_n^jF_n)=\Tr(P_1^jF_1^n),\qquad j=1,2,
\end{equation}
cannot both hold for every sufficiently large $n$.
Here the right insertions are fixed actual base operators,
and the second is genuinely the square of the first.
This last assertion needs no algebraicity, purity,
full-centralizer recovery, or collision specialization.
\end{finding}

\begin{proof}
\par\smallskip\noindent\textbf{The marked central element.}\quad
In the actual algebra set $e=P_1/D$, $f=P_2/D$.
These are idempotents, and
\[
 Z=e+f-ef-fe.
\]
A direct multiplication gives
\[
 Ze=eZ=e-efe,\qquad Zf=fZ=f-fef.
\]
Thus $Z$ and $X$ are central.
In the row-action convention, the actual matrices on $N(a)$
are
\[
 P_1=\begin{pmatrix}D&0\\a&0\end{pmatrix},\qquad
 P_2=\begin{pmatrix}0&a\\0&D\end{pmatrix}.
\]
They give
\[
 Z=(1-a^2/D^2)I.
\]
On both scalar sectors $Z=0$.
Since $Z=z_1$ on $N(\delta)$ and $Z=z_2$ on $N(D-2)$,
\begin{equation}\label{eq:actualXbranches}
 X|_{\epsilon_\pm}=0,\qquad
 X|_{N(\delta)}=0,\qquad
 X|_{N(D-2)}=\Delta I.
\end{equation}
The denominators $D,z_2$ are units at $v=1$ and on
the previously specified base.
Equivalently, if
\[
 B=D(P_1+P_2)-(P_1P_2+P_2P_1),
\]
then
\[
 X=\frac{B((D-1)D-B)}{4(D-1)D^2}.
\]
This displays a regular rational expression in the named
generators. It is not an unspecified observable.

Under the actual coherent algebra model, $X$ is zero
on both scalar factors and the coordinate $x$ on
$M_2(S[x]/(x(x-\Delta)))$.
On the actual module, the regular \fid{F21} splitting
makes \eqref{eq:actualXbranches} an integral action identity.
Its $\Delta$-sector has rank $2\cdot4=8$, and $X(1)=0$
as a matrix, not merely at the level of eigenvalues.

\par\smallskip\noindent\textbf{Three scalar consequences.}\quad
Let $g$ be one algebraic function over $\C(v)$,
evaluated at $b^n$ through places of that same relation.

First, if
\[
 g_{b^n}=C\alpha^n
\]
eventually, where $C,\alpha\ne0$, then for some
$L>0$, $m\in\Z$,
\begin{equation}\label{eq:singleexponentialalgebraic}
 g(v)^L=C^L v^m.
\end{equation}
Indeed, expand its irreducible polynomial relation
$P(b^n,C\alpha^n)=0$.
This is a finite power sum on roots $b^i\alpha^j$.
Vandermonde uniqueness forces two distinct monomials
to have the same root. Their $j$ exponents differ,
so $\alpha^L=b^m$ for some $L,m$.
On $n=Lr$, substitute $v=t^L$ to obtain
$P(t^L,Ct^m)=0$ at infinitely many $t=b^r$.
It is therefore a rational identity.
The irreducible polynomial of $g$ has the root $Ct^m$,
so it divides $Y^L-C^Lv^m$ over $\C(v)$.
This proves \eqref{eq:singleexponentialalgebraic}.
In particular $g$ cannot vanish at any place over $v=1$.

Second, if
\[
 g_{b^n}=\sum_{j=1}^r\alpha_j^n
\]
with genuine multiplicities and a common modulus,
Part~1 gives a common $L>0$ and one integer $m$ such that
$\alpha_j^L=b^m$ for all $j$.
At degrees $Lr'$ the trace is $r b^{mr'}$.
Hence the algebraic function $g^L-r^L v^m$ vanishes
at places over infinitely many distinct parameters.
The nonzero-constant-term argument for its minimal
polynomial, proved in Part~1, forces
\begin{equation}\label{eq:markedcommonradius}
 g(v)^L=r^L v^m.
\end{equation}
Again $g$ cannot vanish above $1$.
This consequence uses no integral-weight assumption
and no continuity of the evaluation sheets.

Third, suppose the base roots have integral weights:
\[
 \alpha_j=b^{w_j}\zeta_j,\qquad
 w_j\in\Z,\quad |\zeta_j|=1.
\]
The mixed-integral trace argument of \fid{F51} gives
\begin{equation}\label{eq:markedintegraltrace}
 g(v)=\sum_{j=1}^r v^{w_j},\qquad g(1)=r.
\end{equation}
To recall the exact deduction, expand
$P(b^n,\sum_j b^{w_jn}\zeta_j^n)$ and group by the
integer power of $b^n$.
Each remaining coefficient is a finite power sum
on unit-modulus roots.
Distinct radii and Vandermonde uniqueness make every
coefficient sequence zero, including at $n=0$.
Those identities are the Laurent coefficients of
$P(v,\sum_jv^{w_j})=0$.
Irreducibility over $\C(v)$ makes this rational root
the algebraic function itself.
Negative weights are allowed in $\C(v)$.
These arguments keep genuine positive multiplicities;
they do not restore eigenvalues canceled in a virtual trace.

\par\smallskip\noindent\textbf{The actual tensor-module gate.}\quad
Since $F_b$ commutes with the actual action,
the rank-$8$ image of $X_b/\Delta(b)$ is invariant
and its restriction is invertible.
Let its eigenvalues be $\alpha_1,\ldots,\alpha_8$,
including algebraic multiplicities.
The algebraic function
\[
 g(v)=\Delta(b)^{-1}\Tr_M(X(v)F(v))
\]
would satisfy
\[
 g_{b^n}=\sum_{j=1}^8\alpha_j^n
\]
under \eqref{eq:markedXretuning}.
But regularity and $X(1)=0$ give $g(1)=0$.
Equation \eqref{eq:markedcommonradius} contradicts this
if the branch spectrum has one common modulus;
\eqref{eq:markedintegraltrace} gives the contradiction
$g(1)=8$ in the integral-mixture case.

A pure Weil spectrum supplies the common-radius hypothesis
at any chosen complex embedding.
No full point-Weil endomorphism algebra has been identified
with the actual algebra in this argument.
The same reasoning works when $X(1)$ is only nilpotent,
provided it commutes with $F(1)$:
their product is then nilpotent and its trace is zero.
For unrestricted multiple radii with nonintegral weights,
neither scalar conclusion is asserted.
For example, on $v=t^2$ the function $g=1+t$ has values
$1+(\sqrt b)^n$ on the positive sheets over $b^n$,
but vanishes at $t=-1$ over $v=1$.
This is a scalar boundary example, not an actual-action
or full-order realization.

\par\smallskip\noindent\textbf{The full relative-centralizer gate.}\quad
Now suppose $\End_F(H)$ is the actual regular order
and the central marking is preserved.
The operator $F$ is central in its full centralizer,
so its component in the matrix-center factor is
$c_0+c_1x$.
Its value on the $x=\Delta$ branch is
$d=c_0+c_1\Delta$.

Over the generic field the actual algebra is split
semisimple, with two scalar and two two-by-two matrix
factors.
A central element is a tuple of scalars, so $F$ is
diagonalizable.
Distinct primitive central blocks have distinct eigenvalues;
otherwise maps between equal-eigenvalue blocks would
enlarge the full centralizer.
Each two-by-two factor is thus the full endomorphism
algebra of a two-dimensional eigenspace.
It follows that
\[
 g(v)=\Tr_H(X(v)F(v))=2\Delta(v)d(v).
\]
At the base,
\[
 \Tr_H(X_bF_b^n)=2\Delta(b)d(b)^n.
\]
The constant $2\Delta(b)$ and the eigenvalue $d(b)$
are nonzero.
If \eqref{eq:markedXretuning} held, the single-exponential
identity \eqref{eq:singleexponentialalgebraic} would make
$g$ a unit at every point over $1$.
Its regular expression $2\Delta d$ instead vanishes there.
This proves the assertion without any weight assumption.

The trace-only model of Part~2 has
$g(v)=2\Delta(v)d(v)$ and $d(2)=8$.
The extra comparison would require
$g(2^n)=2\Delta(2)2^{3n}$.
It fails already at $v=4$, where the exact difference is
\[
 -\frac{197896315512672}{4359848400625}.
\]
Thus that model gives a concrete failure of the marked
axiom while preserving both identity traces.

\par\smallskip\noindent\textbf{Two actual inserted words.}\quad
Finally use the actual relation
\[
 P_n^2=D(b^n)P_n
\]
and the fixed base relation $P_1^2=D(b)P_1$.
If both identities in \eqref{eq:twowordretuning} held,
they would imply
\[
 \bigl(D(b^n)-D(b)\bigr)\Tr(P_1F_1^n)=0.
\]
The function $(v+v^{-1})^2$ is strictly increasing
for $v>1$, so the first factor is nonzero for $n>1$.
Hence $\Tr(P_1F_1^n)$ must vanish eventually.

But $P_1/D(b)$ is an idempotent of rank $40$:
the sectors contribute $20+16+4$.
Its image is $F_1$-invariant because the two commute,
and the restricted $F_1$ is invertible.
If its eigenvalues are $\lambda_i$, with genuine
multiplicities, then
\[
 \Tr(P_1F_1^n)=D(b)\sum_i\lambda_i^n.
\]
Grouped distinct nonzero roots have positive integer
coefficients.
Vandermonde uniqueness forbids their vanishing for
every sufficiently large $n$.
This is the required contradiction.
No commutation hypothesis at the other parameters
was needed for this last calculation.
\end{proof}

\paragraph{Both actual actions remove the weight hypothesis.}
In the actual-module assertion of Part~3, the spectral
hypotheses can be replaced by commutation of $F_b$
with the four named ordinary left operators
$E_V,F_V,E_W,F_W$.
Only commutation at the base is needed.
The regular algebraic family, right-action commutation,
and marked-trace assumptions remain as stated.

\begin{proof}
The complete actual bimodule decomposition in \fid{F43}
identifies the $\Delta$-sector as
\[
 (W_1^V\otimes W_1^W)\otimes N(D(b)-2).
\]
The two $W_1$ factors each have dimension two.
Their labels are $\mathfrak{sl}_2$ highest weights;
the polynomial $GL_2$ labels remain $(2,1)$, with
their determinant twists and total degree three.

On the actual regular basis, each ordinary pair acts by
\[
 Eg_-=g_+,\quad Eg_+=0,\qquad
 Fg_+=g_-,\quad Fg_-=0.
\]
These are the two off-diagonal matrix units, and their
products give both diagonal units.
Thus the two named left actions generate $M_4(\C)$
on the multiplicity factor.

On the right factor, put $a=D(b)-2$.
The established basis change gives
\[
 e=P_1/D=\begin{pmatrix}1&0\\0&0\end{pmatrix},\qquad
 f=P_2/D=
 \begin{pmatrix}1-z_2&1\\z_2(1-z_2)&z_2\end{pmatrix}.
\]
Since $b>1$, $D>4$ and $0<z_2<1$.
The elements
\[
 ef(1-e)=E_{12},\qquad
 (1-e)fe=z_2(1-z_2)E_{21},
\]
together with $e,1-e$, generate $M_2(\C)$.
The joint actual left/right image is therefore
$M_4(\C)\otimes M_2(\C)=M_8(\C)$.
Its commuting $F_b$ is a scalar $\beta I_8$, with
$\beta\ne0$.

Consequently the marked trace would give
\[
 g_{b^n}=8\Delta(b)\beta^n,\qquad
 g(v)=\Tr_M(X(v)F(v)).
\]
The single-exponential lemma
\eqref{eq:singleexponentialalgebraic}
makes $g$ a unit at every point above $v=1$.
The regular actual action instead has $X(1)=0$,
so $g(1)=0$, a contradiction.
\end{proof}

This corollary asks the named left operators themselves
to commute with Frobenius.
It does not apply to coefficient-twisted operators
satisfying covariance rather than commutation, as
distinguished in \fid{F48}.
An unrestricted fractional mixed marked-trace construction
commuting only with the right action can therefore evade
this corollary by not preserving those named ordinary
left morphisms.


\paragraph{The comparison axiom is part of the conclusion.}
The left insertion in \eqref{eq:markedXretuning} is the
actual $X$ at $b^n$; the right insertion is the fixed
$X_b$.
Its trace law is additional data and does not follow
from identity-trace equality.
Although $X$ is a regular expression in named generators,
its parameter-dependent rational coefficients cannot
silently be moved through a fixed-operator comparison.

The two-word obstruction strengthens \fid{F42} to a
finite test under an explicit action and composition
requirement.
It does not exclude a dictionary in which the formal
coefficient $D(v)$ is represented by a graded coefficient
object, and the quadratic relation becomes a tensor
relation with that object.
Such a dictionary does not impose the fixed scalar
$D(b)$ used on the right of \eqref{eq:twowordretuning}.
The separate word objects of \fid{F45}, the coefficient
kernels of \fid{F43}, and the fixed-parameter constructions
of \fid{F50} and \fid{F53} remain outside that obstruction.
The remaining comparison concerns the designated
arithmetic, grading and basis structures.

\subsection{An actual fractional mixed trace model}
The marked-trace gate in \fid{F54} has a genuine
one-place exception when arbitrary mixed radii are allowed.
The following construction retains the actual right module
and two specified trace laws. Its failure of stronger
comparisons can be checked on explicit operators.

\begin{finding}[\fid{F55}, Part 1: a regular one-place model]
\label{f:fractionalescape}
There is one finite localization of the algebraic cover
$v=t^2$ over $\Q(\sqrt2)$, containing $t=-1$ and all points
$t_n=2^{n/2}$, on which the actual rank-$64$ tensor module
has an invertible regular operator $F$ commuting with its
actual right action. With base $t_1=\sqrt2$, it satisfies
for every $n\ge1$
\begin{equation}\label{eq:fractionaltwotraces}
 \Tr F(t_n)=\Tr(F(\sqrt2)^n),\qquad
 \Tr\bigl(X(2^n)F(t_n)\bigr)
       =\Tr\bigl(X(2)F(\sqrt2)^n\bigr),
\end{equation}
where $X$ is the actual marked central element
\eqref{eq:actualmarkedX}.
The family is regular at the chosen collision $t=-1$,
where $v=1$, and retains its order-two collision.
It does not commute with all named ordinary left
operators, does not have full spectral retuning, and
fails the fixed $P_1$ and $P_1^2$ trace comparisons at
every $n>1$. Its operator entries have unavoidable poles
at the other place $t=+1$ above $v=1$.
\end{finding}

\begin{proof}
\par\smallskip\noindent\textbf{The actual summands.}\quad
Use the regular \fid{F21}/\fid{F43} bimodule isomorphism
\begin{align}
 M={}&(W_3\otimes W_3\oplus W_1\otimes W_1)
                          \otimes\epsilon_+
       \ \oplus\ (W_1\otimes W_1)\otimes\epsilon_-
       \notag\\
 &\oplus(W_3\otimes W_1\oplus W_1\otimes W_3)
                          \otimes N(\delta)
       \ \oplus\ (W_1\otimes W_1)\otimes N(D-2).
 \label{eq:fractionalactualmodule}
\end{align}
The right-isotypic multiplicities are $20,4,16,4$,
and their actual dimensions are $20,4,32,8$.
In the final summand the multiplicity space is exactly
$W_1^V\otimes W_1^W$.
Here $W_1$ denotes the $\mathfrak{sl}_2$-highest-weight-one
factor of the polynomial $\GL_2$ representation of
partition $(2,1)$, with its determinant twist retained.

Put
\[
 \delta=v+v^{-1},\qquad D=\delta^2,\qquad
 \Delta=\frac{(D-1)(D-4)}{D^2}.
\]
The element $X$ acts by $\Delta$ on the last rank-eight
summand and by zero on the others; on the actual regular
tensor lattice, $X(1)=0$.
The preceding bimodule comparison is defined over
\[
 S_0=\Q[v,v^{-1},\delta^{-1},(D-1)^{-1},(D-2)^{-1}].
\]
We pull back this actual comparison.
The rank-eight summand is selected by its regular
commutant decomposition, without dividing $X$ by
$\Delta$ at the collision.

\par\smallskip\noindent\textbf{One coefficient cover.}\quad
Set
\[
 s=\sqrt2>0,\qquad K_0=\Q(s),\qquad v=t^2,\qquad
 \eta=\Delta(2)=\frac{189}{625}.
\]
The base is $t=s$, the extension points are $t_n=s^n$,
and the chosen collision is $t=-1$.
Define
\begin{equation}\label{eq:fractionalfunctions}
 g(t)=2(1+t)^2,\qquad
 h(t)=\frac{\eta g(t)}{\Delta(t^2)},\qquad
 \lambda(t)=\frac{h(t)}2-1-2s,\qquad
 u(t)=1+\frac{g(t)-h(t)}{20}.
\end{equation}
The cancellation at $t=-1$ is exact:
\[
 \Delta(t^2)=
 \frac{(t-1)^2(t+1)^2(t^2+1)^2(t^8+t^4+1)}
      {(t^4+1)^4},
\]
\begin{equation}\label{eq:fractionalhcancellation}
 h(t)=\frac{2\eta(t^4+1)^4}
 {(t-1)^2(t^2+1)^2(t^8+t^4+1)}.
\end{equation}
In particular,
\[
\begin{array}{c|ccc}
 t&h(t)&u(t)&\lambda(t)\\ \hline
 s&6+4s&1&2\\
 -1&126/625&6187/6250&-562/625-2s .
\end{array}
\]
At the remaining small extension point,
\[
 h(2)=\frac{1503378}{203125},\qquad
 \lambda(2)=\frac{548564}{203125}-2s\ne0.
\]
The latter is nonzero because its first term is rational
and $s$ is irrational.

Let $R_0$ be the pullback of $S_0$ to $K_0[t,t^{-1}]$,
further localized at $t-1$ and $t^2+1$, and put
\begin{equation}\label{eq:fractionalbase}
 R=R_0[u^{-1},\lambda^{-1}].
\end{equation}
The expression for $h$ belongs to $R_0$, since
\[
 D(t^2)-1=\frac{t^8+t^4+1}{t^4}.
\]
This is one fixed finite localization of the coefficient
cover. Every inherited factor is nonzero at $v=1$
and at real $v>1$. The two new denominator factors
are nonzero at $t=-1$ and at every $t_n>1$.
The displayed values settle $u,\lambda$ at $t=-1$
and $\lambda$ at $n=1,2$.

For $n\ge3$, $t_n\ge2s$ and $0<\Delta(2^n)<1$, so
\[
 \lambda(t_n)>
 \eta(1+2s)^2-1-2s
 =\frac{1076-494s}{625}>0,
\]
where $s<2$ already gives a numerator greater than $88$.
To check $u$ at all $n>1$, $D(v)$ is strictly increasing
for $v>1$, and
\[
 \frac{d}{dD}\left(1-\frac5D+\frac4{D^2}\right)
      =\frac{5D-8}{D^3}>0\qquad(D>4).
\]
Thus $\Delta(2^n)>\eta$, whence
\begin{equation}\label{eq:fractionalmonotonicity}
 h(t_n)<g(t_n),\qquad u(t_n)>1\quad(n>1).
\end{equation}
All required points survive \eqref{eq:fractionalbase}.
At $t=-1$, $v-1=(t-1)(t+1)$ has a simple zero,
so $\Delta(t^2)$ still has a zero of order exactly two.
The actual coherent center and tensor module are
their original pullbacks at this place.

\par\smallskip\noindent\textbf{The actual automorphism.}\quad
On the entire rank-$20$ $\epsilon_+$ block in
\eqref{eq:fractionalactualmodule}, let $F=uI_{20}$.
This includes both the $W_3\otimes W_3$ and
$W_1\otimes W_1$ multiplicity summands of that scalar
right type.
On $\epsilon_-^4\oplus N(\delta)^{16}$, of total
dimension $36$, let $F$ be the identity.
On the final block, in the ordered multiplicity basis
\[
 g_+\otimes g_+,\quad g_+\otimes g_-,\quad
 g_-\otimes g_+,\quad g_-\otimes g_-,
\]
put
\begin{equation}\label{eq:fractionaloperator}
 F=\operatorname{diag}(1,s,s,\lambda(t))\otimes I_2.
\end{equation}
Every entry is regular over $R$, and
\[
 \det F=u^{20}(s^2\lambda)^2=4u^{20}\lambda^2
          \in R^\times.
\]
If $\iota$ is the regular $A$-linear isomorphism
from the original tensor module to
\eqref{eq:fractionalactualmodule}, the operator on the
original tensor module is explicitly
$\iota^{-1}F_{\rm block}\iota$, with the rightmost map
applied first.
Thus the comparison retains the actual named right
generators; it does not substitute a trace-equivalent
module. On the nonscalar blocks $F$ acts on the
multiplicity factor and the right action on $N(a)$.
It therefore commutes with that right action.
No division by $\Delta$ enters its definition.

\par\smallskip\noindent\textbf{The two trace identities.}\quad
On the final rank-eight sector,
\[
 \Tr(F|_{M_2})=2(1+s+s+\lambda)=h.
\]
Consequently
\begin{equation}\label{eq:fractionalalgebraictraces}
 \Tr_M(X(t^2)F(t))=\Delta(t^2)h(t)=\eta g(t),\qquad
 \Tr_M F(t)=20u(t)+36+h(t)=56+g(t).
\end{equation}
At $t=s$ the full base spectrum is
\[
 1^{\,58},\qquad s^{\,4},\qquad 2^{\,2}.
\]
On the marked sector the respective multiplicities
are $2,4,2$, and $X$ has eigenvalue $\eta$.
As $s^2=2$,
\[
 \Tr(F(s)^n)=58+4s^n+2\cdot2^n=56+2(1+s^n)^2,
\]
\[
 \Tr(X(2)F(s)^n)
       =\eta(2+4s^n+2\cdot2^n)=\eta\,2(1+s^n)^2.
\]
Substitution in \eqref{eq:fractionalalgebraictraces}
proves \eqref{eq:fractionaltwotraces} for every $n$.

\par\smallskip\noindent\textbf{Ordinary left action.}\quad
The actual \fid{F43} operator has
$E_V(g_-\otimes g_+)=g_+\otimes g_+$ on the final
multiplicity space. Hence
\[
 (F(s)E_V-E_VF(s))(g_-\otimes g_+)
       =(1-s)(g_+\otimes g_+)\ne0.
\]
Tensor with any nonzero right-factor vector.
This uses composition of actual operators on vectors;
the source's row-matrix/opposite-composition convention
is unchanged.
Thus ordinary left invariance fails explicitly, as
required by the both-actions gate in \fid{F54}.

\par\smallskip\noindent\textbf{Fixed generator words and spectra.}\quad
On $N(a)$, $\Tr P_1=D$.
On $\epsilon_+$ it is $D$ and on $\epsilon_-$ it is zero.
Therefore
\[
 \Tr(P_1(v)F(t))
 =D(v)\left(20u(t)+16+\frac{h(t)}2\right)
 =D(v)\left(36+g(t)-\frac{h(t)}2\right).
\]
At the base,
\[
 \Tr(P_1(2)F(s)^n)
  =D(2)(37+2s^n+2^n)
  =D(2)\left(36+\frac{g(t_n)}2\right).
\]
By \eqref{eq:fractionalmonotonicity} and
$D(2^n)>D(2)$, the first expression at $t_n$
is strictly larger for every $n>1$.
Multiplying by the respective $D$ factors, using
$P_1^2=DP_1$, gives the same strict failure for $P_1^2$.

Full spectral retuning also fails without retaining
any sector labels. The operator $F(s)^n$ has eigenvalue
$1$ with multiplicity exactly $58$.
For $n>1$, $F(t_n)$ has that eigenvalue with multiplicity
at most $40$: the $36$ unchanged dimensions, the two
fixed marked dimensions, and at most two additional
dimensions if $\lambda(t_n)=1$.
Here $u(t_n)>1$ and $s>1$ exclude all other sources.

\par\smallskip\noindent\textbf{The other place above one.}\quad
At $t=+1$, \eqref{eq:fractionalhcancellation} has a pole
of order two: its numerator does not vanish, and its
only vanishing denominator factor is $(t-1)^2$.
Both $u$ and $\lambda$ have poles of order two, with
respective leading terms $-h/20$ and $h/2$.
They are actual eigenvalue functions.
No change of finite free lattice can make $F$ regular
there: a matrix over the valuation ring has integral
eigenvalues, whereas these have negative valuation.
The finite localization explicitly removes $t=+1$
while retaining $t=-1$.
\end{proof}

The construction asserts neither a full-centralizer
identification nor all ordinary left morphisms.
Under $v=q^{1/2}$ with base $q=4$, the marked base
eigenvalues have weights $0,\tfrac12,1$ and different
absolute values. It therefore lies outside both the
common-radius and integral-mixed-weight hypotheses
of the marked-module gate.
It is an algebraic matrix family, not an assertion
of an ordinary geometric mixed object with integral
Deligne weights.
Its regularity at one place is also not a comparison
with VIII's prescribed global integral form or a
continuous choice along the positive real parameter
axis. The ordinary trace and marked trace remain
regular at $t=+1$ through the cancellations in
\eqref{eq:fractionalalgebraictraces}, even though the
actual operator has poles there.

\subsection{Regularity at every collision place}
The preceding exception is excluded by a stronger
regularity comparison. The following operator argument
requires neither commutation with the actual action
nor restrictions on the base spectrum.

\begin{finding}[\fid{F55}, Part 2: the all-places inserted-trace gate]
\label{f:allplacestrace}
Let $\mathcal L/\C(v)$ be a finite coefficient field
for an algebraic matrix family $F(v)$ on the actual
rank-$64$ tensor module.
Assume regular evaluations at $b>1$ and all
sufficiently large $b^n$, with $F_b$ invertible.
If
\[
 \Tr(X_{b^n}F_{b^n})=\Tr(X_bF_b^n)
\]
for all sufficiently large $n$, then $F$ cannot be
regular on the actual pulled-back tensor lattice
at every place of $\mathcal L$ above $v=1$.
Here $X$ is the actual marked central operator,
with $X(1)=0$ on that lattice.

No commutation of $F_b$ with $X_b$, or of $F$ with
the actual right action, is needed.
The result allows arbitrary Jordan form and imposes
no common-radius, weight, or positivity hypothesis.
\end{finding}

\begin{proof}
We prove a general inserted-trace principle.

\par\smallskip\noindent\textbf{Polynomial identities propagate to index zero.}\quad
Let $W$ be a nonzero finite-dimensional complex vector
space, let $F_0$ be invertible and $X_0$ arbitrary, and set
\[
 S_n=\Tr(X_0F_0^n)\qquad(n\ge0).
\]
For any polynomial
$Q(v,Y)=\sum_{i,j}q_{ij}v^iY^j$, if $Q(b^n,S_n)$
vanishes for all sufficiently large $n$, then
\begin{equation}\label{eq:insertedtracezeroindex}
 Q(1,\Tr X_0)=0.
\end{equation}
For the nonzero polynomial case, on the finite direct sum
\[
 W_Q=\bigoplus_{q_{ij}\ne0}W^{\otimes j}
\]
put
\[
 T_Q=\bigoplus_{q_{ij}\ne0}b^iF_0^{\otimes j},
 \qquad
 U_Q=\bigoplus_{q_{ij}\ne0}q_{ij}X_0^{\otimes j},
\]
with zeroth tensor power equal to $\C$ and its
operators equal to $1$.
The operator $T_Q$ is invertible.
Tensor-product trace gives
\[
 z_n:=\Tr(U_QT_Q^n)=Q(b^n,S_n).
\]
No commutation is used: each block product is
$q_{ij}b^{in}(X_0F_0^n)^{\otimes j}$.

Write the characteristic polynomial of $T_Q$ as
$T^r+c_{r-1}T^{r-1}+\cdots+c_0$, with $c_0\ne0$.
Cayley--Hamilton gives
\[
 z_{n+r}+c_{r-1}z_{n+r-1}+\cdots+c_0z_n=0
 \qquad(n\ge0).
\]
Eventual vanishing therefore propagates backwards
one index at a time to $z_0=0$, proving
\eqref{eq:insertedtracezeroindex}.
This includes Jordan terms without choosing a
Jordan basis or asserting effectiveness of their traces.

\par\smallskip\noindent\textbf{A polynomial with prescribed reduction.}\quad
Let $f\in\mathcal L$ be regular with common residue
$c\in\C$ at every place above $v=1$.
Put $O=\C[v]_{(v-1)}$ and let $B$ be its integral closure
in $\mathcal L$.
It is finite over $O$ by the characteristic-zero
finite-normalization theorem
(Stacks Lemma~10.161.8,
\href{https://stacks.math.columbia.edu/tag/032L}{tag~032L}),
and finite free since it is torsion-free over this DVR
(Stacks Lemma~15.22.11(3),
\href{https://stacks.math.columbia.edu/tag/0AUW}{tag~0AUW}).
Its maximal localizations are the DVR places above one,
by lying over and the dimension-one normal-local criterion
(Stacks Lemmas~10.36.17 and~10.119.7,
\href{https://stacks.math.columbia.edu/tag/00GQ}{tag~00GQ}
and
\href{https://stacks.math.columbia.edu/tag/00PD}{tag~00PD}).

Regularity in all these localizations puts $f$ in $B$.
Indeed, if the denominator ideal
\[
 I=\{a\in B:af\in B\}
\]
were proper, choose a maximal ideal $\mathfrak m$
containing it. Membership $f\in B_{\mathfrak m}$
would give $a\notin\mathfrak m$ with $af\in B$,
a contradiction.

The finite $\C$-algebra $\overline B=B/(v-1)B$
has residue $\overline f=c$ at every maximal ideal.
Multiplication by $\overline f-c$ cannot have a
nonzero eigenvalue $\lambda$.
If it did, multiplication by $\overline f-c-\lambda$
would be singular, so this element would be a nonunit
and lie in a maximal ideal.
Its residue there is $-\lambda\ne0$, a contradiction.
Here multiplication is bijective exactly when its
multiplying element is a unit, and all residue fields
of $\overline B$ are $\C$.
Thus multiplication by $\overline f-c$ is nilpotent.

Choose an $O$-basis of $B$ and let $H(v,Y)\in O[Y]$
be the monic characteristic polynomial of multiplication
by $f$. Its degree $d=[\mathcal L:\C(v)]$ is positive,
and
\[
 H(1,Y)=(Y-c)^d.
\]
Cayley--Hamilton applied to $1\in B$ gives $H(v,f)=0$.
Clear its rational coefficient denominators by one
polynomial $a(v)$ with $a(1)\ne0$.
Then $Q=aH\in\C[v,Y]$ satisfies
\begin{equation}\label{eq:allplacespolynomial}
 Q(v,f)=0,\qquad Q(1,Y)=a(1)(Y-c)^d.
\end{equation}

Consequently, whenever the selected evaluations of
this same algebraic element satisfy
\[
 f_{b^n}=\Tr(X_0F_0^n)
\]
for all sufficiently large $n$,
\eqref{eq:insertedtracezeroindex} and
\eqref{eq:allplacespolynomial} force
\begin{equation}\label{eq:allplaceszerothtrace}
 c=\Tr X_0.
\end{equation}

\par\smallskip\noindent\textbf{The actual marked element.}\quad
Apply this principle to
\[
 f(v)=\Tr_M(X(v)F(v)),\qquad X_0=X_b,\quad F_0=F_b.
\]
All-places regularity of $F$ and the actual identity
$X(1)=0$ make $f$ regular with common residue $0$.
But $X_b$ acts by $\Delta(b)$ on its rank-eight
marked sector and by zero on the others, so
\[
 \Tr X_b=8\Delta(b)\ne0\qquad(b>1).
\]
This contradicts \eqref{eq:allplaceszerothtrace}.
Neither invariance of that sector under $F_b$ nor
commutation with the right action was used.
\end{proof}

The same principle covers finite alternating inserted
traces: use a direct-sum base operator and put the
alternating signs into the insertion.
Zero residue at every place then requires
$\sum_j(-1)^j\Tr X_j=0$.
If that zeroth inserted Euler trace is zero,
this obstruction says nothing; it does not turn
a virtual trace into an effective spectrum.

Base invertibility is essential.
For $F_0=0$ and $X_0=I$ on a nonzero space,
the inserted traces vanish for all $n\ge1$ but their
zeroth trace is the dimension.
Actual finite Weil spaces have invertible Frobenius.
If $X$ specializes only to a nilpotent, one must
separately verify that $\Tr(XF)$ has zero residue;
commutation is one sufficient condition in that variant.

The explicit model in Part~1 has poles in its operator
at $t=+1$ and therefore fails precisely the all-places
regularity hypothesis.
Its regularity at $t=-1$, actual right-action
compatibility and two trace laws remain valid.
All-places regularity is extra comparison data;
it is not supplied automatically by an algebraic
coefficient scheme or a perverse realization.
\paragraph{A separate genuine-trace place observation.}
For genuine positive-multiplicity traces, one can also locate a nonzero place without assuming whole-fiber regularity. The earlier argument is retained here because its hypotheses differ from the stronger inserted-trace principle.

Let $f$ be any algebraic function over $\C(v)$
whose selected regular values satisfy
\[
 f_{b^n}=\sum_i m_i\alpha_i^n
 \quad\text{for all sufficiently large }n,\qquad
 b>1,\quad m_i\in\Z_{>0},\quad\alpha_i\ne0.
\]
Put $N=\sum_i m_i$.
By \fid{F54}, Part~1, there are $L>0$ and integers
$a_i$ with $\alpha_i^L=b^{a_i}$.
Let $P(v,Y)$ be the irreducible minimal polynomial of
$f$ over $\C(v)$, and put
\[
 p(t)=\sum_i m_i t^{a_i}.
\]
At the indices $n=Lr$, the selected trace values
are $p(b^r)$. The rational function $P(t^L,p(t))$
vanishes at infinitely many $t=b^r$, away from its
finitely many possible poles. It is identically zero.
The assignments
\[
 v\longmapsto t^L,\qquad f\longmapsto p(t)
\]
therefore give a field embedding
$\C(v,f)\hookrightarrow\C(t)$.
Restrict the $t=1$ place to this field.
Since $t^L-1$ has a simple zero at $1$, it is a place
over $v=1$ with value group $\Z$.
The function $f$ is regular there and has residue
$p(1)=N$.
This proves existence of one such place, not that
every place has this residue.



\subsection{Following the same algebraic branch to the collision}
The one-place construction distinguishes the chosen collision point
from the positive arithmetic branch. A continuous choice of that
same scalar branch gives a complementary obstruction.

\begin{finding}[\fid{F55}, Part 3: continuous genuine trace branches]
\label{f:continuoustracebranch}
Let $b>1$ and let $f:[1,\infty)\to\C$ be continuous.
Suppose one monic irreducible polynomial
$P(v,Y)\in\C(v)[Y]$, of positive $Y$-degree, satisfies
$P(v,f(v))=0$ wherever its coefficients are defined.
Suppose that, for every sufficiently large integer $n$,
\[
 f(b^n)=\sum_i m_i\alpha_i^n,
 \qquad m_i\in\Z_{>0},\quad \alpha_i\in\C^*.
\]
Then there are rational numbers $r_i$ such that
\begin{equation}\label{eq:continuoustracepositive}
 \alpha_i=b^{r_i},\qquad
 f(v)=\sum_i m_i v^{r_i}\quad(v\ge1),\qquad
 f(1)=\sum_i m_i>0.
\end{equation}
All real powers use their positive real branches.

Consequently the actual marked-$X$ trace law of
\fid{F54}, Part~3, is impossible for arbitrary mixed
base spectra if the chosen evaluations of the coefficient
family vary continuously along the positive parameter
axis and are regular at $v=1$ on that same choice.
No weight or full-centralizer condition is needed.
\end{finding}

\begin{proof}
By \fid{F54}, Part~1, there are a common $L>0$ and
integers $a_i$ with
\[
 \alpha_i^L=b^{a_i}.
\]
Put $p(t)=\sum_i m_i t^{a_i}$.
At the indices $n=Lk$ the trace law gives
$f(b^{Lk})=p(b^k)$.
The rational function $P(t^L,p(t))$ therefore has
infinitely many distinct zeros $t=b^k$, away from its
finitely many possible poles, and is identically zero.

The extension $\C(t)/\C(t^L)$ is cyclic Galois, with
automorphisms $t\mapsto\zeta t$, $\zeta^L=1$.
Since $P$ is irreducible and has $p(t)$ as a root,
its roots are among $p(\zeta t)$.
Indeed the orbit polynomial of $p(t)$ has coefficients
fixed by this Galois group, and its irreducible minimal
factor is $P$.
Thus, away from a finite exceptional set, for real $v\ge1$,
\begin{equation}\label{eq:finitepositivebranches}
 f(v)\in\{p(\zeta v^{1/L}):\zeta^L=1\}.
\end{equation}
The inclusion also holds at an exceptional point:
approach through nonexceptional values, pass to a
subsequence with one fixed $\zeta$ from this finite set,
and use continuity of both functions.

Set $p_0(v)=p(v^{1/L})$.
For every $\zeta$ and every positive $v$,
\[
 |p(\zeta v^{1/L})|
 \le\sum_i m_i v^{a_i/L}=p_0(v).
\]
If $p(\zeta v^{1/L})=p_0(v)$, equality in the triangle
inequality, with a positive real sum and strictly
positive coefficients, forces every summand to have
argument zero. Hence $\zeta^{a_i}=1$ for every $i$.
That branch is therefore identically $p_0$.
No genuinely different branch in
\eqref{eq:finitepositivebranches} meets $p_0$ anywhere
on the positive real axis. Negative or zero exponents
do not alter this argument.

The set
\[
 E=\{v\in[1,\infty):f(v)=p_0(v)\}
\]
is closed and contains every sufficiently large
$b^{Lk}$. It is also open: at a point of $E$, the
finitely many genuinely different branch values have
positive distance from $p_0$; continuity preserves
a smaller positive separation nearby, and
\eqref{eq:finitepositivebranches} forces $f$ to remain
on $p_0$.
Connectedness gives $E=[1,\infty)$.

Now $|\alpha_i|=b^{a_i/L}$, and for every sufficiently
large $n$,
\[
 \sum_i m_i\alpha_i^n
   =\sum_i m_i|\alpha_i|^n>0.
\]
Equality in the triangle inequality forces each
$\alpha_i^n$ to be positive real. Applying this at two
consecutive such indices forces each $\alpha_i$ itself
to be positive real. Taking $r_i=a_i/L$ proves
\eqref{eq:continuoustracepositive}.

For the actual module, $X_b$ is $\Delta(b)$ times the
projector onto the rank-eight $N(D-2)^4$ sector.
The assumed right-action commutation makes that sector
invariant under the invertible base operator $F_b$.
The continuous scalar branch
\[
 f(v)=\frac{\Tr_M(X(v)F(v))}{\Delta(b)}
\]
satisfies the irreducible minimal polynomial over
$\C(v)$ of this trace element in the fixed finite
coefficient field, normalized to be monic.
Its marked comparison values are the genuine power
traces of the eight-dimensional restriction of $F_b$.
The proved scalar theorem gives $f(1)=8$.
Regularity on the same branch and the actual matrix
identity $X(1)=0$ give $f(1)=0$, a contradiction.
In fact only commutation of $F_b$ with $X_b$ is needed
for this last trace calculation.
\end{proof}

The continuous choice is additional comparison data.
It is not supplied merely by a finite algebraic cover
or by choosing some place over $v=1$.
In the explicit fractional model, the arithmetic points
have $t=2^{n/2}>0$ over $v=t^2$, while the regular
collision is at $t=-1$.
The normalized marked function $g(t)=2(1+t)^2$ has
$g(1)=8$ and $g(-1)=0$.
Its marked-sector trace $h$ and operator entries have
poles at $t=1$, even though $g$ and the adjusted
ordinary total trace $56+g$ are regular there.
That model does not satisfy this same-branch regularity
condition. Its one-place assertion remains valid.

\paragraph{\fid{F55}, Part 1: a variant retaining the actual formal bar.}
The one-place construction can additionally commute with
an explicit extension of the actual tensor bar.
This variant retains the actual right action and both
trace laws, while still failing all-places regularity
and the stronger left/word comparisons.

\begin{proof}
Keep the actual bimodule decomposition and regular
bar-fixed basis of \fid{F43}, and the constants
$s=\sqrt2$, $\eta=189/625$, $v=t^2$ from Part~1.
Extend the coefficient involution by
\[
 \overline t=t^{-1},\qquad \overline s=s.
\]
This is a stated choice of lift of $v\mapsto v^{-1}$
to the coefficient cover.
Define a second family, distinguished by tildes:
\[
 \widetilde g(t)=\frac{2(1+t)^2}{t}
                  =4+2t+2t^{-1},\qquad
 \widetilde h(t)=\frac{\eta\widetilde g(t)}{\Delta(t^2)},
\]
\begin{equation}\label{eq:barfractionalfunctions}
 \widetilde\lambda(t)=\frac{\widetilde h(t)}2-2-s,\qquad
 \widetilde u(t)=1+\frac{\widetilde g(t)-\widetilde h(t)}{20}.
\end{equation}
All four functions are fixed by this involution:
$\widetilde g(t^{-1})=\widetilde g(t)$ and
$\Delta(v^{-1})=\Delta(v)$.
The cancellation formula is
\[
 \widetilde h(t)=
 \frac{2\eta(t^4+1)^4}
 {t(t-1)^2(t^2+1)^2(t^8+t^4+1)}.
\]
The exact special values are
\[
\begin{array}{c|cccc}
t&\widetilde g&\widetilde h&\widetilde u&
                                   \widetilde\lambda\\ \hline
s&4+3s&4+3s&1&s^{-1}\\
-1&0&-126/625&6313/6250&-1313/625-s .
\end{array}
\]

Use the same $R_0$ as in Part~1 and set
\[
 \widetilde R
   =R_0[\widetilde u^{-1},\widetilde\lambda^{-1}].
\]
It is stable under bar: the inherited factors are
bar invariant, $t-1$ maps to $-(t-1)/t$,
$t^2+1$ to $(t^2+1)/t^2$, and the two new functions
are invariant.
Every inherited denominator is nonzero at $t=-1$
and all $t_n=s^n$.
The displayed values check the additional units
at $-1$ and $s$.
For $n>1$, $\Delta(2^n)>\eta$, so
\[
 \widetilde h(t_n)<\widetilde g(t_n),\qquad
 \widetilde u(t_n)>1.
\]
If $n$ is even, $t_n$ is rational and
$\widetilde\lambda(t_n)\in\Q-s$, hence is nonzero.
For $n=2r+1$ odd, put
\[
 d_n=\Delta(2^n),\qquad c_r=2^r+2^{-r-1}.
\]
Both are rational, and
\[
 \widetilde\lambda(t_n)
 =2\left(\frac{\eta}{d_n}-1\right)
   +s\left(\frac{\eta c_r}{d_n}-1\right).
\]
At $n=1$ this equals $s^{-1}$.
For odd $n>1$, its rational coefficient of $1$
is nonzero because $d_n>\eta$.
Linear independence of $1,s$ proves nonvanishing.
Thus one bar-stable finite localization contains the
chosen collision and every extension point.
The contact order at $t=-1$ is still two.

On the entire rank-$20$ $\epsilon_+$ block use
$\widetilde u I_{20}$, and on the rank-$36$
$\epsilon_-^4\oplus N(\delta)^{16}$ block use identity.
On the final multiplicity factor, in order
$(++,+-,-+,--)$, put
\[
 \widetilde F
   =\operatorname{diag}(1,1,s,\widetilde\lambda)\otimes I_2.
\]
Its determinant is
$2\widetilde u^{20}\widetilde\lambda^2\in\widetilde R^\times$.
As before, transport by
$\iota^{-1}\widetilde F_{\rm block}\iota$ to the
actual tensor lattice.
The operator commutes with the actual right action
because it acts on multiplicity factors.
Its entries and all adapted basis vectors are fixed
by the extended actual bar $\beta$, so
\[
 \beta\widetilde F=\widetilde F\beta
\]
as semilinear maps over $\widetilde R$.
The established bar-compatible basis comparison
retains this equality on the original tensor module.

The marked trace is
$2(1+1+s+\widetilde\lambda)=\widetilde h$.
Hence
\begin{equation}\label{eq:barfractionaltraces}
 \Tr(X(t^2)\widetilde F(t))=\eta\widetilde g(t),\qquad
 \Tr\widetilde F(t)=56+\widetilde g(t).
\end{equation}
At $t=s$, the marked eigenvalues $1,s,s^{-1}$
have multiplicities $4,2,2$; on the full module
the multiplicity of $1$ is $60$.
Thus
\[
 \Tr(\widetilde F(s)^n)
       =60+2s^n+2s^{-n}=56+\widetilde g(t_n),
\]
\[
 \Tr(X(2)\widetilde F(s)^n)
       =\eta(4+2s^n+2s^{-n})=\eta\widetilde g(t_n).
\]
These give both retuning identities for every $n$.

The stronger comparisons still fail explicitly.
On the final multiplicity space,
\[
 (\widetilde F(s)E_V-E_V\widetilde F(s))
       (g_-\otimes g_+)
  =(1-s)(g_+\otimes g_+)\ne0.
\]
The generator calculation gives
\[
 \Tr(P_1(v)\widetilde F(t))
    =D(v)\left(36+\widetilde g(t)
                         -\frac{\widetilde h(t)}2\right),
\]
\[
 \Tr(P_1(2)\widetilde F(s)^n)
    =D(2)\left(38+s^n+s^{-n}\right)
    =D(2)\left(36+\frac{\widetilde g(t_n)}2\right).
\]
At $n>1$ the first is strictly larger, since
$\widetilde h(t_n)<\widetilde g(t_n)$ and $D(2^n)>D(2)$.
Multiplication by the respective $D$ factors proves
the same failure for $P_1^2$.

Unlabelled spectral retuning fails as well:
$\widetilde F(s)^n$ has eigenvalue $1$ of multiplicity
$60$, whereas $\widetilde F(t_n)$ has it with
multiplicity at most $42$ for $n>1$.
The latter count is $36$ unchanged dimensions,
four fixed marked dimensions, and at most two more
if $\widetilde\lambda(t_n)=1$; neither
$\widetilde u(t_n)>1$ nor $s>1$ contributes one.

Finally, $\widetilde h,\widetilde u,
\widetilde\lambda$ have order-two poles at $t=+1$,
with the latter two leading terms
$-\widetilde h/20$ and $\widetilde h/2$.
These are actual eigenvalue poles and cannot be
removed by a finite free lattice change.
The involution fixes $t=+1$ and $t=-1$ separately.
Bar compatibility consequently does not supply
all-places regularity or continuation from the
positive arithmetic sheet to the negative collision.
\end{proof}

At $q=4$ under $v=q^{1/2}$, this variant's marked
weights are $0,\tfrac12,-\tfrac12$.
It is not an integral mixed-geometric or pure-total
realization.
The formal bar covers $t\mapsto t^{-1}$; it is not
asserted to be a Frobenius-equivariant identification
of reciprocal weight lines at one fixed arithmetic
parameter, nor Verdier duality or VIII globalization.
The original Part~1 model and its comparisons are
unchanged.

\subsection{Coherent coefficient kernels and their completion limits}
The coefficient-object alternative in \fid{F54} has an explicit
realization. It retains the actual rank-three word algebra on a
split Grothendieck ring, with natural maps for the presentation
relations. Its completion and adjoint behavior show why this
does not yet select the requested canonical basis.

\begin{finding}[\fid{F56}: coefficient kernels, split $K_0$, and completion]
\label{f:coefficientkernelcategory}
Let $\Lambda=\Z[v,v^{-1}]$ and use the actual positive
rank-three matrices of \fid{F43}.
There is a concrete additive monoidal category
$\mathcal P_{64}$ of finite graded matrix kernels, acting
on $64$ graded coefficient states, such that
\[
 K_0^{\mathrm{split}}(\mathcal P_{64})
   \simeq A_\Lambda^{\mathrm{act}}
   =\Lambda\langle P_1,P_2\rangle
       \subset M_{64}(\Lambda).
\]
It has explicit natural quadratic and degree-five
relation isomorphisms, with coherent associators and
ordered coefficient comparisons.
The algebra on the right is the actual word-algebra
image, not the full ambient matrix algebra.

The category is not closed under direct summands
or under the adjoints of its named generators.
Its Karoubi completion satisfies
\[
 K_0^{\mathrm{split}}\operatorname{Kar}(\mathcal P_{64})
       \simeq\Lambda^{24}\times M_2(\Lambda)^{20},
\]
of rank $104$ over $\Lambda$.
On the faithful six-state version, the completion instead has
\[
 \Lambda^2\times M_2(\Lambda)^2.
\]
Its generic rank is ten, but its localized order is the
split normalization of the actual colliding center.

For fixed $q$, replacing degree $r$ by a specified Weil
line of Frobenius $q^{r/2}$ gives an exact arithmetic
coefficient-object realization, preserving the natural
relation maps at every extension degree.
In odd characteristic, entrywise application of
\fid{F53}'s $J$ gives actual perverse convolution kernels
and maps, with the known weight of every homogeneous piece.
This is a rank-three graded construction, not a selection
of VIII's canonical simple labels or integral forms.
\end{finding}

\begin{proof}
\par\smallskip\noindent\textbf{The concrete matrix-kernel category.}\quad
Let $\mathcal V$ be finite-dimensional $\Z$-graded vector
spaces over a characteristic-zero field, with
degree-preserving maps and ordinary tensor symmetry.
There are no Koszul signs: the grading is formal
quantum degree.
Write $\Bbbk\langle d\rangle$ for a line in degree $d$,
and
\[
 \operatorname{ch}(U)=\sum_d\dim U_d\,v^d.
\]
For a finite state set $I$, define
$\operatorname{Mat}_I(\mathcal V)$ with entrywise maps and
convolution
\begin{equation}\label{eq:gradedmatrixconvolution}
 (K\star L)_{ij}=\bigoplus_h K_{ih}\otimes L_{hj}.
\end{equation}
Its unit is diagonal $\Bbbk$.
It acts on row vectors by
\[
 (T_K X)_j=\bigoplus_i X_i\otimes K_{ij},
 \qquad T_L\circ T_K\simeq T_{K\star L}.
\]

Use the $64$ actual \fid{F43} states: $20$ scalar-plus
states, $4$ scalar-minus states, $16$ copies of $N(\delta)$,
and $4$ copies of $N(D-2)$.
Set
\[
 A_1=\Bbbk\langle-1\rangle\oplus\Bbbk\langle1\rangle,\quad
 A_2=\Bbbk\langle-2\rangle\oplus\Bbbk\langle2\rangle,\quad
 \mathcal D=A_1\otimes A_1.
\]
Their characters are $\delta,D-2,D$.
On a two-state block use
\begin{equation}\label{eq:actualgradedkernels}
 K_1=\begin{pmatrix}\mathcal D&0\\A&0\end{pmatrix},
 \qquad
 K_2=\begin{pmatrix}0&A\\0&\mathcal D\end{pmatrix},
 \qquad A=A_1\ \hbox{or}\ A_2.
\end{equation}
These are the source's actual right-row matrices.
On plus states both kernels are $\mathcal D$; on minus
states they are zero.
The basis is a basis of the actual tensor lattice after
the proved finite localization.
No new equality with the original global tensor-monomial
lattice is asserted.

For any parenthesized convolution word, flattening gives
\[
 \bigoplus_{i_1,\ldots,i_{m-1}}
 (K_1)_{i_0i_1}\otimes\cdots\otimes(K_m)_{i_{m-1}i_m}.
\]
The associator preserves the state sequence and ordered
tensor factors.
Every route around its pentagon therefore has the same
action on every pure tensor in every summand.
The unit identities have the same verification.
Coefficient symmetries permute pure tensor positions.
Adjacent swaps satisfy their square, distant-commutation,
and braid relations because the resulting permutations agree.
This establishes the monoidal coherence used below.

\par\smallskip\noindent\textbf{Quadratic maps and the ordered overlap.}\quad
A two-state $K_i$ resets to state $i$, so
\[
 (K_i\star K_i)_{hi}=(K_i)_{hi}\otimes\mathcal D.
\]
The path identification and coefficient symmetry give
\begin{equation}\label{eq:kernelquadratic}
 K_i\star K_i\simeq K_i\otimes\mathcal D
                   \simeq\mathcal D\otimes K_i.
\end{equation}
These are actual natural maps.
The same description works on plus states and zero
objects on minus states.
For every $m\ge1$, retaining the initial edge and the
chronological list of loops gives
\[
 K_i^{\star m}\simeq
 K_i\otimes\mathcal D^{\otimes(m-1)}.
\]
Flattening makes this independent of convolution
parenthesization.

For precision, use the left-coefficient version of
\eqref{eq:kernelquadratic} and label a triple path
$u\otimes d_1\otimes d_2$.
Reducing the first pair and then the remaining pair
gives $d_1\otimes d_2\otimes u$.
Reducing the last pair, moving its coefficient left,
and then reducing the remaining pair gives
$d_2\otimes d_1\otimes u$.
The comparison is the transposition of the two
$\mathcal D$ factors.
It is nonidentity, witnessed by their distinct
degree-zero labels, but is precisely ordinary
symmetric bookkeeping, not a coherence obstruction.

In general a reduction tree gives a permutation
$\sigma_T$ of the labeled internal $\mathcal D$ factors.
Inductively, local maps only permute those labels and
retain the initial edge.
Compare two normal forms by
$\sigma_{T'}\sigma_T^{-1}$.
All comparison pentagons telescope to the identity,
and the elementary comparisons satisfy the symmetric-group
relations already checked.
This supplies coherent ordered comparisons; it does not
impose strict equality after discarding the labels.

\par\smallskip\noindent\textbf{The actual degree-five relation maps.}\quad
IV, Theorem~11.16, equations~(111)--(112), p.~65, gives
\[
 P_i^2=DP_i,\qquad
 P_1(P_2P_1-a_1^2)(P_2P_1-a_2^2)
 =
 P_2(P_1P_2-a_1^2)(P_1P_2-a_2^2),
\]
with $a_1=\delta$, $a_2=D-2$ \cite{IV}.
Let $V_1=K_1\star K_2\star K_1$ and
$W_1=K_1\star K_2\star K_1\star K_2\star K_1$;
use $V_2,W_2$ for their mirrored words.
Put
\[
 U=A_1^{\otimes2},\quad V=A_2^{\otimes2},\quad
 \mathcal S=U\oplus V,\quad \mathcal T=U\otimes V.
\]
The positive rearrangement has sides
\begin{align}
 L&=W_1\oplus(V_2\otimes\mathcal S)
                         \oplus(K_1\otimes\mathcal T),\nonumber\\
 R&=W_2\oplus(V_1\otimes\mathcal S)
                         \oplus(K_2\otimes\mathcal T).
 \label{eq:kerneldegreefive}
\end{align}
We specify an actual natural isomorphism.

On an $N(A)$ block, every internal alternating step
is the $A$ edge, giving chronological identifications
\[
 V_j\simeq K_j\otimes A^{\otimes2},\qquad
 W_j\simeq K_j\otimes A^{\otimes4}.
\]
Suppress tensor symbols in the following lists.
For $A=A_1$, expand the sides as
\[
 L=(K_1UU,K_2UU,K_2UV,K_1UV),\qquad
 R=(K_2UU,K_1UU,K_1UV,K_2UV).
\]
Map positions by $(1\,2)(3\,4)$, using identity
coefficient maps.
For $A=A_2$, the lists are
\[
 L=(K_1VV,K_2VU,K_2VV,K_1UV),\qquad
 R=(K_2VV,K_1VU,K_1VV,K_2UV).
\]
Send positions $1,2,3,4$ to $3,4,1,2$.
The first and third maps are identities; the other two
interchange the $U,V$ coefficient blocks.
On plus states use the identical expressions
$\mathcal D^5,\mathcal D^3\mathcal S,\mathcal D\mathcal T$,
and on minus states use the unique zero map.

These maps specify every summand and factor permutation.
They are degree preserving and natural.
Exchanging the generator labels gives the inverse:
the displayed permutations reverse, and the opposite
coefficient interchange inverts the first.
The same map on every multiplicity copy commutes with
its matrix units.
Tensoring with any kernel context gives natural relation
maps; disjoint applications commute by the tensor
interchange identity.
Associators and coefficient symmetries have the coherence
already proved.

Thus there is a coherent category/action with explicit
maps for both algebra presentation relations.
No source-prescribed additional higher $2$-relation
among overlapping degree-five maps has been assumed.
An assertion that every unspecified rewriting loop is
identity would be stronger and is not made here.

\par\smallskip\noindent\textbf{The split Grothendieck ring.}\quad
Define $\mathcal P_{64}$ to be the full additive
subcategory whose objects are finite direct sums of
grading shifts of convolution words in $I,K_1,K_2$,
including their isomorphic copies.
Do not close under arbitrary direct summands.
It is monoidal, and use its split exact structure.

Isomorphism classes of finite graded matrix kernels
are determined by entrywise graded dimension.
If those matrices agree, choose an isomorphism in every
homogeneous entry. Fullness puts the resulting
isomorphism in $\mathcal P_{64}$.
The object monoid therefore embeds in
$M_{64}(\mathbb N[v,v^{-1}])$ and is cancellative.
Its image is exactly the nonnegative Laurent
combinations of actual word matrices.

Group completion is injective into $M_{64}(\Lambda)$.
Indeed, equality of two formal differences gives,
after cross-addition, identical dimension matrices,
hence an actual kernel isomorphism and the same
split Grothendieck class.
The image is the $\Z$-span of all shifted words,
which is exactly $A_\Lambda^{\mathrm{act}}$.
Convolution respects multiplication, proving
\[
 K_0^{\mathrm{split}}(\mathcal P_{64})
          \simeq A_\Lambda^{\mathrm{act}}.
\]
All four actual generic simple types occur, so the
actual generic algebra acts faithfully and has
dimension ten.
Removing repeated copies does not change any word
relation. After the established finite localization,
this is the actual order of \fid{F19}.

The qualifier \emph{split} matters.
No abelian or semisimple simple-object decomposition,
idempotent completeness, or prescribed canonical basis
has been asserted for this category.

\par\smallskip\noindent\textbf{Karoubi completion on $64$ states.}\quad
The unit already has the independent state idempotents
in its endomorphism algebra.
They split in the Karoubi completion to give every
$E_{ii}$.
On a two-state block,
$E_{22}\star K_1\star E_{11}$ is the coefficient $A$
in position $(2,1)$.
Splitting a homogeneous basis line and shifting its
degree gives $E_{21}$.
Similarly $K_2$ gives $E_{12}$.
Together with the diagonal objects, these give every
elementary two-by-two kernel and its shifts.
A scalar state supplies its one elementary diagonal kernel.

Every original word is block diagonal with respect to
the $24$ scalar and $20$ two-state copies.
An entrywise idempotent image cannot create a missing
matrix entry.
Conversely every finite graded kernel on these blocks
is a sum of the elementary shifted kernels just constructed.
The ambient fullness of maps therefore gives
\[
 \operatorname{Kar}(\mathcal P_{64})
       \simeq\mathcal V^{24}\times
                    \operatorname{Mat}_2(\mathcal V)^{20}.
\]
Hence its split Grothendieck ring is
\[
 \Lambda^{24}\times M_2(\Lambda)^{20}.
\]
There are $104$ elementary simple-kernel families
up to shift and $64$ simple summands of the unit.
There are no new entries between distinct multiplicity
copies, so the result is not $M_{64}(\Lambda)$.
The resulting positive elementary basis is a basis
of this enlarged algebra, not of the actual word algebra.

\par\smallskip\noindent\textbf{The faithful six-state normalization.}\quad
Retain one copy of each scalar type and each two-state type.
Its word relations are exactly the same, so its
$\mathcal P_6$ has the same split Grothendieck algebra.
The same completion argument gives
\[
 \operatorname{Kar}(\mathcal P_6)
   \simeq\mathcal V^2\times\operatorname{Mat}_2(\mathcal V)^2,
 \qquad
 K_0^{\mathrm{split}}\operatorname{Kar}(\mathcal P_6)
   =\Lambda^2\times M_2(\Lambda)^2.
\]
Its rank is ten, but the relative order has changed.

Work over $R=\mathbb Q[v,v^{-1}]_{(v-1)}$.
The actual center maps to the split branch center by
\[
 R[x]/(x(x-\Delta))\longrightarrow R\times R,\qquad
 f_0+f_1x\longmapsto(f_0,f_0+\Delta f_1).
\]
Its image is precisely the pairs congruent modulo
$\Delta$, and its cokernel is $R/(\Delta)$.
The center cokernel has length two; the four matrix
entries give length eight.
After the regular branch basis changes, the completed
kernel algebra is the split normalization order
\[
 R^2\times M_2(R)^2.
\]
The actual special fiber has center
$\mathbb Q^2\times\mathbb Q[\epsilon]/(\epsilon^2)$
and a four-dimensional matrix radical.
The displayed map kills that radical on specialization
and has six-dimensional image.
The completed order has a reduced split center.
Thus equal generic rank does not retain the collision.

\par\smallskip\noindent\textbf{Adjoints require an enlargement.}\quad
Finite tensor-Hom adjunction gives both left and right
adjoint kernels
\[
 (K^\dagger)_{ij}=(K_{ji})^*.
\]
Graded dual reverses degree.
The chosen $\mathcal D,A_1,A_2$ are self-dual, so
\[
 K_1^\dagger=\begin{pmatrix}\mathcal D&A\\0&0\end{pmatrix},
 \qquad
 K_2^\dagger=\begin{pmatrix}0&0\\A&\mathcal D\end{pmatrix}.
\]
Their supports already show that the named generators
are not self-adjoint, even up to shift.

Neither adjoint belongs to $\mathcal P_{64}$ or
$\mathcal P_6$.
For a nonempty word $w=i_1\cdots i_L$, the rank-one
matrix multiplication of \fid{F43} gives
\[
 P_w=D^{L-1-t}a^t u_{i_1}e_{i_L}^{T},\qquad
 u_1=(D,a)^T,\quad u_2=(a,D)^T,
\]
where $t$ counts adjacent changes.
At $v=1$, $D=4$ and $a=2$.
Both rows of every nonempty word are nonzero; the
identity also has both rows nonzero.
Every nonzero object in the word category is a
positive sum of shifted words.
Its restriction to either two-state block therefore
has both rows nonzero, since shifts evaluate to one
and no cancellation is possible.
But $K_1^\dagger$ has zero second row, and
$K_2^\dagger$ zero first row.
This proves exclusion.

This statement concerns the positive word-object
category. It does not assert that the adjoint operator
lies outside every localization of the actual algebra.
The Karoubi completion contains the transpose-dual
kernels and their elementary summands.

\par\smallskip\noindent\textbf{Arithmetic coefficient objects.}\quad
Fix a prime power $q$, a prime $\ell$ different from its
characteristic, and the algebraic number
$\lambda=\sqrt q>0$ under a specified complex embedding.
Choose its $\ell$-adic realization in a finite
coefficient field $E$.
Let $L$ be the Weil line with geometric Frobenius $\lambda$.
For $E$-linear graded spaces, send
\[
 E\langle r\rangle\longmapsto L^{\otimes r}.
\]
This extends by sums and tensors to an exact symmetric
monoidal functor.
It is fully faithful onto the semisimple category of
these characters: $\lambda^r=\lambda^s$ forces $r=s$,
and Frobenius maps between distinct characters vanish.
The equal-degree maps are the original coefficient maps.

Every specified kernel map is degree preserving,
so it becomes a Frobenius morphism.
The lines are continuous unramified representations:
$\lambda$ is an $\ell$-adic unit, its powers lie in a
compact profinite unit group, and the homomorphism
from $\Z$ extends to its profinite completion.
For every $n$, in the chosen complex embedding,
\begin{equation}\label{eq:gradedcoefficienttrace}
 \Tr\left(F^n\mid\bigoplus_r V_r\otimes L^r\right)
       =\sum_r\dim(V_r)q^{rn/2}.
\end{equation}
Thus $\mathcal D$ becomes
$L^2\oplus E^{\oplus2}\oplus L^{-2}$, and the quantum
relation is a tensor relation with that coefficient object
at every extension degree.
It is not the fixed scalar $D(\sqrt q)$ on one finite space.

Sending degree one instead to $E(-1)$ gives a full
Tate step and the substitution $v=q^n$.
The half-step above is a Weil coefficient line;
it is not asserted to be a geometric Tate object.
Different degrees give mixed weights, and different
words have their own finite trace objects.
This is not the fixed-operator comparison excluded
by \fid{F42} and \fid{F54}.

\par\smallskip\noindent\textbf{Actual perverse coefficient kernels.}\quad
Now assume $q$ is odd.
Use \fid{F53}'s exact fully faithful symmetric strong
monoidal functor $J$ to its perverse Gaussian image,
with specified product $\star_\psi$ and image unit $J(E)$.
Its symmetry is the transported symmetry, including
the established correction to the naive odd-shift swap.

On matrices of image objects define
\[
 (\mathcal K\star\mathcal L)_{ij}
       =\bigoplus_h\mathcal K_{ih}
                      \star_\psi\mathcal L_{hj}.
\]
The unit is diagonal $J(E)$.
Entrywise $J$ after the arithmetic functor gives
\[
 J_I(K\star L)\simeq J_I(K)\star J_I(L).
\]
To check coherence, expand an iterated product by its
intermediate state sequence and use F53's strong-monoidal
comparison on each tensor summand.
Every associator pentagon agrees summand by summand;
direct sums preserve the equality.
Full faithfulness and exactness are entrywise.
The matrix action on the corresponding product of
perverse image categories follows from the same formula.

All constructed quadratic, degree-five, inverse,
tensor-context, and ordered-overlap maps are therefore
actual perverse morphisms with the same comparisons.
This uses the explicit coefficient maps, not an
inference from a Laurent $K_0$ equality.

A degree-$d$ line has Weil weight $d$ in the half-step
convention, and $J(L^d)$ is perverse pure of weight $d$.
In the full Tate convention its weight is $2d$.
An arbitrary coefficient is the specified finite sum
of these pieces; its total need not have one weight.
The split-$K_0$, completion and adjoint calculations
carry to this full image by equivalence.
They do not classify all perverse sheaves in the
ambient constructible category.
\end{proof}

\paragraph{The precise remaining categorical comparison.}
This construction gives actual coherent operations and
maps at rank three, including an all-extension arithmetic
coefficient-object interpretation.
It does not supply compatible positive graded kernels
through all tensor degrees, a source-prescribed higher
categorical presentation, or VIII's global/local lattice and canonical
simple-label comparisons.
The fixed-parameter all-degree realization of \fid{F53}
is a separate statement.

After forgetting arithmetic descent, the lines $L^d$
are ordinary constant coefficient lines and their Gaussian
images have the same underlying geometric rank-one
constituent.
Their degree labels remain arithmetic/grading data,
rather than a new family of geometric cell labels.
The completion and adjoint results show concretely
why the present positive word category is not already
the desired semisimple canonical-basis category.
They do not refute a different categorical construction.

\subsection{The formal bar and the actual rank-three bar}
The absence of adjoints in the positive word category does not
prevent a formal bar. The distinction is the transpose of the
state indices.

\begin{finding}[\fid{F56}, supplement: an entrywise formal bar]
\label{f:formalkernelbar}
On finite graded matrix kernels put
\[
 (\mathsf B K)_{ij}=K_{ij}^*,\qquad
 (U^*)_d=\operatorname{Hom}_{\Bbbk}(U_{-d},\Bbbk).
\]
Entrywise duality is a contravariant strong monoidal involution,
preserving the order of kernel multiplication. It restricts to
$\mathcal P_{64}$ and $\mathcal P_6$, reverses grading shifts,
and fixes their two named generators up to the specified
self-duality maps below. It carries the explicit quadratic,
degree-five and ordered-overlap maps to their inverses under
those identifications.

On the split Grothendieck algebra and its state module, it
induces the actual \fid{F43} bar: $v\mapsto v^{-1}$,
$P_i\mapsto P_i$, and the $64$ state classes are fixed.
This is an entrywise bar, not the transpose-dual adjoint.
\end{finding}

\begin{proof}
\par\smallskip\noindent\textbf{Monoidal and involutive structure.}\quad
The tensorator
\[
 \mu_{K,L}:\mathsf B K\star\mathsf B L
                 \longrightarrow\mathsf B(K\star L)
\]
is, on each intermediate-state summand, the isomorphism
\[
 \phi\otimes\psi\longmapsto
 [u\otimes w\longmapsto\phi(u)\psi(w)].
\]
Finiteness identifies the dual of each direct sum with the
direct sum of its duals. No state index is transposed and no
Koszul sign occurs. Naturality reads
\[
 \mathsf B(f\star g)\mu_{K',L'}
       =\mu_{K,L}(\mathsf Bf\star\mathsf Bg)
\]
for $f:K\to K'$ and $g:L\to L'$.
For a triple product, both coherence routes evaluate as
$\phi(u)\psi(w)\chi(z)$ on the same flattened state path.
The unit map is the scalar evaluation $\Bbbk\simeq\Bbbk^*$
on each diagonal entry. These computations prove all
strong-monoidal associativity and unit conditions.

Let $\eta_K:K\to\mathsf B^2K$ be evaluation,
$\eta_K(u)(\phi)=\phi(u)$, entrywise. The tensorator for
$\mathsf B^2$ is
\[
 \mu^{(2)}_{K,L}
 =\mathsf B(\mu_{K,L})^{-1}
       \mu_{\mathsf B K,\mathsf B L}.
\]
Evaluation on pure tensors gives
\[
 \eta_{K\star L}=\mu^{(2)}_{K,L}(\eta_K\star\eta_L),
 \qquad
 \mathsf B(\eta_K)\eta_{\mathsf B K}
        =\operatorname{id}_{\mathsf B K}.
\]
The second composite sends $\phi$ to the functional with
value $\eta_K(u)(\phi)=\phi(u)$ at $u$.
Thus evaluation is monoidal and satisfies involutive
coherence, not merely an objectwise double-dual isomorphism.

\par\smallskip\noindent\textbf{Specified generator identifications.}\quad
For $r=1,2$, write $e_{r,-},e_{r,+}$ for the homogeneous
basis of $A_r$ in degrees $-r,r$. Give $A_r$ the symmetric
degree-zero form
\[
 b_r(e_{r,-},e_{r,+})=b_r(e_{r,+},e_{r,-})=1,
 \qquad b_r(e_{r,-},e_{r,-})=b_r(e_{r,+},e_{r,+})=0.
\]
Its map $\theta_r:A_r\to A_r^*$ is nondegenerate and obeys
$\mathsf B(\theta_r)\eta_{A_r}=\theta_r$.
Give $\mathcal D=A_1\otimes A_1$ the product form;
its map is $\theta_{\mathcal D}=\mu_{A_1,A_1}
(\theta_1\otimes\theta_1)$.
It pairs the two distinct degree-zero loop labels with
one another, rather than identifying those labels.

Using these maps in every nonzero kernel entry gives
$\theta_i:K_i\simeq\mathsf B K_i$.
Products use product forms on each state path, and sums
use orthogonal direct sums. This specifies symmetric
self-dual identifications on all unshifted word objects.
Since
\[
 \mathsf B(K\langle d\rangle)
     \simeq(\mathsf B K)\langle-d\rangle,
\]
finite sums of shifted words remain in the word category.
Fullness supplies their dual morphisms, so $\mathsf B$
restricts to $\mathcal P_{64}$ and $\mathcal P_6$.
The functor itself is defined entrywise and does not
depend on a chosen word presentation of an object.

\par\smallskip\noindent\textbf{The relation maps.}\quad
Flattened associators and ordinary coefficient permutations
preserve the product pairings. The quadratic path map and
its coefficient flip are therefore isometries, including
the nontrivial transposition of two labeled loop factors
in the triple overlap.
For the degree-five map, on the $A_1$ branch the summand
permutation is $(1\ 2)(3\ 4)$ with identity coefficient maps.
On the $A_2$ branch it sends positions $1,2,3,4$ to
$3,4,1,2$ and interchanges $U,V$ on the two mixed terms.
These are isometries for the product forms on
$U=A_1^{\otimes2}$ and $V=A_2^{\otimes2}$.
The scalar-plus maps are identities and the minus maps
are zero. Hence for every named relation isomorphism
$r:C\to C'$,
\[
 \mathsf B(r)\theta_{C'}r=\theta_C,
 \qquad
 \theta_C^{-1}\mathsf B(r)\theta_{C'}=r^{-1}.
\]
The same argument applies to the specified ordered-overlap
permutations and to unshifted word contexts with these
self-dual identifications. For arbitrary shifted or summed
contexts $H,J$, monoidality instead identifies the dual
of $\operatorname{id}_H\star r\star\operatorname{id}_J$
with $\operatorname{id}_{\mathsf B H}\star r^{-1}
\star\operatorname{id}_{\mathsf B J}$ after identifying only
the two relation objects with their duals. The inverse is
on the corresponding dual contexts; their shifts reverse.
No self-duality of every shifted object or unspecified
further higher relation is inferred.

\par\smallskip\noindent\textbf{The actual bar comparison.}\quad
The character identity is
\[
 \operatorname{ch}(\mathsf B K)_{ij}(v)
       =\operatorname{ch}(K)_{ij}(v^{-1}).
\]
Together with monoidality, this gives the semilinear
ring involution on the actual split Grothendieck image,
fixing $P_1,P_2$. It preserves word order.
On rows put $(\mathsf B_I X)_i=X_i^*$.
The same tensor pairing gives the natural action comparison
\[
 T_{\mathsf B K}(\mathsf B_I X)
        \simeq\mathsf B_I(T_K X).
\]
Its coherence follows on the same state-path tensors;
each degree-zero state class is fixed.

By \fid{F43}, every vector of the chosen actual tensor
basis is fixed by the actual tensor bar. Identifying
state classes with that basis identifies the displayed
coefficientwise involution with that actual semilinear
bar on their full-rank lattice. The already proved,
bar-stable localization makes this the entire actual
tensor lattice there. This is the stated decategorified
comparison; no equality with VIII's unlocalized global
form has been introduced.
\end{proof}

The action adjoint instead has
$(K^\dagger)_{ij}=(K_{ji})^*$.
Its transpose changes the supports of the reset kernels,
so the \fid{F56} proof excluding those adjoints from
the positive word category remains unchanged.
The entrywise bar does not supply adjunction evaluations,
arbitrary summands, or a repair of the Karoubi collision loss.
The chosen bilinear forms are not positive Hermitian
polarizations.

Under the fixed-$q$ coefficient realization it sends a
degree-$d$ Weil line to degree $-d$.
Entrywise \fid{F53} transport gives internal coefficient
duality on the perverse image; it is not an identification
with ambient Verdier duality or the reciprocal-weight dual.
No new weight estimate, designated simple label, global/local
lattice comparison, or all-degree graded construction follows.

\subsection{A negative word trace on an actual rank-four simple module}
The rank-three coefficient kernels do not extend, with the same
strict positivity requirement, to every simple module of the full
rank-four nonstandard Hecke algebra. The obstruction is a trace of
a named word and is therefore independent of the choice of basis.
It concerns an additional simple module outside the two-row quotient.

\begin{finding}[\fid{F57}: an actual rank-four word-trace obstruction]
\label{f:rankfournegativetrace}
Put \(\delta=v+v^{-1}\), \(D=\delta^2\), and
\[
 N=S'\check M_{(3,1)},\qquad
 N^\#=\text{the }\Theta^{\mathrm{op}}\text{-contragredient of }N.
\]
The latter is the additional five-dimensional simple module listed
in IV, \S11.8, p.~67. For the named word
\[
 \omega=P_2P_3P_2P_3P_2P_3P_1P_2P_1P_2P_1
\]
one has
\begin{equation}\label{eq:rankfournegativepolynomial}
\begin{split}
 \operatorname{Tr}_{N^\#}(\omega)=p(D)
 ={}&-D^9+20D^8-163D^7+734D^6-2048D^5\\
    &+3712D^4-4352D^3+3072D^2-1024D.
\end{split}
\end{equation}
Its Laurent coefficients at \(v^{18}\) and \(v^{-18}\) are \(-1\).
The module \(N^\#\) is an actual direct summand of the restriction
of \(M_{(2,1,1)}\otimes M_{(3,1)}\), hence occurs in the actual
nonstandard tensor representation when
\(\dim V\ge3,\ \dim W\ge2\), or with the factors interchanged.

Consequently no basis of this specified generic simple module
makes all three named \(P_i\) matrices coefficientwise nonnegative
Laurent polynomials. In particular there is no effective finite
graded-matrix-kernel realization with this Grothendieck module
and these named generators.

For any fixed \(q>1\), neither
\[
 \pm p\big((q^{n/2}+q^{-n/2})^2\big)
\]
is the power-trace sequence of an invertible endomorphism of a
finite-dimensional vector space for all sufficiently large \(n\).
A common Tate twist does not repair this. A bounded complex
with both cohomological parities does realize the scalar sequence
when \(q\) is a prime power.

The negative \(P\)-word occurs on the extra full-algebra
\(N^\#\) block. It does not exclude a \(P_i\)-only effective
kernel model for the two-row quotient, nor does it contradict
the signed/factored positivity stated in IV.
\end{finding}

\begin{proof}
\par\smallskip\noindent\textbf{The exact source input.}\quad
IV already constructs the two-row tensor basis through rank four:
see \S1.5, p.~13, Conjecture~19.1, pp.~130--131, and
Example~19.3, pp.~133--135.
Its stated coefficient property is signed and factored:
Remark~11.19, p.~67, and Example~19.3, p.~134, allow signs and
\(c_0=[2]^2-4\).
This is distinct from strict nonnegative Laurent matrices.
The canonical cellular basis of the algebra itself remains
unconstructed there at rank four; see \S19.1, p.~136
\cite{IV}.

Use the explicit \(Q_i\) matrices in Figure~22, p.~135,
with \(Q_i=D-P_i\).
This normalization follows from \S11.1, p.~59.
It also avoids propagating the exponent inconsistency between
that section's preliminary \(C'_s,C_s\) expressions and
\S3.1: the displayed matrices themselves satisfy \(Q_i^2=DQ_i\).

Set \(a=D-2,\ c=D-4\). Order the Figure~22 states as its
two upper states from left to right, followed by its three
lower states from left to right. Their descent sets are
\[
 \{2\},\quad\{1,3\},\quad\{1\},\quad\{2\},\quad\{3\}.
\]
Use right row matrices: entry \((r,s)\) is the coefficient of
state \(s\) in the action on state \(r\). The source's descent
and arrow conventions then give
\begin{equation}\label{eq:rankfourfivecellmatrices}
\begin{gathered}
Q_1=\begin{pmatrix}
0&-\delta&c&0&0\\
0&D&0&0&0\\
0&0&D&0&0\\
0&0&a&0&0\\
0&-\delta&0&0&0
\end{pmatrix},\qquad
Q_2=\begin{pmatrix}
D&0&0&0&0\\
-\delta&0&0&0&0\\
0&0&0&a&0\\
0&0&0&D&0\\
a&0&0&-a&0
\end{pmatrix},\\[2mm]
Q_3=\begin{pmatrix}
0&-\delta&0&0&c\\
0&D&0&0&0\\
0&-\delta&0&0&0\\
0&-\delta&0&0&-2\\
0&0&0&0&D
\end{pmatrix}.
\end{gathered}
\end{equation}
They are matrices of the actual \(N=S'\check M_{(3,1)}\)
in Example~19.3, not an abstract representation identified
only by a dimension count.

\par\smallskip\noindent\textbf{An explicit product certificate.}\quad
Write
\[
 w=(1,2,1,2,1,3,2,3,2,3,2),\qquad
 Q_w=Q_1Q_2Q_1Q_2Q_1Q_3Q_2Q_3Q_2Q_3Q_2.
\]
Let
\[
 u=(-2\delta a,D^2,-\delta D,-\delta a,-\delta D)^t,
 \qquad e_2^t=(0,1,0,0,0).
\]
Equation~\eqref{eq:rankfourfivecellmatrices} gives
\(Q_1Q_3=ue_2^t\). After rotating inside the trace, put
\[
 A=Q_2Q_3Q_2Q_3Q_2Q_1Q_2Q_1Q_2.
\]
Successive multiplication of the row \(e_2^t\) by the
displayed matrices gives
\[
 e_2^tA=(r_1,0,0,r_4,0),
\]
where
\[
 r_1=-\delta^5(D^4-9D^3+31D^2-48D+32)
\]
and
\[
 r_4=\delta(D^7-16D^6+109D^5-418D^4
                 +992D^3-1472D^2+1280D-512).
\]
It follows that
\begin{align*}
 \operatorname{Tr}(Q_w)
 &=e_2^tAu=-\delta a(2r_1+r_4)\\
 &=-D(D-2)(D^7-18D^6+127D^5-480D^4\\
 &\hspace{37mm}+1088D^3-1536D^2+1280D-512),
\end{align*}
which expands to \(p(D)\).
For completeness, the only nonzero diagonal entries of
the full product are
\[
\begin{aligned}
(Q_w)_{11}={}&D^{10}-20D^9+180D^8-951D^7+3246D^6\\
&-7440D^5+11520D^4-11712D^3+7168D^2-2048D,\\
(Q_w)_{44}={}&-D^{10}+19D^9-160D^8+788D^7-2512D^6\\
&+5392D^5-7808D^4+7360D^3-4096D^2+1024D.
\end{aligned}
\]
Their sum again gives
\eqref{eq:rankfournegativepolynomial}.

\par\smallskip\noindent\textbf{The source dual and actual occurrence.}\quad
IV, Proposition~11.6, p.~60, supplies \(\Theta(P_i)=Q_i\).
Section~11.4, pp.~62--63, defines \(\#\) using the
anti-involution \(\Theta^{\mathrm{op}}\).
Thus the dual row matrices of \(P_i\) on \(N^\#\) are \(Q_i^t\).
The reversal in
\[
 \omega=P_{w^{\mathrm{rev}}}
\]
is essential to the direct matrix identity
\[
 \rho_{N^\#}(\omega)=(Q_w)^t.
\]
Taking traces proves
\eqref{eq:rankfournegativepolynomial} with the actual
source convention.

The generic direct-summand inclusion
\[
 N\subset S^2M_{(3,1)}
      \subset M_{(3,1)}\otimes M_{(3,1)}
\]
uses IV, (99), p.~63, and the symmetric/exterior splitting.
Apply the exact contragredient functor.
Propositions~11.11--11.12, p.~63, especially (101) and (105),
give
\[
 N^\#\subset
 \operatorname{Res}_{\check H_4}
   \big(M_{(2,1,1)}\otimes M_{(3,1)}\big)
\]
as a direct summand.
Ordinary Schur--Weyl duality puts this restriction in
\((V\otimes W)^{\otimes4}\) for \(\dim V\ge3,\ \dim W\ge2\).
The identification with the actual nonstandard action is
IV, Theorem~12.1, (115), p.~68, and (116)--(117), p.~69.
For dimensions \(3,2\), the ordinary left factors have
dimensions \(3,3\), so this particular isotypic component
contains nine copies of \(N^\#\).
Remark~12.2, p.~68, qualifies the source's \(X/X^*\)
convention. Consistent dualization transposes matrices and
reverses words; it preserves this trace obstruction.

The full-algebra list in IV, \S11.8, p.~67, includes
the additional five-dimensional \(N^\#\).
Theorem~13.4, pp.~74--75, gives the two-row list without it.
The latter squared dimensions sum to
\[
 1+4+9+36+25+4+9+1=89;
\]
the full-algebra dimension is \(89+25=114\).

\par\smallskip\noindent\textbf{Effective-kernel obstruction.}\quad
Expanding in \(v\) gives
\begin{equation}\label{eq:rankfournegativeLaurent}
\begin{split}
p(D)={}&-(v^{18}+v^{-18})+2(v^{16}+v^{-16})
 +4(v^{14}+v^{-14})\\
&+36(v^{12}+v^{-12})+67(v^{10}+v^{-10})
 +176(v^8+v^{-8})\\
&+297(v^6+v^{-6})+476(v^4+v^{-4})
 +657(v^2+v^{-2})+668.
\end{split}
\end{equation}
In particular its two outer coefficients are negative.
If every \(P_i\) had entries in
\(\mathbb N[v,v^{-1}]\), matrix multiplication and trace
would force the same property for every word trace.
Similarity cannot change this trace, so the obstruction
holds for every basis of the specified module.
Effective finite graded-vector-space kernels have exactly
this nonnegative matrix-character property.

There is also a pointwise obstruction:
\[
 p((3+3^{-1})^2)
 =-\frac{99834597608550400}{387420489}<0.
\]
On this displayed specialization, no real basis makes all
three \(P_i\) matrices entrywise nonnegative.
No irreducibility assertion at \(v=3\) is needed.

On a larger module the same conclusion follows if one also
requires an effective kernel projector onto this isotypic
component: its inserted word trace is a positive multiplicity
times \(p(D)\).
Effectivity of that projector is an additional premise;
an arbitrary positive basis of a larger tensor module
need not make its direct-summand projectors effective.

\par\smallskip\noindent\textbf{Finite power traces, parity, and a derived witness.}\quad
At \(v=q^{n/2}\),
\eqref{eq:rankfournegativeLaurent} is a finite exponential
sum with distinct roots \(q^{e/2}\) and the displayed
integer coefficients.
A genuine finite power trace
\(\operatorname{Tr}(F^n;H)\), with \(F\) invertible,
has nonnegative integer multiplicities after equal
eigenvalues are grouped. Jordan blocks do not change this
trace.
Subtracting the two expressions and applying the finite
Vandermonde independence in \fid{F02} contradicts the
coefficient \(-1\) at \(q^9\).
The same proof holds if equality is assumed only for
all sufficiently large \(n\).
The inverse parameter convention gives the same polynomial.

An overall parity shift cannot help:
\(p(D)\) has negative outer coefficients while
\(-p(D)\) has, for example, constant coefficient \(-668\).
A fixed Tate twist multiplies all distinct roots by one
nonzero scalar and leaves these signs unchanged.

There is nevertheless an explicit bounded Tate complex
realizing the scalar trace. For prime-power \(q\), choose
\(\ell\nmid q\) and let
\[
 (b_1,\ldots,b_8)=(657,476,297,176,67,36,4,2),
\]
\[
 E_+=\mathbb Q_\ell^{\,668}
   \oplus\bigoplus_{k=1}^{8}
      \big(\mathbb Q_\ell(-k)\oplus
           \mathbb Q_\ell(k)\big)^{\oplus b_k},
 \qquad
 E_-=\mathbb Q_\ell(-9)\oplus\mathbb Q_\ell(9).
\]
Then \(\dim E_+=4098\), and
\[
 E_+\oplus E_-[-1]
\]
has Frobenius Euler trace \(p(D)\) at every extension degree.
This proves the exact two-parity escape for this scalar
sequence; it supplies no compatible relation maps for
the original word action.
\end{proof}

The calculation is an obstruction to the strict effective
kernel extension on a specified full-algebra simple module.
It does not refute IV's signed/factored positivity, or
exclude alternating derived traces and inserted-operator
traces. On the original two-row \(N\), it is a \(Q\)-word
trace, and \(Q_i=D-P_i\) uses subtraction. Thus the
\(P_i\)-only two-row route remains outside this result.
The source matrices, product diagonals, rank-one reduction,
dual reversal, Laurent expansion and exact specialization
are also checked in the saved exact certificate; the proof
does not infer an all-word claim from a finite search.

\subsection{A normalization that cannot be a finite Frobenius trace}
\begin{finding}[\fid{F26}: rational normalization versus finite Weil spectra]\label{f:normalizedtrace}
Let $b>1$ be real and $f(x)\in\C(x)$. The eventual sequence $f(b^n)$ is
a finite exponential polynomial
\[
 \sum_{j=1}^s p_j(n)\alpha_j^n,
 \qquad p_j\in\C[n],\quad \alpha_j\in\C^*,
\]
if and only if $f$ is a Laurent polynomial.

On the \emph{actual} rank-three tensor module $M$ of \fid{F21}, set
$e_i=P_i/D$, the normalized algebraic idempotents. Then
\[
 \Tr_M(e_1e_2)=24+\frac{16}{D^2},\qquad
 \Tr_M(P_1P_2)=24D^2+16.
\]
For a prime power $q>1$ and either substitution $v=q^{n/2}$ or $v=q^{-n/2}$, the first trace
is not a finite exponential polynomial in $n$, even eventually.
It therefore cannot be the alternating trace of a fixed correspondence
against $F^n$ on any bounded finite-cohomology Weil realization.
This refutes the literal arithmetic word-trace comparison described
below; it does not assert that the GCT sources require that comparison.
\end{finding}
\begin{proof}
If $f=\sum a_jx^j$ is Laurent, then $f(b^n)=\sum a_j(b^j)^n$.
Conversely a finite exponential polynomial satisfies an eventual
constant-coefficient recurrence
\[
 \sum_{j=0}^r c_j f(b^{n+j})=0,\qquad c_0\ne0.
\]
For example its annihilator is a product of powers of the shift
operators $S-\alpha_j$, with nonzero constant term. The rational
function $R(x)=\sum_{j=0}^r c_j f(b^jx)$ vanishes at infinitely many
distinct $b^n$, so $R=0$. If $f$ has a finite nonzero pole $x_0$ of
maximal absolute value, none of $f(b^jx)$ for $j\ge1$ has a pole at
$x_0$: that would require the larger-modulus pole $b^jx_0$ of $f$.
Thus only $c_0f(x)$ has a pole there, a contradiction. A rational
function with poles only at $0$ and infinity is Laurent. Eventual
evaluation is always defined, since a rational function has only
finitely many poles. The same argument applies to $f(b^{-n})$ by
replacing $f(x)$ with $f(x^{-1})$.

In the source row matrices $N(a)$ from \fid{F21},
$\Tr(e_1e_2)=a^2/D^2$. The two scalar sectors contribute $20$ and zero.
Consequently
\[
 \Tr_M(e_1e_2)
 =20+\frac{16}{D}+\frac{4(D-2)^2}{D^2}
 =24+\frac{16}{D^2}.
\]
Multiplication by $D^2$ gives the unnormalized trace. Under the stated
substitutions, $D=q^n+2+q^{-n}$ and the sequence is
\[
 h_n=24+\frac{16q^{-2n}}{(1+q^{-n})^4}.
\]
As a rational function of $x=q^n$ this is
$24+16x^2/(x+1)^4$, which has a genuine pole at $x=-1$ and is not
Laurent. The criterion just proved excludes every finite exponential
polynomial, with arbitrary complex coefficients and Jordan-polynomial
terms allowed.

Finally let $T$ be any fixed endomorphism on a finite-dimensional
cohomology space with invertible Frobenius $F$. Its trace series is
\[
 \sum_{n\ge0}\Tr(TF^n)Z^n
 =\Tr\bigl(T(1-ZF)^{-1}\bigr),
\]
a rational function, equivalently a finite exponential polynomial in
$n$. A finite alternating sum over degrees has the same property.
Commutation of $T$ with $F$ is not needed for this obstruction; for a
Weil endomorphism it holds and only ordinary power sums are needed.
This proves the finite-realization assertion.
\end{proof}

\paragraph{The exact comparison excluded.}
Suppose a proposed finite realization assigns fixed Weil endomorphisms
$E_1,E_2$ to $e_1,e_2$ and requires, for every word $w$ and every
extension degree,
\[
 \sum_i(-1)^i\Tr\bigl(w(E_1,E_2)F^n\mid H^i\bigr)
 =\left.\Tr_M(w(e_1,e_2))\right|_{v=q^{n/2}}.
\]
The word $e_1e_2$ contradicts \fid{F26}; even imposing no idempotency
on the $E_i$ would not help. At one fixed specialization
$v=q^{1/2}$ the algebraic matrices can certainly act on a constant
Weil vector space. The resulting trace against $F^n$ is different
from substituting $v=q^{n/2}$ afresh at each extension degree.

The denominator issue is already visible in the finite Tate ring:
with $\mathbb L=v^2$, $D=\mathbb L+2+\mathbb L^{-1}$ is not a unit
of $\Q[\mathbb L,\mathbb L^{-1}]$. A bounded finite virtual sum of
Tate lines cannot have class $D^{-1}$. In contrast the actual
unnormalized trace $24D^2+16$ is a positive Laurent polynomial and
has a finite graded Tate realization. Individual matrix entries can
likewise be lifted to signed graded classes. Such entrywise lifts do
not construct compatible chain maps, multiplication functors or
categorical splittings of the normalized idempotents.

Localized motive rings, stacks with unbounded cohomology and completed
categories can accommodate rational expressions; they are not excluded
by this theorem. They require an additional finite extraction, or a
trace theorem beyond \fid{F02}/\fid{F24}, for the particular non-Laurent
family above. Indeed its convergent expansion is
\[
 h_n=24+16\sum_{j\ge0}(-1)^j\binom{j+3}{3}q^{-(j+2)n}.
\]
The independent analytic proof in the audit shows genuine generating
series poles at every $q^{j+2}$; the recurrence proof above avoids any
convergence assumption. Clearing denominators before taking traces is
another possible dictionary, but its compatibility with the designated
coefficient and operations is a separate claim.


\subsection{Fixed operators and extension-dependent quantum grading}
Clearing the normalization denominators in \fid{F26} removes
the obstruction for an individual word. It does not produce
one finite Frobenius space supporting all words under the same
extension-dependent substitution.

\begin{finding}[\fid{F42}: a common finite spectrum cannot encode all words]\label{f:commonwords}
On the actual rank-three tensor module of \fid{F21}, put
$D=v^2+2+v^{-2}$. For every $k\ge1$,
\[
 \Tr(P_1^k)=40D^k.
\]
Fix $q>1$ and set $D_n=q^n+2+q^{-n}$.
There is no finite-dimensional complex vector space $V$,
fixed invertible $F\in\End(V)$, and endomorphisms $T_k$ for
an unbounded set of positive integers $k$, such that
\[
 \Tr(T_kF^n)=40D_n^k
\]
for all sufficiently large $n$. The threshold may depend on $k$.
The assertion permits noncommuting $T_k,F$ and Jordan blocks,
and remains valid for finite Euler supertraces. It also remains
valid after multiplying each right-hand side by a nonzero scalar
$a_k$ and a single nonzero exponential factor $\beta_k^n$.

When $q$ is a prime power, the entire actual tensor family at the single
parameter $v_0=q^{1/2}$ has an exact action-compatible realization
by finite Weil objects with identity Frobenius. Its normalized
generators are genuine idempotent endomorphisms. Moreover each
individual sequence $40D_n^k$ has a positive finite Tate
coefficient object, with that object allowed to depend on $k$.
These constructions do not supply the opposite-sign path
comparison required in \fid{F24}.
\end{finding}

\begin{proof}
Use the consistent normalization of \fid{F19}: $c_i^2=\delta c_i$,
$\delta=v+v^{-1}$, and
\[
 P_i=c_i\otimes c_i+(c_i-\delta)\otimes(c_i-\delta).
\]
The two terms are orthogonal and their squares are $D$ times
themselves, so $P_i^2=DP_i$. This uses the symbolic definition
in GCT IV, equation~(93), with the section-3.1 convention
\cite[\S3.1 and Definition~11.1, equation~(93)]{IV},
rather than silently using the inconsistent section-11
assignments already recorded in the source audit.
In the decomposition of \fid{F21}, $P_1$ acts by $D$ on
the scalar positive block of multiplicity 20 and by zero on
the scalar negative block of multiplicity 4. Each of the
two-dimensional blocks has rank one for $P_1$ and trace $D$,
with multiplicities 16 and 4. Hence
\[
 \Tr(P_1^k)=40D^k\qquad(k\ge1).
\]
Both substitutions $v=q^{n/2}$ and $v=q^{-n/2}$ give
\begin{equation}\label{eq:wordpowers}
 40D_n^k
   =40\sum_{j=0}^{2k}\binom{2k}{j}q^{(k-j)n}.
\end{equation}
The $2k+1$ roots in this sum are distinct, and every coefficient
is nonzero.

For fixed $F,T_k$, the generating function beginning at $n=0$ is
\[
 \sum_{n\ge0}\Tr(T_kF^n)T^n
   =\Tr\bigl(T_k(1-TF)^{-1}\bigr).
\]
All its poles are among the inverses of the finitely many
eigenvalues of $F$. A Jordan block can increase a pole order
but cannot add a pole location. Removing finitely many initial
terms changes only a polynomial. On the other hand the
generating function of \eqref{eq:wordpowers} is
\[
 40\sum_{j=0}^{2k}\binom{2k}{j}
      \frac{1}{1-q^{k-j}T}.
\]
Every one of its $2k+1$ poles is genuine: only one summand has
that pole, and its coefficient is nonzero. Thus a common $F$
would require at least $2k+1$ distinct eigenvalues for unbounded
$k$, which is impossible.
For a finite graded object use the union of the even and odd
spectra; for a bounded complex use its finite-dimensional
cohomology. Multiplication by $a_k\beta_k^n$ preserves the
number of distinct nonzero roots, now
$\beta_kq^{k-j}$, proving the normalization variant.
This proof even permits unrelated $T_k$; requiring
$T_k=P^k$ or additional tensor relations cannot evade it.

We give the positive fixed-parameter construction with its
action comparison stated explicitly. Let $q$ be a prime power,
$\ell\nmid q$, and let a finite extension $E/\Q_\ell$
contain $v_0=q^{1/2}$.
Write
\[
 M_{r,q}=(V_q\otimes W_q)^{\otimes r}
\]
for the actual tensor module with the two ordinary Hecke
actions at $v_0$, and let $A_{r,q}$ be the image algebra of its
specified nonstandard generators. An action comparison consists
of finite Weil objects $H_r$, fiber identifications
$\omega(H_r)\simeq M_{r,q}$, and right $A_{r,q}$-actions
by Weil endomorphisms such that:
\begin{enumerate}
\item every named $P_i$ is intertwined by the fiber
identification;
\item the tensor identifications $H_r\otimes H_s\simeq H_{r+s}$
intertwine the local generators on disjoint consecutive blocks,
with the usual tensor associativity coherence.
\end{enumerate}
Extension of the finite field raises Frobenius to its $n$th
power while keeping $v_0$ and these operators fixed.

Give $V_q,W_q$ identity Frobenius and use the induced identity
Frobenius on their tensor powers. Every actual matrix is then
a Weil endomorphism. The local Hecke formula in
GCT IV, \S6.1 \cite[\S6.1, p.~33]{IV}, gives the actual operators, and locality proves
the disjoint-block compatibility. All their algebra relations
are unchanged. This proves the stated action comparison for
the whole actual image-algebra family.
For fixed dimensions two these are the two-row quotient actions;
no faithfulness assertion for the full abstract algebra at
every rank is required. The quantum Hecke operators are extra
morphisms, not an identification with the symmetry maps of the
ordinary Weil category.

Since $D_0=(v_0+v_0^{-1})^2\ne0$, scalar division defines
$e_i=P_i/D_0$ and gives $e_i^2=e_i$. Their images are genuine
Weil subobjects. This divides by a coefficient-field scalar;
it does not tensor-invert the object of \fid{F38}.
The resulting extension trace is
\[
 \Tr\bigl(w(P_i(v_0))F_r^n\bigr),
\]
with fixed operators. For identity Frobenius it is constant
in $n$, not the formal word trace retuned at $q^{n/2}$.
In particular the normalized mixed-word trace of \fid{F26}
is the constant $24+16/D_0^2$. This is an endomorphism trace,
not automatically the Euler multiplicity of an object;
$e_1e_2$ need not be idempotent.

There is also a valid coefficient-object construction for
\eqref{eq:wordpowers}. On a point let
\[
 \mathcal D=E(-1)\oplus E^{\oplus2}\oplus E(1),\qquad
 \mathcal C_k=E^{\oplus40}\otimes\mathcal D^{\otimes k}.
\]
Geometric Frobenius has trace $D_n$ on $\mathcal D$ and
$40D_n^k$ on $\mathcal C_k$ at every extension.
The canonical tensor identification
$\mathcal C_{k+1}\simeq\mathcal D\otimes\mathcal C_k$
realizes multiplication by this coefficient.
Its dimension is $40\cdot4^k$ and its weights range from
$-2k$ to $2k$, so positivity here does not assert purity.

The one-generator quadratic relation itself has a coherent
functor model. On the product of 64 finite Weil-representation
categories, let $\mathcal E_{40}$ project onto 40 chosen
components and set
$\mathcal P=\mathcal D\otimes\mathcal E_{40}$. Then
\[
 \mathcal P\circ\mathcal P
    \simeq\mathcal D\otimes\mathcal P
\]
by tensor associativity and
$\mathcal E_{40}^2=\mathcal E_{40}$, compatibly with further
powers. This models the split action of one generator and its
quadratic relation. It does not identify a simultaneous model
for the second generator, the original integral tensor lattice,
or the designated nonstandard bases.

Finally the action comparison alone cannot imply the main
sign hypothesis. Take a constant rank-$N$ sheaf with identity
Frobenius on affine $d$-space, with any of the preceding
fixed-parameter actions on its fibers. Its tensor and named
operator relations remain exact. Its pointwise trace is $N$,
so the pointwise hypothesis of \fid{F24} holds with $a=0$
and $B=N$, while its total compact trace is
\[
 M_n=Nq^{dn}.
\]
For $d>0$ the visible root $q^d$ exceeds $q^{d/2}$.
The proved \fid{F24} therefore rules out the required
opposite-sign path comparison with sublinear boundary budget
on an unbounded even set for this model. This proves the
claimed insufficiency with an explicit action-compatible
example.
\end{proof}

Thus a common finite operator realization, a fixed-parameter
action, and word-dependent coefficient objects impose different
comparison requirements. The latter objects may grow with the
word; the common-space obstruction does not exclude them.
For the main route one must additionally identify the actual
path complex and its tensor or idempotent construction by a
Weil-compatible comparison, with the required odd shift and a
controlled cone. Equality of fixed-operator traces alone does
not provide that comparison. Nor does fixed-$q$ numerical
positivity imply coefficientwise positivity in a formal
variable: $x^2-x+1$ is positive for every real $x$ but has a
negative coefficient.


\subsection{An actual stack realization of the normalized word trace}
The modern stack formalism changes the finiteness premise of
\fid{F26}. The primary inputs used here are Sun's trace formula
\cite[Theorem~1.1, p.~49; Definition~4.1, p.~84; Theorem~4.2,
p.~85]{SunStacks}, his explicit classifying-stack calculation
\cite[Theorem~4.8 and \S4.8.2, pp.~89--90; Example~7.4,
p.~109]{SunStacks}, and the published Laszlo--Olsson duality and
K\"unneth formulas \cite[\S9.1, Lemmas~9.1.1--9.1.2, p.~198;
Theorem~11.0.14, p.~204]{LO2}. These statements and their applicable
hypotheses were inspected; their full general proofs are cited inputs.

\begin{finding}[\fid{F35}: a convergent stack trace with an unbounded dual]\label{f:stacktrace}
Over $\F_q$ let
\[
 \mathcal Y=[(\A^6\times\mathbb G_m^8)/(\GL_2)^4]
   =\A^6\times\mathbb G_m^8\times(B\GL_2)^4,
\]
with trivial action, and let $\mathcal X$ be the disjoint union
of 24 rational points and 16 copies of $\mathcal Y$.
Then $\mathcal X$ is a finite-type Artin stack with affine stabilizers,
and for every $n\ge1$, with $Q=q^n$,
\begin{equation}\label{eq:stackwordtrace}
 \#\mathcal X(\F_Q)
 =24+\frac{16Q^2}{(Q+1)^4}
 =\sum_i(-1)^i\Tr(F^n\mid H_c^i(\mathcal X_{\bar{\F}_q},\Q_\ell)).
\end{equation}
Here $\#$ means groupoid mass and the series is absolutely convergent.
This realizes the complete scalar sequence of \fid{F26}.

The component $\mathcal Y$ is smooth of dimension $-2$.
Its compact cohomology is unbounded below and satisfies actual duality
\[
 H_c^{-4-i}(\mathcal Y)\simeq H^i(\mathcal Y)^\vee(2).
\]
Its visible compact spectrum is bounded and admissible as in
\fid{F32}; the corresponding ordinary dual spectrum is unbounded.
Thus actual duality holds, while the separate dual upper-bound
hypothesis of \fid{F17}/\fid{F32} fails in this model.
No nonstandard generator action or simultaneous word comparison
is asserted.
\end{finding}

\begin{proof}
Put $U=\A^6\times\mathbb G_m^8$ and $G=(\GL_2)^4$.
The atlas $U\to[U/G]$ is smooth of relative dimension 16, whereas
$\dim U=14$. This gives the stated dimension and smoothness of
$\mathcal Y$. The full $\mathcal X$ is not equidimensional: its
components have dimensions zero and minus two, and $\dim\mathcal X=0$.

A $\GL_2$-torsor over a field is the frame torsor of a rank-two
vector space, hence is trivial. For the trivial action each
$U(\F_Q)$ point therefore contributes $|G(\F_Q)|^{-1}$ to the mass.
Counting linearly independent columns gives
\[
 |\GL_2(\F_Q)|=(Q^2-1)(Q^2-Q)=Q(Q-1)^2(Q+1).
\]
Consequently
\[
 \#\mathcal Y(\F_Q)
 =\frac{Q^6(Q-1)^8}{Q^4(Q-1)^8(Q+1)^4}
 =\frac{Q^2}{(Q+1)^4}.
\]
Sun's Theorem~1.1 gives the asserted convergent trace for the constant
sheaf. Equivalently Theorem~4.2 applies: the constant sheaf is
bounded, geometrically stratifiable, pure of weight zero, and
convergent in Definition~4.1. The structural map is finite type.
The finite-field base also satisfies the excellence, invertibility
of $\ell$, and finite cohomological-dimension hypotheses of the
cited Laszlo--Olsson formalism.

We now compute the actual cohomology, suppressing the geometric
base change and coefficient field in the notation. Define
\[
 \sum_{s\ge0}a_s z^s=(1-z)^{-4}(1-z^2)^{-4}.
\]
For one classifying factor the known polynomial presentation is
$H^*(B\GL_2)=\Q_\ell[c_1,c_2]$, with $\deg c_i=2i$ and
$c_i$ a rational Chern class spanning $\Q_\ell(-i)$.
Sun's Example~7.4 explicitly gives
\[
 H_c^{-8-2s}(B\GL_2)
 =\Q_\ell(s+4)^{\oplus(\lfloor s/2\rfloor+1)},\qquad s\ge0.
\]
The ordinary cohomology of $\mathbb G_m^8$ is the exterior
algebra on eight copies of $\Q_\ell(-1)$ in degree one, while
$\A^6$ has only the unit in ordinary cohomology. K\"unneth gives
\begin{equation}\label{eq:stackordinarygroups}
 H^i(\mathcal Y)=
 \bigoplus_{\substack{0\le b\le8,\ s\ge0\\i=b+2s}}
 \Q_\ell(-(b+s))^{\oplus\binom8b a_s}.
\end{equation}
For a direct compact-support calculation, the affine factor has
$H_c^{12}(\A^6)=\Q_\ell(-6)$, the torus has
$H_c^{16-b}(\mathbb G_m^8)=\Q_\ell(-(8-b))^{\oplus\binom8b}$,
and the four classifying factors give
$H_c^{-32-2s}((B\GL_2)^4)=\Q_\ell(s+16)^{\oplus a_s}$.
The compact-support K\"unneth formula therefore yields
\begin{equation}\label{eq:stackcompactgroups}
 H_c^j(\mathcal Y)=
 \bigoplus_{\substack{0\le b\le8,\ s\ge0\\j=-4-b-2s}}
 \Q_\ell(b+s+2)^{\oplus\binom8b a_s}.
\end{equation}
The sum in each degree is finite. Since $a_s>0$ for every $s$,
these groups are unbounded below. The indicated weight is
$-2(b+s+2)=j-b$, consistent with Sun's affine-stabilizer bound
\cite[Theorem~1.4, p.~50]{SunStacks}. A fixed degree can be mixed:
$H_c^{-6}=\Q_\ell(3)^4\oplus\Q_\ell(4)^{28}$.

For $x=q^{-n}$ the sign in \eqref{eq:stackcompactgroups} is
$(-1)^b$, and the Euler series becomes
\[
 x^2\frac{(1-x)^8}{(1-x)^4(1-x^2)^4}
 =\frac{x^2}{(1+x)^4}.
\]
This grouping is legitimate: for every real $t>0$ the ungrouped
absolute eigenvalue moment is
\[
 \sum_{b,s}\binom8b a_s q^{-t(b+s+2)}
 =\frac{q^{-2t}(1+q^{-t})^4}{(1-q^{-t})^8}<\infty.
\]
The visible coefficient of $q^{-(j+2)}$ is consequently
\[
 \sum_{b=0}^{\min(8,j)}(-1)^b\binom8b a_{j-b}
   =(-1)^j\binom{j+3}{3}.
\]
For instance at $q^{-3}$ the actual degree-$-6$ multiplicity four
and degree-$-5$ multiplicity eight give the signed coefficient
$4-8=-4$. The full $\mathcal X$ therefore has exactly the
admissible visible spectrum
\[
 (1,24),\qquad
 \left(q^{-(j+2)},16(-1)^j\binom{j+3}{3}\right),\quad j\ge0.
\]
Its support and absolute moments meet \fid{F32}'s hypotheses, so
that finding's uniqueness applies in its stated convergence class.

Finally the smooth dualizing complex of $\mathcal Y$ is
$\Q_\ell(-2)[-4]$. The cited duality formula gives
$H_c^{-4-i}=H^i{}^\vee(2)$, also verified term by term by
\eqref{eq:stackordinarygroups}--\eqref{eq:stackcompactgroups}.
A compact root $\alpha=q^{-(b+s+2)}$ is paired with the
ordinary root $q^{-2}/\alpha=q^{b+s}$. After Euler grouping
the latter roots are all $q^j$, with nonzero coefficients
$(-1)^j\binom{j+3}{3}$. Their support is unbounded, and the
ordinary geometric-Frobenius trace terms do not tend to zero.
Thus this dual is outside the independently bounded visible-spectrum
class. For the nonequidimensional $\mathcal X$, Verdier duality
is computed componentwise; no common dimension-minus-two pairing
is imposed on all its components.
\end{proof}

This is a uniform geometric realization of one scalar word-trace
sequence, with unbounded cohomology explaining its compatibility
with \fid{F26}. It supplies neither operators for $e_1,e_2$ nor
their composition law, and no tensor-path sign or nonstandard basis
follows from its groupoid mass. Its failure of the independent dual
bound is a proved property of this model, not an exclusion of every
stack or completed realization.

\subsection{Finite weight windows and a fixed-parameter action}
\begin{finding}[\fid{F41}: finite windows lose the extension-wise normalization]\label{f:finitewindow}
Every graded Weil subquotient of the compact cohomology of the
stack $\mathcal Y$ in \fid{F35}, restricted to a fixed bounded
weight interval, has finite total dimension and bounded
cohomological support. No fixed endomorphism on such an object
has the alternating all-extension trace
\[
 h_n=24+\frac{16q^{2n}}{(q^n+1)^4}
\]
of the actual normalized word in \fid{F26}.

Nevertheless an actual 64-dimensional subquotient of a Tate
twist of $H_c^{-12}(\mathcal Y)$ supports, after choosing an
identification, the entire actual normalized tensor action
at the single parameter $v_0=q^{1/2}$ or $q^{-1/2}$.
This is a transport construction, not a restriction of an action
already constructed on the stack.
\end{finding}
\begin{proof}
The actual groups computed in \fid{F35} are
\[
 H_c^j(\mathcal Y)=
 \bigoplus_{\substack{0\le b\le8,\ s\ge0\\j=-4-b-2s}}
     \Q_\ell(b+s+2)^{\oplus\binom8b a_s},
 \qquad
 \sum_{s\ge0}a_s z^s=(1-z)^{-4}(1-z^2)^{-4}.
\]
Set $r=b+s+2$. Each summand has weight $-2r$ and scalar
Frobenius $q^{-r}$. A bounded weight interval permits only
finitely many integers $r$, and for each one there are only
finitely many pairs $(b,s)$, since $0\le b\le8$ and
$s=r-b-2\ge0$. Every multiplicity is finite. Also
$j=-2r+b$ takes only finitely many values. This proves the
finiteness statement and the same statement for subquotients.

The $q^{-r}$-eigenspaces are canonical in each degree.
Consequently a degree-preserving Frobenius-commuting action,
if supplied, preserves these windows; any invariant
subquotient retains it. This conditional statement does not
supply that action or make every chosen subspace invariant.
For any fixed endomorphism $T$ on a finite graded window,
its alternating trace with $F^n$ has the form
\[
 \sum_r c_r q^{-rn}
\]
with finitely many coefficients. These are the alternating
traces of the diagonal blocks of $T$; even commutation with
Frobenius is unnecessary for this trace formula.
Put $x_n=q^{-n}$. If it agreed with $h_n$ for every
sufficiently large $n$, the Laurent polynomial $\sum_r c_rx^r$
would agree at infinitely many points with
\[
 24+\frac{16x^2}{(1+x)^4}.
\]
Clearing denominators gives a polynomial with infinitely many
roots, hence an identity. But the latter rational function has
a genuine pole at $x=-1$, while a Laurent polynomial does not.
This contradiction persists after any Tate twist. The more
general finite Weil-complex obstruction is \fid{F26}.

For the positive fixed-parameter construction, degree $-12$
and weight $-12$ require
\[
 b+2s=8,\qquad b+s+2=6,
\]
so $b=0,s=4$. The corresponding actual summand is
\[
 \Q_\ell(6)^{85}\subset H_c^{-12}(\mathcal Y),
 \qquad a_4=35+40+10=85.
\]
The three terms count monomials of respective total degrees
$(4,0),(2,1),(0,2)$ in four variables of weight one and four
variables of weight two. Twist this summand by $(-6)$ and
choose a 64-dimensional subspace $W$. It is an effective
pure weight-zero Weil object with Frobenius exactly one.

Choose a finite coefficient extension $E/\Q_\ell$ containing
$v_0=q^{1/2}$, or its inverse. The actual 64-dimensional tensor
representation in \fid{F26} is defined at this parameter because
$D(q)=q+2+q^{-1}\ne0$. Choose an $E$-linear identification
with $W\otimes E$ and transport its specified normalized
generators. These actual endomorphisms obey all the specialized
relations and commute with scalar Frobenius, preserving every
word action at this one parameter.

For $T=e_1e_2$ their extension traces are consequently
$\Tr(TF^n)=h_1$ for every $n$, rather than $h_n$.
Before the twist they would be $q^{-6n}h_1$.
Raising Frobenius to its $n$th power does not reevaluate fixed
matrix coefficients at $q^{n/2}$.
\end{proof}

Thus a finite effective action model at fixed $q$ exists,
whereas retaining the full extension-wise normalized trace
requires different data. The construction gives no canonical
stack correspondence, restriction of a pre-existing stack
action, or nonstandard basis comparison. Its finite-window
obstruction is conditional on the explicitly stated
all-extension trace requirement.


\subsection{An inverse trace need not be an invertible object}
The stack construction in \fid{F35} permits convergent reciprocal
traces. The next calculation identifies the additional obstruction
to interpreting such a reciprocal as a tensor inverse in the usual
derived Weil realization.

\begin{finding}[\fid{F38}: trace inversion and tensor inversion differ]\label{f:traceinverse}
Over $\F_q$, set
\[
 \mathcal Z=\A^1\times\mathbb G_m^2\times B\GL_2,
 \qquad \mathcal W=\mathbb P^1\times\mathcal Z.
\]
For every $Q=q^n$ their groupoid masses are
\[
 \#\mathcal Z(\F_Q)=\frac1{Q+1},\qquad
 \#\mathcal W(\F_Q)=1.
\]
The corresponding compact Euler traces converge absolutely.
Nevertheless $R\Gamma_c(\mathcal W)$ is not the unit: it has
nonzero cohomology in arbitrarily negative degrees, including
$H_c^{-1}(\mathcal W)=\Q_\ell(1)^2$.

In the usual unbounded derived category of Weil vector spaces,
with the algebraic tensor product,
\[
 H=R\Gamma(\mathbb P^1_{\bar{\F}_q},\Q_\ell)
   =\Q_\ell\oplus\Q_\ell(-1)[-2]
\]
has no tensor inverse. More generally, no tensor-invertible object
in this category has all-extension Euler trace
$D(q^n)=q^n+2+q^{-n}$.
Thus a strong monoidal realization with this specified image of
the projective-line object cannot extend through a tensor localization
making that object invertible. No faithfulness assumption is needed.
\end{finding}

\begin{proof}
Both stacks are smooth, of dimensions $-1$ and zero, and have
affine stabilizers. The torsor and column-count calculation in
\fid{F35} gives
\[
 \#\mathcal Z(\F_Q)
 =\frac{Q(Q-1)^2}{|\GL_2(\F_Q)|}
 =\frac1{Q+1}.
\]
Multiplication by $\#\mathbb P^1(\F_Q)=Q+1$ proves the mass
formula for $\mathcal W$.

Let $c_s=\lfloor s/2\rfloor+1$. Using the classifying-stack
calculation \cite[Example~7.4, p.~109]{SunStacks} and the compact
K\"unneth formula \cite[Theorem~11.0.14, p.~204]{LO2}, we obtain
\begin{equation}\label{eq:reciprocalstackgroups}
 H_c^j(\mathcal Z)=
 \bigoplus_{\substack{0\le b\le2,\ s\ge0\\j=-2-b-2s}}
 \Q_\ell(b+s+1)^{\oplus\binom2b c_s}.
\end{equation}
Indeed $\A^1$ contributes degree two and twist $-1$;
$\mathbb G_m^2$ contributes degree $4-b$ and twist $-(2-b)$;
and $B\GL_2$ contributes degree $-8-2s$ and twist $s+4$.
The sum of these degrees and twists is precisely the displayed
formula. The finite-field and finite-type hypotheses are the
same as those verified in \fid{F35}.

Put $K=R\Gamma_c(\mathcal Z)$. The projective bundle formula,
already used in \fid{F15}, gives the displayed decomposition of $H$.
Properness of $\mathbb P^1$ and K\"unneth give
\[
 R\Gamma_c(\mathcal W)=H\otimes K,\qquad
 H_c^j(\mathcal W)=H_c^j(\mathcal Z)
       \oplus H_c^{j-2}(\mathcal Z)(-1).
\]
Hence $H_c^0(\mathcal W)=\Q_\ell$ and
$H_c^{-1}(\mathcal W)=\Q_\ell(1)^2$.
Formula~\eqref{eq:reciprocalstackgroups} also makes the unbounded
negative support explicit. These are actual cohomology direct sums,
so their nonzero summands cannot cancel.

For $x=q^{-n}$ the compact Euler series of $K$ is
\[
 x\frac{(1-x)^2}{(1-x)(1-x^2)}
   =\frac{x}{1+x}
   =\sum_{j\ge0}(-1)^j x^{j+1}.
\]
Every positive real absolute eigenvalue moment converges: if $t>0$,
the ungrouped sum is
\[
 \sum_{b,s}\binom2b c_s q^{-t(b+s+1)}
 =\frac{q^{-t}(1+q^{-t})}{(1-q^{-t})^2}<\infty.
\]
Thus Sun's trace theorem applies, and this explicit convergence
also justifies all regrouping. Tensoring with the finite complex
$H$ preserves absolute convergence. Since its Euler trace is
$1+x^{-1}$, the Euler trace of $H\otimes K$ is exactly one.
After grouping, its only visible root is one with coefficient one;
the remaining actual cohomology cancels between degrees. The
uniqueness in \fid{F32} concerns that visible spectrum and is not
a faithfulness assertion for the trace.

For completeness, over a field every complex, including an
unbounded one, splits as its cohomology with zero differential
plus a contractible complex. Tensoring a contractible complex
remains contractible. With algebraic direct-sum totalization,
the resulting K\"unneth formula is
\[
 H^m(A\otimes B)=\bigoplus_{i+j=m}H^i(A)\otimes H^j(B).
\]
If $A\otimes B$ is the one-dimensional unit, both factors are
nonzero. Two independent total cohomology vectors in either
factor would give at least two in their tensor product.
Consequently each factor has one-dimensional total cohomology,
concentrated in a single degree. Forgetting the Weil action is
exact and conservative, so the same necessity holds for Weil
complexes. Truncation then identifies each factor with a
one-dimensional Weil representation in one degree. Conversely
such a representation, with its shift, has its dual with the
opposite shift as inverse.

In particular $H$ has no inverse. Directly, if $H^a(B)\ne0$,
then $H\otimes B$ has nonzero cohomology in both degrees $a$
and $a+2$. Strong monoidal functors preserve tensor inverses.
Therefore no strong monoidal realization sending a designated
object to this $H$ can factor through a localization in which
that object becomes invertible.

For the nonstandard normalization, under $v=Q^{\pm1/2}$ we have
\[
 D(Q)=(v+v^{-1})^2=Q^{-1}(Q+1)^2.
\]
The actual complexes
\[
 H_D=\Q_\ell(1)\otimes H^{\otimes2},\qquad
 K_D=\Q_\ell(-1)\otimes K^{\otimes2}
\]
have traces $D(Q)$ and $D(Q)^{-1}$, respectively. Their product
has trace one, but $H_D$ has total cohomology dimension four and
cannot have a tensor inverse. This conclusion does not depend
on choosing $H_D$: every invertible Weil complex has trace
$\pm\alpha^n$ for a single nonzero Frobenius scalar $\alpha$.
The minus sign cannot give the positive first and second values
of $D(q^n)$. The plus sign would require $D(q^2)=D(q)^2$, whereas
\[
 D(q)^2-D(q^2)=4(q+q^{-1}+1)>0.
\]
This proves the asserted all-extension obstruction.
\end{proof}

The statement permits formal Euler-class localizations, other
realizations, and completions with a different tensor product or
trace functional. It also permits multiplying an operator by the
numerical scalar $D(q)^{-1}$ at a fixed $q$. Scalar division does
not assert an object-level inverse or a reevaluation at every
$q^n$. The fixed realization and usual derived tensor product
are essential hypotheses of the obstruction.

\subsection{Actual pure classical multiplicity objects via geometric Satake}
\begin{finding}[\fid{F20}: Schur summands have pure multiplicity spaces; their cycle action need not be positive]\label{f:satake}
Let $G=\GL_r/\F_q$ be split, with $\ell\nmid q$, and use the symmetric geometric Satake category with $\overline\Q_\ell$ coefficients. For a dominant partition $\mu$ of length at most $r$, let $\IC_\mu$ be the natural unnormalized Weil intersection complex of its Schubert variety, perverse of weight
\[
 D_\mu=\dim\Gr_\mu=\langle2\rho,\mu\rangle
 =\sum_{i=1}^r(r+1-2i)\mu_i.
\]
For $\nu\vdash d$, a primitive Young idempotent $e_\nu\in\Q[S_d]$ acts, through the parity-corrected symmetric commutativity, on
$K=\IC_\mu^{\star d}$. Its image $P_{\nu,\mu}$ is a pure perverse Weil sheaf of weight $dD_\mu$. The canonical multiplicity Weil module
\[
 M_{\lambda;\nu,\mu}
 =\Hom_{\mathrm{Sat}_{\overline{\F}_q}}
 (\IC_\lambda,P_{\nu,\mu})
\]
is pure of weight $dD_\mu-D_\lambda$ and has dimension
\[
 \dim M_{\lambda;\nu,\mu}
 =[S_\lambda(\overline\Q_\ell^r):
   S_\nu(S_\mu(\overline\Q_\ell^r))].
\]
On the product Grassmannian for $\GL_m\times\GL_n$, apply the Schur idempotent for $\nu\vdash d$ to the $d$th convolution power of
$\IC_{(1),m}\boxtimes\IC_{(1),n}$. Its multiplicity module at
$\IC_{\lambda,m}\boxtimes\IC_{\mu,n}$ has dimension $g(\lambda,\mu,\nu)$ and is pure of weight
\[
 d(m+n-2)-D_{\lambda,m}-D_{\mu,n}.
\]
These are ordinary geometric realizations of the classical coefficients, including rectangular cases within the dimension bounds. They do not by themselves identify the nonstandard quantum bases.

Even in this construction the symmetric-group action on a geometric multiplicity space need not permute a basis. For $G=\GL_2$, $\mu=(1)$ and $d=2$, the determinant multiplicity line is $\overline\Q_\ell(-1)$ and the transposition acts by $-1$. Thus no basis of that line is an $S_2$-permutation basis.
\end{finding}
\begin{proof}
The literature inputs are geometric Satake, its finite-field convolution/symmetry, and ordinary weight theory. We use the published version of Richarz \cite[Theorem~3.1, p.~214; proof pp.~226--227; Proposition~4.1, p.~227; Theorem~5.1, p.~230; Lemma~5.4 and its proof, p.~232]{Richarz}. The tensor equivalence is also Mirkovi\'c--Vilonen, Theorem~7.3 \cite{MV}. These theorems are cited, not re-proved here.

The twisted product of $d$ Schubert varieties is projective and is smooth-locally the product for purposes of its intersection complex. Its natural Weil IC is normalized to restrict to $\overline\Q_\ell[dD_\mu]$ on the dense twisted orbit. With that normalization, rather than only a geometric IC isomorphism, it is the twisted external product of the $\IC_\mu$ and is pure of weight $dD_\mu$; this is the IC identification in the proof of Richarz's Lemma~5.4, together with BBD~5.3.2, second proof (which covers a nonaffine open immersion). Its proper convolution pushforward is pure of that same complex weight (BBD, Remark~5.4.9) and is perverse (Richarz, Theorem~3.1). This proves the assertion about $K$.

The symmetry constructed over $\F_q$, with the component-parity sign specified in the cited proof, gives an $S_d$ action by Weil endomorphisms. In the abelian perverse category the idempotent splits; its image is a direct summand, hence remains pure by BBD~5.3.1. Under the symmetric monoidal Satake equivalence, $\IC_\mu$ corresponds to $S_\mu(\overline\Q_\ell^r)$ and $K$ to its $d$fold tensor power. In Schur--Weyl decomposition, the primitive idempotent $e_\nu$ has rank one on the Specht factor $[\nu]$ and vanishes on the other factors. Its image therefore corresponds to one copy of the Schur functor $S_\nu(S_\mu(\overline\Q_\ell^r))$. A central isotypic idempotent would instead give $\dim[\nu]$ copies; that is why the primitive condition matters.

Here is the Frobenius argument for the multiplicity weight. Geometric semisimplicity gives the canonical evaluation isomorphism
\[
 \bigoplus_\lambda\IC_\lambda\otimes
 \Hom(\IC_\lambda,\overline P_{\nu,\mu})
 \longrightarrow\overline P_{\nu,\mu}.
\]
It is the canonical isotypic decomposition, not an arbitrarily chosen decomposition. Every support and natural $\IC_\lambda$ is defined over the split base $\F_q$, so Frobenius preserves each summand. Give its Hom factor the conjugation Frobenius induced by the two fixed Weil structures. Evaluation is then Weil-equivariant, and the isotypic summand is a pure perverse Weil sheaf of weight $dD_\mu$. This descent/subquotient formalism is BBD~5.1.2; purity of the summand is BBD~5.3.1. On its own smooth open Schubert stratum inside its support, $\IC_\lambda$ is the constant sheaf shifted by $D_\lambda$. Restriction to this open preserves both weight inequalities because $j^*=j^!$ (BBD~5.1.14). The Hom module shifted by $D_\lambda$ is therefore pure of weight $dD_\mu$, which means the module itself has weight $dD_\mu-D_\lambda$ by the shifted convention of BBD~5.1.8. This proves the all-eigenvalue purity statement without assuming arithmetic semisimplicity. The tensor equivalence identifies its dimension with the required multiplicity.

For the product group, external product gives the simple object corresponding to $V\boxtimes W$. Its weight is $(m-1)+(n-1)$. The same Schur-idempotent construction corresponds to $S_\nu(V\otimes W)$ as a $\GL_m\times\GL_n$ representation. The definition of the Kronecker coefficient and the just-proved Hom-weight calculation give the stated formula.

Finally consider $\GL_2$ and two minuscule steps. The twisted convolution variety is a smooth surface. Over the central lattice $t\mathcal O^2$, its fiber consists of the intermediate lattices between $t\mathcal O^2$ and $\mathcal O^2$, hence of lines in $\mathcal O^2/t\mathcal O^2$; it is $\PP^1$. The contribution to the perverse degree-zero stalk is $H^2(\PP^1)=\overline\Q_\ell(-1)$. To identify it as the determinant multiplicity, Satake gives
\[
 V\otimes V=\Sym^2V\oplus\wedge^2V,
\]
with determinant multiplicity one. Equivalently the non-open relevant fiber supplies this unique one-dimensional summand; the IC of the open two-dimensional stratum has no perverse degree-zero stalk contribution at the boundary point by the intermediate-extension condition. The corrected symmetric commutativity corresponds to the ordinary tensor swap, which sends
$e_1\otimes e_2-e_2\otimes e_1$ to its negative. The sole nonzero basis vector of a one-dimensional permutation representation would have to be fixed, so no such basis exists. The symmetric projector kills this line, while the alternating projector keeps it.
\end{proof}
This is a constructive reduction of the ordinary-geometry gate for \emph{classical} coefficients. Its purity uses the existing Weil/BBD theorems; it is not a new application of \fid{F03}, and it does not provide the path-hypersurface identities for a nonstandard canonical coefficient. The remaining problem is more specific: describe the actual projector image, its graded basis and its compatibility with the designated quantum operations, together with effective labels when a counting rule is required. The $\GL_2$ example prevents replacing that task by unsigned orbit-counting on the unprojected component basis.

The construction also provides an exact Grothendieck-framework dictionary. Convolution is a proper constructible pushforward; the Young idempotent is an actual projector in a characteristic-zero symmetric tensor category; its image is an effective object whose $K_0$ class is the Schur operation. The weights record the unnormalized IC shifts, so no half-Tate convention is hidden. The motivating language of projectors occurs in cote~16, scans~35, 37 and~39; the actual theorem here uses the modern Satake equivalence and the precisely cited SGA/BBD weight formalism, not an attribution to those leaves. The explicit symmetric-group action is the additional datum beyond virtual $\lambda$ identities.

\paragraph{\fid{F20}, motivic refinement: scalar Tate Frobenius.}
In the rational motivic Satake category over $\F_q$, these same Schur
projector images have coefficient motives
\[
 \mathbf1(-h)^{\oplus a},\qquad
 h=\frac{dD_\mu-D_\lambda}{2}
   =n(\lambda)-d\,n(\mu),\qquad
 n(\alpha)=\sum_i(i-1)\alpha_i,
\]
where $a=[S_\lambda V:S_\nu(S_\mu V)]$. For the product-group Kronecker
construction the coefficient is
\[
 \mathbf1\bigl(-n(\lambda)-n(\mu)\bigr)^{\oplus g(\lambda,\mu,\nu)}.
\]
Consequently the natural Frobenius in the \emph{specified} Weil
multiplicity spaces is the scalar $q^h$ and has no Jordan part.

\begin{proof}
The input is Richarz--Scholbach \cite[Theorem~4.1; Proposition~4.8;
Proposition~5.9; Definition~6.2; Lemma~6.3; Corollary~6.4;
Lemma~6.5]{RS}. These provide pure intersection motives, weight-exact
convolution, symmetric constraints compatible with the corrected
\(\ell\)-adic realization, and a semisimple abelian compact Satake
category generated by the $\IC_\lambda(n)$ and stable under convolution.
Thus the primitive rational Young idempotent acts motivically and its
image is pure of weight $w=dD_\mu$. Its simple summands have the form
$\IC_\lambda(n)$, of weight $D_\lambda-2n$. Purity forces
$n=(D_\lambda-w)/2=-h$. Compatibility with geometric Satake identifies
their multiplicity with $a$. Restriction of the corresponding isotypic
summand to its open stratum and removal of the shift $D_\lambda$ gives
exactly $\mathbf1(-h)^{\oplus a}$; Lemma~6.3 gives the coefficient-motive
Hom identification. Since $|\lambda|=d|\mu|$ and
$D_\alpha=(r-1)|\alpha|-2n(\alpha)$, the stated integer formula for
$h$ follows. The product-group proof replaces $D_\lambda$ by
$D_{\lambda,m}+D_{\mu,n}$ and $w$ by $d(m+n-2)$, giving
$h=n(\lambda)+n(\mu)$. Finally the realization of $\mathbf1(-h)$ is
$\overline\Q_\ell(-h)$, whose geometric Frobenius is exactly $q^h$.
\end{proof}
Here the coefficient motive is not the untwisted vector space
$\Hom_{\rm mot}(\IC_\lambda,P)$, which can vanish when $h\ne0$.
The relevant multiplicity vector space is
$\Hom_{\rm mot}(\IC_\lambda(-h),P)$; tensoring it with the Tate line
recovers the Weil multiplicity module. This is stronger than purity for
these natural classical objects, and uses the modified dual-group
weight coordinate of \cite[Theorem~6.8]{RS}. It neither chooses a basis
of the coefficient vector space nor computes its rank efficiently.
Arithmetic semisimplicity for arbitrary pure Weil modules is still
false, as in \fid{F05}.

\subsection{A precise positive rule for the Schur-projector gate}
\begin{finding}[\fid{F23}: signed orbits suffice, but need not exist in pure Satake spaces]\label{f:signedorbits}
Let $M$ be a characteristic-zero $S_d$-module with a specified basis
$B$ permuted up to signs. For a label orbit $O$, choose $b_O$, let
$H_O$ be its label stabilizer, and define
$\chi_O(h)$ by $hb_O=\chi_O(h)b_O$. Then
\[
 M\simeq\bigoplus_O\operatorname{Ind}_{H_O}^{S_d}\chi_O,
\]
and
\[
 \dim M^{S_d}=\#\{O:\chi_O=1\},\qquad
 \dim\Hom_{S_d}(\mathrm{sgn},M)
 =\#\{O:\chi_O=\mathrm{sgn}|_{H_O}\}.
\]
For a primitive Young idempotent $e_\nu$,
\begin{equation}\label{eq:orbitprojector}
 \dim e_\nu M
 =\sum_O\dim\Hom_{H_O}(\Res^{S_d}_{H_O}[\nu],\chi_O).
\end{equation}
If every $H_O$ is a Young subgroup $\prod_i S_{a_i}$ and $\chi_O$ is a
product of trivial or sign characters on its blocks, this last sum
counts chains from $\varnothing$ to $\nu$: at block $i$ add a horizontal
$a_i$-strip for the trivial character, or a vertical $a_i$-strip for sign.

These are actual positive formulas for the corresponding Schur
multiplicity whenever the specified signed-basis data exist on
\[
 M_{\lambda,\mu,d}
 =\Hom_{\GL(V)}(S_\lambda V,(S_\mu V)^{\otimes d}).
\]
They do not always exist: for $\GL_2$, $\mu=(1)$, $d=3$ and
$\lambda=(2,1)$ this multiplicity module admits no signed-permutation
basis, although its natural geometric realization is
$\overline\Q_\ell(-1)^2$ with scalar Frobenius.
\end{finding}
\begin{proof}
The sign cocycle implies that $\chi_O:H_O\to\{\pm1\}$ is a character.
The map $k[S_d]\otimes_{k[H_O]}k_{\chi_O}\to kO$ sending
$g\otimes1$ to $gb_O$ is well defined, equivariant and surjective.
Both spaces have dimension $[S_d:H_O]$, so it is an isomorphism.
The distinct orbit spans are a direct sum. Frobenius reciprocity gives
\eqref{eq:orbitprojector}, and the trivial/sign cases have a
one-dimensional Hom precisely when the indicated characters agree.
In $M=\bigoplus_\rho[\rho]\otimes W_\rho$, a primitive idempotent has
rank one on $[\nu]$, hence image dimension $\dim W_\nu$.
A central idempotent would have rank $\dim[\nu]\dim W_\nu$ instead.
When the stabilizer is the stated Young subgroup, the Frobenius
characteristic is a product of $h_{a_i}$ and $e_{a_i}$. The classical
Pieri rule \cite[Theorem~7.10, p.~16]{Pieri} expands this product by exactly
the horizontal/vertical chains just described. This proves the positive
chain formula, without presupposing the desired projector-image basis.

For the counterexample, Schur--Weyl duality gives
\[
 (\overline\Q_\ell^2)^{\otimes3}
 =S_{(3)}V\otimes[(3)]\oplus S_{(2,1)}V\otimes[(2,1)].
\]
Thus the indicated multiplicity is the faithful two-dimensional
standard representation of $S_3$. It is faithful because it is the
sum-zero subspace of the permutation representation on three letters,
whose complementary all-ones line is fixed. A signed-permutation basis
would embed $S_3$ into the group of $2\times2$ signed permutation
matrices, of order $2^2\cdot2!=8$. Lagrange's theorem excludes this.

Here is the geometric identification in full. The threefold minuscule
convolution variety for $\GL_2$ is an iterated $\PP^1$-bundle, smooth of
dimension three. Over the lattice
$t^2\F_q[[t]]e_1\oplus t\F_q[[t]]e_2$, the quotient has dimension
three and multiplication by $t$ is nilpotent of type $(2,1)$. The fiber
parametrizes its full stable flags $L\subset P$. Write
$I=\operatorname{im}N$, $K=\ker N$, of dimensions one and two.
Stability requires $L\subset K$, $I\subset P$, and $N(P)\subset L$.
If $P\ne K$, then $N(P)=I$ and $L=I$; if $P=K$, every line
$L\subset K$ works. Thus the reduced fiber is
\[
 \{L=I,\ P\supset I\}\ \cup\ \{P=K,\ L\subset K\}
 =\PP^1\cup_{\mathrm{pt}}\PP^1.
\]
Both components and the intersection are rational over $\F_q$.
Its $H^2$ is therefore $\overline\Q_\ell(-1)^2$. Since the source
dimension is three and $D_{(2,1)}=1$, proper base change identifies this
group with the degree $-1$ stalk of the pure perverse convolution.
The intermediate extension $\IC_{(3)}$ has no degree
$-\dim\Gr_{(2,1)}=-1$ stalk on this boundary stratum, by its strict
support condition. This stalk is exactly the $\IC_{(2,1)}$ multiplicity
module. Corrected Satake symmetry gives the standard $S_3$ action just
used, while Frobenius is scalar $q$. This proves every assertion about
the example.
\end{proof}
The obstruction is to signs $\{\pm1\}$, not arbitrary scalar-monomial
bases: the standard $S_3$ representation is induced from a nontrivial
character of its cyclic subgroup of order three. Pure Tate Frobenius
does not restrict the phases of this separate commuting group action.
The example also has classical positive tableau formulas; it refutes
only the asserted universal signed-orbit route.

\paragraph{Exact encoding conditions for a counting rule.}
For inputs of binary length $L$, suppose orbit representatives have
unique polynomial-length encodings with polynomial-time validity tests.
For the trivial/sign formula, suppose a polynomial-size stabilizer
generating set and its character values are computable in polynomial
time, including the cost of permutation/sign encoding when $d$ is
binary. The character test on generators makes the orbit label itself
a counted witness. For the mixed-Young formula, additionally require
polynomially many blocks and rows, with polynomial-bit block sizes and
parts, and a polynomial-time procedure computing the block sizes and
trivial/sign types from the orbit label. A witness consists of that label
and its chain of partitions.
All parts are bounded by $d$, so the witness is polynomial length.
For $\alpha\subset\beta$, a horizontal strip is checked by
$\beta_i\ge\alpha_i\ge\beta_{i+1}$; a vertical strip by
$\beta_i-\alpha_i\in\{0,1\}$. Sizes and partitions are checked by
integer arithmetic. These tests count exactly the displayed chains,
giving a conditional \(\#P\) rule. They do not provide a single
polytope, an Ehrhart family or PH1's oracle from purity alone.

\paragraph{A known binary family, with its exact small polytope.}
For $\dim V=2$, Hermite reciprocity and Clebsch--Gordan with swap parity
give the classical formula
\[
 \Sym^d(\Sym^2 V)\simeq\Sym^2(\Sym^dV)
       \simeq\bigoplus_{j=0}^{\lfloor d/2\rfloor}S_{(2d-2j,2j)}V.
\]
The exact sources are Goodman--Wallach, Theorem~4.1.18, p.~191,
and Donley--Kim, Proposition~3.1, equation~(3.4), p.~3, and
Proposition~8.1, equation~(8.2), p.~11 \cite{GW,DK}.
These known theorems are not re-proved here. The passage from $SL_2$
to $GL_2$ fixes total polynomial degree $2d$; in the summand
$S_{(2d-r,r)}$, a highest-weight vector is
$(x_1y_2-x_2y_1)^r x_1^{d-r}x_2^{d-r}$ and swap acts by $(-1)^r$.
For $\lambda=(2d-r,r)$, $0\le r\le d$, the specialized consequence is
\[
 a_\lambda(d[2])=\mathbf1_{2\mid r},\qquad
 a_{N\lambda}((Nd)[2])
 =\#(N P_{d,r}\cap\Z),\quad P_{d,r}=\{r/2\}.
\]
The point is given by $2x\le r$, $-2x\le-r$, so has an explicit
polynomial-bit description and separation test. This is a restricted
fixed-inner-partition Ehrhart formula, not general PH1. On odd-$r$
rays its ordinary support is unsaturated, but the zero odd residue
polynomial is permitted by VI Definition~1.2.4, so the VI saturation
index remains zero, exactly as distinguished in \fid{F18}.


\section{What is reached in the GCT cases requested}\label{sec:rectangles}
We keep three different notions separate: scalar nonnegativity of a representation multiplicity; a positive effective rule for it; and positivity of normalized canonical structure constants. The first is already part of semisimple representation theory and \fid{F12}. The others do not follow from it.

Write $n\times d=(d^n)$, $\lambda\vdash nd$, and
\[
 g_n(\lambda)=g(\lambda,(d^n),(d^n)),\qquad
 a_\lambda(d[n])=[S_\lambda V:\Sym^d(\Sym^nV)].
\]
Let $q^m_\lambda(d[n])$ denote the degree-$d$ multiplicity in the coordinate ring of the orbit closure of $X_{11}^{n-m}\operatorname{per}_m$, in IP's covariant convention. The known inequality is $0\le q^m_\lambda(d[n])\le a_\lambda(d[n])$ (IP, equation~(1.3)).

\subsection{Exact rectangular results already proved}
The attached Ikenmeyer--Panova paper \cite{IP}, matching arXiv 1512.03798v2, states:
\begin{itemize}
\item Theorem~1.4: $q^m_\lambda(d[n])>g_n(\lambda)$ implies $\operatorname{dc}(\operatorname{per}_m)>n$.
\item Theorem~1.6: if $n>3m^4$ and $g_n(\lambda)=0$, then $q^m_\lambda(d[n])=0$.
\item Theorem~1.8: $q^m_\lambda(d[n])>0$ requires $\ell(\lambda)\le m^2$ and $|\bar\lambda|\le md$, where $\bar\lambda$ removes the first row.
\item Theorem~1.10: put
$\mathcal X=\{(1),(1,1),(1^4),(1^6),(2,1),(3,1)\}$.
If $\bar\lambda\in\mathcal X$, then $a_\lambda(d[n])=0$. If $\bar\lambda\notin\mathcal X$ and there is $m\ge3$ with
\[
 \ell(\lambda)\le m^2,\quad |\bar\lambda|\le md,\quad
 d>3m^3,\quad n>3m^4,
\]
then $g_n(\lambda)>0$.
\end{itemize}
These are substantive existing positivity and no-occurrence results. They are not new consequences of our conditional sign theorem. Revised IV, Theorems~1.6/15.21(vii)--(viii), also supplies a positive highest-weight crystal rule when the two restricted factors have at most two rows (and the $\GL_4$ label at most four parts). That includes two-row rectangles, not arbitrary-height rectangles in the asymptotic regime.

\subsection{An obstruction for every nontrivial rectangular tensor square}
\begin{finding}[\fid{F28}: an isolated hook excludes all Young-stabilizer monomial models]\label{f:recthooks}
Let $M$ be a characteristic-zero $S_N$-module. If an internal hook
degree $1\le k\le N-2$ occurs in $M$, while its two neighboring
hook degrees do not, then $M$ has no scalar-monomial basis all of
whose label stabilizers are conjugate to Young subgroups.

Consequently, for \emph{every} rectangle $\rho=(a^b)$ with $a,b\ge2$,
the actual module $[\rho]\otimes[\rho]$ has no such basis, even when
arbitrary nonzero scalar phases are allowed. Thus the Young-stabilizer
version of \fid{F23} cannot describe the unprojected rectangular
Satake multiplicity module, including large-parameter rectangles.
This does not exclude non-Young stabilizers or other positive rules.
\end{finding}
\begin{proof}
For a Young subgroup $H=\prod_{i=1}^r S_{a_i}$ with nonempty blocks,
every one-dimensional character is a product of trivial and sign
characters. Set $s=\sum_{\text{sign blocks }i}(a_i-1)$; singleton
blocks contribute zero under either designation. If $U_N$ is the
standard sum-zero representation, then
\[
 \Res_H U_N=\bigoplus_i U_{a_i}\oplus k^{r-1},\qquad U_1=0.
\]
The classical identification $\wedge^jU_a=[(a-j,1^j)]$ follows here
from the cited Pieri rule: the oriented $j$-subset module
$\wedge^j k^a=\operatorname{Ind}_{S_j\times S_{a-j}}^{S_a}
(\mathrm{sgn}\boxtimes1)$ is the sum of the two adjacent hooks for $1\le j\le a-1$,
and $k^a=1\oplus U_a$ identifies their exterior powers inductively.
In the exterior power of the displayed restriction, trivial on block
$i$ forces degree zero, and sign forces degree $a_i-1$. The remaining
degree is chosen from the trivial $(r-1)$-space. Frobenius reciprocity
therefore gives the exact formula
\[
 [\operatorname{Ind}_H^{S_N}\chi:(N-j,1^j)]
       =\binom{r-1}{j-s}.
\]
Its positive hook support is the consecutive interval
$s,\ldots,s+r-1$, contained in $0,\ldots,N-1$ because
$0\le s\le N-r$.

A scalar-monomial basis decomposes into actual orbit modules
$\operatorname{Ind}_{H_O}^{S_N}\chi_O$, by the orbit-map proof in
\fid{F23}. Some summand must contain the given internal hook. If its
Young subgroup has $r\ge2$ blocks, the interval just computed includes
at least one neighboring hook, contradicting its absence from the
effective direct sum. If $r=1$, the orbit is the trivial or sign
representation of the whole group, whose hook degree is $0$ or
$N-1$, not internal. This proves the first assertion. There is no
cancellation in an effective representation decomposition.

The exact source input for rectangles is Ikenmeyer--Panova
\cite[Proposition~6.4, pp.~15--16]{IP}. With $m=\min(a,b)$ and $N=ab$,
it states
\begin{equation}\label{eq:recthookproduct}
 \sum_{j=0}^{N-1}g((N-j,1^j),\rho,\rho)t^j
       =\prod_{i=2}^{m}(1+t^{2i-1}).
\end{equation}
Transposition of both rectangular Specht factors justifies the
symmetric shorter-side convention. For $N\ge6$, degree three is
internal. The product has coefficient one at degree three and zero
at degrees two and four. Apply the isolated-hook assertion.
The only smaller rectangle is $(2,2)$ of size four. Its square is
$[(4)]\oplus[(2,2)]\oplus[(1^4)]$ by the already checked character
calculation, or Manivel's Theorem~1. Its hook degrees are $0,3$.
Every Young orbit with at least two blocks would force consecutive
hooks; all remaining whole-group orbits are one-dimensional and
cannot supply $[(2,2)]$. This proves the edge case as well.
\end{proof}

The actual multiplicity identification is explicit: for $\dim U=\dim V=b$,
Schur--Weyl duality gives
\[
 \Hom_{\GL(U)\times\GL(V)}
 (\det^aU\boxtimes\det^aV,(U\otimes V)^{\otimes ab})
 \simeq[(a^b)]\otimes[(a^b)]
\]
with its diagonal symmetric-group action. Its primitive Schur
projectors give precisely the rectangular Kronecker coefficients.
The obstruction thus concerns the requested actual family, not an
unrelated test representation. A virtual alternating induction formula
can cancel the forbidden hooks, but would not be the positive orbit
decomposition considered here.

The scope is sharp already at size four. The permutation module on
the three pairings of four letters is $1\oplus[(2,2)]$, so
$[(2,2)]^{\otimes2}$ is that module plus a sign line: it does have
a signed-permutation model. A pairing stabilizer is
$S_2\wr S_2$, of order eight, whereas a Young subgroup of $S_4$ has
order $24,6,4,2$ or $1$. The stabilizer is not Young.

\paragraph{A working ordinary monomial model in the known height-two family.}
Manivel \cite[Theorem~1, p.~1; \S\S2--3, pp.~2--4]{Manivel} proves
that $[(d,d)]^{\otimes2}$ is multiplicity free, with constituents the
partitions of $2d$ having at most four even parts, or exactly four odd
parts. In fact, for dimensions $(\dim U,\dim V,\dim W)=(2,2,4)$,
his proof identifies the highest-weight invariant algebra
\[
 \C[U\otimes V\otimes W]^{SL(U)\times SL(V)\times N_W}
 =\C[f_1,f_2,f_3,f_4]
\]
in the source's dual-polynomial convention. The four $W$-partition
labels are $(1^4)$, $(2)$, $(2,2)$ and $(2,2,2)$, with degrees $4,2,4,6$.
These are known source results, not consequences newly obtained from
purity. A monomial with label $\nu$, padded to four parts, necessarily
has exponents
\[
 a_1=\nu_4,\quad a_2=(\nu_1-\nu_2)/2,\quad
 a_3=(\nu_2-\nu_3)/2,\quad a_4=(\nu_3-\nu_4)/2.
\]
They are nonnegative by the partition inequalities, and integral
exactly when the padded parts have common parity. The two determinant
exponents are $d=2a_1+a_2+2a_3+3a_4=|\nu|/2$.
Thus this is an actual known monomial basis with a verified
multiplicity comparison, in a different representation-theoretic model.

Equivalently, for $\ell(\nu)\le4$, let
\[
 P_\nu=\{((\nu_1-\nu_4)/2,(\nu_2-\nu_4)/2,(\nu_3-\nu_4)/2)\}.
\]
Its coordinates are $(a_2+a_3+a_4,a_3+a_4,a_4)$, an integral
unimodular change from the last three exponents. The exact restricted
Ehrhart formula is
\[
 g(N\nu,(Nd,Nd),(Nd,Nd))=\#(NP_\nu\cap\Z^3).
\]
Use the empty polytope for length greater than four. The point has
denominator at most two and a polynomial-bit description and
separation test. It is the source's actual exponent rule, not a
point assigned after consulting a zero/nonzero oracle. Its torus
monomial basis, for $d\ge2$, is not a scalar-monomial $S_{2d}$ Young-orbit basis;
the latter is exactly what \fid{F28} excludes.


\subsection{Rectangular hooks in ordinary group cohomology}
This is an application of the known rectangular hook formula and
a known split-group motive decomposition, rather than a new
Kronecker rule.

\begin{finding}[\fid{F33}: a group-scheme realization with mixed degrees]\label{f:hookgeometry}
Let $n,d\ge1$, $N=nd$, $m=\min(n,d)$ and $0\le k<N$.
For $G=SL_{m,\F_q}$ and $\ell\nmid q$,
\[
 g((N-k,1^k),(d^n),(d^n))
   =\dim_{\Q_\ell}H^k(G_{\bar{\F}_q},\Q_\ell).
\]
More precisely there is an actual Weil decomposition
\begin{equation}\label{eq:hookweil}
 H^k(G_{\bar{\F}_q},\Q_\ell)=
 \bigoplus_{\substack{I\subseteq\{2,\ldots,m\}\\
                     \sum_{i\in I}(2i-1)=k}}
       \Q_\ell\left(-\sum_{i\in I}i\right).
\end{equation}
The degree can be mixed: for $m=8$,
\[
 H^{15}(SL_{8,\bar{\F}_q},\Q_\ell)
       =\Q_\ell(-8)\oplus\Q_\ell(-9).
\]
This realizes the numerical coefficients in ordinary cohomology;
it does not identify a nonstandard action or canonical basis.
\end{finding}

\begin{proof}
The known formula in IP, Proposition~6.4, is
\[
 \sum_{k=0}^{N-1}g((N-k,1^k),(d^n),(d^n))t^k
       =\prod_{i=2}^m(1+t^{2i-1}).
\]
For the geometric input, Merkurjev gives the canonical split
special-linear-group motive decomposition over a perfect field
\cite[\S5, Corollary~5.5, p.~8; notation on p.~7]{MerkurjevSL}.
In his notation a sequence
$1\le j_1<\cdots<j_a\le m-1$ contributes
$\mathbf1(d+a)[2d+a]$, with $d=\sum j_r$.
Reindexing by $I=\{j_1+1,\ldots,j_a+1\}$ and
$s_I=\sum_{i\in I}i$ gives
\[
 M(SL_m)_{\Q}
       \simeq\bigoplus_{I\subseteq\{2,\ldots,m\}}
                 \mathbf1(s_I)[2s_I-|I|].
\]
The displayed known motive theorem is invoked, not reproved.
Finite fields satisfy its perfect-field hypothesis.
Under ordinary cohomological $\ell$-adic realization, the summand
$\mathbf1(s)[a]$ contributes $\Q_\ell(-s)$ in degree $a$.
The decomposition is defined over $\F_q$, so Frobenius on that
summand is exactly $q^s$. This proves \eqref{eq:hookweil}.
Taking dimensions yields the product above, hence the required
coefficient identity.

For a summand in degree $k=2s_I-|I|$, its weight is
$2s_I=k+|I|$. Different subset cardinalities can therefore give
different weights in the same degree. At $m=8,k=15$, the only
subsets are $\{8\}$ and $\{2,3,4\}$. A two-element subset has
even degree, while four or more elements have minimum degree
$3+5+7+9=24$. Among three-element subsets the minimum 15 is
attained only by $\{2,3,4\}$. Their sums are eight and nine,
giving the stated Tate summands, weights 16 and 18, and traces
$q^{8r}+q^{9r}$ over $\F_{q^r}$.
\end{proof}

The corresponding hook coefficient is two for every rectangle
whose minimum side is eight. A single Tate twist cannot make the
two displayed weights equal. This is consistent with \fid{F20},
whose pure Satake multiplicity motive is a different realization
of the same integer. The affine group $G$ supplies ordinary
cohomology; no assertion of projectivity or purity of each
cohomological degree is needed.


\subsection{The open-orbit method has a precise higher-height limit}
\begin{finding}[\fid{F29}: the same flag-orbit mechanism stops at height two]\label{f:sphericity}
Let $n\ge3$, $E=\C^n\otimes\C^n$, $N=n^2$, and
$G=SL_n\times SL_n$. For every $2\le r\le N$, $G$ has no
open orbit on $\operatorname{Flag}(1,\ldots,r;E)$.
Equivalently, $G\times B_r$ has no open orbit on $E\otimes\C^r$,
where $B_r$ is a Borel subgroup. For $r=2$ the generic orbit has
codimension $n-2$. The same open-orbit proof underlying Manivel's
height-two formula therefore does not extend to these heights,
even after restricting the output length to two.
\end{finding}
\begin{proof}
On the dense locus of injections $(\C^r)^*\to E$, quotienting
by the appropriate dual Borel records the flag of successive images.
For each flag the compatible ordered injections form one free Borel
orbit. This is the usual frame bundle of the filtered tautological
bundle, locally a principal Borel bundle. It proves the equivalence
of the two open-orbit questions. The partial flag dimension is
\[
 rN-r(r+1)/2,\qquad \dim G=2(N-1).
\]
At $r=3$ the difference is $N-4>0$; successive increments of the
flag dimension are $N-r-1\ge0$ for $3\le r<N$. Thus no open
orbit exists for $r\ge3$. At full length the deficit is
$(N-1)(N-4)/2$.

For $r=2$, view $E$ as matrices, with action $X\mapsto gXh^t$.
The invertible first lines form a dense open orbit in $\mathbf P(E)$:
over $\C$ a scalar can normalize any invertible matrix's determinant.
A generic flag can therefore be written
$\C I\subset\C I+\C B$, with $B$ regular semisimple.
Its stabilizer satisfies $gh^t=cI$, $c^n=1$, and
$gBg^{-1}=aI+bB$, $b\ne0$. The affine transformation
$x\mapsto a+bx$ must permute the distinct eigenvalues of $B$.
Only finitely many such affine transformations exist, because the
images of two distinct eigenvalues determine $a,b$. The identity
component therefore has $c=1,a=0,b=1$ and is the centralizer of
$B$ in $SL_n$, a torus of dimension $n-1$. The generic orbit has
dimension $2N-n-1$ in a flag variety of dimension $2N-3$,
giving codimension $n-2>0$. Any open orbit would meet this dense
regular-semisimple locus, which is impossible.
\end{proof}
For $r=1$ the invertible matrix lines do form an open orbit, so the
boundary is real. The precise source premise being tested is Manivel
\cite[Lemma~1 and Corollary~1, \S3, p.~3]{Manivel}. Failure
of this open-orbit premise excludes its multiplicity-free algebra
mechanism; it does not exclude positive bases, useful generators,
or counting rules for multiplicity-bearing invariant algebras.

\subsection{An actual primitive motive for two-row rectangular coefficients}
The coefficient identity used in this subsection is known:
if $N=nd$, $0\le r\le\lfloor N/2\rfloor$, and $p_r(n,d)$ counts
partitions of $r$ in an $n\times d$ rectangle, then
\begin{equation}\label{eq:rectprimitiveknown}
 g_r:=g((N-r,r),(d^n),(d^n))=p_r(n,d)-p_{r-1}(n,d),
 \qquad p_{-1}=0.
\end{equation}
For $r\ge1$ the inspected primary locator is Pak--Panova,
\emph{Bounds on certain classes of Kronecker and $q$-binomial
coefficients}, author manuscript dated 13 October 2016,
Lemma~1.3, p.~2 \cite{PPBounds}; the source preserves the earlier
attributions. For $r=0$ the identity is the multiplicity one of the
trivial representation in a Specht tensor square. We apply this
known identity and classical Lefschetz theory; no historical
novelty is claimed for the underlying geometric interpretation.

\begin{finding}[\fid{F31}: a primitive Grassmannian coefficient object]\label{f:grassprimitive}
Let $n,d\ge1$, $N=nd$, $0\le r\le\lfloor N/2\rfloor$, and
$X=\Gr(n,n+d)$ over $\F_q$. Write $L=c_1(\mathcal O_X(1))$
for the Pl\"ucker class in $H^2(X_{\bar{\F}_q},\Q_\ell(1))$,
where $\ell\nmid q$. The actual Weil-module cokernel
\begin{equation}\label{eq:rectprimitivecoker}
 C_r=\operatorname{coker}\bigl(
 L:H^{2r-2}(X_{\bar{\F}_q},\Q_\ell)(-1)
       \longrightarrow H^{2r}(X_{\bar{\F}_q},\Q_\ell)\bigr)
 \simeq\Q_\ell(-r)^{\oplus g_r}
\end{equation}
is pure of weight $2r$, with scalar geometric Frobenius $q^r$;
the source is zero for $r=0$.

For $V_j=H^{2j}(X_{\bar{\F}_q},\Q_\ell)(j)$ and $r\ge1$, put
\[
 A_r=L^{N-2r+1}:V_r\longrightarrow V_{N-r+1},\qquad
 B_r=L^{N-2r+2}:V_{r-1}\longrightarrow V_{N-r+1}.
\]
Then $B_r$ is invertible and the Frobenius-equivariant projector
\begin{equation}\label{eq:rectprimitiveprojector}
 E_r=\id_{V_r}-L B_r^{-1}A_r
\end{equation}
has image $\ker A_r$ and kernel $L V_{r-1}$. For $r=0$ set
$E_0=\id$. Untwisting identifies its image with $C_r$.
It is the realization of an actual rational Chow-motive summand
$\mathsf P_r\simeq(\mathbb L^r)^{\oplus g_r}$, where
$\mathbb L$ denotes the Lefschetz motive. Its matrix in the
Schubert basis need not be nonnegative or integral.
\end{finding}

\begin{proof}
We first specify the arithmetic Tate decomposition. Let $\mathbb G_m$
act on the ambient space with weights $1,2,\ldots,n+d$.
The fixed $n$-planes are the rational coordinate planes indexed by
$I=\{i_1<\cdots<i_n\}$. The tangent weights are $j-i$ for
$i\in I$, $j\notin I$, so the positive-weight dimension is
\[
 a_I=\sum_{s=1}^n(d+s-i_s).
\]
The sequence $\lambda_s=d+s-i_s$ is a partition in the
$n\times d$ rectangle, and the inverse assignment is
$i_s=d+s-\lambda_s$. Thus $a_I=j$ occurs exactly $p_j(n,d)$
times. The motivic Bia\l ynicki--Birula decomposition in
Brosnan \cite[Theorems~3.3 and~3.5, author PDF pp.~5--6]{Brosnan}
applies to this smooth projective variety and split action.
Its fixed pieces are rational points, hence it gives, over $\F_q$,
\[
 h(X)\simeq\bigoplus_{j=0}^{N}(\mathbb L^j)^{\oplus p_j(n,d)},
 \qquad
 H^{2j}(X_{\bar{\F}_q},\Q_\ell)
    \simeq\Q_\ell(-j)^{\oplus p_j(n,d)},\quad H^{2j+1}=0.
\]
In particular Frobenius on $V_j$ is the identity. This actual
Tate decomposition supplies arithmetic splitting; purity alone
would not exclude a unipotent arithmetic extension.

The Pl\"ucker bundle is ample and defined over $\F_q$, and
$X$ is smooth projective of dimension $N$. BBD Theorem~5.4.10,
p.~144, applies to $X\to\Spec\F_q$ and the pure perverse
sheaf $\Q_\ell[N]$. Hard Lefschetz gives the stated isomorphism
$B_r$. Since $A_rL=B_r$, the map $L:V_{r-1}\to V_r$ is injective.
Set $T_r=L B_r^{-1}A_r$. Then
\[
 T_r^2=L B_r^{-1}(A_rL)B_r^{-1}A_r=T_r,\qquad
 A_r E_r=0,\quad E_r L=0,\quad
 E_r|_{\ker A_r}=\id.
\]
Also $E_rv=0$ implies $v=T_rv\in L V_{r-1}$.
Consequently
\[
 V_r=L V_{r-1}\oplus\ker A_r.
\]
The quotient map restricts to an isomorphism from the primitive
kernel. Its dimension is $p_r-p_{r-1}=g_r$ by the cited identity.
All maps are Frobenius equivariant, including $B_r^{-1}$, so
untwisting gives \eqref{eq:rectprimitivecoker}. More explicitly
the primitive subspace in the original degree is
\[
 \ker\bigl(L^{N-2r+1}:H^{2r}(X_{\bar{\F}_q},\Q_\ell)
 \longrightarrow
 H^{2N-2r+2}(X_{\bar{\F}_q},\Q_\ell)(N-2r+1)\bigr).
\]
The quotient-to-kernel splitting is canonical once $L$ is fixed.
The twist $(-1)$ in \eqref{eq:rectprimitivecoker} is necessary:
the rational divisor class is invariant in $H^2(1)$.

For the motivic assertion use rational coefficients in the displayed
Tate decomposition. Morphisms between sums of a fixed Tate motive
are rational matrices, and degree-zero morphisms between distinct
Tate powers vanish. After the displayed degree pieces are normalized
by their Tate powers, multiplication by $L$ and its powers are such
matrices, with rational entries given by algebraic correspondences.
Hard Lefschetz makes $B_r$'s determinant nonzero, so its inverse
is a rational matrix and an actual morphism in this Tate category.
Thus $E_r$ is an algebraic idempotent; the idempotent completion
gives its motive image with the stated rank. Brosnan's \S2,
author PDF pp.~3--4, specifies the correspondence and projector
category used here. This proves algebraicity in the cellular Tate
case; it does not assume the general standard conjecture on
Lefschetz inverses.

Finally take $X=\Gr(2,4)$ and $r=2$. In its integral Schubert
cycle lattice the add-one-box Pieri rule gives
\[
 L\sigma_1=\sigma_2+\sigma_{11},\qquad
 L\sigma_2=L\sigma_{11}=\sigma_{21}.
\]
These are the ordinary Schubert specializations of the Pieri rule
\cite[Theorem~7.10, p.~16]{Pieri}. Here $A_2=(1,1)$, $B_2=(2)$,
and
\[
 E_2=\frac12
 \begin{pmatrix}1&-1\\-1&1\end{pmatrix}.
\]
Its image is the primitive Tate line. The lattices
$\Z(\sigma_2+\sigma_{11})$ and
$\Z(\sigma_2-\sigma_{11})$ have sum of index two in
$\Z\sigma_2\oplus\Z\sigma_{11}$. Thus the prescribed
primitive-versus-Lefschetz decomposition does not split this
integral lattice, and its rational projector has negative entries.
This does not prohibit every integral projector onto the primitive
line: $\bigl(\begin{smallmatrix}1&0\\-1&0\end{smallmatrix}\bigr)$
is an integral idempotent with the same image and a different kernel.
\end{proof}

The two rectangular factors in \fid{F31} have arbitrary height and
width. Its restriction is the two-row output $(nd-r,r)$: it covers
single-part padded tails, not IP's general multipart tails. Within
IP's padded-tail range $|\bar\lambda|\le md$, it applies when
$r\le\min(md,\lfloor nd/2\rfloor)$; its own construction requires
no IP size threshold. The Satake coefficient in \fid{F20} for the
same two rectangles, using $\GL_n\times\GL_n$, has weight
$2nd(n-1)$, whereas \fid{F31} has weight $2r$. A Tate twist can
align weights, but their equal ranks do not choose an isomorphism,
an operation-compatible basis, or the nonstandard path comparison.
This finding uses classical hard Lefschetz and a known coefficient
identity; it does not derive that identity or purity anew from the
sign-only path argument.


% Required bibliography keys: PPS (arXiv:2511.02312v3, Lemma 1.3, p.2)
% and KM (Kahle--Michalek, EJC 25(1) (2018), P1.41).
\subsection{A positive plethysm ray whose required dimension exceeds the compact bound}
The compact input convention in the supplied VI is explicit: zero parts
are suppressed, and height is the number of positive parts
\cite[\S1.1, p.~9; opening of \S3.4, p.~53]{VI}.
In particular, the source deliberately omits the trailing zeros of an
outer weight even when the ambient group has very large rank.
Its stretching operation fixes the inner weight
\cite[equation~(1.5), p.~20]{VI}.
These conventions matter for the following elementary deduction from a
classical two-row identity.

\begin{finding}[\fid{F34}: a necessary dimension and an encoding obstruction]\label{f:binarydimension}
Fix $k\ge3$, take $H=\GL_2$, inner weight $\mu=(k)$,
outer weight $\lambda=(1)$, and output weight $\pi=(k-1,1)$.
Their stretching multiplicity is
\[
 c_k(N)=a^{((k-1)N,N)}_{(N),(k)}
       =[\Sym^N(\Sym^k\C^2):S_{((k-1)N,N)}\C^2].
\]
For every $N\ge0$,
\[
 c_k(N)=
 \#\{(x_2,\ldots,x_k)\in\Z_{\ge0}^{k-1}:
                          2x_2+\cdots+kx_k=N\}.
\]
It is the Ehrhart function of the rational simplex
\[
 P_k=\{x\in\R_{\ge0}^{k-1}:2x_2+\cdots+kx_k=1\},
\]
whose affine dimension is $k-2$. Every fixed bounded polytope with this
Ehrhart function has dimension at least $k-2$.

Consequently the uniform dimension bound in the supplied
VI, Hypothesis~1.6.4(2), p.~23,
\[
 m=\operatorname{poly}
   (\langle\mu\rangle,\operatorname{ht}(\lambda),
                         \operatorname{ht}(\pi)),
\]
is false under its stated compact binary input convention.
Moreover, the plethysm counting function with all nonzero parts of
all three partitions encoded in binary is not a $\#P$ function:
already this family can require exponentially many output bits
relative to the input length.
These assertions do not apply to a unary weight convention or to
a bound polynomial in the numerical $k$.
\end{finding}
\begin{proof}
At $N=0$ both counts are one, so assume $N\ge1$.
The classical two-row specialization, restated in
\cite[Lemma~1.3, p.~2]{PPS}, gives
\[
 c_k(N)=p_N(N,k)-p_{N-1}(N,k),
\]
where $p_r(N,k)$ counts partitions of $r$ in the $k$ by $N$ rectangle.
For $r=N,N-1$ the width bound is automatic. After conjugation, the
remaining condition is that all parts have size at most $k$. Thus
\[
 c_k(N)
 =[z^N](1-z)\prod_{i=1}^k(1-z^i)^{-1}
 =[z^N]\prod_{i=2}^k(1-z^i)^{-1},
\]
which proves the positive count and the asserted simplex formula.

The weight calculation also checks which plethysm stretch is involved.
A monomial of $\Sym^N(\Sym^k\C^2)$ has nonnegative multiplicities
$a_0,\ldots,a_k$ satisfying $\sum a_i=N$ and second-coordinate weight
$\sum ia_i=r$. At $r=N$ or $N-1$, the determined value of $a_0$ is
automatically nonnegative. The GL$_2$ highest-weight multiplicity is
the difference of these adjacent weight dimensions, since
$N\le kN/2$. The inner $k$ is fixed, exactly as in VI's definition.

Put $L_k=\operatorname{lcm}(2,\ldots,k)$. For $N=L_ks$, the map
\[
 (y_2,\ldots,y_k)\longmapsto
       ((L_k/2)y_2,\ldots,(L_k/k)y_k)
\]
injects the nonnegative solutions of $\sum y_i=s$ into the displayed
count. Hence
\[
 c_k(L_ks)\ge\binom{s+k-2}{k-2}.
\]
Conversely, choosing $x_2,\ldots,x_{k-1}$ determines $x_k$ if it
exists, so $c_k(N)\le(N+1)^{k-2}$.

For any fixed bounded polytope $Q$ of affine dimension $d$,
$\#(NQ\cap\Z^m)=O_Q(N^d)$. To see this without a further asymptotic
theorem, choose $d$ coordinate functions injective on the direction
space of its affine hull. They inject the lattice points of $NQ$
into an integer box with side lengths $O_Q(N)$.
Comparison along $N=L_ks$ therefore forces $d\ge k-2$.
The displayed simplex has this dimension and attains the bound.

Here $\langle\mu\rangle=O(\log k)$,
$\operatorname{ht}(\lambda)=1$ and
$\operatorname{ht}(\pi)=2$, while the original group $H$ has fixed
rank two. Any ambient dimension for the claimed PH1 polytope is at
least its affine dimension $k-2$, which eventually exceeds every
fixed polynomial in these compact parameters. The omitted trailing
zeros in the $\GL_{k+1}$ outer weight cannot be charged to the input:
VI explicitly excludes them.

For the output-length assertion choose an integer $m\ge2$ and set
$k=2m$, $N=m(m+1)$. In the positive count choose independently
\[
 x_{2i}\in\{0,1\}\quad(2\le i\le m),\qquad
 x_j=0\quad(j\text{ odd}),
\]
and set
\[
 x_2=\frac{N-\sum_{i=2}^m2i\,x_{2i}}2.
\]
This is a nonnegative integer because the maximum subtracted sum
is $m(m+1)-2=N-2$. Distinct choices give distinct solutions, so
\[
 a^{((2m-1)N,N)}_{(N),(2m)}\ge2^{m-1}.
\]
The three partitions together have four nonzero parts, each with
$O(\log m)$ bits. A $\#P$ function on inputs of length $L$ has at
most $2^{p(L)}$ accepting witnesses for some fixed polynomial $p$.
The lower bound $2^{m-1}$ on these $O(\log m)$-bit inputs contradicts
that requirement for every fixed $p$.
\end{proof}

\paragraph{Known precedent and the scope of this deduction.}
The universal ordinary Ehrhart-existence claim in the supplied VI
was already refuted by Kahle--Micha{\l}ek
\cite[Theorem~5 and Remark~6, p.~4]{KM}.
Their fixed-small-parameter example has outer weight $(4)$, inner
weight $(3)$ and output weight $(7,5)$; its stretching function
violates Ehrhart reciprocity. Their p.~2 explicitly identifies the
question with VI, Hypothesis~1.6.4. Their pp.~3--4 distinguish the
later inhomogeneous formulation, which their result does not
dispose of in every dimension.
We cite that prior result for the historical failure of the old
existence assertion. Finding~\fid{F34} isolates the additional
compact-dimension issue on a family that \emph{does} have an exact
positive simplex model; no priority claim is made.

In the effective representation ring the numbers $c_k(N)$ are actual
Schur multiplicities, and their ordinary pure Satake realizations
are compatible with \fid{F20}/\fid{F31}. Purity does not change the
growth degree of the multiplicity. An encoding that includes the
numerical rank $\dim\Sym^k\C^2=k+1$ can accommodate this simplex,
but is different from VI's compact input. The associated rectangular
Kronecker input explicitly lists $k$ rows, so its input length is
already at least $k$; the binary plethysm output-length obstruction
does not apply to that encoding.


% Required bibliography keys: PPS (arXiv:2511.02312v3) and KM.
\subsubsection{The actual polynomial highest-weight algebra on the same ray}
The positive rule in \fid{F34} can be realized inside its ordinary
highest-weight algebra. The following direct slice calculation is a
constructive supplement to that finding, without a priority claim for
the classical method.

\paragraph{Claim.}
For $k\ge2$ put $V=\Sym^k\C^2$ and choose its divided basis
\[
 v_i=x^{k-i}y^i/i!,\qquad 0\le i\le k.
\]
In the free symmetric algebra $A=\Sym(V)=\C[v_0,\ldots,v_k]$,
let $D(v_0)=0$, $D(v_i)=v_{i-1}$, and assign
$\deg v_i=1$, $\operatorname{ind}v_i=i$.
The ray algebra
\[
 B=\bigoplus_{N\ge0}
       \bigl(A_{\deg=N,\,\operatorname{ind}=N}\cap\ker D\bigr)
\]
is a polynomial ring on $k-1$ explicit generators $f_2,\ldots,f_k$
of degrees $2,\ldots,k$. Its actual degree-$N$ monomial basis is
indexed by $2a_2+\cdots+ka_k=N$, and
\[
 \dim B_N=[\Sym^N(\Sym^k\C^2):S_{((k-1)N,N)}\C^2]=c_k(N).
\]
The symbols $v_i$ are free algebra generators here; their displayed
binary forms specify a basis and action, rather than an identification
of $A$ with a subalgebra of $\C[x,y]$.

\begin{proof}
The raising operator $x\partial_y$ acts by $Dv_i=v_{i-1}$ in
the divided basis. A degree-$N$, index-$r$ monomial has torus
weight $(kN-r,r)$. In characteristic zero, the vectors of this
weight killed by the raising operator are its highest-weight
vectors. Since $k\ge2$, the weight $((k-1)N,N)$ is dominant,
and ordinary highest-weight decomposition proves the multiplicity
identity. This also agrees with the classical two-row specialization
in \cite[Lemma~1.3, p.~2]{PPS} used in \fid{F34}.

In $A[v_0^{-1}]$ put $t=v_1/v_0$, so $Dt=1$, and define
\[
 f_i=v_0^{i-1}\sum_{j=0}^i\frac{(-t)^j}{j!}v_{i-j},
       \qquad 2\le i\le k.
\]
These elements belong to $A$. Indeed the terms with $j\le i-1$
are
\[
 \frac{(-1)^j}{j!}v_0^{i-1-j}v_1^jv_{i-j},
\]
and the last term is $(-1)^iv_1^i/i!$.
Each has degree and index $i$. Differentiating the sum, the terms
from $D((-t)^j/j!)$ cancel those from $D(v_{i-j})$ after shifting
the index; the remaining endpoint has $D(v_0)=0$.
Thus $Df_i=0$. For example,
\[
 f_2=v_0v_2-\tfrac12v_1^2,\qquad
 f_3=v_0^2v_3-v_0v_1v_2+\tfrac13v_1^3.
\]

The formula for $f_i$ is triangular, with leading term
$v_0^{i-1}v_i$ and remaining terms involving only earlier generators.
The leading coefficient is a unit in the localization. Solving
successively for $v_1=tv_0,v_2,\ldots,v_k$ therefore gives an
invertible change of free coordinates,
\[
 A[v_0^{-1}]=\C[v_0,v_0^{-1},t,f_2,\ldots,f_k].
\]
In these coordinates $D=\partial/\partial t$. Characteristic zero
then gives
\[
 A[v_0^{-1}]^D=\C[v_0,v_0^{-1},f_2,\ldots,f_k].
\]
This assertion concerns the explicitly computed localization; no
description of the entire unlocalized invariant ring is assumed.

Every monomial $v_0^a\prod_{i=2}^k f_i^{a_i}$ in this Laurent
polynomial ring has degree minus index equal to $a$.
Algebraic independence and homogeneity therefore imply that its
component with degree equal to index has $a=0$ and is exactly
$\C[f_2,\ldots,f_k]$. All the $f_i$ already belong to $A$, so
intersecting with $A$ loses nothing. Hence
\begin{equation}\label{eq:actualrayalgebra}
 B=\C[f_2,\ldots,f_k],\qquad \deg f_i=i.
\end{equation}
Its displayed monomials form the asserted actual highest-weight
basis. The construction is defined over $\Q$ and then base changed
to $\C$; all divided bases, factorials and the derivative-kernel
argument use characteristic zero.
\end{proof}

This gives multiplication constants equal to one in this particular
ordinary algebra. Its $k-1$ generators and their formulas have size
polynomial in the numerical $k$, not in $\log k$. The construction
therefore retains both the exact positive simplex and the compact
dimension obstruction in \fid{F34}. It does not identify this
monomial basis with VIII's nonstandard multiplication or
comultiplication basis, and asserts no theta or Weil realization.


\subsection{Short witnesses for the actual marked rectangular tree rule}
Pak--Panova--Swanson \cite[Theorem~4.5, p.~10]{PPS} give a positive
marked-KOH-tree rule for
\[
 g((nk-r,r),(n^k),(n^k)),\qquad
 n,k\ge1,\quad 0\le r\le nk/2.
\]
Their rule is a known result. The following encoding deduction makes
its witness size explicit when the rectangular partitions are supplied
as lists of $k$ binary-encoded parts.

\begin{finding}[\fid{F37}: a bijective compression of marked KOH witnesses]\label{f:kohcompression}
Every source marked KOH tree has a unique compressed encoding of length
\[
 O\!\left(k^2\log(k+1)+k\log(nk+1)\right)
\]
with a polynomial-time verifier that does not expand long unary chains.
The compression preserves the source objects bijectively.
Consequently the equal-rectangle, two-row-output Kronecker counting
problem is in $\#P$ under the explicit-row, binary-part input convention.
No bound polynomial in $\log k+\log n+\log(r+1)$ is asserted for a
rectangle specified only by two binary integers, and no Ehrhart PH1
conclusion follows.
\end{finding}
\begin{proof}
We first state the source rules needed for the encoding
\cite[Definition~3.2, p.~6; Definition~4.4, p.~10]{PPS}.
A node has a label $(\mu,a,b)$ with $a\ge0$, $b\ge1$, $\mu\vdash b$.
Leaves have $b=1$, $\mu=(1)$. A nonleaf has one child for every
distinct part $j$ of $\mu$, in increasing $j$ order, with type
\[
 a_j=(a+2)j-2Q_j(\mu),\qquad b_j=m_j(\mu),\qquad
 Q_j(\mu)=\sum_s\min(j,\mu_s).
\]
All child $a_j$ must be nonnegative. The root type is $(n,k)$.
For depth-first leaf labels $a_1,\ldots,a_t$, put
$\sigma=nk-\sum_i a_i$. The leaf marks satisfy
\[
 0=u_1\le\cdots\le u_t=r-\sigma/2,\qquad
 u_{i+1}-u_i\le
 \min\left(\sum_{j\le i}a_j-2u_i,\ a_{i+1}\right).
\]

Call a nonleaf a repeat node when $\mu=(1^b)$, $b\ge2$.
It has exactly one child, with type
\[
 (a,b)\longmapsto(a-2(b-1),b).
\]
Every other nonleaf obeys
\[
 \sum_j b_j=\ell(\mu)\le b-1.
\]
Contract each maximal consecutive run of repeat nodes, storing its
length $h$ in binary before the next surviving node. The entire run
is the update
\[
 (a,b)\longmapsto(a-2h(b-1),b).
\]
Its endpoint being nonnegative guarantees every intermediate
$a$-value is nonnegative. Retain the surviving partition labels,
ordered children and all leaf marks. A leaf has incoming repeat
count zero, because repeats preserve $b\ge2$.

This compression is bijective. Starting from the root type, a
compressed record first performs the displayed run update, then
requires its surviving partition to be different from $(1^b)$
unless it is a leaf. Its ordered child types are determined by the
source formula. Expanding a count $h$ inserts exactly $h$ nodes
labeled $(1^b)$ with the uniquely determined successive $a$-values.
The surviving node is nonrepeat, so the inserted run is maximal.
Compression and expansion are inverse operations. Unary contraction
does not change depth-first leaf order, any leaf label, or any mark;
all marking inequalities are therefore preserved.

Let the compressed tree have $H$ internal vertices and $t$ leaves.
The $b$-labels on each contracted edge are unchanged. Summing the
positive losses at internal vertices telescopes:
\[
 H\le
 \sum_{v\ {\rm internal}}
       \left(b_v-\sum_{w\ {\rm child\ of}\ v}b_w\right)
 =k-t.
\]
Thus there are at most $k$ surviving vertices. Every surviving
partition has at most $k$ parts, each at most $k$, accounting for
$O(k^2\log(k+1))$ bits.

To bound all numerical labels, set
$\beta(\mu)=\sum_s(s-1)\mu_s$. Decreasing row lengths give
\[
 \sum_jm_j(\mu)Q_j(\mu)
 =\sum_{r,s}\min(\mu_r,\mu_s)
 =b+2\beta(\mu).
\]
Hence
\[
 \sum_j a_jb_j=ab-4\beta(\mu)\le ab.
\]
In a valid tree all these terms are nonnegative, so every vertex
satisfies $0\le ab\le nk$. Each $a$-label is at most $nk$, and
every contracted repeat count is at most $a/(2(b-1))\le nk$.
The same identity gives $\sigma\ge0$. Every mark lies between
zero and $r\le nk/2$. At most $k$ repeat counts and $k$ marks
therefore take $O(k\log(nk+1))$ bits. Tree punctuation and
partition-length fields fit in the stated bound as well.

Use one canonical preorder serialization, with canonical integer
fields, decreasing partition parts, and the fixed child and leaf
orders. The verifier checks the number of surviving nodes, the
partition conditions, the root and recursive child types, the
binary run update, zero incoming repeat counts at leaves, and
then all leaf marking inequalities.
It can reject any derived type violating $0\le ab\le nk$;
the preceding identity proves that no valid source object is lost.
All work uses polynomially many operations on integers of the
stated bitlengths, and no run is expanded.

If the explicit-row input has length $L$, then $L\ge k$ and it
includes the binary lengths of $n$ and $r$. The witness bound and
the verification time are polynomial in $L$. Canonical serialization,
with required zero padding to a fixed polynomial length if desired,
gives exactly one accepting string for every compressed object.
The bijection and the cited PPS counting theorem identify the
number of accepting strings with the Kronecker coefficient.
\end{proof}

\paragraph{A literal tree can be exponentially long.}
Take $n=2s$, $k=2$, $r=2s$. At $b=2$ the only source partitions
are $(1,1)$ and $(2)$. After $j$ repeats of $(1,1)$, choosing $(2)$
ends in one leaf labeled $4s-4j$, so $\sigma=4j$.
Its sole mark is zero, forcing $0=r-\sigma/2=2s-2j$.
Thus the unique marked tree has $s$ repeat vertices, one final
$(2)$-vertex and one leaf. It has $s+2$ vertices on an input of
$O(\log s)$ bits. Its compressed encoding has two vertices and
one binary repeat count $s$. This is why local checks alone would
not have justified the binary witness bound.

\paragraph{Encoding and polyhedral limits.}
The related plethysm inputs $(n)$ and $(k)$ can both have one part,
so a bound polynomial in the numerical $k$ is not automatically
polynomial in their compact input length. Finding~\fid{F34} gives
an explicit output-length obstruction to that stronger assertion.
The present rectangular input instead lists $k$ rows.

Nor does a short witness rule imply one ordinary dilation polytope.
Here the tree shape and its partition labels are discrete choices,
and the local relations include fixed offsets. Indeed the known
Kahle--Micha{\l}ek counterexample
\cite[Theorem~5 and Remark~6, p.~4; rectangular identification, p.~6]{KM}
already supplies an equal-rectangle, two-row-output ray that is
not an ordinary Ehrhart function. The positive marked-tree rule
and the verified compression still apply to it.
This finding settles a restricted effective-label gate; it supplies
no designated nonstandard canonical basis, theta functor,
structure-constant positivity, or orbit-closure obstruction.


\subsection{Occurrence versus multiplicity}
Let $\Omega_n=\overline{\GL_{n^2}\cdot\det_n}$, $Z_{n,m}$ the padded-permanent orbit closure, and $K_n(\lambda)$ the degree-$d$ multiplicity in $\C[\Omega_n]$. B\"urgisser--Ikenmeyer--Panova \cite[Theorem 1.4]{BIP}, in the attached arXiv 1604.06431v3, proves for $n\ge m^{25}$:
\[
 q^m_\lambda(d[n])>0\quad\Longrightarrow\quad K_n(\lambda)>0.
\]
Theorem~1.5 is the numbering in an older version, not this attachment. The paper's \S1(b) distinguishes
\[
 K_n(\lambda)\le sk_n(\lambda)\le g_n(\lambda),
\]
where $sk_n$ is the symmetric rectangular coefficient controlling the determinant \emph{orbit}, not automatically its closure.

If $Z_{n,m}\subseteq\Omega_n$, restriction is a graded equivariant surjection
$\C[\Omega_n]\twoheadrightarrow\C[Z_{n,m}]$. Semisimplicity of the characteristic-zero representation in each finite-dimensional degree implies
\[
 q^m_\lambda(d[n])\le K_n(\lambda).
\]
Thus a strict reverse multiplicity inequality remains a valid obstruction, as BIP \S1(c) explains. A theorem about nonzero support does not settle it. The ordinary Kronecker upper bound supplies a sufficient stronger inequality, not an equality with closure multiplicities.

\paragraph{Consequence of the present work.}
\fid{F03}--\fid{F04} establish purity only for data satisfying their trace and sign hypotheses. No such data have been identified here for the normalized coefficients in VIII, the full rectangular region, or the closure multiplicity comparison. Therefore no new general positivity or multiplicity obstruction is derived from them. This is the exact limit of the proved transfer, not a restatement of numerical nonnegativity as a new theorem. The permanent--determinant problem concerns algebraic complexity classes; an additional argument would be required before making a Boolean $P\ne NP$ claim.

\part{The findings in Grothendieck's language}
\section{An exact dictionary, not an attribution of new theorems}\label{sec:groth}
The new deductions above are restated here in the requested frameworks. A framework locator says where the language and operations live; it does not claim that a theorem about GCT occurs in Grothendieck's manuscripts. Where no archival proof or identification was verified, we say so explicitly.

\subsection{Schemes and group schemes over the integers}
The classical objects are group schemes $\GL_{r,\Z}$ and their polynomial representations on finite locally free modules. Tensor products and integral Schur constructions give maps such as
\[
 \GL(V)\times\GL(W)\longrightarrow\GL(V\otimes W),\qquad
 \GL(V)\longrightarrow\GL(S_\mu V).
\]
Representations of an affine group scheme are comodules of its \emph{commutative} coordinate Hopf algebra (SGA~3, I, \S4.2 and 4.7.2). Representations of diagonalizable groups correspond to graded modules (I, 4.7.3). These descriptions carry \eqref{eq:plethysm}, \fid{F11} and \fid{F13} into group-scheme language. Cote~90, scan~17, gives an authentic manuscript anchor for a flat affine group over a base, a group action, and equivariant quasi-coherent modules.

Over characteristic zero the representation categories of the reductive groups here are semisimple, so their actual classes form the positive monoid in the irreducible basis of $K_0$. This is the precise content of scalar nonnegativity. Over $\Z$ or in characteristic $p$, do not replace a filtration by a direct sum of simple modules. A divided Young symmetrizer is only an idempotent when its denominators are invertible; integral Schur functors must be constructed integrally. For the integral Specht-model representations used in the Kronecker and plethysm instances of \fid{F12}, tensor the free integral lattices with $\Z[1/|H|]$. The averaging idempotent then has a finite locally free direct-summand image, an effective $K_0$ class whose complex rank is the stated multiplicity. More generally this construction works over any coefficient ring carrying the representation and in which $|H|$ is invertible; an arbitrary complex representation need not descend to $\Z[1/|H|]$.

The nonstandard bialgebra preserves some comodule language but is not automatically a commutative affine group scheme. \fid{F11}'s matrix calculation proves why a reduction modulo $p$ supplies no automatic absolute Frobenius. This is a limitation of that construction, not a statement against separately defined quantum Frobenius maps.

\subsection{Representation rings, $\lambda$-rings and the $\gamma$-filtration}
For locally free sheaves, $\lambda^i[E]=[\wedge^iE]$ and Schur operations are the associated symmetric-function operations. On virtual classes the identity $\lambda_t(x-y)=\lambda_t(x)/\lambda_t(y)$ belongs to the group-completed ring. SGA~6, Expos\'e~V, \S\S2--3, contains the $\lambda$ and $\gamma$ formalism; cote~133, scan~18, explicitly displays symmetric/exterior generating functions. Cote~97, scan~4, concerns Chern classes of representations. These are the appropriate anchors for the formal calculation in \fid{F11}.

The representation ring example is an exact dictionary, not an analogy: $z$ is the standard one-dimensional comodule of $\mathbb G_m$; $z-1$ is a rank-zero virtual class; its exterior-power series gives $\lambda^2(z-1)=1-z$. The augmentation ideal and its $\gamma$-filtration retain characteristic-class information but introduce no order turning this virtual class into an actual representation. In particular the equality of two expressions in a $\lambda$-ring cannot replace the effective-simple-basis hypothesis of \fid{F07}.

For singular schemes, one must distinguish $K_0$ of vector bundles or perfect complexes from $G_0$ of coherent sheaves. The GCT graded pieces themselves are finite-dimensional representations, so their representation-$K_0$ classes are well-defined without identifying $K_0(X)$ and $G_0(X)$ on a singular orbit closure. Coherent Euler pushforward may involve alternating higher cohomology; an actual section module is more data than a virtual Euler class.

\subsection{Weil realizations and virtual versus effective traces}
In the trace formalism, a bounded complex of finite-dimensional Weil representations yields
\[
 Z(K,T)=\prod_i\det(1-TF\mid H^i(K))^{(-1)^{i+1}},\qquad
 \log Z(K,T)=\sum_{n\ge1}\Tr(F^n\mid K)\frac{T^n}{n}.
\]
The trace is an alternating trace. \fid{F02} bounds the surviving finite spectrum after equal roots have been grouped. \fid{F05}(b) says that $K_0$ and $Z(K,T)$ can erase identical spectra in adjacent degrees. An effective parity statement prevents this erasure; that is the exact role of (S2).

The geometric dictionary of \S\ref{sec:tracegeom} is the proof-carrying transfer of \fid{F01}--\fid{F04} and \fid{F17}. Character orthogonality is expressed by Artin--Schreier sheaves; the external tensor product realizes powers; compactly supported direct image realizes the sum; one-degree support fixes the alternating sign; Tate twist supplies the factor $q^n$; duality supplies the reciprocal eigenvalues. Framework locators are SGA~4, XVII, \S3 (compact supports), XVIII, 3.2.5 (smooth duality); SGA~5, III, 4.7 (Lefschetz--Verdier); and SGA~$4\frac12$, \emph{Rapport}, 3.2 (the trace formula). The exact ordinary weight estimates are Weil~II, 3.3.1 and 3.3.4--3.3.6. These general theorems are cited background; this run does not re-prove their machinery.

For a scheme over $\F_q$ and a closed point $x$ with residue field $\F_{q^{\deg x}}$, a Weil sheaf is pure of weight $w$ if every eigenvalue $\alpha$ at $x$ has
\[
 |\sigma(\alpha)|=q^{(\deg x)w/2}
\]
for every relevant complex embedding. For complexes the weight conditions are degree-shifted and involve duality; purity of a complex does not assert that every singular stalk has exact weight equal to its degree. The KL theorem uses the stronger stalkwise even/Tate statement \eqref{eq:KL}. This distinction translates \fid{F06} and prevents a false purity-to-coefficient inference.

\subsection{Orbit closures, Hilbert/Quot, and motives}
\fid{F15}'s kernel construction is the affine scheme-theoretic-image construction of EGA~I (1960), \S9.5. Its explicit $f=p$ example is already a group-scheme example over $\Spec\Z$. Projective closures with fixed Hilbert polynomial and coherent quotients belong to FGA, Bourbaki Expos\'e~221, \S3, Theorems~3.1--3.2, pp.~259--260. Fixing the Hilbert polynomial uses the chosen polarization, and a Quot functor specifies the coherent sheaf being quotiented. They do not turn coherent section multiplicities into intersection-cohomology multiplicities by representability alone.

For smooth projective schemes with algebraic correspondences, pure motives and their realizations provide another language for duality and projectors. The Tate realization in \fid{F16} is legitimate. By contrast a general singular orbit-closure intersection complex is not automatically a Chow motive furnished by SGA~4/5/6, and a quantum coordinate algebra acquires no such motive from its name. Cote~16, scans~35, 37 and~39, supplies a historical anchor for Lefschetz operators and primitive projectors; scans~47--49 discuss conditional correspondence/motive consequences, with conjectural letters interpreted in the letter's own notation. Their algebraic-projector and index assumptions cannot be replaced by the Euler sign $(-1)^i$.

\subsection{The constructive models and their exact comparison limits}
\paragraph{\fid{F18}: weight lattices and rational semigroups.}
Regard $g(\nu,(d,d),(d,d))$ as the coefficient of
$\det^d\boxtimes\det^d$ in the restriction of $S_\nu(V\otimes W)$,
where $V,W$ have rank two. This is an equality in the representation
$K_0$ of the group scheme $\GL(V)\times\GL(W)$ after passage to
characteristic zero. Characters of a diagonalizable group correspond
to gradings (SGA~3, I, 4.7.3). Integral difference-constraint arrays
would supply a particular saturated graded semigroup. The flooring
and dilation proof of \fid{F18} transfers unchanged to that weight
lattice. Its parity example requires rational-lattice or congruence
data instead. SGA~6's formal $\lambda$ identities do not supply those
data; cote~133, scans~18--21, is a vocabulary anchor rather than a
positive-rule theorem.

\paragraph{\fid{F19}, \fid{F21}, \fid{F25}: coherent algebra and its actual module.}
Here there is a specified scheme over a localization of
$\Spec\Z[v,v^{-1}]$: the finite flat center $S$ of \fid{F19}, with
the explicitly trivial bundle $E$ of ranks $1,1,2$.
The noncommutative algebra is $\Gamma(S,\End E)$, and the actual
right tensor module is $\Gamma(S,\mathcal Hom(E,\mathcal F))$ for
the displayed coherent branch sheaf $\mathcal F$. Thus the module
dictionary is an explicit action-preserving isomorphism, not just
an equality of characters. The matrices in \fid{F21} carry the proof
over to this sheaf language. Equivariant quasi-coherent modules are
the framework of cote~90, scan~17, and SGA~3, I, 4.7.2.

The branch modules $L_i$ are coherent but are not asserted to be
locally free or perfect over the singular center. Their natural
Grothendieck group is $G_0(C)$; one must not silently apply vector-bundle
$\lambda$ operations from $K_0(C)$ to them. The periodic resolution
in \fid{F25} computes their actual coherent derived Hom. Its
$R/(\Delta)$ records an order-two deformation invisible in the
constant coefficient realization. The regular algebra basis in that
finding is unrelated to any assertion that this singular coherent
category is semisimple. Its scalar $\lambda$ has compact-support motive
\[
 \mathbb L^{-2}(\mathbb L-1)^2(1+\mathbb L+\mathbb L^2)
\]
in the virtual Tate ring. The trace equality follows from
$\mathbb G_m^2\times\PP^2$ and the trace formula; its mixed signs are
cohomological Euler signs, not dimensions of one pure degree.

\paragraph{\fid{F22}: the \'etale topos and the structure sheaf retain different data.}
The finite maps $\Spec C_n\to\Spec C_m$ are universal
homeomorphisms. EGA~IV, 18.1.2, in the precise modern formulation
Stacks 04DZ/03SI, transports the entire ordinary coefficient-sheaf
topos over the trait. The proof in \fid{F22} explicitly identifies
the constant nearby-cycle triangle with
$\Q_\ell\to\Q_\ell^2\to\Q_\ell$; all three objects have scalar
weight-zero Frobenius. In contrast the conductor and crossed coherent
$\operatorname{Ext}^1$ retain $m$. This is an exact instance of the
difference between a coefficient topos and a ringed coherent space.
It does not forbid adding a separately specified nilpotent action to
a Weil sheaf. No manuscript statement is claimed to prove this
nonstandard deformation comparison.

\paragraph{\fid{F20}: effective Schur operations and Tate coefficient motives.}
On a split affine Grassmannian, the primitive Young projector is an
actual idempotent in a rational symmetric tensor category. Its image
is an effective object, and its class is the specified Schur operation
on the locally defined representation. The modern motivic Satake
theorems cited in \fid{F20} identify its coefficient motive as
$\mathbf1(-h)^{\oplus a}$, not merely an alternating virtual class.
The ordinary coefficient realization has scalar Frobenius $q^h$.
This supplies an actual motive in these classical cases; it does not
assert the existence of a Chow motive for every singular orbit closure.
The source of the motivic theorem is Richarz--Scholbach, not SGA~6
or cote~16. The latter's scans~35,37,39 supply the historical language
of algebraic projectors. Canonical bases and efficient rank labels
remain separate information even when the motive is Tate.

\paragraph{\fid{F23}: induction in the effective representation monoid.}
A signed orbit is the representation of a constant finite group
scheme induced from a one-dimensional stabilizer representation.
Over $\Z[1/d!]$ with the relevant integral lattices, finite-group
averaging splits invariant projectors; over characteristic zero
the representation category is semisimple. The direct-sum equality
of \fid{F23} is effective in representation $K_0$, and its Young
stabilizer specialization is the ordinary symmetric/exterior Pieri
operation. The proof by orbit modules and restriction of Schur
functors is unchanged in this language. The pure Tate $\GL_2$
example shows why the $S_d$ correspondence action is additional to
the weight and motive. Cotes~90/17 and~133/19--21 give the respective
module and finite-set-operation anchors; neither provides the missing
universal signed basis.

\paragraph{\fid{F24}: a signed class can still retain a largest eigenvalue.}
For a bounded Weil complex, group equal roots in its alternating trace.
Its visible multiplicities are integers, possibly of both signs.
The external tensor product maps this class to its power in the
group ring of eigenvalues. The augmentation-and-Vandermonde argument
in \fid{F24} proves that an even power retains a positive top-radius
coefficient. This transfers the proof exactly to the triangulated
$K_0$ trace formalism of SGA~4/5, without requiring one input degree.
The negative visible path-output spectrum and any reciprocal visible
duality are still hypotheses. A zero virtual class may erase actual
cohomology, and finite extensions can erase a signed top spectrum;
the two countermodels specify those limits.

\paragraph{\fid{F26}--\fid{F27}: localization, bounded complexes, and finite traces.}
Finite Weil cohomology gives rational resolvents and zeta factors
$\det(1-TF)$, including their alternating products. This is the
finite determinant mechanism underlying the trace formalism; the
non-Laurent rational grading function in \fid{F26} cannot be realized
by such a fixed finite trace in all extensions. Its recurrence proof
is independent of any particular choice of category. The generic
Fredholm determinant formulas visible in cote~1, scans~77--78,
do not identify it with a finite Weil object or a $\Xi$-function.

Bounded complexes in the Hecke category do realize the signed Euler
class $P$. Their four basis coordinates prove the one-parity
obstruction in \fid{F27}. The rank-one bimodule comes from the
finite flat correspondence
$\Spec A\times_{\Spec A^s}\Spec A$. Separately, the coefficient
Hecke algebra $H=R[c]/(c^2-\delta c)$ normalizes to $R^2$
(and its tensor square to $R^4$); this separates its character
points and yields the other positive cone.
The actual four-point Tate construction retains the rank-one action
after the stated localization. Algebraic localization, coherent
normalization, and a trace-compatible tensor realization are three
specified operations here; none is inferred from the others.


\subsection{Finding-by-finding transfer table}
{\small\begin{longtable}{p{.12\linewidth}p{.48\linewidth}p{.32\linewidth}}
\toprule IDs&Exact Grothendieck-framework form&Locator/anchor and limit\\\midrule\endhead
F01--F04, F17&Trace of external tensor powers on the path hypersurface; parity of $R\Gamma_c$; Tate duality; optional finite base extension for each $k$. The calculations of \S\ref{sec:main} apply after the trace-formula identification.&SGA~4 XVII \S3, XVIII 3.2.5; SGA~5 III 4.7; Weil~II 3.3.1. Cote~16 provides only the distinct projector/duality framework, not this theorem.\\
F05&The zero virtual class $[\Q_\ell(-1)]-[\Q_\ell(-1)]$ has invisible large eigenvalues. A pure arithmetic Weil module can be a nontrivial extension.&Triangulated $K_0$; BBD 5.3.8 and its following example, pp.~138--139. No archival proof of the new examples is asserted.\\
F06--F07&Effective graded simple classes yield nonnegative multiplicities; a positive phase and Lefschetz maps are additional. Euler classes alone do not suffice.&BBD 5.3.2, 5.4.5, 5.4.10; cote~16, scans~35, 37, 39 for operator language.\\
F08&The rank vector bundles on two points form a special $\lambda$-ring; an effective integral basis need not be its simple positive basis.&SGA~6 V \S\S2--3; cote~133, scan~18. The entire counterexample is proved here.\\
F09--F10&An algebra inclusion need not preserve the effective monoid of a split $K_0$ categorification; exclusive monomial pivots give a sufficient alternative.&Split Grothendieck groups, effective objects, grading shifts; cote~90, scan~17 for equivariant-module context. We do not attribute the nonstandard inclusion claim to the manuscripts.\\
F11&Virtual exterior powers fail effectivity; noncommutative $p$th powers fail additivity, so do not define a scheme Frobenius.&SGA~3 I 4.2, 4.7.2--3; SGA~6 V \S\S2--3; cotes~97/4 and~133/18.\\
F12--F13&For the integral Specht models, invariant projectors are direct summands over $\Z[1/|H|]$; positive restriction matrices compose in representation $K_0$.&SGA~3 I 4.7.2--3; cote~90/17. Actual theta/GCT functor is still absent.\\
F14&Frobenius-stable subquotients inherit a weight; endomorphism spectra are ratios.&Weil representations; BBD 5.3.1 for perverse subquotient weight bounds. Realization/stability is still required.\\
F15&Kernel-defined orbit schemes over $\Z$; coherent polarization and linearization retained in $K_0/G_0$; finite-field base change is separate.&EGA~I \S9.5; FGA~221, 3.1--3.2; cote~90/17 for group actions.\\
F16&The centered polynomial of $\mathbf1\oplus\mathbb L^{\oplus2}\oplus\mathbb L^2$ realizes the rank-one product coefficient; its action comparison fails.&Projective bundle/K\"unneth calculations; cote~16 correspondence language. Pure Tate pieces, not an ungraded pure weight-zero module.\\
F18&Homogeneous weight-lattice difference arrays force ordinary support saturation; a rational half-point retains parity.&SGA~3 I4.7.3 for gradings; exact proof F18; VI Definition1.2.4 is a distinct residuewise notion.\\
F19, F21, F25&A finite flat center, a vector bundle endomorphism algebra and the actual coherent right tensor module; coherent Ext measures the collision.&Cote90/17; affine scheme/module dictionary; explicit complete proofs in these findings. Branch modules belong to coherent $G_0$.\\
F20&Primitive Schur projectors on convolution have effective pure Tate coefficient motives with fixed twists.&Richarz--Scholbach 4.1,4.8,5.9,6.2--6.5; BBD weight formalism; cote16 projector vocabulary only.\\
F22&Universal homeomorphisms preserve the ordinary coefficient topos while conductor and coherent Ext vary.&EGA~IV18.1.2; Stacks04DZ,03SI,0A4E.\\
F23, F28&Effective induced representation classes yield orbit/Pieri formulas when their stabilizer hypotheses hold; an isolated rectangular hook excludes that Young-orbit model.&SGA~3 representations; cote133/19--21 only a finite-set-operation anchor. IP Proposition6.4 supplies the hook input.\\
F24&Even external tensor powers of visible signed Euler spectra retain a positive largest-radius coefficient.&SGA~4/5 trace formalism; full proof F24, with powered-spectrum visibility and visible duality.\\
F26--F27&Finite Weil resolvents are rational; non-Laurent extension evaluations are excluded. Bounded Hecke complexes and normalized character schemes have different effective cones.&Split/triangulated $K_0$; finite normalization; cote1/77--78 gives generic Fredholm vocabulary, not this finite trace identity.\\\bottomrule
\end{longtable}}

\paragraph{\fid{F28}: rectangular projectors in representation $K_0$.}
The unprojected multiplicity module is the actual constant-group
representation $[(a^b)]\otimes[(a^b)]$. A proposed monomial orbit
basis would be an effective sum of induced one-dimensional Young
subgroup representations. Exterior powers of the standard
representation compute their consecutive hook supports. The
isolated hook supplied by IP Proposition~6.4 contradicts such an
effective sum; this proof carries verbatim to representation $K_0$.
A virtual sum may cancel those hooks and is a different assertion.
The known Manivel algebra instead has a torus-graded monomial basis,
with the displayed determinant-character comparison. This belongs
to the affine-group and weight-grading framework of SGA~3, I,
4.7.2--3 and cote~90, scan~17. Cote~133's operations on finite sets
do not identify these two different positive structures.


\paragraph{\fid{F29}: flag schemes and orbit dimensions.}
The action is that of the split group scheme $SL_n\times SL_n$ on
the partial flag scheme of its tensor product representation,
base-changed to $\C$. These flag schemes are iterated relative
Grassmannians, hence instances of the locally free Quot framework
in FGA, Expos\'e~221, \S3, Theorems~3.1--3.2. Their filtered
frame bundles give the Borel torsors used in the proof; the group
actions belong to SGA~3, I, \S4. The complete dimension and
stabilizer calculation proves only the absence of an open orbit for
this action. Cote~90, scan~17, anchors equivariant affine-module
language, not an attribution of this orbit theorem to the archive.

\paragraph{\fid{F30}: vector bundles and coherent commutants under specialization.}
In the language of schemes, $M_R$ is a vector bundle on $\Spec R$
with an action of a finite coherent algebra. Images of the regular
idempotents above are coherent direct summands and hence vector
subbundles. The regular mixing is an automorphism of the underlying
bundle; it does not commute with the whole algebra action. The coherent
commutant is the kernel of the finitely many linear commutation
equations for the generators and is the scalar sheaf in this case.
Its formation does not commute with the special fiber: that fiber
has a second classical projector. This is an explicit base-change
phenomenon for coherent endomorphisms, compatible with Grothendieck's
scheme and module language. It supplies neither a Frobenius action nor
a weight statement, and purity alone does not decide which algebra
equivariance a proposed comparison should retain.

The affine module/sheaf equivalence is
\cite[Lemma~26.7.5, tag~01IB]{StacksAffine}; its historical
EGA~I framework is Theorem~1.3.7, p.~86, Corollary~1.3.8,
p.~87, and Theorem~1.4.1, pp.~90--92. Equivariant modules use
SGA~3, I, 4.7.2--3; cote~90,
scan~17, is the audited archive anchor. Every needed idempotent
image, scalar commutant and special-fiber calculation is proved
in \fid{F30}; no general base-change theorem is asserted without
its hypotheses.

\paragraph{\fid{F31} in Grothendieck's language.}
The integral Grassmannian is the projective Quot scheme of
rank-$d$ locally free quotients of $\mathcal O_{\Z}^{n+d}$,
with its universal sequence and Pl\"ucker line. FGA, Expos\'e~221,
\S3, Theorems~3.1--3.2, pp.~259--260, supplies the representability
framework. Its split $\GL_{n+d}$ action and the cocharacter used
above belong to SGA~3, I, 4.7.2--3. The cokernel is an actual
Weil object and the projector an actual rational correspondence,
with effective class $g_r[\mathbb L^r]$. The inverse exists here
because of the full matrix Hom spaces of the Tate motive, while
the integral Schubert cone need not be preserved. Cote~16,
scans~35, 37 and~39, gives the audited Lefschetz/primitive-projector
vocabulary; these leaves are not asserted to contain this
Kronecker identification or its proof.


\paragraph{\fid{F32}: completed Tate classes and the limits of visible duality.}
For a countable direct sum of Tate pieces whose Euler trace has a
convergent nonnegative absolute moment, the theorem specifies exactly
the discrete completion being used. It is not an element of ordinary
finite $K_0$ merely by notation. Tensor multiplication is legitimate
because the absolute-moment estimate proves convergence and finite
grouping at each nonzero root. The formal binomial expansion of
$(1+\mathbb L^{-1})^{-4}$ in the actual trace has such a realization
at the level of convergent Euler spectra; its inverse Tate degrees
are unbounded. Dualization sends them to unbounded positive degrees,
outside this bounded-spectrum class. Thus an effective or signed
completed class still needs a specified duality and trace functor.
The finite trace framework is SGA~4\,$\frac12$, \emph{Rapport},
3.2, and SGA~5, III, 4.7. Cote~1, scans~77--78, records the
audited generic Fredholm trace/determinant vocabulary; neither it
nor finite $K_0$ supplies this countable extension. Its convergence
and reciprocal-finiteness proof is given in \fid{F32} itself.

\paragraph{\fid{F35} in Grothendieck's language.}
The ordinary split group schemes and quotient presentation are
explicit. Positive-dimensional stabilizers are retained, so this
construction uses the Artin-stack extension of the six operations,
not the bounded constructible pushforward theorem for schemes.
Its rational mass and infinite alternating Tate trace are connected
by the precisely cited Laszlo--Olsson/Sun formalism. The effective
cohomology groups in \eqref{eq:stackcompactgroups} have positive
dimensions, while their grouped Euler coefficients have both signs.
The duality exists but does not preserve the chosen convergent
geometric-Frobenius trace class. The audited manuscript determinant
and projector vocabulary is only a framework anchor; no specific
stack, nonstandard action, or proof here is attributed to those leaves.

\paragraph{\fid{F38} in Grothendieck's language.}
The projective-line decomposition and the stack's actual Tate
cohomology distinguish a positive object from its virtual Euler
class. A rational mass can invert an Euler trace while leaving
unbounded nonzero cohomology. Thus localization of a coefficient
ring or of a convergent trace algebra does not itself provide an
invertible geometric object. The split $\GL_2$ group scheme,
projective bundle formula, and modern stack six operations are
the precise frameworks used here. No archive leaf is claimed to
contain this construction or to identify it with a nonstandard
canonical basis. A different categorical localization requires a
specified new realization and its own comparison theorem.

\paragraph{\fid{F39} in Grothendieck's language.}
The singular piece is the actual affine scheme
$\Spec k[\epsilon]/\epsilon^2$, and matrix Morita equivalence
expresses its coherent category without changing the defect.
Smoothness here asks for a perfect diagonal bimodule, a perfect-complex
condition distinct from finite-dimensionality or ordinary $\ell$-adic
purity. The displayed resolution proves its failure directly. EGA~I's
affine scheme language and the perfect-complex framework of SGA~6
provide the classical dictionary; the smooth proper dg and
noncommutative-motive extension is explicitly the modern cited theory.
The audited cote~16, scans~48--49, concerns conditional semisimplicity
and Weil statements in its own motivic framework; it is not a source
for this dg construction. A new realization and basis comparison is
required before its Frobenius can address the GCT coefficients.


\paragraph{\fid{F34}: effective representation classes and growth of graded algebras.}
The actual highest-weight multiplicities are ranks of effective Schur
summands, in the characteristic-zero representation $K_0$ of $\GL_2$.
Their exact positive simplex has dimension $k-2$; purity cannot alter
that growth degree. The polynomial bound fails only for the specified
compact encoding. Locally free Schur operations and their virtual
extensions use the SGA~3 representation and SGA~6 $\lambda$-ring
frameworks. Cote~133, scans~18--21, anchors those operations, not the
computational bound. The dimension and output-length arguments are
proved explicitly in \fid{F34}; they are not archival theorems.

\paragraph{\fid{F37}: effective multiplicity objects with compressed labels.}
The known rectangular identity identifies the marked-tree count with
the rank of the primitive Tate summand constructed in \fid{F31}.
Compression is a bijection on the finite indexing sets for each input,
so it changes neither that effective representation class nor its rank.
It gives a polynomial witness encoding under the listed-row convention;
it does not identify the trees with a canonical cycle basis or make their
recursion into a theta correspondence. Cote~133, scans~19--21, is a
framework anchor for finite-set operations, and cote~16, scans~35,37,39,
for primitive projectors. The exact counting input is PPS, Theorem~4.5;
the new encoding argument is the proof of \fid{F37} itself.



\paragraph{\fid{F36}: coherent Frobenius, normalization and descent.}

This is a calculation of finite maps, coherent modules, relative
Frobenius and derived tensor products on an explicit singular scheme.
The relative map from collision order $mp$ to order $m$ is a finite
universal homeomorphism, like those in \fid{F22}; its coherent cokernel
is nevertheless nonzero. Ordinary constant \'etale traces lose this
length. The normalization has unit coherent Frobenius and preserves
the original algebra action after restriction, while pushforward
through the collision does not preserve the unit property.


The affine module/sheaf language is
\cite[Lemma~26.7.5, tag~01IB]{StacksAffine}, within the
historical EGA~I framework of Theorem~1.3.7, p.~86,
Corollary~1.3.8, p.~87, and Theorem~1.4.1, pp.~90--92; the action and
equivariant-module language is SGA~3, I, 4.7.2--3, with the audited
cote~90, scan~17, as its archival anchor. The derived calculations
are the explicit resolutions in the proof, not an invocation of
flat base change at the singular point. No archival identification
of coherent absolute Frobenius with a nonstandard Weil realization
is asserted.


\paragraph{\fid{F34}, constructive supplement: the actual graded ray scheme.}
The actual scheme $\Spec B$ is affine $(k-1)$-space with positive
grading $2,\ldots,k$; $\operatorname{Proj}B$ is the corresponding
weighted projective scheme. The degree-$N$ graded coordinate piece
is the highest-weight multiplicity space just proved, so its Hilbert
series is $\prod_{i=2}^k(1-z^i)^{-1}$. The exponent simplex has
dimension $k-2$. These are exact statements about the affine
coordinate ring and its grading, without an identification of
coherent Hilbert multiplicities with arbitrary constructible
cohomology. The ordinary affine and graded-scheme dictionary applies;
no arithmetic weight, designated nonstandard basis, or archival
claim is supplied by this coordinate computation alone.

The affine dictionary is \cite[Lemma~26.7.5, tag~01IB]{StacksAffine},
within EGA~I's Theorem~1.3.7, p.~86, Corollary~1.3.8,
p.~87, and Theorem~1.4.1, pp.~90--92; the highest-weight
representation dictionary uses SGA~3, I, \S4. The separate source anchors for Schur operations are recorded above; no complexity bound is ascribed to those texts.


\paragraph{\fid{F40} in Grothendieck's language.}
Affine endomorphism algebras, idempotent corners and perfect
diagonals give the precise coherent dictionary. The explicit
finite diagonal resolution proves smoothness of a larger category,
while its projective objects retain the singular corner and its
actual module. This uses the perfect-complex language associated
with SGA~6 and the explicitly cited modern smooth dg and
noncommutative-motive theory. The exceptional decomposition
explains the unit Frobenius spectrum without turning the corner
subcategory into an admissible motivic summand. No unpublished
archive leaf is claimed to contain this construction. The relative
order is supplied by \fid{F44}; realization of the designated
morphisms and comparison of the required bases remain separate tasks.

The affine module equivalence is
\cite[Lemma~26.7.5, tag~01IB]{StacksAffine}.
Perfect complexes use SGA~6, I, Definition~4.7, pp.~121--122,
and Example~4.8, p.~122, with the modern bounded-projective criterion
\cite[Definition~15.76.1, tag~0657]{StacksAffine}.
The modern perfect-diagonal smoothness definition is
\cite[\S5.1, pp.~13--14]{TabuadaWeil}. Cote~16, scans~48--49, is a vocabulary anchor for conditional motivic semisimplicity and Weil statements, not for this smooth dg construction.


\paragraph{\fid{F33} in Grothendieck's language.}
The geometric object is the ordinary split group scheme $SL_m$.
Its rational Tate-motive decomposition supplies actual Weil
summands and scheme Frobenius, so the equality of dimensions
has a precise realization. The relevant frameworks are SGA~3's
split group schemes and the explicitly cited modern motive
theorem. The generator degree $2i-1$ and weight $2i$ remain
different gradings. No manuscript leaf is claimed to contain
this hook formula or a nonstandard basis comparison.


The group-scheme language is SGA~3, I, \S4; cote~90, scan~17, is the audited group-action anchor. The motive decomposition itself is Merkurjev, Corollary~5.5, and not an assertion from the archive.


\paragraph{\fid{F41} in Grothendieck's language.}
The effective Tate summands and their weights are those of the
actual stack calculation in \fid{F35}. Eigenspace restriction
belongs to ordinary Weil linear algebra, and a finite weight
window restores the finiteness premise of \fid{F26}.
The chosen fixed-parameter action demonstrates that a geometric
coefficient object and an extension-wise parameter comparison
are distinct requirements. No action on the whole stack, or
archive attribution for such an action, is asserted.

The finite trace framework is SGA~4\,$\frac12$, \emph{Rapport}, 3.2, and SGA~5, III, 4.7. The infinite stack input uses the precisely cited modern extension; no bounded scheme theorem is applied to the unbounded object.


\paragraph{\fid{F42} in Grothendieck's language.}
The finite-resolvent obstruction concerns one finite realization
and its fixed Frobenius, whereas the Tate coefficient objects
use tensor powers with growing spectra. The fixed-parameter
construction is an exact action in the category of point Weil
objects, with named generators and tensor compatibility; it is
not an extension-wise quantum grading comparison. Scalar
division, an idempotent image, and inversion of a tensor object
are distinct operations here. The remaining path comparison
must preserve the actual operations and control its boundary
complex; no such comparison is attributed to an archive leaf.

The exact trace-framework locators are SGA~4\,$\frac12$, \emph{Rapport}, 3.2, and SGA~5, III, 4.7; the tensor/action language uses SGA~3, I, 4.7.2--3. Cote~90, scan~17, supplies the audited action vocabulary, not the comparison axiom.


\paragraph{\fid{F43} in Grothendieck's language.}
Over $\Spec S$ this is an explicit rank-64 vector bundle with
actual commuting algebra actions and a verified change of its
integral lattice. The full commutant is a product of matrix
algebras acting on its actual multiplicity bundles.
Separately, the finite graded kernels can be described by
$\mathbb G_m$-graded vector spaces on 64 points; their
Grothendieck group records the Laurent matrices and their
composition records paths. Closing that semisimple kernel
construction under every summand loses the coherent branch-gluing
condition, as the displayed specialization test proves.
The coordinate and $K_0$ descriptions make the positive action
concrete without identifying a formal quantum degree with a Weil
weight or specifying the dual FRT multiplication across degrees.
The actual bar and ordinary left action are now supplied in this
example; the designated nonstandard lattice, globalization and
geometric comparison remain separate requirements.

The equivariant-module and grading framework is SGA~3, I,
4.7.2--3, with cote~90, scan~17; the virtual $K_0$ and $\lambda$
framework is SGA~6, V, \S\S2--3 and cote~133, scan~18.
The displayed lattice, commutant and summand computations provide
the proof of this new application, rather than an archival attribution.


\paragraph{\fid{F44} in Grothendieck's language.}
The affine module/sheaf equivalence, precisely
Stacks Lemma~26.7.5, tag~01IB, translates the actual
corner lattice into coherent algebra on the affine
center. Perfect complexes mean bounded finite
projective complexes, as in Stacks Definition~15.76.1,
tag~0657; the modern dg smoothness criterion is
the perfect diagonal
\cite[\S5.1, pp.~13--14]{TabuadaWeil}.
The explicit relative diagonal and coherent derived
Hom prove the smooth repair and retain the order of
collision, while the nonadmissible corner prevents
an unwarranted motivic direct-summand assertion.
The Frobenius computation separates extension of
scalars over the endomorphism algebra from
linearization over the base. SGA~6, I, Definition~4.7, pp.~121--122, and
Example~4.8, p.~122, supply the historical perfect-complex framework, not a
source for this noncommutative construction.
The unpublished archive leaves are framework
anchors only; none is claimed to contain this
relative algebra, its action comparison, or its
Frobenius obstruction.

The representation language is SGA~3, I, 4.7.2--3 and
cote~90, scan~17. Cote~16, scans~48--49, provides the
audited conditional semisimplicity/Weil vocabulary, not
a theorem about this relative noncommutative algebra.



\paragraph{\fid{F45} in Grothendieck's language.}
The trace-formula framework is SGA~5, III, 4.7; the data
are smooth projective schemes, rational algebraic
correspondences, a projector, integer Tate twists, a centered
realization and the explicitly chosen ample line bundle
$\mathcal O(1,\ldots,1)$. Its invariant Chern class gives the
Lefschetz isomorphisms on the projector image. Milne's projective
bundle calculation and BBD~5.1.8/5.4.10 supply the exact classical
cohomological framework used above. Cote~16, scans~35, 37 and~39,
is the audited historical anchor for primitive-projector and
Lefschetz language \cite{Archive}; the present polarization and
construction are not attributed to those leaves.
The resulting centered complex is pure of weight zero, with
degree-$i$ cohomology of weight $i$, and its total ungraded
realization has several weights. The equality with the actual
word-trace polynomial is proved for all finite extensions, while
the summand test states exactly which extra compatibility fails
for the indicated positive $\mathrm{SL}_2$-kernel framework.


\paragraph{\fid{F46} in Grothendieck's language.}
A finite flat divisor on $\mathbf P^1_S$ and its section
sheaves supply actual coherent coefficient objects with
Cartier perfect resolutions. Their endomorphisms, module
actions and crossed extensions preserve the specified
family and its collision order. Multiplication by a
coordinate gives separately chosen linear Weil data after
coefficient specialization; it is not the arithmetic
Frobenius of that coefficient scheme. This distinction
allows an exact pure End/Hom realization while keeping the
required grading, globalization and path comparisons as
additional conditions. No unpublished archive leaf is
claimed to contain the Cayley construction or these
nonstandard comparisons.

The affine module/sheaf dictionary is EGA~I, Theorem~1.3.7,
p.~86, Corollary~1.3.8, p.~87, and Theorem~1.4.1,
pp.~90--92. Perfect complexes use SGA~6, I, Definition~4.7,
pp.~121--122, and Example~4.8, p.~122. The equivariant-module
framework is SGA~3, I, 4.7.2--3, with cote~90, scan~17,
as the audited archive anchor. The eigenvalue and complex-weight
conventions are BBD~5.1.2 and~5.1.8. These locators supply
the framework; the displayed resolutions and matrix-unit
arguments prove this particular comparison.

\paragraph{\fid{F47} in Grothendieck's language.}
The finite-flat scheme $\Spec C\to\Spec R$ has an algebra trace,
and $C\to C^\vee$, $a\mapsto(b\mapsto\Tr(ab))$, is a morphism
of finite locally free sheaves. Its determinant ideal is
$(\Delta^2)$ and its special-fiber kernel is the nilpotent
direction, by the displayed complete computation. Equivariant
modules use SGA~3, I, 4.7.2--3 and cote~90, scan~17.
The $K_0$ assertion applies the precisely cited modern theorem
on finite based rings; it is stronger than the virtual
$\lambda$ identities in SGA~6, V, \S\S2--3 and cote~133,
scan~18. For a proposed motivic polarization, the exact
additional assumptions are a positive-definite realization and
self-adjoint correspondences for the two named generators.
Under these assumptions the radical is necessarily killed.
Cote~16, scans~35,37,39 and~48--49, anchors the primitive-projector,
and conditional semisimplicity vocabulary. No
existence of such GCT correspondences or polarization is
attributed to those leaves. Neither a cohomological sign nor
Weil purity alone supplies this positive-adjoint structure.



\paragraph{\fid{F48} in Grothendieck's language.}
The tensor realization starts from the actual bialgebra
coaction and its forgetful fiber functor. The precise
coaction and equivariant-module framework is SGA~3,
I, 4.7.2--3, with cote~90, scan~17 as the audited archive
anchor. The representation-ring and $\lambda$-operation
vocabulary for ordinary coefficient objects is in
SGA~6, V, \S\S2--3, with cote~133, scan~18 as the audited anchor;
it does not identify the nonstandard exterior relations
with the symmetric constraint of the Weil target.
Purity here uses the eigenvalue and complex conventions
of BBD~5.1.2 and~5.1.8.
The displayed low-degree commutants, rational recurrence
and pole argument, and cyclic-module center calculation
prove the specific additional restrictions.
These sources supply their frameworks, not an unpublished
nonstandard tensor or arithmetic comparison.


\paragraph{\fid{F49} in Grothendieck's language.}
The original lattice is a vector bundle over $\Spec R$.
The rational idempotent $p$ extends and splits that bundle
over the open set $h\ne0$, but its explicit $1/h$ coefficient
prevents extension as an endomorphism of this bundle over
the whole parameter scheme. The precise affine dictionary is
EGA~I, Theorem~1.3.7, p.~86, Corollary~1.3.8, p.~87, and
Theorem~1.4.1, pp.~90--92.
The perfect direct-summand framework is SGA~6, I,
Definition~4.7, pp.~121--122, and Example~4.8, p.~122;
cote~90, scan~17 is an audited equivariant-module anchor.
None of those sources is claimed to contain this projector
calculation or the comparison with VIII.
The coherent model retains the actual framed bar and projector;
descent of its full specified localized projector system to
the designated global form fails under the actual tensor
comparison. Changing the integral form or restricting the
operator algebra remains possible. Purity of the separately
chosen Weil coefficients does not supply that comparison.


\paragraph{\fid{F50} in Grothendieck's language.}
The concrete schemes are an affine line, the path quadric,
its quadratic pencil, a finite etale degeneracy divisor,
and a smooth projective double cover. Artin--Schreier
and Kummer torsors supply the actual coefficient sheaves.
Compact-support base change, Kunneth, duality and Leray
prove the single-degree statement; the trace formula and
Tate twist give the opposite sign and factor $q^n$.
The exact primary inputs are SGA~$4\frac12$,
\emph{Sommes trig.}, 1.9.2, 1.18 and 4.2--4.4, and
the Milne locators displayed above; the trace framework
also agrees with SGA~5, III, 4.7 \cite{SGA45,Milne,SGA5}.
The fixed-parameter action language is that of equivariant
modules, SGA~3, I, 4.7.2--3, and cote~90, scan~17
\cite{SGA3,Archive}.

The normalized construction gives a fully faithful
degree-zero perverse realization of the actual
\fid{F46} End/Hom data, while remaining a module functor
for constant coefficients rather than a strong monoidal
functor for the ordinary tensor product. Finding~\fid{F53}
supplies a coherent twisted convolution for which it is strong monoidal. It retains arithmetic nilpotents and does not
create distinct geometric simple labels. Its graded
version retains the actual \fid{F43} coefficient kernels
with the stated $K_0$ and coherence qualifications.
The opposite-sign comparison is proved; the remaining
pointwise and dual requirements are exactly
\eqref{eq:gausspointwisegate}--\eqref{eq:gausscommonweight}.
A further advance would require replacing the constant
unknown coefficient factor by a varying local system
with an independently established pointwise bound and
the required actual operation comparisons.


\paragraph{\fid{F51} in Grothendieck's language.}
The coefficient extension is the algebraic field
extension explicitly allowed in VIII, \S2.2.
The module and action interpretation uses EGA~I,
Theorem~1.3.7, p.~86, Corollary~1.3.8, p.~87, and
Theorem~1.4.1, pp.~90--92; the coaction framework is
SGA~3, I, 4.7.2--3, with cote~90, scan~17 as the
audited anchor.
The pure-weight convention is BBD~5.1.2 and~5.1.8.
These frameworks do not assert the arithmetic
comparison: the displayed polynomial, radius and
Vandermonde argument proves its particular restriction.
No unpublished archive leaf is claimed to contain
this algebraic trace obstruction.

\paragraph{\fid{F52} in Grothendieck's language.}
The original order and its first-order overorder
define different coherent algebra sheaves on the
parameter trait. Their abstract special fibers agree,
while their embedded center lattices differ, with
the displayed nonzero cokernel lengths.
The affine module/sheaf framework is EGA~I,
Theorem~1.3.7, p.~86, Corollary~1.3.8, p.~87,
and Theorem~1.4.1, pp.~90--92.
The action comparison uses SGA~3, I, 4.7.2--3,
with cote~90, scan~17 as the audited framework anchor.
The finite coefficient extension is the type
allowed in VIII, \S2.2, p.~7.
The finite-cover, flat-kernel, valuation and
saturation arguments have been supplied explicitly
above; no unverified archive statement is substituted
for them.
The precise additional gate is preservation of the
whole relative order under the arithmetic comparison,
not merely its generic or special isomorphism type.

\paragraph{\fid{F53} in Grothendieck's language.}
The scheme is the additive line with its group law.
Convolution is compact pushforward along addition;
the normalized constant object is its dualizing
complex $E(1)[2]$ with compact trace $E$.
The origin Gysin map compares the ambient delta unit
with its idempotent image unit.
These are the actual Gysin, compact-trace, Kunneth
and duality constructions at the exact Milne locators
given in the proof, within the SGA compact-support
and trace frameworks used in \fid{F50}.

The Gaussian gauge tensors with a specified
Artin--Schreier line and an invertible constant
complex. The resulting kernel is the sheaf for
$2xy$, with an explicitly transported associator
and the symmetry correction required by its odd
shifts. No unspecified nonstandard sheaf theory
is needed for this operation.

The source remains the actual bialgebra/comodule
family. Its module/sheaf language is EGA~I,
Theorem~1.3.7, Corollary~1.3.8 and Theorem~1.4.1;
its coaction language is SGA~3, I, 4.7.2--3,
with cote~90, scan~17 as the audited archive anchor
\cite{EGA,SGA3,Archive}.
Those are framework references, not sources for
this particular convolution deduction.
The construction supplies coherent operations and
actual maps through all degrees together with the
original opposite-sign path cohomology.
The designated grading, full-source-Hom, simple-label,
canonical-basis and independent pointwise-weight
requirements remain distinct.

\paragraph{\fid{F54}, Part 1, in Grothendieck's language.}
The group algebra $\C[\Lambda]$ is the coordinate
ring of a split torus with character lattice
$\Lambda$. The element $B$ is a character, and
$U$ is the character of a genuine direct sum of
weight lines. Its algebraicity over the one
character forces its support onto that character's
rational line; the lattice index gives the common
power. This is the character/coaction language of
SGA~3, I, 4.7.2--3, with the Laurent and derivative
argument supplied explicitly here.

Under $K_0\operatorname{Rep}(T)=\Z[\Lambda]$, the
element $U$ belongs to the effective cone: its
coefficients are actual weight-line multiplicities.
Virtual classes can cancel them, exactly the step
excluded by the proof.
SGA~6, V, \S\S2--3, and cote~133, scan~18, supply
the previously audited $\lambda/\sigma$ framework
vocabulary \cite{SGA6,Archive}; no $\gamma$-filtration
positivity is inferred.

The algebraic parameter extension is the format
allowed by VIII, \S2.2, p.~7.
The regular matrix and centralizer language uses
EGA~I, Theorem~1.3.7, p.~86, Corollary~1.3.8,
p.~87, and Theorem~1.4.1, pp.~90--92.
Cote~90, scan~17, is only the audited action-framework
anchor \cite{VIII,EGA,SGA3,Archive}.
No manuscript leaf is claimed to contain the new
trace theorem, and no unspecified geometric
Frobenius comparison is being assumed.

\paragraph{\fid{F54}, Part 2, in Grothendieck's language.}
This is a finite family of modules over the actual
coherent center algebra, equipped with a chosen invertible
linear operator and an exact invariant-Hom action comparison.
The nilpotent center survives the collision even when
the two unmarked identity traces fail to determine it.
The affine module/sheaf framework is EGA~I,
Theorem~1.3.7, p.~86, Corollary~1.3.8, p.~87, and
Theorem~1.4.1, pp.~90--92.
The specified action is in the language of SGA~3, I,
4.7.2--3, with cote~90, scan~17 as the audited action
anchor \cite{EGA,SGA3,Archive}.
These framework references are distinct from the
explicit End/Hom and trace proofs supplied here.

\paragraph{\fid{F54}, Part 3, in Grothendieck's language.}
The marked insertion belongs to the actual coherent
algebra action, rather than being reconstructed from an
unmarked scalar trace.
Its trace-retuning law is a specified additional
arithmetic comparison.
The operator/representation language is SGA~3, I,
4.7.2--3, with cote~90, scan~17 as the audited framework
anchor; the underlying coherent module language is
EGA~I, Theorem~1.3.7, p.~86, Corollary~1.3.8, p.~87,
and Theorem~1.4.1, pp.~90--92
\cite{SGA3,EGA,Archive}.
A changing coefficient object belongs to the tensor
and Grothendieck-group side of the comparison, as in
SGA~6, V, \S\S2--3 and cote~133, scan~18.
No eigenvalue or purity theorem is inferred merely
from the existence of a coherent module.
The both-actions corollary is a statement about the centralizer of specified representation morphisms: the actual left and right operators generate the full endomorphism algebra on the marked sector. It requires those left operators themselves to be Weil morphisms, rather than maps twisted by coefficient lines. The same SGA~3 and cote~90 action dictionary carries its entire matrix proof.
The finite-power, scalar-function and two-word
obstructions have all been proved explicitly above.

\paragraph{\fid{F55}, Part 1, in Grothendieck's language.}
The constructive model is a finite coefficient cover $v=t^2$, followed by an explicitly specified finite localization, carrying the actual pulled-back coherent algebra, tensor module and right-module automorphism.
The affine module/sheaf dictionary is EGA~I, Theorem~1.3.7, p.~86, Corollary~1.3.8, p.~87, and Theorem~1.4.1, pp.~90--92.
The representation and commuting-morphism language is SGA~3, I, 4.7.2--3, with cote~90, scan~17 as the audited framework anchor \cite{EGA,SGA3,Archive}.
The two places above $v=1$ and the operator pole distinguishing their regularity claims are explicit algebraic data.
Neither the coherent model nor its two trace identities chooses the designated canonical basis, global integral form, or target positivity inequalities.

\paragraph{\fid{F55}, Part 2, strengthened, in Grothendieck's language.}
The normalization $B$ is a finite free algebra over
the local parameter ring, and multiplication by $f$
is an actual endomorphism whose characteristic
polynomial records the special-fiber scalar reduction.
The underlying affine module language is EGA~I,
Theorem~1.3.7, p.~86, Corollary~1.3.8, p.~87,
and Theorem~1.4.1, pp.~90--92 \cite{EGA}.
The exact normalization, freeness and place statements
are the primary Stacks lemmas cited in the proof.

The polynomial trace identity is represented by
specified finite direct sums and tensor powers
of endomorphisms. The alternating extension uses
an inserted Euler trace, not an effective-spectrum
or $\gamma$-positivity inference.
The established representation/Grothendieck framework
anchors remain SGA~3, I, 4.7.2--3 and
SGA~6, V, \S\S2--3, with cote~90, scan~17 and
cote~133, scan~18 \cite{SGA3,SGA6,Archive}.
The proof concerns this added all-places regularity
comparison with the actual lattice; it does not
identify a canonical basis or supply the independent
pointwise coefficient estimates.
The separately retained genuine-trace place observation uses the effective torus-character sum of F54. Its positive multiplicities use the SGA~6 coefficient-character language and cote~133, scan~18 already specified above; the strengthened all-places argument itself does not require effectivity.

\paragraph{\fid{F55}, Part 3, in Grothendieck's language.}
The coefficient functions lie on a finite cover of the
parameter torus. The family, action and marked
endomorphism are affine coherent-module data in
EGA~I, Theorem~1.3.7, p.~86, Corollary~1.3.8, p.~87,
and Theorem~1.4.1, pp.~90--92.
The representation and commuting-morphism dictionary
is SGA~3, I, 4.7.2--3, with cote~90, scan~17 as the
audited action anchor \cite{EGA,SGA3,Archive}.
After $v=t^L$, the genuine trace is an effective sum
of torus characters $t^{a_i}$.
Its positive multiplicities are the effective cone in
$K_0$ of torus representations; the tensor-character
language is SGA~6, V, \S\S2--3 and cote~133, scan~18
\cite{SGA6,Archive}.
This is not a positivity assertion for the
$\gamma$-filtration or a conjectural canonical basis.
The continuous branch over the positive real locus
is extra comparison data on the complex points of
the cover; the EGA module/sheaf equivalence does not
provide it. The full separation and phase-removal
proof was supplied above, rather than attributed to
an archival theorem. A virtual alternating trace
does not acquire the required positive multiplicities
merely by having a specified numerical dimension.

\paragraph{\fid{F56} in Grothendieck's language.}
The effective object monoid, its split group completion,
and its Karoubi completion are distinct constructions.
Here cancellation and the exact algebra image are proved,
and the completion changes the positive generators and
the relative order.
SGA~6, V, \S\S2--3, with cote~133, scan~18, provides the
audited representation/$\lambda$/$\sigma$ vocabulary;
it applies here to the symmetric coefficient category.
No $\lambda$-ring structure is asserted on the
noncommutative split $K_0(\mathcal P_{64})$, and
no $\gamma$-filtration positivity is inferred
\cite{SGA6,Archive}.

The actual coherent algebra and its action use the
affine module framework of EGA~I, Theorem~1.3.7, p.~86,
Corollary~1.3.8, p.~87, and Theorem~1.4.1, pp.~90--92,
and the coaction language of SGA~3, I, 4.7.2--3,
with cote~90, scan~17 \cite{EGA,SGA3,Archive}.
The normalization formula and all length computations
have been proved explicitly.
The perverse realization uses the actual Gaussian
convolution and primary compact-trace formalism
already supplied in \fid{F53}.
None of these framework identifications makes the
elementary matrix-kernel basis into the designated
nonstandard canonical basis.

\paragraph{\fid{F55}, the bar-compatible variant, in Grothendieck's language.}
The coefficient cover carries a specified involution
lifting inversion on the parameter torus, and the actual
coherent module automorphism is equivariant for this
involution and the actual tensor bar.
The affine module framework is EGA~I, Theorem~1.3.7,
p.~86, Corollary~1.3.8, p.~87, and Theorem~1.4.1,
pp.~90--92.
The representation/coaction framework is SGA~3, I,
4.7.2--3, with cote~90, scan~17 as the audited anchor
\cite{EGA,SGA3,Archive}.
The involution on the algebraic coefficient cover is
explicit additional equivariance data.
It does not select VIII's prescribed global integral
form or convert this mixed algebraic spectrum into
an integral-weight arithmetic realization.
The actual bar-fixed adapted basis is the proved
input from \fid{F43}; the invariant operator family
and all unit and trace identities are proved above.

\paragraph{\fid{F56}'s formal bar in Grothendieck's language.}
Finite graded coefficient spaces are split-torus
representations; their dual reverses the character lattice.
Evaluation and product pairings provide the actual maps
proved above. The representation framework is SGA~3,
I, 4.7.2--3, with cote~90, scan~17 \cite{SGA3,Archive}.
SGA~6, V, \S\S2--3 and cote~133, scan~18 supply the
coefficient representation/$\lambda$/$\sigma$ vocabulary,
not a $\lambda$-ring assertion for the noncommutative
word algebra \cite{SGA6,Archive}.
The actual localized module comparison uses EGA~I,
Theorem~1.3.7, p.~86, Corollary~1.3.8, p.~87, and
Theorem~1.4.1, pp.~90--92 \cite{EGA}.
This entrywise contravariant duality is distinct from
the transpose action adjoint, Verdier duality, and the
designated global-basis question. The new coherence and
inverse-map proofs were supplied explicitly; they are
not attributed to an archival theorem.

\paragraph{\fid{F57}'s trace obstruction in Grothendieck's language.}
Finite graded coefficient spaces are representations of
a split torus. Their effective character cone is
\(\mathbb N[v,v^{-1}]\); duality reverses the character,
and tensor product multiplies characters.
The representation framework is SGA~3,
I, 4.7.2--3, with cote~90, scan~17 \cite{SGA3,Archive}.
This is action language, not a new scheme-group realization
of the noncommutative algebra.
SGA~6, V, \S\S2--3 and cote~133, scan~18 supply the
representation and \(\lambda/\sigma\) vocabulary
\cite{SGA6,Archive}; no \(\lambda\)-ring structure on the
noncommutative word algebra is asserted.
The actual module, its dual, and its endomorphism matrices
are affine module data in EGA~I, Theorem~1.3.7, p.~86,
Corollary~1.3.8, p.~87, and Theorem~1.4.1,
pp.~90--92 \cite{EGA}.
The new calculation shows that the specified word trace
lies outside the effective cone with either overall sign.
It remains a valid virtual character, as the explicit
bounded Tate complex demonstrates. Its underlying bounded
coefficient complex fits SGA~6, I, Definition~4.7,
pp.~121--122, and Example~4.8, p.~122 \cite{SGA6};
the Frobenius action is additional specified coefficient data.
The grading-to-Frobenius comparison and an effective
summand projector remain extra hypotheses.
This is an effectivity
and comparison obstruction, not a failure of a
Riemann-hypothesis theorem or a source claim that every
nonstandard action matrix was strictly positive.

\section{Archival audit and the Fable leads}\label{sec:archive}
The original Montpellier PDFs add a cover. For the documents inspected here, a bottom-stamped archive scan $s$ is one-based PDF page $s+1$. Catalogue titles and scanned contents are distinguished. The following checks were visual readings of the specified leaves, not certifications of full manuscript proofs.

{\small\begin{longtable}{p{.19\linewidth}p{.36\linewidth}p{.37\linewidth}}
\toprule Cote and scan&Verified content&Consequence for this investigation\\\midrule\endhead
90, scan 17 (PDF 18)&Flat affine group over a base; group action; quasi-coherent equivariant modules.&Authentic anchor for the representation/group-scheme dictionary.\\
97, scan 4 (PDF 5)&Chern classes of regular representations and products over characters.&Characteristic-class anchor, not a plethysm-positivity theorem.\\
133, scan 18 (PDF 19)&Symmetric and exterior generating operations, including $\sigma_t(E)\lambda_{-t}(E)=1$.&Direct language anchor for the $\lambda$-ring calculation.\\
133, scans 19--21 (PDF 20--22)&Virtual finite $G$-sets and operations; scan 20 explicitly names plethysm.&These finite-set operations must not be identified with a proved representation-theoretic positive rule.\\
16, scans 35, 37, 39 (PDF 36, 38, 40)&Handwritten Lefschetz/primitive-summand/projector formulas, including hyperplane-section comparisons.&Relevant standard-conjecture framework. No identification with Alp\"oge's path sign is proved.\\
16, scans 34, 36, 38, 40&Reused typescript backs; scan 34 starts Karoubi's \emph{Cohomologie des cat\'egories de Banach}.&The Fable range ``34--40'' must be refined; these backs are not Grothendieck's handwritten Hodge-index proof.\\
16, scans 47--49 (PDF 48--50)&Letter to Coates: correspondences/duality; on scan 48, $C(X)$ plus $I(X\times X)$ implies correspondence-ring semisimplicity; on scan 49, deduction of Weil from $C$ and $A$, in the letter's notation.&Conditional standard-conjecture statements. The semisimplicity passage is scan 48, not scan 50.\\
16, scan 50 (PDF 51)&Separate handwritten endomorphism-algebra lemma.&Not the location of the preceding semisimplicity sentence.\\
1, archival scans 77--78 (PDF78--79), handwritten folios79--80&Exterior-power/trace-norm estimates; scan78 has a Fredholm determinant product and trace sum.&Generic operator-determinant framework is visible; the heavily revised manuscript proof is not certified.\\
1, archival scans109--110 (PDF110--111), handwritten folios111--112&Entire functions, order/genus and canonical products.&No $\Xi$ identity or RH-equivalent positive operator was verified.\\
1, literal archival scans79--80 and111--112 (PDF80--81,112--113)&Corrected determinant calculations, convex/supporting functions, and entire-function or positive-sequence convergence material.&Both numbering interpretations were checked; these are not interchangeable leaves.\\\bottomrule
\end{longtable}}
Cote~133's catalogue title concerns Witt groups, quadratic forms and formal groups; ``plethysm'' is a content lead, not its official title. Cote~60 was found in the catalogue and inspected at thumbnail level, but no precise theorem from it is used.

The Fable suggestions are treated as follows. Suggestion~3's norm estimate is proved and extended in \fid{F01}; its literal positive-Fourier-operator identification fails by \fid{F05}(e). Suggestion~4's identification of the cohomological sign with a polarized evanescent intersection form remains unproved: parity in a trace formula and a positive-definite form are different structures. Suggestion~8 is made precise by \fid{F13}, conditional on an actual functor and endpoint identification. No assertion about the unprovided earlier runs or their theorem IDs is imported.

Some published-source checks have a narrower scope than the archive checks.
The original EGA~I affine-module statements and their proof
contexts were read, and printed pp.~86, 87 and~90 were visually
checked: Theorem~1.3.7, Corollary~1.3.8 and Theorem~1.4.1
(the last proof continues through p.~92). Stacks Lemma~26.7.5,
tag~01IB, supplies a separately read modern formulation.
EGA~I, \S9.5, is reserved for scheme-theoretic images.
The original SGA~6, I, Definition~4.7, pp.~121--122, and
Example~4.8, p.~122, were visually read from the facsimile;
the primary Thomason locator check
\cite[\S3.1, pp.~8--9]{ThomasonPerfect} is retained as
corroboration. Expos\'e~V, \S\S2--3, remains background for
the $\lambda/\gamma$ formalism, not a fresh original-source
proof check. The specific locator V, Lemma~2.2.1, p.~315, is independently cited by Altman--Kleiman \cite[\S2.10]{AK}; it is a locator check, not a re-proof of SGA~6. The modern transcription consulted as a finding aid is not used as technical authority. All elementary $\lambda$ identities used in the paper have their displayed direct proofs.

\part{Unproved bridges, attempted routes, and current state}
\section{What a sufficient nonstandard package would have to contain}\label{sec:gates}
There are no unchecked affirmative ``new theorems'' in this part. The following is a sufficient research specification, not a quotation from Mulmuley and not a proved construction.

\begin{enumerate}
\item \gate{R: realization and normalization.} For each relevant nonstandard coefficient construct a finite-field Weil object or a defined categorical replacement, with explicit Frobenius, trace formula, degrees and weights. Identify $v$ with the correct grading, and specify signs, half-powers and rational rescalings.
\item \gate{S: spectral sign and duality.} Verify the pointwise input bound, the finite-spectrum identification, the output sign for unbounded $k$, and controlled background terms. F24 allows signed visible inputs at even $k$; its finite-extension version requires retained top-radius visibility. Supply either the Tate self-duality of \fid{F03} or the dual system of \fid{F17}. These inputs imply the stated purity theorem by the checked proof.
\item \gate{E: effective basis and operations.} Identify the normalized coefficient with a graded multiplicity of actual simple objects and prove compatibility with multiplication/comultiplication. Purity alone does not give this. The fixed inherited ambient cone is ruled out by \fid{F09}; F27 restores the actual rank-one character action in a different localized cone. Neither identifies VIII's basis.
\item \gate{L: Lefschetz and defects.} Prove the appropriate Lefschetz maps for unimodality. In general plethysm, define the nonflatness correction and prove the quantitative error bounds of the intended version of VIII 2.9/3.6. A spectral error bound such as \fid{F04} is not automatically that coefficient-vector error bound.
\item \gate{C: canonical and comparison data.} Prove basis existence/globalization, cell and quasi-cell claims, and the required specialization. For general plethysm retain VII's duality/reciprocity hypotheses; for Kronecker do not re-open the already proved revised-IV duality.
\item \gate{P: effective labels under a stated encoding.} Specify the input convention before asking for polynomial witnesses. Finding~\fid{F37} settles the equal-rectangle, two-row-output witness problem with explicitly listed rows. Finding~\fid{F34} excludes a general \(\#P\) claim under fully compact binary plethysm encoding. The attached homogeneous PH1 is already refuted by the cited Kahle--Micha{\l}ek theorem; it is not an open requirement to be proved. Any replacement using inhomogeneous polytopes or another counting model needs its own exact formulation, size bounds and algorithm. SH or the appropriate index bound remains separate.
\item \gate{O: actual obstructions.} Establish a strict multiplicity comparison for the orbit closures in an appropriate asymptotic family. Positive formulas and efficient decisions do not themselves produce that inequality.
\end{enumerate}
The spectral theorems remove an estimate once the explicit trace and sign hypotheses are supplied. F32 extends the upper-radius argument to a specified countable completion, but bounded visible reciprocal duality then forces finiteness. F19--F21 and F27 construct specified coherent or ordinary motivic realizations. F36 constructs action-compatible coherent Frobenius models and computes their descent defect. F35 closes the scalar stack-realization question for the actual normalized word trace; its dual does not satisfy the needed second bound. F38 excludes the proposed object-level tensor inverse in the usual derived Weil realization, and F39 excludes the direct nonsmooth collision category from the smooth proper noncommutative Weil framework. F40 repairs that gate by an explicit smooth proper corner category, retaining the exact specialized action; its whole-category TP spectrum is pure but does not yet realize the designated coefficient morphisms. These results narrow particular realizations; the designated based operations and arithmetic comparison remain open.  F43 supplies the actual positive joint action, full regular
commutant and tensor bar in rank three. F44 gives the relative
smooth corner and a precise Frobenius tradeoff. F45 supplies
projective word-dependent motives and hard Lefschetz, while proving
an obstruction when a specified commutant projector is retained.
F46 realizes the actual algebra, both actions and collision
extensions by pure Weil End/Hom data and coherent projective-line
objects; its chosen Frobenius fails the specified extension law
and constant-coefficient integrality comparison. F47 proves the
positive-adjoint radical obstruction. Mere existence of a finite
pure rank-three coefficient realization is therefore no longer
the gate; operation-specific sign comparisons, arithmetic grading and
designated basis compatibility remain. F48 supplies an actual all-degree faithful tensor family but excludes the stronger full point-Weil centralizer comparison. Its rational spectral theorem classifies the retuned characteristic polynomial and excludes both the single-weight generic and regular mixed-weight collision comparisons. F51 extends the integral-weight results to algebraic coefficient fields. Under full spectral retuning, F52 removes the weight restriction when the actual whole relative order is required to be the full regular centralizer: its second-order contact contradicts the first-order contact of distinct retuned algebraic eigencharacters. F49 proves that an actual positive commutant projector cannot descend to VIII's designated global form. F50 supplies the normalized perverse realization and actual opposite-sign path cohomology for effective finite coefficients, including the actual fixed-parameter End/Hom data and the graded word objects. In that specified construction the remaining pointwise and reciprocal estimates are exactly the desired spectral bounds; arbitrary inserted operator traces and a graded comparison through all tensor degrees remain separate; F56 supplies explicit maps for both actual rank-three presentation relations. F53 supplies coherent geometric tensor operations for the actual fixed-parameter family through all degrees, with precisely specified units and symmetry. Its identity-Frobenius character is still the dimension character, so it does not identify the designated simple basis or improve the weight estimate. F54 removes the integral-weight restriction from the algebraic common-radius identity-trace obstruction, constructs a genuine mixed-radius two-trace model retaining the whole actual order, and proves stronger marked-action limits. One regular central insertion excludes the specified pure or integrally mixed trace comparison on the actual tensor module without requiring full endomorphism recovery. Requiring the base Frobenius also to commute with the four named ordinary left operators removes that weight restriction: the marked rank-eight sector is jointly irreducible for the actual two actions. F55 constructs an actual fractional mixed model preserving the right action and both identity and marked traces at one selected collision place; its operator has poles at the other place. A second such family also commutes with an explicit extension of the actual tensor bar; this still does not supply all-places regularity. Marked retuning is excluded with arbitrary Jordan form and without any action-commutation or positivity assumption if the family is regular at every collision place; the same continuous positive-parameter branch also excludes it for genuine base-sector traces. These are explicit extra comparison hypotheses. The two fixed generator words $P_1$ and $P_1^2$ already exclude a composition-preserving parameter-retuned trace law. F56 now constructs the rank-three graded coefficient-object alternative with actual quadratic and degree-five maps, exact split Grothendieck algebra, and a perverse arithmetic lift. Its Karoubi completion enlarges the actual64-state algebra to rank104; the faithful six-state completion keeps generic rank10 but loses the collision contact. Its named adjoints also require enlargement. Entrywise coefficient duality nevertheless supplies a coherent formal bar, matching the actual localized tensor bar on split $K_0$ and the state classes, and inverting the specified relation maps; it is distinct from the transpose action adjoint. F57 excludes extending the strict effective-kernel construction simple by simple to the full rank-four algebra: an actual five-dimensional extra simple has a negative named-word trace. This extra simple is absent from the two-row quotient, and signed/factored positivity is not contradicted. Compatible graded kernels for larger whole modules or for the two-row quotient, a source-prescribed higher categorical presentation and designated simple/global labels remain. Word-dependent realizations remain another alternative. Replacing all seven by the single phrase ``nonstandard RH implies positivity'' hides independent claims.

For completeness, an effective-label claim has a simple formal implication: if the labels $\alpha$ in \eqref{eq:comparison} and witnesses for $m_\lambda^\alpha,n_\pi^\alpha$ have polynomial-length unique encodings and polynomial-time tests, then triples consisting of a label and the two witnesses count the sum of products. A nondeterministic machine guesses and checks the triple, giving a \(\#P\) function. This is just the definition of \(\#P\); neither the encodings nor a replacement for the refuted PH1 follows from the abstract multiplicity equality. No complexity-class collapse or lower bound is asserted.

\section{Approaches attempted and their exact outcomes}
{\small\begin{longtable}{p{.23\linewidth}p{.40\linewidth}p{.29\linewidth}}
\toprule Approach&Proved outcome or failure point&Next unresolved input\\\midrule\endhead
Literal source formulation&No finite-field nonstandard RH statement in the supplied versions; actual signed/near-positive targets located.&Choose a precise R/E target.\\
Path orthogonality and repeated Parseval&F01 gives the bound, but the constant trace example defeats purity without spectral signs.&S1/S2 realization.\\
Tensor powers and inverse roots&F02--F04 remove the fixed exponent loss; F17 allows separate finite extensions and a dual system.&Prove the concentration/error and duality hypotheses.\\
Fourier-positive operator&The bare Fourier-after-nonnegative-multiplication claim is false, F05(e).&A different operator and a proved identity to the coefficient.\\
Ordinary ambient purity&Stable subquotients inherit purity, F14. The actual inherited positive Hecke cone has only scalar image, F09.&A new effective realization/heart or different character map.\\
Change of rank-one positive model&F16 loses the actual action; F27 restores it in a localized four-character mixed Tate cone.&Identify the designated basis and operations; the repaired object is not pure of one weight.\\
Positive supports/crystals&Monomial pivots suffice, F10. Global positive disjoint supports and local disjointness alone fail.&Global exclusive pivots or another valid basis argument.\\
Orbit/Hilbert/Quot geometry&Integral models and ordinary IC purity exist; polarized coherent multiplicities are not determined by underlying Weil cohomology, F15.&A coherent-to-constructible multiplicity comparison.\\
Determinants and Gram forms&Positive projectors already express the ranks, F12; chosen-basis coefficient positivity fails in F08.&Image basis, independence and operation compatibility.\\
Theta and interlacing chains&F13 gives path/polytope counts for specified exact functors.&Identify the actual theta chain and endpoints with GCT restriction.\\
$\lambda$-rings and motives&Formal operations and Tate examples translate exactly; virtual effectivity and algebraic projectors remain extra.&Effective classes and the required realization, not only ring identities.\\\bottomrule
\end{longtable}}

\section{States of the five requested tasks}
\begin{enumerate}
\item \textbf{Statement needed.} The source-location and target audit is completed as A01 and \S\ref{sec:targets}. The demand for exact nonstandard Frobenius objects/weights cannot be met from those papers because they are not specified. The replacement package in \S\ref{sec:gates} is explicitly our sufficient specification.
\item \textbf{Template.} The standard geometric dictionary and every missing role are stated in \S\ref{sec:KL}. F19/F21 supply the actual rank-three coherent algebra and tensor module; F43 supplies both actions, bar and a positive joint basis; F46 supplies exact pure Weil End/Hom and coherent projective-line objects. F20 supplies ordinary pure Tate classical coefficients and F45 supplies projective word motives. The designated arithmetic grading and based realization remain open at R/E.
\item \textbf{Sign-only mechanism.} F01--F04, F17, F24 and the admissible countable extension F32 are proved and checked; omitted-hypothesis versions are refuted by the displayed countermodels. F26 excludes the finite normalized all-extension word dictionary. F35 supplies a genuine stack scalar realization with explicit cohomology, while F38 separates its reciprocal traces from actual tensor inverses. F43--F46 now retain specified actual operators through positive, projective, relative smooth and pure coefficient constructions. F50 supplies an exact normalized perverse realization and opposite-sign path cohomology for their effective coefficient objects. F53 adds coherent geometric tensor operations carrying the actual all-degree maps. The pointwise and reciprocal estimates remain equivalent to the desired coefficient weights, so an independently bounded varying local system remains to be constructed. F56 gives actual rank-three graded coefficient kernels with coherent presentation maps and formal bar, but its completion loses the original collision order. F54--F55 give the precise fixed-word, marked-trace, all-places and continuous-branch comparison obstructions, with explicit mixed countermodels outside those hypotheses. Designated arithmetic and basis comparisons remain separate. Thus the general transfer remains open at R/S/E, with explicit realizations and sign geometry now supplied in the stated modes.
\item \textbf{Positivity reached.} F31 gives actual primitive Tate motives for equal rectangles with two-row output. F37 gives polynomially verifiable compressed witnesses for the known marked-tree rule under the listed-row encoding. F34 proves the separate compact dimension and output-length obstructions. The supplied homogeneous PH1 is refuted by a cited prior theorem, including an equal-rectangle example. General designated canonical-basis positivity and strict orbit-closure multiplicity inequalities remain open at E/L/C/O; replacement effective rules require the revised P gate.
\item \textbf{Theta positivity.} F12--F13 and F23 state the exact projector/path consequences; F18 and F28 delimit integral-interlacing and Young-orbit versions. The theta/GCT functor and basis identification remain open. F08 refutes the generic determinant-to-basis-positivity shortcut.
\end{enumerate}

Research is continuing on the remaining gates under the scheduled-review rule. The present checkpoint is not a final return. No result here proves or refutes the intended full nonstandard positivity program, or $P\ne NP$.

\appendix
\section{Verification and reproducibility}
The independent spectral audit reconstructed F01, reproved F02 by Ces\`aro orthogonality rather than poles, and checked the two-limit argument of F03. The initial all-conjugate duality gap was recorded and repaired before the final theorem. Later full-paper reviews checked each proof, including the actual Hecke-cone example, the positive-pivot countermodels, the categorical and theta statements, and the scheme/base-change examples. Exact dates, reviewer runs, model identity, exposed effort information and check scopes are in \texttt{CHECK\_EVENTS.tsv}.

The file \texttt{check\_path\_bound.py} checks 54 small scalar/vector examples, including complex weights, for the exact path identity and the estimate. These floating-point checks are supplementary; the proof is F01. The optional exact Murnaghan--Nakayama script and character certificate compute a small rectangular coefficient and check character orthogonality; they are not used in any proof here and are not a Lean certificate. The source manifest records the unchanged attachments by SHA-256. Source reads and whole-proof checks are distinguished from cited background theorems and from formal verification.

The dated ledger records actual changes of gates. Scheduled checkpoints are due only at 45-minute intervals; scheduled progress reviews begin only after approximately three hours. The metadata records which such deadlines, if any, occurred; no review or checkpoint is invented to imply elapsed work that did not occur.

\begin{thebibliography}{99}
\bibitem{Alpoge} L.~Alp\"oge, \emph{Remark}, supplied four-page note, 20 November 2026; \url{https://alpo.ge/remark.pdf}. Main locators: p.~3 footnote~3 and steps~3--5; p.~4 steps~7--8. Supplied formalization entry points: \texttt{ALPOGE\_Main.lean}, \texttt{ALPOGE\_Statements.lean}; scan record \texttt{ALPOGE\_LEAN\_SCAN.tsv}. Results accepted as stipulated; missing imported development not executed.
\bibitem{IV} J.~Blasiak, K.~D.~Mulmuley and M.~Sohoni, \emph{Geometric Complexity Theory IV: nonstandard quantum group for the Kronecker problem}, arXiv:cs/0703110v4, 2013. \url{https://arxiv.org/abs/cs/0703110v4}. Exact statement locators are given in the text; attached source \texttt{Apr13KroneckerGCT4.tex}.
\bibitem{VI} K.~D.~Mulmuley, \emph{Geometric Complexity Theory VI: the flip via saturated and positive integer programming in representation theory and algebraic geometry}, arXiv:0704.0229v4, 2009. \url{https://arxiv.org/abs/0704.0229v4}.
\bibitem{VII} K.~D.~Mulmuley, \emph{Geometric Complexity Theory VII: Nonstandard quantum group for the plethysm problem}, arXiv:0709.0749v2, 2008. \url{https://arxiv.org/abs/0709.0749v2}.
\bibitem{VIII} K.~D.~Mulmuley, \emph{Geometric Complexity Theory VIII: On canonical bases for the nonstandard quantum groups}, arXiv:0709.0751v2, 2008. \url{https://arxiv.org/abs/0709.0751v2}.
\bibitem{IP} C.~Ikenmeyer and G.~Panova, \emph{Rectangular Kronecker coefficients and plethysms in geometric complexity theory}, arXiv:1512.03798v2, 2017; especially Theorems~1.4, 1.6, 1.8, 1.10. \url{https://arxiv.org/abs/1512.03798v2}.
\bibitem{BIP} P.~B\"urgisser, C.~Ikenmeyer and G.~Panova, \emph{No occurrence obstructions in geometric complexity theory}, arXiv:1604.06431v3, 2018; Theorem~1.4, \S1(c), Proposition~3.3, equation~(4.1), Proposition~4.5, Theorem~4.7. \url{https://arxiv.org/abs/1604.06431v3}.
\bibitem{KL} D.~Kazhdan and G.~Lusztig, \emph{Schubert varieties and Poincar\'e duality}, in \emph{Geometry of the Laplace Operator}, Proc. Sympos. Pure Math.~36, 1980, pp.~185--203; Theorem~4.3. The exact normalized form used here is checked through Haines's primary proof below.
\bibitem{Haines} T.~J.~Haines, \emph{A proof of the Kazhdan--Lusztig purity theorem via the decomposition theorem of BBD}; Theorems~2.0.1, 5.0.4, Proposition~3.0.2, \S4. \url{https://www.math.umd.edu/~tjh/KL_purity1.pdf}.
\bibitem{WeilII} P.~Deligne, \emph{La conjecture de Weil II}, Publ. Math. IH\'ES~52 (1980), pp.~137--252; Definitions~1.2.2 and~6.2.2--6.2.4, Theorem~3.3.1 and Corollaries~3.3.4--3.3.6. \url{https://www.numdam.org/item/PMIHES_1980__52__137_0/}.
\bibitem{BBD} A.~Beilinson, J.~Bernstein and P.~Deligne, \emph{Faisceaux pervers}, Ast\'erisque~100 (1982). Exact locators: 5.1.2, 5.1.8, 5.1.14, 5.3.1--5.3.2 (including the second proof), 5.3.8 and the following example, 5.4.4--5.4.5, Remark~5.4.9 (p.~143), 5.4.10 (p.~144), 6.2.5. \url{https://www.numdam.org/item/AST_1982__100__1_0/}.
\bibitem{Molev} A.~I.~Molev, \emph{Gelfand--Tsetlin bases for classical Lie algebras}, Theorem~2.1, equation~(2.2), patterns~(2.4); author-hosted preprint, printed pp.~7--9. \url{https://www.maths.usyd.edu.au/u/pubs/publist/preprints/2002/molev-18.pdf}.
\bibitem{Milne} J.~S.~Milne, \emph{Lectures on \'Etale Cohomology}, v2.21, 2013; Proposition~14.2, p.~90; Theorem~16.1, p.~108; Theorem~16.7 and preceding adjunction, p.~111; Example~16.3, p.~109; Theorem~17.10, p.~116; Proposition~18.4, pp.~121--122; Theorem~22.1, Remarks~22.3(b),(d) and Proposition~22.5, pp.~135--137; Theorem~23.2, pp.~139--140; Theorem~24.1(a),(b), p.~144. \url{https://www.jmilne.org/math/CourseNotes/LEC.pdf}.
\bibitem{EGA} A.~Grothendieck and J.~Dieudonn\'e, \emph{\'El\'ements de g\'eom\'etrie alg\'ebrique I: Le langage des sch\'emas}, Publ. Math. IH\'ES~4 (1960), Theorem~1.3.7, p.~86, Corollary~1.3.8, p.~87, and Theorem~1.4.1, pp.~90--92 (affine module/sheaf framework; original scan and proof contexts checked) and \S9.5 (scheme-theoretic images, with the affine instance proved in F15). \url{https://www.numdam.org/item/PMIHES_1960__4__5_0/}.
\bibitem{SGA3} M.~Demazure and A.~Grothendieck (eds.), \emph{Sch\'emas en groupes (SGA 3)}, Expos\'e~I by M.~Demazure, \S4.2, 4.7.2--4.7.3. Checked in the Gille--Polo re-edition, respecting explicitly marked editorial additions. \url{https://wstein.org/sga/SGA3/Exp1-20jan2010.pdf}.
\bibitem{SGA4} M.~Artin, A.~Grothendieck and J.-L.~Verdier (eds.), \emph{Th\'eorie des topos et cohomologie \'etale des sch\'emas (SGA 4)}, Expos\'es~XVII, \S3, and XVIII, 3.2.5. Background framework references.
\bibitem{SGA45} P.~Deligne, \emph{Cohomologie \'etale (SGA $4\frac12$)}, LNM~569, 1977; \emph{Rapport sur la formule des traces}, 3.2 (background trace-formula reference); \emph{Applications de la formule des traces aux sommes trigonom\'etriques}, 1.7--1.9, especially (1.9.2); 1.18(b)--(d), pp.~178--179; Propositions~4.2--4.3, pp.~196--197; Remark~4.4, p.~197. The latter original source statements and Gaussian proof were read in this run. \url{https://publications.ias.edu/sites/default/files/Number32.pdf}.
\bibitem{SGA5} A.~Grothendieck (ed.), \emph{Cohomologie \(\ell\)-adique et fonctions L (SGA 5)}, LNM~589, 1977; Expos\'e~III, 4.7. Background Lefschetz--Verdier reference.
\bibitem{SGA6} P.~Berthelot, A.~Grothendieck and L.~Illusie (eds.), \emph{Th\'eorie des intersections et th\'eor\`eme de Riemann--Roch (SGA 6)}, LNM~225, 1971; Expos\'e~I by L.~Illusie, Definition~4.7, pp.~121--122, and Example~4.8, p.~122 (perfect complexes; original facsimile visually checked); Expos\'e~V by P.~Berthelot, \S\S2--3 ($\lambda/\gamma$ operations; carried as background). Facsimile pages: \url{https://wstein.org/sga/sga/6/6t_121.jpeg} and \url{https://wstein.org/sga/sga/6/6t_122.jpeg}.
\bibitem{FGA} A.~Grothendieck, \emph{Techniques de construction et th\'eor\`emes d'existence en g\'eom\'etrie alg\'ebrique IV: les sch\'emas de Hilbert}, S\'eminaire Bourbaki, Expos\'e~221 (1961), \S3, Theorems~3.1--3.2, pp.~259--260. \url{https://www.numdam.org/article/SB_1960-1961__6__249_0.pdf}.
\bibitem{AK} A.~B.~Altman and S.~L.~Kleiman, \emph{Foundations of the theory of Fano schemes}, Compositio Math.~34 (1977), pp.~3--47; \S2.10, used only to corroborate the SGA~6 locator. \url{https://www.numdam.org/item/CM_1977__34_1_3_0/}.
\bibitem{Archive} Universit\'e de Montpellier, \emph{Archives Grothendieck}, cotes~1, 16, 90, 97, 133, scans specified in \S\ref{sec:archive}. Catalogue: \url{https://grothendieck.umontpellier.fr/archives-grothendieck/}; original files \url{https://grothendieck.umontpellier.fr/16.pdf}, \url{https://grothendieck.umontpellier.fr/90.pdf}, \url{https://grothendieck.umontpellier.fr/97.pdf}, \url{https://grothendieck.umontpellier.fr/133.pdf}. Cote~1 was checked at eight selected leaves, with both scan and handwritten-folio numbering as recorded in the text.
\bibitem{Fable} Claude Fable 5.2, effort max as reported in its supplied file, \emph{FABLE\_PROGRAM\_SUGGESTIONS\_20261120.md} and \emph{\ldots\_EXPLAINED.md}, suggestions~3, 4 and~8. Unchecked leads, not cited theorems.
\bibitem{Richarz} T.~Richarz, \emph{A new approach to the geometric Satake equivalence}, Documenta Mathematica 19 (2014), 209--246. Theorem~3.1, Proposition~4.1, Theorem~5.1, Lemma~5.4, and their proofs; the numbering differs from arXiv 1207.5314v2. \url{https://ems.press/content/serial-article-files/26224?nt=1}.
\bibitem{MV} I.~Mirkovi\'c and K.~Vilonen, \emph{Geometric Langlands duality and representations of algebraic groups over commutative rings}, Annals of Mathematics 166 (2007), 95--143; arXiv math/0401222v5 includes the later Appendix~B correction. Theorem~7.3, Lemma~7.1, Corollary~7.4. \url{https://arxiv.org/abs/math/0401222}.
\bibitem{StacksNil} The Stacks Project, Theorem~59.45.2, tag~04DZ, and Proposition~59.45.4, tag~03SI, topological invariance of the \'etale site/topos; Lemma~59.91.9, tag~0A4E, finite base change. \url{https://stacks.math.columbia.edu/tag/03SI}.
\bibitem{RS} T.~Richarz and J.~Scholbach, \emph{The motivic Satake equivalence}, Math. Ann.~380 (2021), 1595--1653; arXiv:1909.08322v3. Theorem~4.1, Propositions~4.8 and~5.9, Definition~6.2, Lemma~6.3, Corollary~6.4, Lemma~6.5, Theorem~6.8. \url{https://arxiv.org/abs/1909.08322}.
\bibitem{GW} R.~Goodman and N.~R.~Wallach, \emph{Symmetry, Representations, and Invariants}, author-hosted Chapter~4, Theorem~4.1.18 (Hermite reciprocity), p.~191. \url{https://mathweb.ucsd.edu/~nwallach/chapter4.pdf}.
\bibitem{DK} R.~W.~Donley, Jr., and W.~G.~Kim, \emph{A rational theory of Clebsch--Gordan coefficients}, arXiv:1707.03022v2; Proposition~3.1, equation~(3.4), p.~3; Proposition~8.1, equation~(8.2), p.~11. \url{https://arxiv.org/abs/1707.03022}.
\bibitem{Pieri} A.~Postnikov, \emph{MIT 18.315 lecture notes}, scribed by C.~Lian, 2014; Theorem~7.10, p.~16. Used for the classical horizontal and vertical Pieri rules. \url{https://web.mit.edu/clian/www/315_notes.pdf}.
\bibitem{EW} B.~Elias and G.~Williamson, \emph{The Hodge theory of Soergel bimodules}, Annals of Mathematics~180 (2014), 1089--1136; \S3.5, Theorem~3.5, pp.~1107--1108; Theorem~1.1, p.~1093. \url{https://annals.math.princeton.edu/wp-content/uploads/annals-v180-n3-p06-p.pdf}.
\bibitem{Manivel} L.~Manivel, \emph{A note on certain Kronecker coefficients}, arXiv:0809.3710v1 (2008); Theorem~1, p.~1; \S3, pp.~3--4, for the four-generator invariant algebra. \url{https://arxiv.org/abs/0809.3710}.

\bibitem{BlasiakRep} J.~Blasiak, \emph{Representation theory of the nonstandard Hecke algebra}, arXiv:1201.2209v2 (2012); Theorem~5.1, p.~14; Proposition~5.2, pp.~14--15; Definition~6.1 and Remark~6.2, p.~23. PDF equation~(25) has a different number in the converted HTML. \url{https://arxiv.org/abs/1201.2209v2}.
\bibitem{PPBounds} I.~Pak and G.~Panova, \emph{Bounds on certain classes of Kronecker and $q$-binomial coefficients}, author manuscript dated 13 October 2016; Lemma~1.3, p.~2. \url{https://www.math.ucla.edu/~pak/papers/QBC23.pdf}.
\bibitem{Brosnan} P.~Brosnan, \emph{On motivic decompositions arising from the method of Bia\l ynicki--Birula}, Invent. Math.~161 (2005), 91--111; inspected author PDF \S2, pp.~3--4, Theorems~3.3 and~3.5, pp.~5--6. \url{https://math.umd.edu/~pbrosnan/Papers/bb.pdf}.


\bibitem{KM} T.~Kahle and M.~Micha{\l}ek, \emph{Plethysm and lattice point counting}, Electronic Journal of Combinatorics~25(1) (2018), P1.41; Theorem~5 and Remark~6, p.~4; formulation history, pp.~2--4; rectangular identification, p.~6. \url{https://www.combinatorics.org/ojs/index.php/eljc/article/download/v25i1p41/pdf/}.
\bibitem{PPS} I.~Pak, G.~Panova and J.~P.~Swanson, \emph{A combinatorial interpretation for certain plethysm and Kronecker coefficients}, arXiv:2511.02312v3, uploaded 27 July 2026, title-page date 28 July 2026; Lemma~1.3, p.~2; \S2.2, p.~5; Definition~3.2 and Proposition~3.4, p.~6; Lemma~4.2, p.~9; Definition~4.4 and Theorem~4.5, p.~10. \url{https://arxiv.org/abs/2511.02312v3}.
\bibitem{SunStacks} S.~Sun, \emph{L-series of Artin stacks over finite fields}, Algebra \& Number Theory~6 (2012), 47--122; Theorem~1.1, p.~49; Definition~4.1, p.~84; Theorem~4.2, p.~85; Theorem~4.8 and \S4.8.2, pp.~89--90; Example~7.4, p.~109. \url{https://msp.org/ant/2012/6-1/ant-v6-n1-p03-s.pdf}.
\bibitem{LO2} Y.~Laszlo and M.~Olsson, \emph{The six operations for sheaves on Artin stacks II: Adic coefficients}, Publ. Math. IH\'ES~107 (2008), 169--210; \S9.1, Lemmas~9.1.1--9.1.2, p.~198; Theorem~11.0.14, p.~204. \url{https://numdam.org/item/10.1007/s10240-008-0012-5.pdf}.
\bibitem{TabuadaWeil} G.~Tabuada, \emph{Noncommutative Weil conjecture}, inspected version arXiv:1808.00950v7 (2019); \S1, pp.~2--3, Theorem~1.5; \S5.1, pp.~13--14; \S6.1, pp.~15--16; \S7, equations~(7.1)--(7.4), p.~17; Example~12.3, p.~24; Theorem~2.12, p.~8, and proof equation~(12.6), p.~26; \S6.2, equation~(6.9), p.~16. \url{https://arxiv.org/abs/1808.00950v7}.

\bibitem{MerkurjevSL} A.~S.~Merkurjev, \emph{Motivic decomposition of certain special linear groups}, inspected author PDF, \S5, Corollary~5.3 and notation, p.~7; Corollary~5.5, p.~8. \url{https://www.math.uni-bielefeld.de/LAG/man/530.pdf}.
\bibitem{EGNO2009} P.~Etingof, S.~Gelaki, D.~Nikshych and V.~Ostrik, \emph{Tensor Categories}, MIT 18.769 lecture notes (2009); Definition~1.42.2 and Proposition~1.42.4, pp.~76--77; Corollary~1.43.5, p.~80. \url{https://math.mit.edu/~etingof/tenscat1.pdf}.
\bibitem{StacksAffine} The Stacks Project, Lemma~26.7.5, tag~01IB (quasi-coherent sheaves on affines), and Definition~15.76.1, tag~0657 (perfect complexes). Primary statements read on 21 November 2026. \url{https://stacks.math.columbia.edu/tag/01IB}; \url{https://stacks.math.columbia.edu/tag/0657}.

\bibitem{ThomasonPerfect} R.~W.~Thomason, \emph{The classification of triangulated subcategories}, Compositio Math.~105 (1997), 1--27; \S3.1, pp.~8--9, corroborating SGA~6, I, 2.1 and 4.7 and II, 2.2.8. Used here only for the precise perfect-complex framework locator. \url{https://doi.org/10.1023/A:1017932514274}.

\bibitem{WeilI} P.~Deligne, \emph{La conjecture de Weil. I},
Publ. Math. IH\'ES~43 (1974), 273--307; Theorem~(1.6), p.~276,
and its derivation from Theorem~(1.7), pp.~276--277.
\url{https://www.numdam.org/item/PMIHES_1974__43__273_0.pdf}.

\bibitem{StacksRH} The Stacks Project, Lemma~53.12.2, tag~0C1D (Riemann--Hurwitz), and Lemma~53.12.4, tag~0C1F (different and tame ramification). Exact statements and the Section~53.12 derivation were read on 21 November 2026. \url{https://stacks.math.columbia.edu/tag/0C1D}; \url{https://stacks.math.columbia.edu/tag/0C1F}.

\bibitem{Run} This run, \texttt{GCT\_NSRH\_20261121}; findings A01 and F01--F57, with agent source memoranda, proof reviews, computations and \texttt{CHECK\_EVENTS.tsv} in the final package. This is the only current-run identifier used.
\end{thebibliography}
\end{document}
